% GENERATED FILE -- DO NOT EDIT.
%
% Produced by harness/assemble_contract.py from module.yaml and claims/*.tex.
% Edit the claim files and regenerate; `--check` fails if this file is stale.
\documentclass[11pt,a4paper]{article}
\pdftrailerid{885ca663e2d8fd768f4e5ab492db590b}
%% --- 1. Mathematics (order matters: physics after braket wins \ket, \tr) ---
\usepackage{amsmath}
\usepackage{amssymb}
\usepackage{amsfonts}
\usepackage{amsthm}
\usepackage{mathtools}
\usepackage{dsfont}
\usepackage{braket}
\usepackage{physics}
%% --- 2. Page furniture ------------------------------------------------------
\usepackage{xcolor}
\usepackage{array}    % >{...} column prefixes in the wide tables of some claims
\usepackage{enumitem}
\usepackage{graphicx}
\usepackage{booktabs}
\usepackage{longtable}
\usepackage{xltabular}              % front- and back-matter tables across pages
\usepackage{needspace}
\usepackage{xurl}
\usepackage[noautomatic]{imakeidx}  % defined-terms index; latexmk runs makeindex
\makeindex[title=Index of defined terms,intoc]
\usepackage{tcolorbox}
\tcbuselibrary{skins,breakable}
%% --- 3. Links and navigation (hyperref late; bookmark, refcount after it) ---
\usepackage{hyperref}
\usepackage{bookmark}
\usepackage{refcount}
%% --- 4. Colours: every distinction also has a shape (rule, caps, size, weight)
\definecolor{nnGrey}{gray}{0.40}      % non-normative text
\definecolor{metaGrey}{gray}{0.30}    % the generated meta line
\definecolor{techGrey}{gray}{0.45}    % heading of a technical claim
\definecolor{mainRule}{RGB}{0,110,140}
\definecolor{reqAccent}{RGB}{165,70,0}
\definecolor{linkBlue}{RGB}{0,62,150}
\definecolor{extLink}{RGB}{120,40,140}  % links into another PDF
\hypersetup{colorlinks=true, linkcolor=linkBlue, citecolor=linkBlue,
  urlcolor=linkBlue, filecolor=extLink, bookmarksnumbered=true,
  bookmarksopen=true, bookmarksopenlevel=1, pdfdisplaydoctitle=true}
% A link into another PDF (\href to a .pdf, \hyperref{file}..., xr-hyper) is a
% relative URI, <file>#nameddest=<destination>, in the file-link colour:
% browser viewers (Edge, Chrome) ignore hyperref's default GoToR action.
\makeatletter
\renewcommand*{\hyper@linkfile}[3]{% anchor text, file name, destination
  \begingroup
    \ifx\\#3\\\def\CQ@target{#2}\else\def\CQ@target{#2\hyper@hash nameddest=#3}\fi
    \let\@urlcolor\@filecolor
    \hyper@linkurl{#1}{\CQ@target}%
  \endgroup}
\makeatother
%% --- 5. The mapped Unicode set (harness/contract_navigation.py, UNICODE_TEX)
\DeclareUnicodeCharacter{03B1}{\ensuremath{\alpha}}
\DeclareUnicodeCharacter{03B2}{\ensuremath{\beta}}
\DeclareUnicodeCharacter{03B3}{\ensuremath{\gamma}}
\DeclareUnicodeCharacter{03B4}{\ensuremath{\delta}}
\DeclareUnicodeCharacter{03B5}{\ensuremath{\varepsilon}}
\DeclareUnicodeCharacter{03F5}{\ensuremath{\epsilon}}
\DeclareUnicodeCharacter{03B6}{\ensuremath{\zeta}}
\DeclareUnicodeCharacter{03B7}{\ensuremath{\eta}}
\DeclareUnicodeCharacter{03B8}{\ensuremath{\theta}}
\DeclareUnicodeCharacter{03D1}{\ensuremath{\vartheta}}
\DeclareUnicodeCharacter{03B9}{\ensuremath{\iota}}
\DeclareUnicodeCharacter{03BA}{\ensuremath{\kappa}}
\DeclareUnicodeCharacter{03BB}{\ensuremath{\lambda}}
\DeclareUnicodeCharacter{03BC}{\ensuremath{\mu}}
\DeclareUnicodeCharacter{03BD}{\ensuremath{\nu}}
\DeclareUnicodeCharacter{03BE}{\ensuremath{\xi}}
\DeclareUnicodeCharacter{03C0}{\ensuremath{\pi}}
\DeclareUnicodeCharacter{03C1}{\ensuremath{\rho}}
\DeclareUnicodeCharacter{03F1}{\ensuremath{\varrho}}
\DeclareUnicodeCharacter{03C3}{\ensuremath{\sigma}}
\DeclareUnicodeCharacter{03C2}{\ensuremath{\varsigma}}
\DeclareUnicodeCharacter{03C4}{\ensuremath{\tau}}
\DeclareUnicodeCharacter{03C5}{\ensuremath{\upsilon}}
\DeclareUnicodeCharacter{03C6}{\ensuremath{\varphi}}
\DeclareUnicodeCharacter{03D5}{\ensuremath{\phi}}
\DeclareUnicodeCharacter{03C7}{\ensuremath{\chi}}
\DeclareUnicodeCharacter{03C8}{\ensuremath{\psi}}
\DeclareUnicodeCharacter{03C9}{\ensuremath{\omega}}
\DeclareUnicodeCharacter{0393}{\ensuremath{\Gamma}}
\DeclareUnicodeCharacter{0394}{\ensuremath{\Delta}}
\DeclareUnicodeCharacter{0398}{\ensuremath{\Theta}}
\DeclareUnicodeCharacter{039B}{\ensuremath{\Lambda}}
\DeclareUnicodeCharacter{039E}{\ensuremath{\Xi}}
\DeclareUnicodeCharacter{03A0}{\ensuremath{\Pi}}
\DeclareUnicodeCharacter{03A3}{\ensuremath{\Sigma}}
\DeclareUnicodeCharacter{03A5}{\ensuremath{\Upsilon}}
\DeclareUnicodeCharacter{03A6}{\ensuremath{\Phi}}
\DeclareUnicodeCharacter{03A8}{\ensuremath{\Psi}}
\DeclareUnicodeCharacter{03A9}{\ensuremath{\Omega}}
\DeclareUnicodeCharacter{00E1}{\'a}
\DeclareUnicodeCharacter{00E0}{\`a}
\DeclareUnicodeCharacter{00E2}{\^a}
\DeclareUnicodeCharacter{00E4}{\"a}
\DeclareUnicodeCharacter{00E5}{\aa{}}
\DeclareUnicodeCharacter{00E9}{\'e}
\DeclareUnicodeCharacter{00E8}{\`e}
\DeclareUnicodeCharacter{00EA}{\^e}
\DeclareUnicodeCharacter{00EB}{\"e}
\DeclareUnicodeCharacter{00ED}{\'{\i}}
\DeclareUnicodeCharacter{00EF}{\"{\i}}
\DeclareUnicodeCharacter{00F3}{\'o}
\DeclareUnicodeCharacter{00F2}{\`o}
\DeclareUnicodeCharacter{00F4}{\^o}
\DeclareUnicodeCharacter{00F6}{\"o}
\DeclareUnicodeCharacter{00F8}{\o{}}
\DeclareUnicodeCharacter{00FA}{\'u}
\DeclareUnicodeCharacter{00F9}{\`u}
\DeclareUnicodeCharacter{00FB}{\^u}
\DeclareUnicodeCharacter{00FC}{\"u}
\DeclareUnicodeCharacter{00F1}{\~n}
\DeclareUnicodeCharacter{00E7}{\c{c}}
\DeclareUnicodeCharacter{00DF}{\ss{}}
\DeclareUnicodeCharacter{00C4}{\"A}
\DeclareUnicodeCharacter{00C5}{\AA{}}
\DeclareUnicodeCharacter{00C9}{\'E}
\DeclareUnicodeCharacter{00D6}{\"O}
\DeclareUnicodeCharacter{00DC}{\"U}
\DeclareUnicodeCharacter{2013}{\textendash{}}
\DeclareUnicodeCharacter{2014}{\textemdash{}}
\DeclareUnicodeCharacter{2018}{`}
\DeclareUnicodeCharacter{2019}{'}
\DeclareUnicodeCharacter{201C}{``}
\DeclareUnicodeCharacter{201D}{''}
\DeclareUnicodeCharacter{2026}{\ldots{}}
\DeclareUnicodeCharacter{00A0}{~}
\DeclareUnicodeCharacter{00A7}{\S{}}
\DeclareUnicodeCharacter{2264}{\ensuremath{\le}}
\DeclareUnicodeCharacter{2265}{\ensuremath{\ge}}
\DeclareUnicodeCharacter{2260}{\ensuremath{\neq}}
\DeclareUnicodeCharacter{2248}{\ensuremath{\approx}}
\DeclareUnicodeCharacter{2261}{\ensuremath{\equiv}}
\DeclareUnicodeCharacter{223C}{\ensuremath{\sim}}
\DeclareUnicodeCharacter{2243}{\ensuremath{\simeq}}
\DeclareUnicodeCharacter{2245}{\ensuremath{\cong}}
\DeclareUnicodeCharacter{221D}{\ensuremath{\propto}}
\DeclareUnicodeCharacter{226A}{\ensuremath{\ll}}
\DeclareUnicodeCharacter{226B}{\ensuremath{\gg}}
\DeclareUnicodeCharacter{2192}{\ensuremath{\to}}
\DeclareUnicodeCharacter{2190}{\ensuremath{\leftarrow}}
\DeclareUnicodeCharacter{2194}{\ensuremath{\leftrightarrow}}
\DeclareUnicodeCharacter{21D2}{\ensuremath{\Rightarrow}}
\DeclareUnicodeCharacter{21D0}{\ensuremath{\Leftarrow}}
\DeclareUnicodeCharacter{21D4}{\ensuremath{\Leftrightarrow}}
\DeclareUnicodeCharacter{21A6}{\ensuremath{\mapsto}}
\DeclareUnicodeCharacter{2208}{\ensuremath{\in}}
\DeclareUnicodeCharacter{2209}{\ensuremath{\notin}}
\DeclareUnicodeCharacter{2286}{\ensuremath{\subseteq}}
\DeclareUnicodeCharacter{2282}{\ensuremath{\subset}}
\DeclareUnicodeCharacter{2287}{\ensuremath{\supseteq}}
\DeclareUnicodeCharacter{2283}{\ensuremath{\supset}}
\DeclareUnicodeCharacter{2229}{\ensuremath{\cap}}
\DeclareUnicodeCharacter{222A}{\ensuremath{\cup}}
\DeclareUnicodeCharacter{2216}{\ensuremath{\setminus}}
\DeclareUnicodeCharacter{2205}{\ensuremath{\emptyset}}
\DeclareUnicodeCharacter{2200}{\ensuremath{\forall}}
\DeclareUnicodeCharacter{2203}{\ensuremath{\exists}}
\DeclareUnicodeCharacter{2115}{\ensuremath{\mathbb{N}}}
\DeclareUnicodeCharacter{2124}{\ensuremath{\mathbb{Z}}}
\DeclareUnicodeCharacter{211A}{\ensuremath{\mathbb{Q}}}
\DeclareUnicodeCharacter{211D}{\ensuremath{\mathbb{R}}}
\DeclareUnicodeCharacter{2102}{\ensuremath{\mathbb{C}}}
\DeclareUnicodeCharacter{00B7}{\ensuremath{\cdot}}
\DeclareUnicodeCharacter{00D7}{\ensuremath{\times}}
\DeclareUnicodeCharacter{2212}{\ensuremath{-}}
\DeclareUnicodeCharacter{00B1}{\ensuremath{\pm}}
\DeclareUnicodeCharacter{2213}{\ensuremath{\mp}}
\DeclareUnicodeCharacter{221A}{\ensuremath{\surd}}
\DeclareUnicodeCharacter{2211}{\ensuremath{\sum}}
\DeclareUnicodeCharacter{220F}{\ensuremath{\prod}}
\DeclareUnicodeCharacter{222B}{\ensuremath{\int}}
\DeclareUnicodeCharacter{221E}{\ensuremath{\infty}}
\DeclareUnicodeCharacter{2202}{\ensuremath{\partial}}
\DeclareUnicodeCharacter{2207}{\ensuremath{\nabla}}
\DeclareUnicodeCharacter{2297}{\ensuremath{\otimes}}
\DeclareUnicodeCharacter{2295}{\ensuremath{\oplus}}
\DeclareUnicodeCharacter{2218}{\ensuremath{\circ}}
\DeclareUnicodeCharacter{2308}{\ensuremath{\lceil}}
\DeclareUnicodeCharacter{2309}{\ensuremath{\rceil}}
\DeclareUnicodeCharacter{230A}{\ensuremath{\lfloor}}
\DeclareUnicodeCharacter{230B}{\ensuremath{\rfloor}}
\DeclareUnicodeCharacter{27E8}{\ensuremath{\langle}}
\DeclareUnicodeCharacter{27E9}{\ensuremath{\rangle}}
\DeclareUnicodeCharacter{2016}{\ensuremath{\|}}
\DeclareUnicodeCharacter{2113}{\ensuremath{\ell}}
\DeclareUnicodeCharacter{00B0}{\ensuremath{^{\circ}}}
\DeclareUnicodeCharacter{00B2}{\ensuremath{^{2}}}
\DeclareUnicodeCharacter{00B3}{\ensuremath{^{3}}}
\DeclareUnicodeCharacter{00B9}{\ensuremath{^{1}}}
\DeclareUnicodeCharacter{2080}{\ensuremath{_{0}}}
\DeclareUnicodeCharacter{2081}{\ensuremath{_{1}}}
\DeclareUnicodeCharacter{2082}{\ensuremath{_{2}}}
\DeclareUnicodeCharacter{207F}{\ensuremath{^{n}}}
\DeclareUnicodeCharacter{2096}{\ensuremath{_{k}}}
%% --- 6. Numbering: one counter per section -----------------------
\numberwithin{equation}{section}
\setcounter{tocdepth}{1}
%% --- 7. The project macro set, then the claim environments -----------------
% spec/macros.tex, read relative to this module's directory (latexmk -cd).
% =============================================================================
%  spec/macros.tex -- the one macro set every generated contract inputs.
%
%  Two parts: part A, the macros spec/notation.md catalogues, and part B,
%  the macros local to one module, grouped by module.
%  harness/assemble_contract.py emits % =============================================================================
%  spec/macros.tex -- the one macro set every generated contract inputs.
%
%  Two parts: part A, the macros spec/notation.md catalogues, and part B,
%  the macros local to one module, grouped by module.
%  harness/assemble_contract.py emits % =============================================================================
%  spec/macros.tex -- the one macro set every generated contract inputs.
%
%  Two parts: part A, the macros spec/notation.md catalogues, and part B,
%  the macros local to one module, grouped by module.
%  harness/assemble_contract.py emits \input{../../macros} for it, so the
%  file is read from spec/modules/<id>/.
%
%  RULES
%    * A macro that spec/notation.md catalogues is a verbatim copy of it
%      (part A). Changing one is a notation change: bump spec/notation.md.
%    * Part B holds the module-local macros. A macro defined by two modules
%      with the same expansion appears once, under the first module (in the
%      order below) that defines it. Two macros with the same name and
%      DIFFERENT expansions are both kept under a module prefix, and the
%      claims write the prefixed names:
%        \Chain  -> \redChain (reductions, \mathfrak{C}: the reduction chain)
%                   \pcChain  (protocol-compiler, \mathfrak{R}: the two-segment
%                             chain; \Chsrc and \Chrec are built on it)
%        \Nmsm   -> \redNmsm  (reductions, \mathcal{N}: a POVM measurement
%                             channel)
%                   \pcNmsm   (\mathcal{N}_{\mathsf A}: Alice's measurement
%                             channel, written by
%                             transmitter/thm:certificate-sufficiency,
%                             transmitter/rem:xreg and
%                             reductions/lem:source-transfer)
%    * Deliberately omitted: a \provenance{tag}{details} macro. Provenance
%      is a header field of every claim file (spec/claim-format.md section
%      3); a claim body that calls such a macro fails to compile, which is
%      the intended signal.
%
%  REQUIRED PACKAGES, loaded by the assembler's preamble before this file:
%    amsmath, amssymb, amsfonts, mathtools, dsfont, braket, physics, xcolor
%    (braket before physics, so that physics wins for \ket, \bra, \tr, ...).
%
%  This file contains only definitions: no \usepackage, no document body.
% =============================================================================


% =============================================================================
%  PART A -- notation macros, the ones spec/notation.md catalogues, each a
%  verbatim copy of its catalogue row; grouped under numbered headings.
% =============================================================================

% -----------------------------------------------------------------------------
% 1. State spaces and operator sets (a "system label" -> spaces of operators)
% -----------------------------------------------------------------------------
\newcommand{\Hs}[1]{\mathcal{H}(#1)}        % Hilbert space of a system label
\newcommand{\Lin}[1]{\mathcal{L}(#1)}       % bounded linear operators on H(#1)
\newcommand{\Linn}[2]{\mathcal{L}(#1,#2)}   % linear maps H(#1) -> H(#2)
\newcommand{\Pos}[1]{\mathrm{Pos}(#1)}      % positive semidefinite operators on H(#1)
\newcommand{\Density}{\mathrm{D}}           % density operators: \Density(A) = D(A)
\newcommand{\Seq}[1]{S_{=}(#1)}             % normalized states (trace = 1)
\newcommand{\Sleq}[1]{S_{\leq}(#1)}         % subnormalized states (0 <= trace <= 1)
\newcommand{\id}{\mathds{1}}                % identity operator
\newcommand{\idon}[1]{\mathds{1}_{#1}}      % identity on a named system, e.g. \idon{C}

% DEPRECATED: applies \mathcal to both arguments itself, so it takes bare
% letters, not system labels. Prefer \Linn; do not use in new contracts.
\providecommand{\Lop}[2]{\mathcal{L}(\mathcal{#1}, \mathcal{#2})}


% -----------------------------------------------------------------------------
% 2. Composite systems
% -----------------------------------------------------------------------------
\newcommand{\dsum}{\oplus}                  % direct sum of system labels
% Tensor product: use \otimes directly.
% n-fold direct sum of an operator space: \Lin{A}^{\oplus n}


% -----------------------------------------------------------------------------
% 3. Events, conditioning and classical registers (for STATES)
%    Convention:  \wedge  -> subnormalized "and-event" component, e.g. \rho_{A\wedge x}
%                 |        -> normalized conditional state,        e.g. \rho_{A|x}
% (Write these directly in subscripts; no extra macros needed.)
% -----------------------------------------------------------------------------


% -----------------------------------------------------------------------------
% 4. Photon-number subspaces / direct-sum blocks
%    Convention:  operators carry the restriction "on top", in brackets;
%                 states carry it as a \wedge subscript (see Sec. 3).
% -----------------------------------------------------------------------------
\newcommand{\Pnum}[1]{P_{#1}}               % photon-number projector (system from context)
\newcommand{\Pnumon}[2]{P_{#1}^{#2}}        % with explicit system superscript, e.g. P_n^C
\newcommand{\blk}[2]{#1^{(#2)}}             % operator block, e.g. \blk{M}{\leqN} = M^{(<=N)}
\newcommand{\leqN}{\leq N}                  % use as \blk{M}{\leqN}
\newcommand{\gtN}{> N}                      % use as \blk{M}{\gtN}
\newcommand{\eqn}{= n}                      % use as \blk{M}{\eqn}


% -----------------------------------------------------------------------------
% 5. System labels used in the protocol
% -----------------------------------------------------------------------------
\newcommand{\Chan}{C}                       % quantum-channel (Eve) input to Bob
\newcommand{\Prb}{P}                        % probe input
\newcommand{\Flag}{{\mathrm{Flg}}}          % flag system (EXISTS ONLY AFTER squashing)
\newcommand{\Loss}{{\mathrm{Lss}}}          % loss system
\newcommand{\Vc}{{\mathrm{Vcm}}}            % vacuum system
\newcommand{\Nout}{N_{\mathrm{out}}}
\newcommand{\Nin}{N_{\mathrm{in}}}


% -----------------------------------------------------------------------------
% 6. Alphabets
% -----------------------------------------------------------------------------
\newcommand{\Yset}{\mathcal{Y}}             % Bob's (coarse-grained) detection alphabet
\newcommand{\Yraw}{\mathcal{Y}_{\mathrm{raw}}}% raw click patterns {0,1}^4
\newcommand{\Yscset}{\mathcal{Y}_{\singlec}}% single-click outcomes
\newcommand{\Ymcset}{\mathcal{Y}_{\mc}}     % multi-click outcomes
\newcommand{\Tset}{\mathcal{T}_{\mathrm{set}}}% private setting set (the probe settings in the optical instance)
\newcommand{\Xset}{\mathcal{X}}             % Alice's alphabet
\newcommand{\Dset}{\{H,V,+,-\}}             % polarization / detector labels
\newcommand{\pass}{\top}                    % pass sifting


% -----------------------------------------------------------------------------
% 7. Mode / outcome abbreviations
% -----------------------------------------------------------------------------
\newcommand{\on}{{\mathrm{on}}}
\newcommand{\off}{{\mathrm{off}}}
\newcommand{\dc}{{\mathrm{dc}}}             % dark count
\newcommand{\sq}{\mathrm{sq}}               % squashed
\newcommand{\mc}{\mathrm{mc}}               % multi-click
\newcommand{\nc}{\mathrm{nc}}               % no-click
\newcommand{\singlec}{\mathrm{sc}}          % single-click
\newcommand{\vac}{\mathrm{vac}}             % vacuum
\newcommand{\kout}{{k_\mathrm{out}}}
\newcommand{\kin}{{k_\mathrm{in}}}


% -----------------------------------------------------------------------------
% 8. Maps, channels and miscellaneous operators
% -----------------------------------------------------------------------------
\newcommand{\Prob}{\mathrm{Prob}}           % probability (use \Prob[\cdot] or \Prob(\cdot))
\newcommand{\lmin}{\lambda_{\min}}          % minimum eigenvalue
\newcommand{\Wflag}{W_{\Flag}}              % flag-space weight
\newcommand{\transpose}{\mathsf{T}}
\newcommand{\innerp}[2]{\left\langle #1 \vert #2 \right\rangle}
\renewcommand{\vec}[1]{\boldsymbol{#1}}     % NB: overrides the LaTeX arrow accent


% -----------------------------------------------------------------------------
% 9. Key-rate numerics
% -----------------------------------------------------------------------------
\newcommand{\cG}{\mathcal{G}}               % protocol map
\newcommand{\cZ}{\mathcal{Z}}               % pinching map
\newcommand{\Amap}{\mathcal{A}_{\mathrm{ann}}} % announcement + measurement map
\newcommand{\Amapnb}{\mathcal{A}_{\mathrm{ann}|\mathrm{nb}}}
\newcommand{\simpl}{}                       % empty suffix for simplified block maps
\newcommand{\acSet}{\mathcal{S}_{\mathrm{acc}}} % acceptance set
\newcommand{\testE}{\mathcal{E}_{\mathrm{test}}} % test-statistics map
\newcommand{\ecLeak}{\mathrm{leak}^{\mathrm{EC}}_{\mathrm{obs}}}
\newcommand{\pgen}{p_{\mathrm{gen}}}
\newcommand{\askeyr}{R_{\infty}}             % asymptotic key rate
\newcommand{\ppass}{p_{\mathrm{pass}}}       % sifting pass probability


% -----------------------------------------------------------------------------
% 10. Finite-size analysis / Renyi entropies
% -----------------------------------------------------------------------------

% --- Round types: test vs. key generation (register T, alphabet \Fset) --------
\newcommand{\tst}{\mathrm{test}}            % test round
\newcommand{\gen}{\mathrm{gen}}             % key-generation round
\newcommand{\Fset}{\mathcal{T}}             % round-type alphabet {test, gen} (not \Tset)
\newcommand{\Chat}{\widehat{C}}             % acceptance-test register
\newcommand{\Chatset}{\widehat{\mathcal{C}}}% its alphabet

% --- Basis / key extraction functions ----------------------------------------
\newcommand{\bA}{b_{A}}                     % Alice's basis-extraction function
\newcommand{\bB}{b_{B}}                     % Bob's basis-extraction function
\newcommand{\kB}{k_{B}}                     % Bob's key-extraction function

% --- Maps --------------------------------------------------------------------
\newcommand{\Mprot}{\mathcal{M}}            % single-round protocol map
\newcommand{\Phimap}{\boldsymbol{\Phi}}     % test-statistics map
\newcommand{\Pdeph}{\mathcal{P}}            % dephasing (pinching) map on \Blreg

% --- Isometries --------------------------------------------------------------
\newcommand{\Vmsm}{V_{\mathrm{msm}}}        % measurement isometry
\newcommand{\Vann}{V_{\mathrm{ann}}}        % announcement isometry
\newcommand{\VI}{V_{I}}                     % Eve's copy of the announcements
\newcommand{\VChat}{V_{\Chat}}              % acceptance-data isometry
\newcommand{\VR}{V_{R}}                     % sifting + key map -> raw key R
\newcommand{\VZ}{V_{Z}}                     % raw key -> key register Z
\newcommand{\Wprot}{W}                      % full protocol isometry

% --- Renyi entropies ---------------------------------------------------------
% SECURITY-CRITICAL: see spec/notation.md §4. Never change an arrow to
% make an inequality work out; never convert silently between the sandwiched
% (tilde) and Petz (overline) families.
\newcommand{\Hsandup}{\widetilde{H}^{\uparrow}_{\alpha}}   % sandwiched, up
\newcommand{\Hsanddown}{\widetilde{H}^{\downarrow}_{\alpha}}% sandwiched, down
\newcommand{\Hpetzup}{\overline{H}^{\uparrow}_{\alpha}}    % Petz, up
\newcommand{\Hpetzdown}{\overline{H}^{\downarrow}_{\alpha}}% Petz, down
\newcommand{\hup}{h^{\uparrow}_{\alpha}}                   % single-round rate
\newcommand{\hqkd}{h^{\uparrow,\mathrm{QKD}}_{\alpha}}     % single-round QKD rate
\newcommand{\hlow}{\underline{h}_{\alpha}}                 % lower bound on \hqkd
% --- spec/notation.md 3.5 and 4.5 -------------------------------------------
% \hlb  is the first-order linearization lower bound on \hlow: the headline
%       key rate (the `first-order` tier of foundations/def:tiers). NOT \hlow
%       itself.
% \tnum is the instance's one declared numerical tolerance. Its
%       near-collision with the declared marginal \tau_A is disclosed in
%       spec/notation.md 3.5; bare \tau is reserved and denotes neither.
\newcommand{\hlb}{\underline{h}^{\mathrm{FO}}_{\alpha}}    % first-order lower bound on \hlow
\newcommand{\tnum}{\tau_{\mathrm{num}}}                    % declared numerical tolerance


% =============================================================================
%  PART B -- module-local macros, grouped by module. Deprecated aliases,
%  defined so that a text written with them compiles, used in no new contract: \Lop,
%  \testE, \Pdeph, \pass, \ecLeak (spec/notation.md section 3.9).
% =============================================================================

% --- meat-core -----------------------------------------------------------------
\newcommand{\SecSet}{\mathbb{S}}            % sector index set of Q
\newcommand{\Qsys}{Q}                       % declared sector-carrying measured system
\newcommand{\Vset}{\mathcal{V}}             % declared joint outcome alphabet of the Q-POVM
\newcommand{\CAset}{\mathcal{C}_{A}}        % Alice announcement alphabet
\newcommand{\CBset}{\mathcal{C}_{B}}        % Bob announcement alphabet
\newcommand{\Party}{\mathfrak{P}}           % declared party partition
\newcommand{\Oset}[1]{\mathcal{O}_{#1}}     % readout alphabet of party #1
\newcommand{\rd}[1]{\pi_{#1}}               % readout map of party #1
\newcommand{\Sadm}{\mathcal{S}_{A}}         % admissible Alice subspace
\newcommand{\nomarg}{\mathsf{none}}         % "no trusted marginal declared"
\newcommand{\Zset}{\mathcal{Z}_{\perp}}     % raw-key alphabet incl. \perp
\newcommand{\Zval}{\mathcal{Z}}             % key alphabet without \perp
\newcommand{\Etil}{\widetilde{E}}           % single-round purifying register
\newcommand{\Ehat}{\widehat{E}}             % Eve's n-round register
\newcommand{\OmAT}{\Omega_{\mathrm{AT}}}    % acceptance event
\newcommand{\epa}{\varepsilon_{\mathrm{PA}}}
\newcommand{\eev}{\varepsilon_{\mathrm{EV}}}
\newcommand{\lEC}{\lambda_{\mathrm{EC}}}
\newcommand{\Decl}{\mathfrak{D}}            % protocol declaration
\newcommand{\CList}{\mathfrak{L}}           % declared constraint list
\newcommand{\PinS}{\mathcal{P}_{\SecSet}}   % sector pinching on Q
\newcommand{\Feas}{\mathcal{W}}             % feasible set of block states
\newcommand{\Chkey}{\Chatset_{\mathrm{key}}}% key-live acceptance values
\newcommand{\Kann}{\mathcal{K}}             % announcement-class alphabet
\newcommand{\hold}{\mathsf{h}}              % declared key holder (A or B)
\newcommand{\sfD}{D}                        % classical KL divergence (base 2)
\newcommand{\Dtwo}{D}                       % quantum relative entropy (base 2)
\newcommand{\Dsand}{\widetilde{D}_{\alpha}} % sandwiched Renyi divergence
\newcommand{\idmap}{\mathrm{id}}            % identity channel
\newcommand{\Aev}{\mathcal{A}}              % Eve's per-round attack channel
\newcommand{\AttSet}{\mathfrak{A}}          % licensed attack-channel set
\newcommand{\Chacc}{\Chatset_{\mathrm{acc}}}% acceptance support of \acSet
\newcommand{\epsstar}{\varepsilon^{\star}}      % distance to an exact POVM
\newcommand{\dfH}{\Delta_{\mathrm{herm}}}       % Hermiticity defect
\newcommand{\dfP}{\Delta_{\mathrm{pos}}}        % negative part
\newcommand{\dfC}{\Delta_{\mathrm{cpl}}}        % completeness defect
\newcommand{\lmoq}{g}                       % acceptance-side gradient block
\newcommand{\lmoJ}{G}                       % feasible-side gradient block
\newcommand{\bacc}{\beta_{\mathrm{acc}}}    % certified acceptance-side minimum
\newcommand{\Itset}{\mathfrak{I}}           % finite index set of iterates (C22)

% --- reductions ----------------------------------------------------------------
\newcommand{\Bsys}{B}                        % physical (pre-squashing) system
\newcommand{\Tcl}[1]{\mathcal{T}(#1)}        % trace-class operators on H(#1)
\newcommand{\Ins}{\mathcal{I}}               % generalized-instrument branch map
\newcommand{\redNmsm}{\mathcal{N}}              % POVM measurement channel
\newcommand{\Eprot}{\mathcal{E}}             % n-round protocol channel
\newcommand{\Grnd}{\mathsf{G}}               % one round map of a protocol
\newcommand{\Kfin}{\mathsf{K}}               % end-of-protocol processing
\newcommand{\Crit}{\mathsf{sec}}             % value of a security criterion
\newcommand{\Pflag}{\Pi_{\Flag}}             % flag-sector projector on Q
\newcommand{\Pleq}{\Pi_{\leqN}}              % preserved-block projector
\newcommand{\Pgt}{\Pi_{\gtN}}                % above-cutoff projector
\newcommand{\Sph}{S^{\mathrm{ph}}}           % family of physical inputs
\newcommand{\Srch}{S^{\mathrm{rch}}}         % reachable squashed states
\newcommand{\Ssq}{S^{\sq}}                   % squashed admissible set
\newcommand{\Ffss}{F^{\mathrm{FSS}}}         % flag-state-squasher target
\newcommand{\Fwp}{F^{\mathrm{WP}}}           % multi-bound WPFSS target
\newcommand{\Mwp}{M^{\mathrm{WP}}}           % renormalized above-cutoff POVM
\newcommand{\Mdec}{{M'}^{\mathrm{ph}}}       % declared physical family on B
\newcommand{\Ppvm}{P^{\Flag}}                % flag PVM
\newcommand{\Lmp}{{\Lambda_{\mathrm{mp}}}}   % measure-and-prepare channel
\newcommand{\Lwp}{{\Lambda^{\mathrm{WP}}}}   % multi-bound WPFSS channel
\newcommand{\Lfss}{{\Lambda^{\mathrm{FSS}}}} % flag-state-squasher channel
\newcommand{\lmax}{\lambda_{\max}}           % maximum eigenvalue
\newcommand{\Lnoise}{\widetilde{\Lambda}}    % flag-noise channel
\newcommand{\cg}{\varpi}                     % declared coarse-graining map
\newcommand{\MCK}{\mathsf{MC}}               % multiclick raw-pattern set
\newcommand{\redChain}{\mathfrak{C}}            % reduction chain
\newcommand{\Lnk}{\ell}                      % a single typed link
\newcommand{\dsq}{\preceq}                      % the squashing relation
\newcommand{\opn}[1]{\left\lVert #1\right\rVert}           % operator (spectral) norm
\newcommand{\mxn}[1]{\left\lVert #1\right\rVert_{\max}}    % entrywise max modulus
\newcommand{\Fld}{\mathbb{F}}                   % exact ordered substrate field
\newcommand{\Fldi}{\mathbb{F}_{i}}              % its Gaussian extension F + iF
\newcommand{\lat}[1]{\big\lfloor #1\big\rfloor_{\kappa}} % round DOWN onto the lattice

% --- protocol-compiler ---------------------------------------------------------
\newcommand{\Spec}{\mathfrak{S}}             % protocol specification (compiler input)
\newcommand{\pcChain}{\mathfrak{R}}            % reduction chain (two typed segments)
\newcommand{\Chsrc}{\pcChain_{\mathrm{src}}}   % its Alice-side source-map segment
\newcommand{\Chrec}{\pcChain_{\mathrm{rec}}}   % its receiver-side squashing segment
\newcommand{\rvirt}{\rho^{\mathrm{v}}}       % virtual (PRE-source-map) signal state
\newcommand{\Fph}{F^{\mathrm{ph}}}           % physical (pre-reduction) POVM element
\newcommand{\Vph}{\mathcal{V}^{\mathrm{ph}}} % physical outcome alphabet
\newcommand{\Tpub}{\mathcal{T}_{\mathrm{pub}}}% public-choice alphabet
\newcommand{\Uset}{\mathcal{U}}              % residual outcome alphabet
\newcommand{\idxm}{\mathsf{idx}}             % label-to-index map
\newcommand{\Asig}{A'}                       % transmitted (signal) system -- the REAL one
\newcommand{\Avirt}{A''}                     % virtual transmitted system (source descriptor of Xi_1)
\newcommand{\Abar}{\bar{A}}                  % Alice's measured register
\newcommand{\Ahat}{\hat{A}}                  % Alice's shield register
\newcommand{\pcNmsm}{\mathcal{N}_{\mathsf A}}  % Alice's measurement channel (see the RULES block)
\newcommand{\Phip}{\Phi^{+}}                 % maximally entangled vector label
\newcommand{\Smap}{\mathfrak{s}}             % the emitted sector map (fraktur: a compile-time object)

% --- receiver ------------------------------------------------------------------
\newcommand{\Modes}{\mathbb{M}}              % declared ordered mode set
\newcommand{\Inp}{\mathbb{I}}                % declared input modes
\newcommand{\InpC}{\mathbb{I}_{\Chan}}       % channel input modes
\newcommand{\InpS}{\mathbb{I}_{\Prb}}        % probe (source) input modes
\newcommand{\Env}{\mathbb{E}}                % dilation (environment) modes
\newcommand{\Port}[1]{\mathbb{P}_{#1}}       % mode set of detector port #1
\newcommand{\Emd}[1]{E_{#1}}                 % coordinate projector on modes #1
\newcommand{\Gmat}{G}                        % instrument (transfer) matrix
\newcommand{\Fock}{\mathfrak{F}}             % Fock / second-quantization functor
\newcommand{\Symn}[1]{\mathrm{Sym}^{#1}}     % symmetric tensor power
\newcommand{\normord}[1]{{:}\mspace{2mu}#1\mspace{2mu}{:}}   % normal ordering
\newcommand{\rk}{\operatorname{rank}}        % rank (physics.sty owns \rank)
\newcommand{\SCK}{\mathsf{SC}}               % single-click pattern set
\newcommand{\Recv}{\mathfrak{R}}             % receiver instrument tuple
\newcommand{\Cat}{\mathfrak{C}}              % component catalogue
\newcommand{\Grph}{\mathfrak{E}}             % declared element graph
\newcommand{\defic}{\mathbb{L}}              % loss deficiency operator
\newcommand{\Pvac}[1]{\Pi_{\vac}(#1)}        % joint no-click element on ports #1
\newcommand{\pmc}{p_{\mc}}                   % uniform high-photon mc weight
\newcommand{\Bnd}{\beta}                     % the scalar bound function beta(n)
\newcommand{\Rot}{R}                         % planar rotation matrix
\newcommand{\Wplate}{W}                      % retarder (wave-plate) matrix
\newcommand{\Bsp}{B}                         % beamsplitter core matrix

% --- statistics ----------------------------------------------------------------
\newcommand{\honest}{p^{\mathrm{hon}}}          % honest per-round distribution
\newcommand{\ecomp}{\varepsilon_{\mathrm{comp}}}% completeness budget
\newcommand{\simplex}[1]{\Delta(#1)}            % probability simplex
\newcommand{\Chatcon}{\Chatset_{\mathrm{con}}}  % constrained outcomes
\newcommand{\Ireg}[3]{\mathcal{I}_{#1}(#2,#3)}  % regularized incomplete beta
\newcommand{\Ncount}[1]{\mathsf{N}_{#1}}        % per-outcome count
\newcommand{\qlo}{\mathbf{q}^{-}}               % box lower bounds
\newcommand{\qhi}{\mathbf{q}^{+}}               % box upper bounds
 for it, so the
%  file is read from spec/modules/<id>/.
%
%  RULES
%    * A macro that spec/notation.md catalogues is a verbatim copy of it
%      (part A). Changing one is a notation change: bump spec/notation.md.
%    * Part B holds the module-local macros. A macro defined by two modules
%      with the same expansion appears once, under the first module (in the
%      order below) that defines it. Two macros with the same name and
%      DIFFERENT expansions are both kept under a module prefix, and the
%      claims write the prefixed names:
%        \Chain  -> \redChain (reductions, \mathfrak{C}: the reduction chain)
%                   \pcChain  (protocol-compiler, \mathfrak{R}: the two-segment
%                             chain; \Chsrc and \Chrec are built on it)
%        \Nmsm   -> \redNmsm  (reductions, \mathcal{N}: a POVM measurement
%                             channel)
%                   \pcNmsm   (\mathcal{N}_{\mathsf A}: Alice's measurement
%                             channel, written by
%                             transmitter/thm:certificate-sufficiency,
%                             transmitter/rem:xreg and
%                             reductions/lem:source-transfer)
%    * Deliberately omitted: a \provenance{tag}{details} macro. Provenance
%      is a header field of every claim file (spec/claim-format.md section
%      3); a claim body that calls such a macro fails to compile, which is
%      the intended signal.
%
%  REQUIRED PACKAGES, loaded by the assembler's preamble before this file:
%    amsmath, amssymb, amsfonts, mathtools, dsfont, braket, physics, xcolor
%    (braket before physics, so that physics wins for \ket, \bra, \tr, ...).
%
%  This file contains only definitions: no \usepackage, no document body.
% =============================================================================


% =============================================================================
%  PART A -- notation macros, the ones spec/notation.md catalogues, each a
%  verbatim copy of its catalogue row; grouped under numbered headings.
% =============================================================================

% -----------------------------------------------------------------------------
% 1. State spaces and operator sets (a "system label" -> spaces of operators)
% -----------------------------------------------------------------------------
\newcommand{\Hs}[1]{\mathcal{H}(#1)}        % Hilbert space of a system label
\newcommand{\Lin}[1]{\mathcal{L}(#1)}       % bounded linear operators on H(#1)
\newcommand{\Linn}[2]{\mathcal{L}(#1,#2)}   % linear maps H(#1) -> H(#2)
\newcommand{\Pos}[1]{\mathrm{Pos}(#1)}      % positive semidefinite operators on H(#1)
\newcommand{\Density}{\mathrm{D}}           % density operators: \Density(A) = D(A)
\newcommand{\Seq}[1]{S_{=}(#1)}             % normalized states (trace = 1)
\newcommand{\Sleq}[1]{S_{\leq}(#1)}         % subnormalized states (0 <= trace <= 1)
\newcommand{\id}{\mathds{1}}                % identity operator
\newcommand{\idon}[1]{\mathds{1}_{#1}}      % identity on a named system, e.g. \idon{C}

% DEPRECATED: applies \mathcal to both arguments itself, so it takes bare
% letters, not system labels. Prefer \Linn; do not use in new contracts.
\providecommand{\Lop}[2]{\mathcal{L}(\mathcal{#1}, \mathcal{#2})}


% -----------------------------------------------------------------------------
% 2. Composite systems
% -----------------------------------------------------------------------------
\newcommand{\dsum}{\oplus}                  % direct sum of system labels
% Tensor product: use \otimes directly.
% n-fold direct sum of an operator space: \Lin{A}^{\oplus n}


% -----------------------------------------------------------------------------
% 3. Events, conditioning and classical registers (for STATES)
%    Convention:  \wedge  -> subnormalized "and-event" component, e.g. \rho_{A\wedge x}
%                 |        -> normalized conditional state,        e.g. \rho_{A|x}
% (Write these directly in subscripts; no extra macros needed.)
% -----------------------------------------------------------------------------


% -----------------------------------------------------------------------------
% 4. Photon-number subspaces / direct-sum blocks
%    Convention:  operators carry the restriction "on top", in brackets;
%                 states carry it as a \wedge subscript (see Sec. 3).
% -----------------------------------------------------------------------------
\newcommand{\Pnum}[1]{P_{#1}}               % photon-number projector (system from context)
\newcommand{\Pnumon}[2]{P_{#1}^{#2}}        % with explicit system superscript, e.g. P_n^C
\newcommand{\blk}[2]{#1^{(#2)}}             % operator block, e.g. \blk{M}{\leqN} = M^{(<=N)}
\newcommand{\leqN}{\leq N}                  % use as \blk{M}{\leqN}
\newcommand{\gtN}{> N}                      % use as \blk{M}{\gtN}
\newcommand{\eqn}{= n}                      % use as \blk{M}{\eqn}


% -----------------------------------------------------------------------------
% 5. System labels used in the protocol
% -----------------------------------------------------------------------------
\newcommand{\Chan}{C}                       % quantum-channel (Eve) input to Bob
\newcommand{\Prb}{P}                        % probe input
\newcommand{\Flag}{{\mathrm{Flg}}}          % flag system (EXISTS ONLY AFTER squashing)
\newcommand{\Loss}{{\mathrm{Lss}}}          % loss system
\newcommand{\Vc}{{\mathrm{Vcm}}}            % vacuum system
\newcommand{\Nout}{N_{\mathrm{out}}}
\newcommand{\Nin}{N_{\mathrm{in}}}


% -----------------------------------------------------------------------------
% 6. Alphabets
% -----------------------------------------------------------------------------
\newcommand{\Yset}{\mathcal{Y}}             % Bob's (coarse-grained) detection alphabet
\newcommand{\Yraw}{\mathcal{Y}_{\mathrm{raw}}}% raw click patterns {0,1}^4
\newcommand{\Yscset}{\mathcal{Y}_{\singlec}}% single-click outcomes
\newcommand{\Ymcset}{\mathcal{Y}_{\mc}}     % multi-click outcomes
\newcommand{\Tset}{\mathcal{T}_{\mathrm{set}}}% private setting set (the probe settings in the optical instance)
\newcommand{\Xset}{\mathcal{X}}             % Alice's alphabet
\newcommand{\Dset}{\{H,V,+,-\}}             % polarization / detector labels
\newcommand{\pass}{\top}                    % pass sifting


% -----------------------------------------------------------------------------
% 7. Mode / outcome abbreviations
% -----------------------------------------------------------------------------
\newcommand{\on}{{\mathrm{on}}}
\newcommand{\off}{{\mathrm{off}}}
\newcommand{\dc}{{\mathrm{dc}}}             % dark count
\newcommand{\sq}{\mathrm{sq}}               % squashed
\newcommand{\mc}{\mathrm{mc}}               % multi-click
\newcommand{\nc}{\mathrm{nc}}               % no-click
\newcommand{\singlec}{\mathrm{sc}}          % single-click
\newcommand{\vac}{\mathrm{vac}}             % vacuum
\newcommand{\kout}{{k_\mathrm{out}}}
\newcommand{\kin}{{k_\mathrm{in}}}


% -----------------------------------------------------------------------------
% 8. Maps, channels and miscellaneous operators
% -----------------------------------------------------------------------------
\newcommand{\Prob}{\mathrm{Prob}}           % probability (use \Prob[\cdot] or \Prob(\cdot))
\newcommand{\lmin}{\lambda_{\min}}          % minimum eigenvalue
\newcommand{\Wflag}{W_{\Flag}}              % flag-space weight
\newcommand{\transpose}{\mathsf{T}}
\newcommand{\innerp}[2]{\left\langle #1 \vert #2 \right\rangle}
\renewcommand{\vec}[1]{\boldsymbol{#1}}     % NB: overrides the LaTeX arrow accent


% -----------------------------------------------------------------------------
% 9. Key-rate numerics
% -----------------------------------------------------------------------------
\newcommand{\cG}{\mathcal{G}}               % protocol map
\newcommand{\cZ}{\mathcal{Z}}               % pinching map
\newcommand{\Amap}{\mathcal{A}_{\mathrm{ann}}} % announcement + measurement map
\newcommand{\Amapnb}{\mathcal{A}_{\mathrm{ann}|\mathrm{nb}}}
\newcommand{\simpl}{}                       % empty suffix for simplified block maps
\newcommand{\acSet}{\mathcal{S}_{\mathrm{acc}}} % acceptance set
\newcommand{\testE}{\mathcal{E}_{\mathrm{test}}} % test-statistics map
\newcommand{\ecLeak}{\mathrm{leak}^{\mathrm{EC}}_{\mathrm{obs}}}
\newcommand{\pgen}{p_{\mathrm{gen}}}
\newcommand{\askeyr}{R_{\infty}}             % asymptotic key rate
\newcommand{\ppass}{p_{\mathrm{pass}}}       % sifting pass probability


% -----------------------------------------------------------------------------
% 10. Finite-size analysis / Renyi entropies
% -----------------------------------------------------------------------------

% --- Round types: test vs. key generation (register T, alphabet \Fset) --------
\newcommand{\tst}{\mathrm{test}}            % test round
\newcommand{\gen}{\mathrm{gen}}             % key-generation round
\newcommand{\Fset}{\mathcal{T}}             % round-type alphabet {test, gen} (not \Tset)
\newcommand{\Chat}{\widehat{C}}             % acceptance-test register
\newcommand{\Chatset}{\widehat{\mathcal{C}}}% its alphabet

% --- Basis / key extraction functions ----------------------------------------
\newcommand{\bA}{b_{A}}                     % Alice's basis-extraction function
\newcommand{\bB}{b_{B}}                     % Bob's basis-extraction function
\newcommand{\kB}{k_{B}}                     % Bob's key-extraction function

% --- Maps --------------------------------------------------------------------
\newcommand{\Mprot}{\mathcal{M}}            % single-round protocol map
\newcommand{\Phimap}{\boldsymbol{\Phi}}     % test-statistics map
\newcommand{\Pdeph}{\mathcal{P}}            % dephasing (pinching) map on \Blreg

% --- Isometries --------------------------------------------------------------
\newcommand{\Vmsm}{V_{\mathrm{msm}}}        % measurement isometry
\newcommand{\Vann}{V_{\mathrm{ann}}}        % announcement isometry
\newcommand{\VI}{V_{I}}                     % Eve's copy of the announcements
\newcommand{\VChat}{V_{\Chat}}              % acceptance-data isometry
\newcommand{\VR}{V_{R}}                     % sifting + key map -> raw key R
\newcommand{\VZ}{V_{Z}}                     % raw key -> key register Z
\newcommand{\Wprot}{W}                      % full protocol isometry

% --- Renyi entropies ---------------------------------------------------------
% SECURITY-CRITICAL: see spec/notation.md §4. Never change an arrow to
% make an inequality work out; never convert silently between the sandwiched
% (tilde) and Petz (overline) families.
\newcommand{\Hsandup}{\widetilde{H}^{\uparrow}_{\alpha}}   % sandwiched, up
\newcommand{\Hsanddown}{\widetilde{H}^{\downarrow}_{\alpha}}% sandwiched, down
\newcommand{\Hpetzup}{\overline{H}^{\uparrow}_{\alpha}}    % Petz, up
\newcommand{\Hpetzdown}{\overline{H}^{\downarrow}_{\alpha}}% Petz, down
\newcommand{\hup}{h^{\uparrow}_{\alpha}}                   % single-round rate
\newcommand{\hqkd}{h^{\uparrow,\mathrm{QKD}}_{\alpha}}     % single-round QKD rate
\newcommand{\hlow}{\underline{h}_{\alpha}}                 % lower bound on \hqkd
% --- spec/notation.md 3.5 and 4.5 -------------------------------------------
% \hlb  is the first-order linearization lower bound on \hlow: the headline
%       key rate (the `first-order` tier of foundations/def:tiers). NOT \hlow
%       itself.
% \tnum is the instance's one declared numerical tolerance. Its
%       near-collision with the declared marginal \tau_A is disclosed in
%       spec/notation.md 3.5; bare \tau is reserved and denotes neither.
\newcommand{\hlb}{\underline{h}^{\mathrm{FO}}_{\alpha}}    % first-order lower bound on \hlow
\newcommand{\tnum}{\tau_{\mathrm{num}}}                    % declared numerical tolerance


% =============================================================================
%  PART B -- module-local macros, grouped by module. Deprecated aliases,
%  defined so that a text written with them compiles, used in no new contract: \Lop,
%  \testE, \Pdeph, \pass, \ecLeak (spec/notation.md section 3.9).
% =============================================================================

% --- meat-core -----------------------------------------------------------------
\newcommand{\SecSet}{\mathbb{S}}            % sector index set of Q
\newcommand{\Qsys}{Q}                       % declared sector-carrying measured system
\newcommand{\Vset}{\mathcal{V}}             % declared joint outcome alphabet of the Q-POVM
\newcommand{\CAset}{\mathcal{C}_{A}}        % Alice announcement alphabet
\newcommand{\CBset}{\mathcal{C}_{B}}        % Bob announcement alphabet
\newcommand{\Party}{\mathfrak{P}}           % declared party partition
\newcommand{\Oset}[1]{\mathcal{O}_{#1}}     % readout alphabet of party #1
\newcommand{\rd}[1]{\pi_{#1}}               % readout map of party #1
\newcommand{\Sadm}{\mathcal{S}_{A}}         % admissible Alice subspace
\newcommand{\nomarg}{\mathsf{none}}         % "no trusted marginal declared"
\newcommand{\Zset}{\mathcal{Z}_{\perp}}     % raw-key alphabet incl. \perp
\newcommand{\Zval}{\mathcal{Z}}             % key alphabet without \perp
\newcommand{\Etil}{\widetilde{E}}           % single-round purifying register
\newcommand{\Ehat}{\widehat{E}}             % Eve's n-round register
\newcommand{\OmAT}{\Omega_{\mathrm{AT}}}    % acceptance event
\newcommand{\epa}{\varepsilon_{\mathrm{PA}}}
\newcommand{\eev}{\varepsilon_{\mathrm{EV}}}
\newcommand{\lEC}{\lambda_{\mathrm{EC}}}
\newcommand{\Decl}{\mathfrak{D}}            % protocol declaration
\newcommand{\CList}{\mathfrak{L}}           % declared constraint list
\newcommand{\PinS}{\mathcal{P}_{\SecSet}}   % sector pinching on Q
\newcommand{\Feas}{\mathcal{W}}             % feasible set of block states
\newcommand{\Chkey}{\Chatset_{\mathrm{key}}}% key-live acceptance values
\newcommand{\Kann}{\mathcal{K}}             % announcement-class alphabet
\newcommand{\hold}{\mathsf{h}}              % declared key holder (A or B)
\newcommand{\sfD}{D}                        % classical KL divergence (base 2)
\newcommand{\Dtwo}{D}                       % quantum relative entropy (base 2)
\newcommand{\Dsand}{\widetilde{D}_{\alpha}} % sandwiched Renyi divergence
\newcommand{\idmap}{\mathrm{id}}            % identity channel
\newcommand{\Aev}{\mathcal{A}}              % Eve's per-round attack channel
\newcommand{\AttSet}{\mathfrak{A}}          % licensed attack-channel set
\newcommand{\Chacc}{\Chatset_{\mathrm{acc}}}% acceptance support of \acSet
\newcommand{\epsstar}{\varepsilon^{\star}}      % distance to an exact POVM
\newcommand{\dfH}{\Delta_{\mathrm{herm}}}       % Hermiticity defect
\newcommand{\dfP}{\Delta_{\mathrm{pos}}}        % negative part
\newcommand{\dfC}{\Delta_{\mathrm{cpl}}}        % completeness defect
\newcommand{\lmoq}{g}                       % acceptance-side gradient block
\newcommand{\lmoJ}{G}                       % feasible-side gradient block
\newcommand{\bacc}{\beta_{\mathrm{acc}}}    % certified acceptance-side minimum
\newcommand{\Itset}{\mathfrak{I}}           % finite index set of iterates (C22)

% --- reductions ----------------------------------------------------------------
\newcommand{\Bsys}{B}                        % physical (pre-squashing) system
\newcommand{\Tcl}[1]{\mathcal{T}(#1)}        % trace-class operators on H(#1)
\newcommand{\Ins}{\mathcal{I}}               % generalized-instrument branch map
\newcommand{\redNmsm}{\mathcal{N}}              % POVM measurement channel
\newcommand{\Eprot}{\mathcal{E}}             % n-round protocol channel
\newcommand{\Grnd}{\mathsf{G}}               % one round map of a protocol
\newcommand{\Kfin}{\mathsf{K}}               % end-of-protocol processing
\newcommand{\Crit}{\mathsf{sec}}             % value of a security criterion
\newcommand{\Pflag}{\Pi_{\Flag}}             % flag-sector projector on Q
\newcommand{\Pleq}{\Pi_{\leqN}}              % preserved-block projector
\newcommand{\Pgt}{\Pi_{\gtN}}                % above-cutoff projector
\newcommand{\Sph}{S^{\mathrm{ph}}}           % family of physical inputs
\newcommand{\Srch}{S^{\mathrm{rch}}}         % reachable squashed states
\newcommand{\Ssq}{S^{\sq}}                   % squashed admissible set
\newcommand{\Ffss}{F^{\mathrm{FSS}}}         % flag-state-squasher target
\newcommand{\Fwp}{F^{\mathrm{WP}}}           % multi-bound WPFSS target
\newcommand{\Mwp}{M^{\mathrm{WP}}}           % renormalized above-cutoff POVM
\newcommand{\Mdec}{{M'}^{\mathrm{ph}}}       % declared physical family on B
\newcommand{\Ppvm}{P^{\Flag}}                % flag PVM
\newcommand{\Lmp}{{\Lambda_{\mathrm{mp}}}}   % measure-and-prepare channel
\newcommand{\Lwp}{{\Lambda^{\mathrm{WP}}}}   % multi-bound WPFSS channel
\newcommand{\Lfss}{{\Lambda^{\mathrm{FSS}}}} % flag-state-squasher channel
\newcommand{\lmax}{\lambda_{\max}}           % maximum eigenvalue
\newcommand{\Lnoise}{\widetilde{\Lambda}}    % flag-noise channel
\newcommand{\cg}{\varpi}                     % declared coarse-graining map
\newcommand{\MCK}{\mathsf{MC}}               % multiclick raw-pattern set
\newcommand{\redChain}{\mathfrak{C}}            % reduction chain
\newcommand{\Lnk}{\ell}                      % a single typed link
\newcommand{\dsq}{\preceq}                      % the squashing relation
\newcommand{\opn}[1]{\left\lVert #1\right\rVert}           % operator (spectral) norm
\newcommand{\mxn}[1]{\left\lVert #1\right\rVert_{\max}}    % entrywise max modulus
\newcommand{\Fld}{\mathbb{F}}                   % exact ordered substrate field
\newcommand{\Fldi}{\mathbb{F}_{i}}              % its Gaussian extension F + iF
\newcommand{\lat}[1]{\big\lfloor #1\big\rfloor_{\kappa}} % round DOWN onto the lattice

% --- protocol-compiler ---------------------------------------------------------
\newcommand{\Spec}{\mathfrak{S}}             % protocol specification (compiler input)
\newcommand{\pcChain}{\mathfrak{R}}            % reduction chain (two typed segments)
\newcommand{\Chsrc}{\pcChain_{\mathrm{src}}}   % its Alice-side source-map segment
\newcommand{\Chrec}{\pcChain_{\mathrm{rec}}}   % its receiver-side squashing segment
\newcommand{\rvirt}{\rho^{\mathrm{v}}}       % virtual (PRE-source-map) signal state
\newcommand{\Fph}{F^{\mathrm{ph}}}           % physical (pre-reduction) POVM element
\newcommand{\Vph}{\mathcal{V}^{\mathrm{ph}}} % physical outcome alphabet
\newcommand{\Tpub}{\mathcal{T}_{\mathrm{pub}}}% public-choice alphabet
\newcommand{\Uset}{\mathcal{U}}              % residual outcome alphabet
\newcommand{\idxm}{\mathsf{idx}}             % label-to-index map
\newcommand{\Asig}{A'}                       % transmitted (signal) system -- the REAL one
\newcommand{\Avirt}{A''}                     % virtual transmitted system (source descriptor of Xi_1)
\newcommand{\Abar}{\bar{A}}                  % Alice's measured register
\newcommand{\Ahat}{\hat{A}}                  % Alice's shield register
\newcommand{\pcNmsm}{\mathcal{N}_{\mathsf A}}  % Alice's measurement channel (see the RULES block)
\newcommand{\Phip}{\Phi^{+}}                 % maximally entangled vector label
\newcommand{\Smap}{\mathfrak{s}}             % the emitted sector map (fraktur: a compile-time object)

% --- receiver ------------------------------------------------------------------
\newcommand{\Modes}{\mathbb{M}}              % declared ordered mode set
\newcommand{\Inp}{\mathbb{I}}                % declared input modes
\newcommand{\InpC}{\mathbb{I}_{\Chan}}       % channel input modes
\newcommand{\InpS}{\mathbb{I}_{\Prb}}        % probe (source) input modes
\newcommand{\Env}{\mathbb{E}}                % dilation (environment) modes
\newcommand{\Port}[1]{\mathbb{P}_{#1}}       % mode set of detector port #1
\newcommand{\Emd}[1]{E_{#1}}                 % coordinate projector on modes #1
\newcommand{\Gmat}{G}                        % instrument (transfer) matrix
\newcommand{\Fock}{\mathfrak{F}}             % Fock / second-quantization functor
\newcommand{\Symn}[1]{\mathrm{Sym}^{#1}}     % symmetric tensor power
\newcommand{\normord}[1]{{:}\mspace{2mu}#1\mspace{2mu}{:}}   % normal ordering
\newcommand{\rk}{\operatorname{rank}}        % rank (physics.sty owns \rank)
\newcommand{\SCK}{\mathsf{SC}}               % single-click pattern set
\newcommand{\Recv}{\mathfrak{R}}             % receiver instrument tuple
\newcommand{\Cat}{\mathfrak{C}}              % component catalogue
\newcommand{\Grph}{\mathfrak{E}}             % declared element graph
\newcommand{\defic}{\mathbb{L}}              % loss deficiency operator
\newcommand{\Pvac}[1]{\Pi_{\vac}(#1)}        % joint no-click element on ports #1
\newcommand{\pmc}{p_{\mc}}                   % uniform high-photon mc weight
\newcommand{\Bnd}{\beta}                     % the scalar bound function beta(n)
\newcommand{\Rot}{R}                         % planar rotation matrix
\newcommand{\Wplate}{W}                      % retarder (wave-plate) matrix
\newcommand{\Bsp}{B}                         % beamsplitter core matrix

% --- statistics ----------------------------------------------------------------
\newcommand{\honest}{p^{\mathrm{hon}}}          % honest per-round distribution
\newcommand{\ecomp}{\varepsilon_{\mathrm{comp}}}% completeness budget
\newcommand{\simplex}[1]{\Delta(#1)}            % probability simplex
\newcommand{\Chatcon}{\Chatset_{\mathrm{con}}}  % constrained outcomes
\newcommand{\Ireg}[3]{\mathcal{I}_{#1}(#2,#3)}  % regularized incomplete beta
\newcommand{\Ncount}[1]{\mathsf{N}_{#1}}        % per-outcome count
\newcommand{\qlo}{\mathbf{q}^{-}}               % box lower bounds
\newcommand{\qhi}{\mathbf{q}^{+}}               % box upper bounds
 for it, so the
%  file is read from spec/modules/<id>/.
%
%  RULES
%    * A macro that spec/notation.md catalogues is a verbatim copy of it
%      (part A). Changing one is a notation change: bump spec/notation.md.
%    * Part B holds the module-local macros. A macro defined by two modules
%      with the same expansion appears once, under the first module (in the
%      order below) that defines it. Two macros with the same name and
%      DIFFERENT expansions are both kept under a module prefix, and the
%      claims write the prefixed names:
%        \Chain  -> \redChain (reductions, \mathfrak{C}: the reduction chain)
%                   \pcChain  (protocol-compiler, \mathfrak{R}: the two-segment
%                             chain; \Chsrc and \Chrec are built on it)
%        \Nmsm   -> \redNmsm  (reductions, \mathcal{N}: a POVM measurement
%                             channel)
%                   \pcNmsm   (\mathcal{N}_{\mathsf A}: Alice's measurement
%                             channel, written by
%                             transmitter/thm:certificate-sufficiency,
%                             transmitter/rem:xreg and
%                             reductions/lem:source-transfer)
%    * Deliberately omitted: a \provenance{tag}{details} macro. Provenance
%      is a header field of every claim file (spec/claim-format.md section
%      3); a claim body that calls such a macro fails to compile, which is
%      the intended signal.
%
%  REQUIRED PACKAGES, loaded by the assembler's preamble before this file:
%    amsmath, amssymb, amsfonts, mathtools, dsfont, braket, physics, xcolor
%    (braket before physics, so that physics wins for \ket, \bra, \tr, ...).
%
%  This file contains only definitions: no \usepackage, no document body.
% =============================================================================


% =============================================================================
%  PART A -- notation macros, the ones spec/notation.md catalogues, each a
%  verbatim copy of its catalogue row; grouped under numbered headings.
% =============================================================================

% -----------------------------------------------------------------------------
% 1. State spaces and operator sets (a "system label" -> spaces of operators)
% -----------------------------------------------------------------------------
\newcommand{\Hs}[1]{\mathcal{H}(#1)}        % Hilbert space of a system label
\newcommand{\Lin}[1]{\mathcal{L}(#1)}       % bounded linear operators on H(#1)
\newcommand{\Linn}[2]{\mathcal{L}(#1,#2)}   % linear maps H(#1) -> H(#2)
\newcommand{\Pos}[1]{\mathrm{Pos}(#1)}      % positive semidefinite operators on H(#1)
\newcommand{\Density}{\mathrm{D}}           % density operators: \Density(A) = D(A)
\newcommand{\Seq}[1]{S_{=}(#1)}             % normalized states (trace = 1)
\newcommand{\Sleq}[1]{S_{\leq}(#1)}         % subnormalized states (0 <= trace <= 1)
\newcommand{\id}{\mathds{1}}                % identity operator
\newcommand{\idon}[1]{\mathds{1}_{#1}}      % identity on a named system, e.g. \idon{C}

% DEPRECATED: applies \mathcal to both arguments itself, so it takes bare
% letters, not system labels. Prefer \Linn; do not use in new contracts.
\providecommand{\Lop}[2]{\mathcal{L}(\mathcal{#1}, \mathcal{#2})}


% -----------------------------------------------------------------------------
% 2. Composite systems
% -----------------------------------------------------------------------------
\newcommand{\dsum}{\oplus}                  % direct sum of system labels
% Tensor product: use \otimes directly.
% n-fold direct sum of an operator space: \Lin{A}^{\oplus n}


% -----------------------------------------------------------------------------
% 3. Events, conditioning and classical registers (for STATES)
%    Convention:  \wedge  -> subnormalized "and-event" component, e.g. \rho_{A\wedge x}
%                 |        -> normalized conditional state,        e.g. \rho_{A|x}
% (Write these directly in subscripts; no extra macros needed.)
% -----------------------------------------------------------------------------


% -----------------------------------------------------------------------------
% 4. Photon-number subspaces / direct-sum blocks
%    Convention:  operators carry the restriction "on top", in brackets;
%                 states carry it as a \wedge subscript (see Sec. 3).
% -----------------------------------------------------------------------------
\newcommand{\Pnum}[1]{P_{#1}}               % photon-number projector (system from context)
\newcommand{\Pnumon}[2]{P_{#1}^{#2}}        % with explicit system superscript, e.g. P_n^C
\newcommand{\blk}[2]{#1^{(#2)}}             % operator block, e.g. \blk{M}{\leqN} = M^{(<=N)}
\newcommand{\leqN}{\leq N}                  % use as \blk{M}{\leqN}
\newcommand{\gtN}{> N}                      % use as \blk{M}{\gtN}
\newcommand{\eqn}{= n}                      % use as \blk{M}{\eqn}


% -----------------------------------------------------------------------------
% 5. System labels used in the protocol
% -----------------------------------------------------------------------------
\newcommand{\Chan}{C}                       % quantum-channel (Eve) input to Bob
\newcommand{\Prb}{P}                        % probe input
\newcommand{\Flag}{{\mathrm{Flg}}}          % flag system (EXISTS ONLY AFTER squashing)
\newcommand{\Loss}{{\mathrm{Lss}}}          % loss system
\newcommand{\Vc}{{\mathrm{Vcm}}}            % vacuum system
\newcommand{\Nout}{N_{\mathrm{out}}}
\newcommand{\Nin}{N_{\mathrm{in}}}


% -----------------------------------------------------------------------------
% 6. Alphabets
% -----------------------------------------------------------------------------
\newcommand{\Yset}{\mathcal{Y}}             % Bob's (coarse-grained) detection alphabet
\newcommand{\Yraw}{\mathcal{Y}_{\mathrm{raw}}}% raw click patterns {0,1}^4
\newcommand{\Yscset}{\mathcal{Y}_{\singlec}}% single-click outcomes
\newcommand{\Ymcset}{\mathcal{Y}_{\mc}}     % multi-click outcomes
\newcommand{\Tset}{\mathcal{T}_{\mathrm{set}}}% private setting set (the probe settings in the optical instance)
\newcommand{\Xset}{\mathcal{X}}             % Alice's alphabet
\newcommand{\Dset}{\{H,V,+,-\}}             % polarization / detector labels
\newcommand{\pass}{\top}                    % pass sifting


% -----------------------------------------------------------------------------
% 7. Mode / outcome abbreviations
% -----------------------------------------------------------------------------
\newcommand{\on}{{\mathrm{on}}}
\newcommand{\off}{{\mathrm{off}}}
\newcommand{\dc}{{\mathrm{dc}}}             % dark count
\newcommand{\sq}{\mathrm{sq}}               % squashed
\newcommand{\mc}{\mathrm{mc}}               % multi-click
\newcommand{\nc}{\mathrm{nc}}               % no-click
\newcommand{\singlec}{\mathrm{sc}}          % single-click
\newcommand{\vac}{\mathrm{vac}}             % vacuum
\newcommand{\kout}{{k_\mathrm{out}}}
\newcommand{\kin}{{k_\mathrm{in}}}


% -----------------------------------------------------------------------------
% 8. Maps, channels and miscellaneous operators
% -----------------------------------------------------------------------------
\newcommand{\Prob}{\mathrm{Prob}}           % probability (use \Prob[\cdot] or \Prob(\cdot))
\newcommand{\lmin}{\lambda_{\min}}          % minimum eigenvalue
\newcommand{\Wflag}{W_{\Flag}}              % flag-space weight
\newcommand{\transpose}{\mathsf{T}}
\newcommand{\innerp}[2]{\left\langle #1 \vert #2 \right\rangle}
\renewcommand{\vec}[1]{\boldsymbol{#1}}     % NB: overrides the LaTeX arrow accent


% -----------------------------------------------------------------------------
% 9. Key-rate numerics
% -----------------------------------------------------------------------------
\newcommand{\cG}{\mathcal{G}}               % protocol map
\newcommand{\cZ}{\mathcal{Z}}               % pinching map
\newcommand{\Amap}{\mathcal{A}_{\mathrm{ann}}} % announcement + measurement map
\newcommand{\Amapnb}{\mathcal{A}_{\mathrm{ann}|\mathrm{nb}}}
\newcommand{\simpl}{}                       % empty suffix for simplified block maps
\newcommand{\acSet}{\mathcal{S}_{\mathrm{acc}}} % acceptance set
\newcommand{\testE}{\mathcal{E}_{\mathrm{test}}} % test-statistics map
\newcommand{\ecLeak}{\mathrm{leak}^{\mathrm{EC}}_{\mathrm{obs}}}
\newcommand{\pgen}{p_{\mathrm{gen}}}
\newcommand{\askeyr}{R_{\infty}}             % asymptotic key rate
\newcommand{\ppass}{p_{\mathrm{pass}}}       % sifting pass probability


% -----------------------------------------------------------------------------
% 10. Finite-size analysis / Renyi entropies
% -----------------------------------------------------------------------------

% --- Round types: test vs. key generation (register T, alphabet \Fset) --------
\newcommand{\tst}{\mathrm{test}}            % test round
\newcommand{\gen}{\mathrm{gen}}             % key-generation round
\newcommand{\Fset}{\mathcal{T}}             % round-type alphabet {test, gen} (not \Tset)
\newcommand{\Chat}{\widehat{C}}             % acceptance-test register
\newcommand{\Chatset}{\widehat{\mathcal{C}}}% its alphabet

% --- Basis / key extraction functions ----------------------------------------
\newcommand{\bA}{b_{A}}                     % Alice's basis-extraction function
\newcommand{\bB}{b_{B}}                     % Bob's basis-extraction function
\newcommand{\kB}{k_{B}}                     % Bob's key-extraction function

% --- Maps --------------------------------------------------------------------
\newcommand{\Mprot}{\mathcal{M}}            % single-round protocol map
\newcommand{\Phimap}{\boldsymbol{\Phi}}     % test-statistics map
\newcommand{\Pdeph}{\mathcal{P}}            % dephasing (pinching) map on \Blreg

% --- Isometries --------------------------------------------------------------
\newcommand{\Vmsm}{V_{\mathrm{msm}}}        % measurement isometry
\newcommand{\Vann}{V_{\mathrm{ann}}}        % announcement isometry
\newcommand{\VI}{V_{I}}                     % Eve's copy of the announcements
\newcommand{\VChat}{V_{\Chat}}              % acceptance-data isometry
\newcommand{\VR}{V_{R}}                     % sifting + key map -> raw key R
\newcommand{\VZ}{V_{Z}}                     % raw key -> key register Z
\newcommand{\Wprot}{W}                      % full protocol isometry

% --- Renyi entropies ---------------------------------------------------------
% SECURITY-CRITICAL: see spec/notation.md §4. Never change an arrow to
% make an inequality work out; never convert silently between the sandwiched
% (tilde) and Petz (overline) families.
\newcommand{\Hsandup}{\widetilde{H}^{\uparrow}_{\alpha}}   % sandwiched, up
\newcommand{\Hsanddown}{\widetilde{H}^{\downarrow}_{\alpha}}% sandwiched, down
\newcommand{\Hpetzup}{\overline{H}^{\uparrow}_{\alpha}}    % Petz, up
\newcommand{\Hpetzdown}{\overline{H}^{\downarrow}_{\alpha}}% Petz, down
\newcommand{\hup}{h^{\uparrow}_{\alpha}}                   % single-round rate
\newcommand{\hqkd}{h^{\uparrow,\mathrm{QKD}}_{\alpha}}     % single-round QKD rate
\newcommand{\hlow}{\underline{h}_{\alpha}}                 % lower bound on \hqkd
% --- spec/notation.md 3.5 and 4.5 -------------------------------------------
% \hlb  is the first-order linearization lower bound on \hlow: the headline
%       key rate (the `first-order` tier of foundations/def:tiers). NOT \hlow
%       itself.
% \tnum is the instance's one declared numerical tolerance. Its
%       near-collision with the declared marginal \tau_A is disclosed in
%       spec/notation.md 3.5; bare \tau is reserved and denotes neither.
\newcommand{\hlb}{\underline{h}^{\mathrm{FO}}_{\alpha}}    % first-order lower bound on \hlow
\newcommand{\tnum}{\tau_{\mathrm{num}}}                    % declared numerical tolerance


% =============================================================================
%  PART B -- module-local macros, grouped by module. Deprecated aliases,
%  defined so that a text written with them compiles, used in no new contract: \Lop,
%  \testE, \Pdeph, \pass, \ecLeak (spec/notation.md section 3.9).
% =============================================================================

% --- meat-core -----------------------------------------------------------------
\newcommand{\SecSet}{\mathbb{S}}            % sector index set of Q
\newcommand{\Qsys}{Q}                       % declared sector-carrying measured system
\newcommand{\Vset}{\mathcal{V}}             % declared joint outcome alphabet of the Q-POVM
\newcommand{\CAset}{\mathcal{C}_{A}}        % Alice announcement alphabet
\newcommand{\CBset}{\mathcal{C}_{B}}        % Bob announcement alphabet
\newcommand{\Party}{\mathfrak{P}}           % declared party partition
\newcommand{\Oset}[1]{\mathcal{O}_{#1}}     % readout alphabet of party #1
\newcommand{\rd}[1]{\pi_{#1}}               % readout map of party #1
\newcommand{\Sadm}{\mathcal{S}_{A}}         % admissible Alice subspace
\newcommand{\nomarg}{\mathsf{none}}         % "no trusted marginal declared"
\newcommand{\Zset}{\mathcal{Z}_{\perp}}     % raw-key alphabet incl. \perp
\newcommand{\Zval}{\mathcal{Z}}             % key alphabet without \perp
\newcommand{\Etil}{\widetilde{E}}           % single-round purifying register
\newcommand{\Ehat}{\widehat{E}}             % Eve's n-round register
\newcommand{\OmAT}{\Omega_{\mathrm{AT}}}    % acceptance event
\newcommand{\epa}{\varepsilon_{\mathrm{PA}}}
\newcommand{\eev}{\varepsilon_{\mathrm{EV}}}
\newcommand{\lEC}{\lambda_{\mathrm{EC}}}
\newcommand{\Decl}{\mathfrak{D}}            % protocol declaration
\newcommand{\CList}{\mathfrak{L}}           % declared constraint list
\newcommand{\PinS}{\mathcal{P}_{\SecSet}}   % sector pinching on Q
\newcommand{\Feas}{\mathcal{W}}             % feasible set of block states
\newcommand{\Chkey}{\Chatset_{\mathrm{key}}}% key-live acceptance values
\newcommand{\Kann}{\mathcal{K}}             % announcement-class alphabet
\newcommand{\hold}{\mathsf{h}}              % declared key holder (A or B)
\newcommand{\sfD}{D}                        % classical KL divergence (base 2)
\newcommand{\Dtwo}{D}                       % quantum relative entropy (base 2)
\newcommand{\Dsand}{\widetilde{D}_{\alpha}} % sandwiched Renyi divergence
\newcommand{\idmap}{\mathrm{id}}            % identity channel
\newcommand{\Aev}{\mathcal{A}}              % Eve's per-round attack channel
\newcommand{\AttSet}{\mathfrak{A}}          % licensed attack-channel set
\newcommand{\Chacc}{\Chatset_{\mathrm{acc}}}% acceptance support of \acSet
\newcommand{\epsstar}{\varepsilon^{\star}}      % distance to an exact POVM
\newcommand{\dfH}{\Delta_{\mathrm{herm}}}       % Hermiticity defect
\newcommand{\dfP}{\Delta_{\mathrm{pos}}}        % negative part
\newcommand{\dfC}{\Delta_{\mathrm{cpl}}}        % completeness defect
\newcommand{\lmoq}{g}                       % acceptance-side gradient block
\newcommand{\lmoJ}{G}                       % feasible-side gradient block
\newcommand{\bacc}{\beta_{\mathrm{acc}}}    % certified acceptance-side minimum
\newcommand{\Itset}{\mathfrak{I}}           % finite index set of iterates (C22)

% --- reductions ----------------------------------------------------------------
\newcommand{\Bsys}{B}                        % physical (pre-squashing) system
\newcommand{\Tcl}[1]{\mathcal{T}(#1)}        % trace-class operators on H(#1)
\newcommand{\Ins}{\mathcal{I}}               % generalized-instrument branch map
\newcommand{\redNmsm}{\mathcal{N}}              % POVM measurement channel
\newcommand{\Eprot}{\mathcal{E}}             % n-round protocol channel
\newcommand{\Grnd}{\mathsf{G}}               % one round map of a protocol
\newcommand{\Kfin}{\mathsf{K}}               % end-of-protocol processing
\newcommand{\Crit}{\mathsf{sec}}             % value of a security criterion
\newcommand{\Pflag}{\Pi_{\Flag}}             % flag-sector projector on Q
\newcommand{\Pleq}{\Pi_{\leqN}}              % preserved-block projector
\newcommand{\Pgt}{\Pi_{\gtN}}                % above-cutoff projector
\newcommand{\Sph}{S^{\mathrm{ph}}}           % family of physical inputs
\newcommand{\Srch}{S^{\mathrm{rch}}}         % reachable squashed states
\newcommand{\Ssq}{S^{\sq}}                   % squashed admissible set
\newcommand{\Ffss}{F^{\mathrm{FSS}}}         % flag-state-squasher target
\newcommand{\Fwp}{F^{\mathrm{WP}}}           % multi-bound WPFSS target
\newcommand{\Mwp}{M^{\mathrm{WP}}}           % renormalized above-cutoff POVM
\newcommand{\Mdec}{{M'}^{\mathrm{ph}}}       % declared physical family on B
\newcommand{\Ppvm}{P^{\Flag}}                % flag PVM
\newcommand{\Lmp}{{\Lambda_{\mathrm{mp}}}}   % measure-and-prepare channel
\newcommand{\Lwp}{{\Lambda^{\mathrm{WP}}}}   % multi-bound WPFSS channel
\newcommand{\Lfss}{{\Lambda^{\mathrm{FSS}}}} % flag-state-squasher channel
\newcommand{\lmax}{\lambda_{\max}}           % maximum eigenvalue
\newcommand{\Lnoise}{\widetilde{\Lambda}}    % flag-noise channel
\newcommand{\cg}{\varpi}                     % declared coarse-graining map
\newcommand{\MCK}{\mathsf{MC}}               % multiclick raw-pattern set
\newcommand{\redChain}{\mathfrak{C}}            % reduction chain
\newcommand{\Lnk}{\ell}                      % a single typed link
\newcommand{\dsq}{\preceq}                      % the squashing relation
\newcommand{\opn}[1]{\left\lVert #1\right\rVert}           % operator (spectral) norm
\newcommand{\mxn}[1]{\left\lVert #1\right\rVert_{\max}}    % entrywise max modulus
\newcommand{\Fld}{\mathbb{F}}                   % exact ordered substrate field
\newcommand{\Fldi}{\mathbb{F}_{i}}              % its Gaussian extension F + iF
\newcommand{\lat}[1]{\big\lfloor #1\big\rfloor_{\kappa}} % round DOWN onto the lattice

% --- protocol-compiler ---------------------------------------------------------
\newcommand{\Spec}{\mathfrak{S}}             % protocol specification (compiler input)
\newcommand{\pcChain}{\mathfrak{R}}            % reduction chain (two typed segments)
\newcommand{\Chsrc}{\pcChain_{\mathrm{src}}}   % its Alice-side source-map segment
\newcommand{\Chrec}{\pcChain_{\mathrm{rec}}}   % its receiver-side squashing segment
\newcommand{\rvirt}{\rho^{\mathrm{v}}}       % virtual (PRE-source-map) signal state
\newcommand{\Fph}{F^{\mathrm{ph}}}           % physical (pre-reduction) POVM element
\newcommand{\Vph}{\mathcal{V}^{\mathrm{ph}}} % physical outcome alphabet
\newcommand{\Tpub}{\mathcal{T}_{\mathrm{pub}}}% public-choice alphabet
\newcommand{\Uset}{\mathcal{U}}              % residual outcome alphabet
\newcommand{\idxm}{\mathsf{idx}}             % label-to-index map
\newcommand{\Asig}{A'}                       % transmitted (signal) system -- the REAL one
\newcommand{\Avirt}{A''}                     % virtual transmitted system (source descriptor of Xi_1)
\newcommand{\Abar}{\bar{A}}                  % Alice's measured register
\newcommand{\Ahat}{\hat{A}}                  % Alice's shield register
\newcommand{\pcNmsm}{\mathcal{N}_{\mathsf A}}  % Alice's measurement channel (see the RULES block)
\newcommand{\Phip}{\Phi^{+}}                 % maximally entangled vector label
\newcommand{\Smap}{\mathfrak{s}}             % the emitted sector map (fraktur: a compile-time object)

% --- receiver ------------------------------------------------------------------
\newcommand{\Modes}{\mathbb{M}}              % declared ordered mode set
\newcommand{\Inp}{\mathbb{I}}                % declared input modes
\newcommand{\InpC}{\mathbb{I}_{\Chan}}       % channel input modes
\newcommand{\InpS}{\mathbb{I}_{\Prb}}        % probe (source) input modes
\newcommand{\Env}{\mathbb{E}}                % dilation (environment) modes
\newcommand{\Port}[1]{\mathbb{P}_{#1}}       % mode set of detector port #1
\newcommand{\Emd}[1]{E_{#1}}                 % coordinate projector on modes #1
\newcommand{\Gmat}{G}                        % instrument (transfer) matrix
\newcommand{\Fock}{\mathfrak{F}}             % Fock / second-quantization functor
\newcommand{\Symn}[1]{\mathrm{Sym}^{#1}}     % symmetric tensor power
\newcommand{\normord}[1]{{:}\mspace{2mu}#1\mspace{2mu}{:}}   % normal ordering
\newcommand{\rk}{\operatorname{rank}}        % rank (physics.sty owns \rank)
\newcommand{\SCK}{\mathsf{SC}}               % single-click pattern set
\newcommand{\Recv}{\mathfrak{R}}             % receiver instrument tuple
\newcommand{\Cat}{\mathfrak{C}}              % component catalogue
\newcommand{\Grph}{\mathfrak{E}}             % declared element graph
\newcommand{\defic}{\mathbb{L}}              % loss deficiency operator
\newcommand{\Pvac}[1]{\Pi_{\vac}(#1)}        % joint no-click element on ports #1
\newcommand{\pmc}{p_{\mc}}                   % uniform high-photon mc weight
\newcommand{\Bnd}{\beta}                     % the scalar bound function beta(n)
\newcommand{\Rot}{R}                         % planar rotation matrix
\newcommand{\Wplate}{W}                      % retarder (wave-plate) matrix
\newcommand{\Bsp}{B}                         % beamsplitter core matrix

% --- statistics ----------------------------------------------------------------
\newcommand{\honest}{p^{\mathrm{hon}}}          % honest per-round distribution
\newcommand{\ecomp}{\varepsilon_{\mathrm{comp}}}% completeness budget
\newcommand{\simplex}[1]{\Delta(#1)}            % probability simplex
\newcommand{\Chatcon}{\Chatset_{\mathrm{con}}}  % constrained outcomes
\newcommand{\Ireg}[3]{\mathcal{I}_{#1}(#2,#3)}  % regularized incomplete beta
\newcommand{\Ncount}[1]{\mathsf{N}_{#1}}        % per-outcome count
\newcommand{\qlo}{\mathbf{q}^{-}}               % box lower bounds
\newcommand{\qhi}{\mathbf{q}^{+}}               % box upper bounds

\makeatletter
\newcommand*{\ClaimHeadFont}{\bfseries}   % switched by prominence
\newcommand*{\ClaimNoteFont}{\mdseries}
\newcommand*{\ClaimShortId}{}
\newcommand*{\SetShortId}[1]{\renewcommand*{\ClaimShortId}{#1}}
\newcommand*{\RestatedName}{}
% Upright bodies throughout: italic is reserved for non-normative text.
% \@endparpenalty\@M keeps the meta line on the statement's page.
\newtheoremstyle{claim}{\topsep}{\topsep}{\normalfont\@endparpenalty\@M}{}%
  {\ClaimHeadFont}{.}{.5em}{\thmname{#1}\thmnumber{ #2}\thmnote{ {\ClaimNoteFont(#3)}}}
\newtheoremstyle{claimass}{\topsep}{\topsep}{\normalfont\@endparpenalty\@M}{}%
  {\ClaimHeadFont}{.}{.5em}{\thmname{#1}\thmnumber{ #2}%
  \ifx\ClaimShortId\@empty\else\ [\textsf{\ClaimShortId}]\fi%
  \thmnote{ {\ClaimNoteFont(#3)}}}
\newtheoremstyle{claimremark}{.5\topsep}{.5\topsep}{\small\@endparpenalty\@M}{}%
  {\itshape}{.}{.5em}{\thmname{#1}\thmnumber{ #2}\thmnote{ {\upshape(#3)}}}
\newtheoremstyle{claimreq}{\topsep}{\topsep}{\normalfont\@endparpenalty\@M}{}%
  {\color{reqAccent}\scshape}{.}{.5em}%
  {\thmname{#1}\thmnumber{ #2}\thmnote{ {\upshape(#3)}}}
\makeatother
\theoremstyle{claim}
\newtheorem{theorem}{Theorem}[section]
\newtheorem{lemma}[theorem]{Lemma}
\newtheorem{proposition}[theorem]{Proposition}
\newtheorem{corollary}[theorem]{Corollary}
\newtheorem{fact}[theorem]{Fact}
\newtheorem{definition}[theorem]{Definition}
\newtheorem*{restatedclaim}{\RestatedName}
\theoremstyle{claimass}
\newtheorem{assumption}[theorem]{Assumption}
\newtheorem*{restatedassumption}{\RestatedName}
\theoremstyle{claimremark}
\newtheorem{remark}[theorem]{Remark}
\newtheorem{example}[theorem]{Example}
\theoremstyle{claimreq}
\newtheorem{requirement}[theorem]{Requirement}
\newtheorem*{restatedrequirement}{\RestatedName}
%% --- 8. Prominence: main results a margin rule and a bold heading ----------
\newtcolorbox{mainclaim}{enhanced jigsaw, breakable, blanker,
  borderline west={1.5pt}{-8pt}{mainRule}, left=0pt, right=0pt, top=0pt, bottom=0pt,
  boxsep=0pt, before skip=\topsep, after skip=\topsep, parbox=false,
  before upper={\renewcommand*{\ClaimNoteFont}{\bfseries}}}
\newenvironment{technicalclaim}{\renewcommand*{\ClaimHeadFont}{\mdseries\color{techGrey}}%
  \renewcommand*{\ClaimNoteFont}{}}{}
%% --- 9. Non-normative text, the meta line, navigation ------------------------
\makeatletter
% Tie what follows to this paragraph: no page break, no list penalty before
% the next theorem head (the \@nobreak route a section heading uses).
\newcommand{\keepwithnext}{\par\nobreak\vspace{3pt}\@nobreaktrue
  \everypar{\if@nobreak\@nobreakfalse\clubpenalty\@M\else\clubpenalty\@clubpenalty\everypar{}\fi}}
\newcommand{\claimsummary}[2]{% #1 claim id (named destination), #2 summary
  \par\addvspace{\bigskipamount}\Needspace{6\baselineskip}%
  {\small\itshape\color{nnGrey}\noindent\hypertarget{claim-#1}{}#2\par}\keepwithnext}
\newcommand{\claimanchor}[1]{% a claim without a summary: the destination alone
  \par\addvspace{\bigskipamount}\Needspace{6\baselineskip}%
  \phantomsection\hypertarget{claim-#1}{}\keepwithnext}
\newcommand{\claimleadin}[1]{{\color{nnGrey}\noindent#1\par}\keepwithnext}
\newcommand{\claimmeta}[1]{\par\nobreak\vspace{2pt}%
  {\footnotesize\color{metaGrey}\raggedright\interlinepenalty\@M\noindent#1\par}%
  \addvspace{\medskipamount}}
% After an assumption: the sentence generated from its verification field; after
% any claim with one: its plain-text reading (user_statement). Both
% non-normative, both grey.
\newcommand{\claimverification}[1]{\par\nobreak\vspace{2pt}%
  {\footnotesize\color{metaGrey}\noindent#1\par}\nobreak}
\newcommand{\claimreading}[1]{\par\nobreak\vspace{2pt}%
  {\footnotesize\color{nnGrey}\noindent\emph{Reading:}\enspace#1\par}\nobreak}
\makeatother
\newcommand{\claimbookmark}[4]{% #1 id, #2 kind word, #3 label, #4 plain title
  \bookmark[dest=claim-#1,level=2]{#2 \getrefnumber{#3} #4}}
\newcommand{\metasep}{\unskip\enspace\textperiodcentered\enspace\ignorespaces}
\newcommand{\metalabel}[1]{\emph{#1}\enspace}
\newcommand{\sectionintro}[1]{{\color{nnGrey}\noindent#1\par}\smallskip}
\newcommand{\sectionimports}[1]{{\footnotesize\color{metaGrey}\raggedright\noindent\metalabel{Imports}#1\par}\smallskip}
\newenvironment{overview}{\par\color{nnGrey}\itshape\noindent\ignorespaces}{\par}
% References: the whole "Def. 2.1" is the link; ids are the fallback text of
% every link into another PDF (a relative URI to a named destination).
\newcommand{\mref}[2]{\hyperref[#2]{#1~\ref*{#2}}}
\newcommand{\aref}[2]{\hyperref[#2]{\textsf{#1}}}
\DeclareUrlCommand{\claimid}{\urlstyle{sf}\def\UrlBreaks{\do\/\do\:\do\-\do\.\do\_}%
  \def\UrlBigBreaks{}\def\UrlNoBreaks{}}
\newcommand{\extref}[2]{\hyperref{#1}{}{claim-#2}{\claimid{#2}}}
\newcommand{\extaref}[3]{\hyperref{#1}{}{claim-#2}{\textsf{#3}}}
\newcommand{\coderef}[2]{\href{#1}{\claimid{#2}}}
% Claim-body commands (spec/claim-format.md section 4): \xref{<full id>} prints
% the kind, title and id of a claim of another contract, linked to its stable
% destination in that contract's PDF (one \DeclareXref per cited claim below);
% \term[gloss]{phrase} marks a defined phrase for the index (the gloss feeds
% the Terms table); \tref[phrase]{text} is a use of a defined term. The
% assembler rewrites both in the body into \termdef (the definition's named
% destination term-<slug>) and \termref / \exttermref (a link to it, in this
% or another contract's PDF); the fallbacks below serve text it copies
% elsewhere, such as the digest. \iface is an interface key.
\makeatletter
\newcommand{\DeclareXref}[4]{% id, kind word, title, contract PDF
  \@namedef{xref@#1}{\hyperref{#4}{}{claim-#1}{#2 ``#3'' ({\small\claimid{#1}})}}}
\newcommand{\xref}[1]{\@ifundefined{xref@#1}%
  {\textbf{[unresolved \nolinkurl{#1}]}\@latex@warning{Unresolved xref #1}}%
  {\@nameuse{xref@#1}}}
\makeatother
\DeclareUrlCommand{\iface}{\urlstyle{tt}}
\newcommand{\term}[2][]{\emph{#2}\index{#2}}
\newcommand{\tref}[2][]{#2}
\newcommand{\termdef}[2]{\hypertarget{#1}{}\emph{#2}\index{#2}}
\newcommand{\termref}[2]{\hyperlink{#1}{#2}}
\newcommand{\exttermref}[3]{\hyperref{#1}{}{#2}{#3}}
\hypersetup{pdftitle={Transmitter source model (transmitter) -- contract 0.13.0}}

% Claims of other contracts cited here (\xref, [[...]]).
\DeclareXref{meat-core/def:declaration}{Definition}{The protocol declaration as the MEAT core reads it}{../meat-core/contract.pdf}
\DeclareXref{meat-core/def:roundmodel}{Definition}{Round model and licensed attack sets}{../meat-core/contract.pdf}
\DeclareXref{meat-core/lem:product-marginal}{Lemma}{Product marginal of Alice's retained systems}{../meat-core/contract.pdf}
\DeclareXref{meat-core/rem:modes-da1}{Remark}{The two source modes on the program}{../meat-core/contract.pdf}
\DeclareXref{protocol-compiler/ass:link}{Assumption}{Soundness of every reduction link}{../protocol-compiler/contract.pdf}
\DeclareXref{protocol-compiler/def:chaincond}{Definition}{Interface conditions on a chain}{../protocol-compiler/contract.pdf}
\DeclareXref{protocol-compiler/def:declaration}{Definition}{Protocol declaration}{../protocol-compiler/contract.pdf}
\DeclareXref{protocol-compiler/def:eligible}{Definition}{Key-length eligibility}{../protocol-compiler/contract.pdf}
\DeclareXref{protocol-compiler/prop:conservative}{Proposition}{Dropping a trusted marginal enlarges the feasible set}{../protocol-compiler/contract.pdf}
\DeclareXref{protocol-compiler/thm:chain-consumption}{Theorem}{Decided identities and named obligations of a consumed chain}{../protocol-compiler/contract.pdf}
\DeclareXref{reductions/def:chain}{Definition}{Reduction chain}{../reductions/contract.pdf}
\DeclareXref{reductions/def:link}{Definition}{Typed link}{../reductions/contract.pdf}
\DeclareXref{reductions/lem:precomposition}{Lemma}{Licensed admissibility survives precomposition with a source channel}{../reductions/contract.pdf}
\DeclareXref{reductions/lem:source-transfer}{Lemma}{Exactness of the source map in trusted-source mode}{../reductions/contract.pdf}
\DeclareXref{reductions/thm:two-sided}{Theorem}{The two-sided reduction in trusted-source mode}{../reductions/contract.pdf}

\begin{document}
\begin{center}
{\LARGE\bfseries Transmitter source model (transmitter)}\\[4pt]
{\large Contract, version 0.13.0}
\end{center}
\begin{overview}
The transmitter module describes Alice's side of a protocol: what she prepares, the system she retains, and whether her source is trusted. Its input is a \termref{term-source-specification}{source specification}, with an ordered alphabet, a retained system carrying a \termref{term-sector-structure}{sector structure}, her measurement and the signal ensemble, declared as binary64 matrices or through a submodule's \termref{term-source-model}{source model} and connected by \termref{term-source-link}{source links} to the \termref{term-virtual-ensemble}{virtual ensemble} on which the certificate is written. In \termref{term-trusted-source-mode}{trusted-source mode} a \termref{term-source-replacement-certificate}{source-replacement certificate}, declared as a pure projector or as positive sector blocks or defined from the virtual ensemble, describes the preparation as a measurement of Alice's retained system. The module's main results are these: \hyperlink{claim-transmitter/thm:certificate-sufficiency}{Theorem~\ref*{thm:certificate-sufficiency}} shows that the certificate's algebraic conditions are, in either declared form, finitely many exact decisions and that her measurement commutes with every attack that leaves her retained system and outcome alone, which retention and the sequential round structure make every attack of a round; \hyperlink{claim-transmitter/lem:ensemble-presentation}{Lemma~\ref*{lem:ensemble-presentation}} constructs from the virtual ensemble, given as sector blocks, a certificate with pure blocks whose defining identities hold with no error term; and \hyperlink{claim-transmitter/lem:matched-pinching}{Lemma~\ref*{lem:matched-pinching}} states what a block certificate implies for the states an attack reaches. Every condition the module decides is decided exactly on the declared entries, while the retention of Alice's system remains an attestation that the module records and never verifies. The module emits an artefact: the declared data with the \termref{term-source-segment}{source segment}, a licence record for a verified certificate, which lets a consumer fix Alice's marginal, and an assumption ledger naming the assumptions the artefact rests on. Its consumers are the protocol compiler, the reductions module and the key-rate engine of meat-core; the module itself states no key rate, builds no optimisation problem and restricts no feasible set. Its decoy submodule declares a weak-coherent-pulse source and the tagging link that makes it finite.
\end{overview}
\medskip
\noindent{\small\textcolor{nnGrey}{\itshape Grey text}, here and throughout, is non-normative: it explains, and nothing in it is assumed or proved.}

\smallskip
\noindent{\small The \href{https://github.com/fraunhofersg/certiqs-key/blob/v0.1.0/spec/modules/transmitter/design.md}{design notes} (non-normative) explain why the module is built as it is and draw its processing flowchart. The dependencies between the sections are drawn in \href{graphs/overview.svg}{the section dependency matrix}, and \href{https://github.com/fraunhofersg/certiqs-key/blob/v0.1.0/spec/modules/transmitter/graphs/README.md}{the graph index} opens every section's dependency graph. The \href{digest.pdf}{statements digest} holds every definition, assumption, result and requirement without the proofs. \href{https://github.com/fraunhofersg/certiqs-key/blob/v0.1.0/harness/contract-reading.md}{How to read a contract} explains this layout and its links.}

{\small\tableofcontents}

\section*{Main results}\addcontentsline{toc}{section}{Main results}
\sectionintro{The module's most important results, those its overview names as what it establishes, with every assumption each rests on, directly or through the claims it uses. Each links to its statement and proof; the statements digest holds them in full.}
\par\noindent\hyperref[lem:ensemble-presentation]{\textbf{Lemma~\ref*{lem:ensemble-presentation}}}\enspace Certificate defined by an ensemble presentation\par\nobreak
{\small\itshape\color{nnGrey}\noindent A virtual ensemble in sector blocks defines, by canonical purification and a square-root measurement, a block certificate with pure blocks that exactly reproduces Alice's marginal and, under that measurement, the ensemble. Only her marginal, a transposed sum of declared blocks, is formed as numbers.\par}\nobreak
\claimmeta{\metalabel{Rests on}no assumption}
\par\noindent\hyperref[thm:certificate-sufficiency]{\textbf{Theorem~\ref*{thm:certificate-sufficiency}}}\enspace Checkability and sufficiency of the source-replacement certificate\par\nobreak
{\small\itshape\color{nnGrey}\noindent In any form (pure projector, sector blocks, ensemble presentation) the certificate's algebraic conditions are finitely many exact decisions making Alice's measurement reproduce the virtual ensemble. Under retention and sequential rounds she may measure before or after any attack: source replacement.\par}\nobreak
\claimmeta{\metalabel{Rests on}\aref{A-SEQ}{ass:seq}, \aref{A-SRC}{ass:src}}
\par\noindent\hyperref[lem:matched-pinching]{\textbf{Lemma~\ref*{lem:matched-pinching}}}\enspace Matched pinching of a block certificate\par\nobreak
{\small\itshape\color{nnGrey}\noindent A block certificate is unchanged by pinching Alice's sectors or the virtual signal's sectors paired with them, so every state an admitted attack reaches from it is block diagonal with Alice's declared marginal. An operator merely constrained to have that marginal need not be block diagonal.\par}\nobreak
\claimmeta{\metalabel{Rests on}\aref{A-SRC}{ass:src}}
\par\noindent\hyperref[prop:tagging]{\textbf{Proposition~\ref*{prop:tagging}}}\enspace Tagging channel for weights at or below the Poisson weights\par\nobreak
{\small\itshape\color{nnGrey}\noindent One channel, independent of the label, maps each tagged virtual state to its phase-randomised weak coherent state whenever the declared sector weights lie at or below the Poisson weights. The weights therefore need only a certified inequality, never an exact Poisson weight.\par}\nobreak
\claimmeta{\metalabel{Rests on}no assumption}
\par\noindent\hyperref[cor:emitted-decoy]{\textbf{Corollary~\ref*{cor:emitted-decoy}}}\enspace The emitted decoy declaration\par\nobreak
{\small\itshape\color{nnGrey}\noindent The emitted decoy declaration has an exact certificate against the tagged ensemble, its source identity by the tagging channel, and block-diagonal attacked states. Exactness is decided on products of declared weights and formed dyad entries; lowering a weight gives another tagged declaration.\par}\nobreak
\claimmeta{\metalabel{Rests on}\aref{A-SEQ}{ass:seq}, \aref{A-SRC}{ass:src}, \aref{A-WCP}{ass:wcp}}

\section*{Assumptions at a glance}\addcontentsline{toc}{section}{Assumptions at a glance}
\sectionintro{Every assumption this contract states or its claims assume, with its category and the claims of this contract that assume it directly. An assumption of another contract links into that contract's PDF.}
{\footnotesize\begin{xltabular}{\linewidth}{@{}>{\raggedright\arraybackslash}p{1.6cm}>{\raggedright\arraybackslash}p{1.9cm}>{\raggedright\arraybackslash}X>{\raggedright\arraybackslash}p{2.3cm}>{\raggedright\arraybackslash}p{2.8cm}@{}}
\toprule
Assumption & Category & Title & Stated as & Assumed directly by\\
\midrule
\endhead
\aref{A-ENS}{ass:ens} & source & Emission of the specified ensemble & \mref{Assumption}{ass:ens} & \mref{Fact}{fact:sr}\\
\addlinespace
\aref{A-SRC}{ass:src} & source & Retention of the system A in trusted-source mode & \mref{Assumption}{ass:src} & \mref{Thm.}{thm:certificate-sufficiency}, \mref{Lem.}{lem:matched-pinching}, \mref{Cor.}{cor:emitted-decoy}\\
\addlinespace
\aref{A-WCP}{ass:wcp} & source & Emission of the declared weak-coherent-pulse states & \mref{Assumption}{ass:wcp} & \mref{Cor.}{cor:emitted-decoy}\\
\addlinespace
\aref{A-SEQ}{ass:seq} & protocol structure & Sequential round structure & \mref{Assumption}{ass:seq} & \mref{Fact}{fact:sr}, \mref{Thm.}{thm:certificate-sufficiency}, \mref{Rem.}{rem:xreg}, \mref{Cor.}{cor:emitted-decoy}\\
\bottomrule
\end{xltabular}}

\section*{Notation}\addcontentsline{toc}{section}{Notation}
\sectionintro{The symbols of the repository's notation catalogue (\href{../../notation.md}{spec/notation.md}) that this contract uses, in the catalogue's order, with their meaning.}
{\small\begin{xltabular}{\linewidth}{@{}>{\raggedright\arraybackslash}p{3.6cm}>{\raggedright\arraybackslash}X@{}}
\toprule
Symbol & Meaning\\
\midrule
\endhead
$\Hs{A}$ & Hilbert space of the system label $A$\\
$\Lin{A}$ & Bounded linear operators on $\mathcal{H}(A)$\\
$\Pos{A}$ & Positive semidefinite operators\\
$\Seq{A}$ & Normalized states\\
$\id$, $\idon{C}$ & Identity operator (contextual / named)\\
$\idmap$ & Identity channel\\
$\Xset$ & Outcome alphabet of the measurement on $A$ \textemdash{} Alice's under the default partition\\
$\CList$ & Declared constraint list $\{(\hat O_k,\bowtie_k,\Theta_k)\}$\\
$\Dset$ & Polarization / detector labels\\
$\on$, $\off$, $\dc$, $\sq$, $\mc$, $\nc$, $\singlec$, $\vac$ & Mode / outcome abbreviations\\
$\tst$, $\gen$ & Round-type values\\
$\transpose$ & Transpose, $X^{\mathsf T}$\\
\bottomrule
\end{xltabular}}

\section*{Terms}\addcontentsline{toc}{section}{Terms}
\sectionintro{The phrases this contract defines, and the terms of other contracts it uses, each with a one-line gloss and the definition that fixes it. In the text, a linked term opens its definition.}
{\small\begin{xltabular}{\linewidth}{@{}>{\raggedright\arraybackslash}p{3.6cm}>{\raggedright\arraybackslash}X>{\raggedright\arraybackslash}p{2.6cm}@{}}
\toprule
Term & Meaning & Defined in\\
\midrule
\endhead
\exttermref{../foundations/contract.pdf}{term-binary64-numbers}{binary64 numbers} & the finite values of the IEEE 754 binary64 format, each a dyadic rational & \extref{../foundations/contract.pdf}{foundations/def:binary64}\\
\addlinespace[2pt]
\termref{term-block-form}{block form} & The certificate form declared as positive blocks, one per sector, on matched sector structures of Alice's system and the virtual system. & \mref{Definition}{def:srcert}\\
\addlinespace[2pt]
\termref{term-declared-descriptor-family}{declared descriptor family} & The family of source descriptors for which a source link's channel identity is asserted. & \mref{Definition}{def:source-link}\\
\addlinespace[2pt]
\exttermref{../foundations/contract.pdf}{term-declared-order}{declared order} & a total order on a finite alphabet, fixed together with the alphabet & \extref{../foundations/contract.pdf}{foundations/def:declared-order}\\
\addlinespace[2pt]
\termref{term-ensemble-presentation}{ensemble presentation} & The certificate form that declares only the weighted virtual ensemble in sector blocks and defines the certificate from it. & \mref{Definition}{def:srcert}\\
\addlinespace[2pt]
\termref{term-entropy-cutoff}{entropy cutoff} & The largest photon number whose sector of the virtual source stays quantum & \mref{Definition}{def:tagged-source}\\
\addlinespace[2pt]
\termref{term-intensities}{intensities} & The declared mean photon numbers of the weak coherent pulses, non-negative reals & \mref{Definition}{def:wcp-model}\\
\addlinespace[2pt]
\termref{term-interface-conditions}{interface conditions} & The conditions (S1) to (S4) a source segment satisfies: descriptor match, index preservation, the source identity, and the slots each link carries. & \mref{Definition}{def:source-segment}\\
\addlinespace[2pt]
\termref{term-matched-subspace}{matched subspace} & The direct sum, over the sectors, of the tensor products of matching sectors of Alice's system and the virtual system. & \mref{Definition}{def:srcert}\\
\addlinespace[2pt]
\termref{term-obligation-record}{obligation record} & What a theorem-supplied source link carries in place of a decided identity: its kind, the soundness claim, the assumptions the obligation carries, and the family member its descriptors are. & \mref{Definition}{def:source-link}\\
\addlinespace[2pt]
\termref{term-one-sector-structure}{one-sector structure} & The default sector structure, whose single sector is the whole retained system. & \mref{Definition}{def:source-spec}\\
\addlinespace[2pt]
\termref{term-poisson-weight}{Poisson weight} & The probability that a phase-randomised coherent pulse of a given intensity carries a given number of photons & \mref{Definition}{def:wcp-model}\\
\addlinespace[2pt]
\termref{term-polarisation-flag}{polarisation flag} & A one-dimensional sector for one photon number between the two cutoffs and one polarisation, shared by all intensities & \mref{Definition}{def:tagged-source}\\
\addlinespace[2pt]
\termref{term-polarisations}{polarisations} & The declared polarisation labels, each standing for a unit vector of the one-photon space that is an object of the model & \mref{Definition}{def:wcp-model}\\
\addlinespace[2pt]
\termref{term-pure-form}{pure form} & The certificate form declared as a rank-one projector, on the one-sector structure. & \mref{Definition}{def:srcert}\\
\addlinespace[2pt]
\termref{term-retained-register-sector-map}{retained-register sector map} & An optional matched pair of sector structures, on Alice's retained system and on the descriptor the certificate is written against. & \mref{Definition}{def:source-link}\\
\addlinespace[2pt]
\termref{term-round-type}{round type} & The third component of a letter, marking a test round or a generation round & \mref{Definition}{def:wcp-model}\\
\addlinespace[2pt]
\exttermref{../foundations/contract.pdf}{term-rounds-nowhere}{rounds nowhere} & said of an accumulation every one of whose binary64 additions is exact & \extref{../foundations/contract.pdf}{foundations/def:declared-order}\\
\addlinespace[2pt]
\termref{term-sector-index-set}{sector index set} & The ordered sectors of the virtual source, those that some letter puts positive weight on & \mref{Definition}{def:tagged-source}\\
\addlinespace[2pt]
\termref{term-sector-structure}{sector structure} & A decomposition of Alice's retained system into orthogonal sectors, each spanned by declared basis vectors. & \mref{Definition}{def:source-spec}\\
\addlinespace[2pt]
\termref{term-sound}{sound} & Said of a source link whose channel maps every source member to its target member at every member of its descriptor family. & \mref{Definition}{def:source-link}\\
\addlinespace[2pt]
\termref{term-soundness-registry}{soundness registry} & The closed list, emitted with every artefact, of admitted source-link kinds with their soundness claims and carried assumptions. & \mref{Definition}{def:source-segment}\\
\addlinespace[2pt]
\termref{term-source-cutoff}{source cutoff} & The largest photon number that keeps a sector of its own in the virtual source & \mref{Definition}{def:tagged-source}\\
\addlinespace[2pt]
\termref{term-source-descriptor}{source descriptor} & The transmitted system and signal ensemble a source link starts from. & \mref{Definition}{def:source-link}\\
\addlinespace[2pt]
\termref{term-source-link}{source link} & A link of a reduction chain's source stage that maps one Alice-side ensemble to another with the same labels and weights. & \mref{Definition}{def:source-link}\\
\addlinespace[2pt]
\termref{term-source-model}{source model} & A submodule's named definition of the signal ensemble as mathematical objects, possibly infinite dimensional, rather than as matrices. & \mref{Definition}{def:source-spec}\\
\addlinespace[2pt]
\termref{term-source-segment}{source segment} & The ordered, possibly empty list of source links leading from the virtual ensemble to the specified one. & \mref{Definition}{def:source-segment}\\
\addlinespace[2pt]
\termref{term-source-specification}{source specification} & The declared description of Alice's side of a protocol: alphabet, retained system with sectors, measurement, signal ensemble and source mode. & \mref{Definition}{def:source-spec}\\
\addlinespace[2pt]
\termref{term-source-replacement-certificate}{source-replacement certificate} & An operator on Alice's retained system and the virtual transmitted system whose marginal on her side is her declared marginal and which, measured on her side by her POVM, leaves the weighted virtual members, with retention attested. & \mref{Definition}{def:srcert}\\
\addlinespace[2pt]
\termref{term-specified-signal-ensemble}{specified signal ensemble} & The weights and signal states, one per letter, that Alice's source is specified to emit. & \mref{Definition}{def:source-spec}\\
\addlinespace[2pt]
\termref{term-tagged-source-declaration}{tagged source declaration} & A finite virtual source over a weak-coherent-pulse model: two cutoffs, sector weights at most the Poisson weights, flags and one tagging link & \mref{Definition}{def:tagged-source}\\
\addlinespace[2pt]
\termref{term-tail-flag}{tail flag} & A one-dimensional sector for one polarisation and one intensity, carrying the weight the declared sectors leave out & \mref{Definition}{def:tagged-source}\\
\addlinespace[2pt]
\termref{term-tail-weight}{tail weight} & One minus the sum of an intensity's declared sector weights, the weight its tail flags carry & \mref{Definition}{def:tagged-source}\\
\addlinespace[2pt]
\termref{term-target-descriptor}{target descriptor} & The transmitted system and signal ensemble a source link produces. & \mref{Definition}{def:source-link}\\
\addlinespace[2pt]
\termref{term-theorem-constructor}{theorem constructor} & A named claim of a submodule that supplies a source link's channel for each member of a declared family of descriptors. & \mref{Definition}{def:source-link}\\
\addlinespace[2pt]
\termref{term-transmitted-system}{transmitted system} & The system Alice sends into the channel, carrying the specified signal ensemble. & \mref{Definition}{def:source-spec}\\
\addlinespace[2pt]
\termref{term-transported-constraint-entries}{transported constraint entries} & Entries of a consumer's constraint list carried across a link, each with the certificate case that discharges it. & \mref{Definition}{def:source-link}\\
\addlinespace[2pt]
\termref{term-trusted-source-mode}{trusted-source mode} & The source mode in which Alice's retained system never reaches the adversary, with a declared marginal and, above dimension one, a certificate. & \mref{Definition}{def:source-spec}\\
\addlinespace[2pt]
\termref{term-untrusted-source-mode}{untrusted-source mode} & The source mode with no certificate and no marginal, in which the adversary may have prepared Alice's retained system. & \mref{Definition}{def:source-spec}\\
\addlinespace[2pt]
\termref{term-virtual-ensemble}{virtual ensemble} & The finite-dimensional ensemble the source segment starts from and the certificate is written against. & \mref{Definition}{def:source-segment}\\
\addlinespace[2pt]
\termref{term-virtual-transmitted-system}{virtual transmitted system} & The transmitted system of the virtual ensemble, where the source segment starts. & \mref{Definition}{def:source-segment}\\
\addlinespace[2pt]
\termref{term-weak-coherent-pulse-model}{weak-coherent-pulse model} & A source model that declares phase-randomised weak coherent states on a two-mode Fock space as mathematical objects, not as matrices & \mref{Definition}{def:wcp-model}\\
\addlinespace[2pt]
\termref{term-weighted-members}{weighted members} & The positive operators, weight times state, one per letter, through which an ensemble is declared and read. & \mref{Definition}{def:source-spec}\\
\bottomrule
\end{xltabular}}

\clearpage
\part{Mathematics}

\section{Source specifications and source modes}\label{group-transmitter/def:source-spec}
\sectionintro{Alice's side of a protocol is fixed by what she declares and by the mode in which she declares it. \hyperlink{claim-transmitter/def:source-spec}{Definition~\ref*{def:source-spec}} lists the declared data and the number model in which every condition is decided exactly. In \termref{term-trusted-source-mode}{trusted-source mode} her retained system never reaches the adversary and a certificate ties it to the signal ensemble; in \termref{term-untrusted-source-mode}{untrusted-source mode} nothing does. \hyperlink{claim-transmitter/rem:source-mode}{Remark~\ref*{rem:source-mode}} says what each mode asserts, why a declared marginal needs an attestation, and why untrusted-source mode is always the safe choice.}

% claim transmitter/def:source-spec
\claimsummary{transmitter/def:source-spec}{Alice's alphabet, retained system with sectors, measurement and signal ensemble are declared, with binary64 numbers and exact decisions, or defined from declared data or a source model. Trusted-source mode adds a marginal and, above dimension one, a certificate; untrusted-source mode has neither.}
\claimbookmark{transmitter/def:source-spec}{Definition}{def:source-spec}{Source specification}
\claimleadin{%
Alice's side of a protocol is described by what she prepares and what she
keeps. The definition lists the data a specification declares, the
conditions on them, and the objects that are defined from declared data
rather than declared themselves; its closing number model says how declared
data are given and how every condition is decided. In trusted-source mode the specification asserts that Alice's retained system
$A$ never reaches the adversary, and a \termref{term-source-replacement-certificate}{source-replacement certificate}
(Definition~\ref{def:srcert}) describes her preparation as her measurement
$\{M_x\}$ of $A$. In untrusted-source mode no certificate ties $A$ or
$\{M_x\}$ to the signal ensemble: $\{M_x\}$ is Alice's measurement of a system
the adversary may have prepared, and the mode is the conservative
relaxation of trusted-source mode (Remark~\ref{rem:source-mode}). The sector
structure on $A$ is what lets a certificate be declared in \termref{term-block-form}{block form}
(Definition~\ref{def:srcert}), which a source that first draws a sector and
then prepares a state in it needs (Remark~\ref{rem:certificate-forms}).
At $d_A=1$ the two modes coincide (Remark~\ref{rem:da1}).
}
\begin{definition}[Source specification]\label{def:source-spec}
A \termdef{term-source-specification}{source specification} $\Spec_{\mathrm{src}}$ consists of the
following data. A finite non-empty alphabet $\Xset$ with a declared order. A
finite-dimensional system $A$ with a declared computational basis,
$d_A\coloneqq\dim\Hs{A}$, and a declared \termdef{term-sector-structure}{sector structure} on it: an
ordered finite index set $\mathbb{T}$ and a decomposition
$\Hs{A}=\bigoplus_{t\in\mathbb{T}}\Hs{A_t}$ into subspaces each spanned by
vectors of the declared basis, $d_{A_t}\coloneqq\dim\Hs{A_t}\ge1$, with
projectors $\Pi^{A}_t$; the \termdef{term-one-sector-structure}{one-sector structure}, $\mathbb{T}$ a
single index and $A_t=A$, is the default. A POVM $\{M_x\}_{x\in\Xset}$ on
$A$ presented as a block list against that structure,
$M_x=\bigoplus_{t}M^{(t)}_x$ with $M^{(t)}_x\in\Lin{A_t}$, each block
Hermitian and positive semidefinite, and
$\sum_{x\in\Xset}M^{(t)}_x=\idon{A_t}$ for every $t$ with the sum taken in
the declared order. The \termdef{term-specified-signal-ensemble}{specified signal ensemble}
$\{p_x,\rho_x\}_{x\in\Xset}$ on the \termdef{term-transmitted-system}{transmitted system} $\Asig$, with
$p_x\ge0$, $\sum_xp_x=1$ and $\rho_x$ a state on $\Asig$. With one sector
the block list $M_x$ is the matrix $M_x$, and every condition on its blocks
is the same condition on that matrix.

\emph{How the specified ensemble is declared.} In one of two ways.
(a) \emph{As matrices}, for a finite-dimensional $\Asig$: by its
\termdef{term-weighted-members}{weighted members} $\sigma_x\in\Pos{\Asig}$, with
$p_x\coloneqq\tr\sigma_x$, $\sum_xp_x=1$ in the declared order, and
$\rho_x\coloneqq\sigma_x/p_x$ where $p_x>0$. A declaration by the pair
$(p_x,\rho_x)$, $\rho_x\in\Seq{\Asig}$, is the declaration of the weighted
member $\sigma_x=p_x\rho_x$ with the product formed exactly, and it requires
$\rho_x\in\Seq{\Asig}$ of \emph{every} letter, a letter with $p_x=0$
included. Every condition of this module that reads an admitted ensemble
reads it through its weighted members only.
(b) \emph{By a} \termdef{term-source-model}{source model}, for a transmitted system $\Asig$ whose
Hilbert space is separable, possibly infinite dimensional: by a named
definition of a submodule of this module that fixes the weights $p_x$ and the
states $\rho_x$ as mathematical objects and not as matrices, together with
the physical-source assumption which that submodule declares in the place of
Assumption~\ref{ass:ens}. A specification declared by a source model is
presented with a non-empty \termref{term-source-segment}{source segment}
(Definition~\ref{def:source-segment}): what a consumer receives is never the
model but the \termref{term-virtual-ensemble}{virtual ensemble}, which is finite dimensional.

\emph{The two modes.} The specification is declared in one of two modes. In
\termdef{term-trusted-source-mode}{trusted-source mode} it additionally carries a marginal
$\tau_A\in\Seq{A}$, presented as a block list
$\tau_A=\bigoplus_t\tau^{(t)}_A$ with $\tau^{(t)}_A\in\Pos{A_t}$ and
$\sum_t\tr\tau^{(t)}_A=1$, and, at $d_A\ge2$, a source-replacement
certificate in the sense of Definition~\ref{def:srcert}, written against the
\emph{virtual} ensemble $\{p_x,\rvirt_x\}$ on the \termref{term-virtual-transmitted-system}{virtual transmitted system}
$\Avirt$: the finite-dimensional ensemble the source segment starts from, whose
channels produce the specified ensemble from it
(Definition~\ref{def:source-segment}\,(S1), (S3)); when the source segment is
empty, $\Avirt=\Asig$ and $\rvirt_x=\rho_x$. The certificate takes one of the
three forms of Definition~\ref{def:srcert}: a rank-one projector (the \termref{term-pure-form}{pure
form}), positive sector blocks (the block form), or an \termref{term-ensemble-presentation}{ensemble presentation},
which declares only the weighted virtual ensemble, in blocks along a matched
sector structure of $\Avirt$, and defines the certificate from it. A sector structure with more than one sector is declared only in
trusted-source mode and only together with a certificate in block form or an
ensemble presentation. \emph{With an ensemble presentation} the specification
declares neither $A$, its basis and its sector structure, nor $\{M_x\}$, nor
$\tau_A$: they are constituents of the specification that \eqref{eq:enspres}
defines from the declared virtual ensemble, the conditions this definition
places on them hold by Lemma~\ref{lem:ensemble-presentation}\,(i), and of
them only $\tau_A$ is formed as numbers, that is, computed entry by entry from
the declared blocks (Lemma~\ref{lem:ensemble-presentation}\,(iii));
$\{M_x\}$ and the certificate are defined and never computed. At $d_A=1$ no
certificate is demanded and $\tau_A=[1]$ is the only marginal, so that there
the two modes impose the same condition. In
\termdef{term-untrusted-source-mode}{untrusted-source mode} the specification carries no certificate and no
marginal, written
$\tau_A=\nomarg$; the specification with
$d_A=1$ and a one-letter $\Xset$ is its degenerate case, in which Alice
prepares nothing that this module describes.

\emph{What the sector structure is not.} It is a declaration about
$\{M_x\}$ and $\tau_A$, which are block diagonal by their presentation. It
is no statement about any operator on $A$ together with a further system,
and no consumer may read it as one: that a joint state is block diagonal
along $\mathbb{T}$ is either proved of the states a certificate reaches
(Lemma~\ref{lem:matched-pinching}) or licensed by a theorem of the consuming
module that restricts its own feasible set, under that module's own
preconditions, and is never inferred from the presentation of $\tau_A$ or
$\{M_x\}$.

\emph{The number model.} Every declared entry is a binary64 number, hence a
dyadic rational, and is admissible exactly as declared. A source model
(case (b)) declares no entries: its weights and states are mathematical
objects, and what is decided of such a specification is decided on the
entries it does declare, those of its virtual ensemble among them. Every
condition above is decided \emph{exactly} on the declared entries: a
quantity that a condition compares --- a sum in the declared order, a
product, a partial
trace --- is \emph{formed} in exact arithmetic on the declared entries as
well as compared exactly, so that the identity decided is the identity
written and rounding enters nowhere. That the operands are binary64 is not
what makes a decision exact; a product of two binary64 numbers need not be
one. Positive semidefiniteness is decided by an exact decision on the
declared entries and never by a floating-point eigenvalue.
\end{definition}
\claimmeta{\metalabel{Used by}\mref{Rem.}{rem:source-mode}, \mref{Def.}{def:source-segment}, \mref{Fact}{fact:sr}, \mref{Def.}{def:srcert}, \mref{Lem.}{lem:ensemble-presentation}, \mref{Thm.}{thm:certificate-sufficiency}, \mref{Cor.}{cor:noncanonical}, \mref{Lem.}{lem:matched-pinching}, \mref{Rem.}{rem:da1}, \mref{Ass.}{ass:ens}, \mref{Ass.}{ass:src}, \mref{Def.}{def:wcp-model}, \mref{Cor.}{cor:emitted-decoy}, \extref{../meat-core/contract.pdf}{meat-core/lem:product-marginal}, \extref{../protocol-compiler/contract.pdf}{protocol-compiler/def:compile}, \extref{../protocol-compiler/contract.pdf}{protocol-compiler/def:spec}, \extref{../protocol-compiler/contract.pdf}{protocol-compiler/thm:gates}, \extref{../reductions/contract.pdf}{reductions/lem:source-transfer}
\metasep newly derived
\metasep implemented
\metasep \metalabel{Code}\coderef{https://github.com/fraunhofersg/certiqs-key/blob/v0.1.0/src/certiqs_key/transmitter/artifact.py}{certiqs_key.transmitter.artifact:SourceSpecification}}

% claim transmitter/rem:source-mode
\claimsummary{transmitter/rem:source-mode}{Only trusted-source mode asserts that Alice's retained system never reaches the adversary. A declared marginal deletes attacks at d\_A \ensuremath{\ge} 2 and needs an attestation there; zero never means no marginal, and the maximally mixed marginal means none only at d\_A = 1, where it is [1].}
\claimbookmark{transmitter/rem:source-mode}{Remark}{rem:source-mode}{The two source modes}
\begin{remark}[The two source modes]\label{rem:source-mode}
\emph{(a) What each mode asserts.} In \emph{trusted-source mode}
($\tau_A\in\Seq{A}$) the specification asserts, through
Assumption~\ref{ass:src}, that Alice's register $A$ never passes through
the adversary's hands, so that in every round the joint state of $A$, the
transmitted system and the adversary's registers has $A$-marginal $\tau_A$.
That is the content of the source-replacement step, and the published
analyses state it in those words:
\cite[Eqs.~(36)--(39)]{winick2018reliable} adds marginal constraints to
the entanglement-based statistics constraints \emph{for prepare-and-measure
protocols}, on the stated ground that ``system $A$ is inaccessible to Eve,
and hence [$\rho_A$] is fixed, independent of Eve's attack'';
\cite[Rem.~1]{kamin2025renyi} imposes
$\tr_{B^n}[\sigma_{A^nB^n}]=\tau_A^{\otimes n}$ ``because we assume Eve has
no access to Alice's lab''. That $n$-round product form of the condition is
the consuming engine's \xref{meat-core/lem:product-marginal}, which derives
it from Assumption~\ref{ass:src}, the independence of the source's rounds
(Assumption~\ref{ass:ens} or the premise a \termref{term-source-model}{source model} declares in
its place) and Assumption~\ref{ass:seq}, together with that engine's
premises on private local randomness, a sealed laboratory and the finite
dimension of the adversary's registers. In \emph{untrusted-source mode}
($\tau_A=\nomarg$) the specification asserts nothing about $A$: the
adversary may prepare $A$ jointly with the transmitted system, and no
marginal condition is asserted at all. $\{M_x\}$ is then Alice's
measurement of whatever system $A$ reaches her, and no condition relates
$A$ or $\{M_x\}$ to the signal ensemble, the certificate that would do so
being absent; the consuming engine's round model accordingly lets the
attack supply $A$ (\xref{meat-core/def:roundmodel}). A protocol with no
trusted source can also be declared at $d_A=1$, with Alice's measured system
placed on the receiver side (Remark~\ref{rem:da1}).

\emph{(b) The mode is not a free choice.} A declared $\tau_A$ \emph{deletes}
from a consumer's feasible set every state whose $A$-marginal differs,
hence every attack that would have produced one. That is a different
\emph{kind} of input from a constraint on a register the adversary holds,
and a more demanding one to justify: such a constraint restricts what the
adversary's channel may do to a register she holds, whereas $\tau_A$
asserts that a register never reached her at all. Declaring one at
$d_A\ge2$ without the attestation of Assumption~\ref{ass:src} overstates the
key in the unsafe direction; at $d_A=1$ it deletes nothing
(Remark~\ref{rem:da1}).
Untrusted-source mode is always available and always more conservative:
every state feasible with the marginal declared is feasible without it, so
dropping the marginal loses rate, never gains any, and can never turn an
unsound declaration into a sound-looking one. That inclusion is proved
where the feasible sets are defined, in \xref{meat-core/rem:modes-da1}\,(a)
and \xref{protocol-compiler/prop:conservative}.

\emph{(c) How ``no marginal'' is expressed.} The way to say ``no marginal''
is $\tau_A=\nomarg$. The value $\tau_A=0$ is inadmissible, for two
independent reasons. First, trusted-source mode requires
$\tau_A\in\Seq{A}$, a normalized state, which $\tau_A=0$ fails. Second,
even with normalization dropped it would express the opposite of leaving
the marginal free: a state $\omega$ on $A$ and the transmitted system with
$\omega_A=0$ has $\tr\omega=\tr\omega_A=0$ and hence $\omega=0$, so the
condition admits no state at all. Nor, at $d_A\ge2$, is
$\tau_A=\idon{A}/d_A$ a neutral way to say that nothing is known about the
source: it fixes the marginal to the maximally mixed state and excludes
every adversarial source that sends anything else. At $d_A=1$ it is $[1]$,
the only state on $A$, and imposes the same condition as $\nomarg$
(Remark~\ref{rem:da1}).
\end{remark}
\claimmeta{\metalabel{Uses}\mref{Def.}{def:source-spec}
\metasep newly derived
\metasep accepted}

\section{Source links and the source segment}\label{group-transmitter/def:source-link}
\sectionintro{A security proof may analyse a replacement of the protocol, reached by a reduction chain whose source stage, the \termref{term-source-segment}{source segment}, changes what Alice emits and whose receiver stage changes the measurement. A \termref{term-source-link}{source link} maps one Alice-side ensemble to another with the same labels and weights, and a segment of links runs from the \termref{term-virtual-ensemble}{virtual ensemble}, on which the certificate is written, to the specified one; each channel is declared data or a named theorem of a submodule, and in the second case the link carries an \termref{term-obligation-record}{obligation record}. Because each link applies one channel to the \termref{term-transmitted-system}{transmitted system} for every label, the adversary could apply it herself, so, in \termref{term-trusted-source-mode}{trusted-source mode}, security of the virtual protocol carries over to the specified one (\xref{reductions/lem:source-transfer}). \hyperlink{claim-transmitter/def:source-segment}{Definition~\ref*{def:source-segment}} states which link kinds are admitted and what a consumer decides of a presented segment; the one admitted kind is the decoy submodule's tagging link (\hyperlink{claim-transmitter/decoy/prop:tagging}{Proposition~\ref*{prop:tagging}}).}

% claim transmitter/def:source-link
\claimsummary{transmitter/def:source-link}{A source link maps one Alice-side signal ensemble to another with the same labels and weights, through a channel given as data or by a named theorem. A theorem-supplied channel carries an obligation record, whose registry entry a consumer decides and whose family membership it does not.}
\claimbookmark{transmitter/def:source-link}{Definition}{def:source-link}{Source link}
\claimleadin{%
A source link replaces Alice's signal ensemble by another one, carrying the
labels and the weights through unchanged; along a source segment the links
run from the virtual ensemble a certificate is written against to the
specified ensemble (Definition~\ref{def:source-segment}\,(S1), (S3)). Besides
its channel, the tuple records the family of descriptors the channel's
identity is asserted for, which is what makes a link sound rather than merely
well typed. The one kind admitted is the decoy submodule's tagging kind, whose
channel is supplied by a theorem (Proposition~\ref{prop:tagging}); the kind
\textsf{source-map} that this module declares has no soundness claim and is
refused. How source links compose with the receiver-stage links of a chain is
stated by the module that composes chains; which links are admitted at all is
Definition~\ref{def:source-segment}.
}
\begin{definition}[Source link]\label{def:source-link}
A \termdef{term-source-link}{source link} $\Xi$ is a typed link of a reduction chain
(\xref{reductions/def:chain}) at the stage
\textsf{source}: a tuple
\begin{equation}\label{eq:source-link}
  \Xi=\Big(\textsf{kind},\
  \big(A^{\mathrm{in}},\{\sigma^{\mathrm{in}}_x\}_{x\in\Xset}\big),\
  \big(A^{\mathrm{out}},\{\sigma^{\mathrm{out}}_x\}_{x\in\Xset}\big),\
  \Gamma_{\Xi},\ \mathcal{D}_{\Xi},\ \Smap_{\Xi},\
  \CList_{\Xi},\ \textsf{licence},\ \textsf{status}\Big)
\end{equation}
whose slots are the following.

\emph{(a) The kind.} \textsf{kind} is a source-stage link kind. The
vocabulary is the kinds this module declares, of which there is exactly one,
\textsf{source-map}, together with the kinds a submodule of this module
declares in its own interface. A submodule that declares a kind supplies its
soundness claim, and a kind for which no claim of this repository proves the
link's identity is declared and refused
(Definition~\ref{def:source-segment}, ``Which kinds are admitted''), as
\textsf{source-map} is. An admitted kind has an entry in the soundness
registry of that definition, and a kind without one is not admitted whatever
a link declares.

\emph{(b) The two descriptors.} The \termdef{term-source-descriptor}{source descriptor}
$(A^{\mathrm{in}},\{\sigma^{\mathrm{in}}_x\})$ is the \termref{term-transmitted-system}{transmitted system} and
signal ensemble the link starts from and the \termdef{term-target-descriptor}{target descriptor}
$(A^{\mathrm{out}},\{\sigma^{\mathrm{out}}_x\})$ the one it produces, both
indexed by the same alphabet $\Xset$ in the same declared order, and both
read through their \termref{term-weighted-members}{weighted members} $\sigma_x=p_x\rho_x$, as every ensemble
of this module is (Definition~\ref{def:source-spec}): $\sigma^{\mathrm{in}}_x$ and
$\sigma^{\mathrm{out}}_x$ are positive semidefinite with
$\tr\sigma^{\mathrm{in}}_x=\tr\sigma^{\mathrm{out}}_x=p_x$, which is what
carrying the weights untouched means.

\emph{(c) The channel witness, in one of two ways.} $\Gamma_{\Xi}$ is a
channel from $A^{\mathrm{in}}$ to $A^{\mathrm{out}}$, supplied either
\emph{as data} or \emph{as a theorem constructor}.
\begin{enumerate}[label=(c\arabic*),leftmargin=2.6em]
\item \emph{As data}: both systems are finite dimensional and the channel is
declared, so that the identity
$\Gamma_{\Xi}(\sigma^{\mathrm{in}}_x)=\sigma^{\mathrm{out}}_x$ for every
$x\in\Xset$ is an exact matrix identity on declared numbers and is decided as
such (Definition~\ref{def:source-segment}\,(S3)).
\item \emph{As a} \termdef{term-theorem-constructor}{theorem constructor}: the witness is supplied by a
named claim of a submodule of this module, which constructs one channel for
each member of a declared family of descriptors and proves the identity of
(c1) for it. The identity is then \emph{not} decided on numbers: the link
carries an \termdef{term-obligation-record}{obligation record} of four components, the triple
\begin{equation}\label{eq:obligation}
  \big(\textsf{kind},\ \text{the id of the soundness claim},\
  \mathcal{A}(\textsf{kind})\big)
\end{equation}
and, as the fourth, the member of the declared family the link's descriptors
are. The third component is the assumptions the obligation carries, which
Definition~\ref{def:source-segment}\,(A1) defines as the soundness claim's
assumption closure together with the physical-source assumption the
submodule declares for the \termref{term-source-model}{source model} the claim is stated over, an
assumption the soundness claim itself does \emph{not} assume. The obligation
record travels with the link and, where the link carries a licence record, in
it. \emph{Its first three components are decided}, against the soundness
registry of the artefact that carries the link
(Definition~\ref{def:source-segment}\,(A1), (A2)), by a comparison of
declared strings. The fourth is not: which member of the family the
descriptors are is asserted by the module that emits the link and re-proved
by no consumer, and no verdict on a number is carried, there being none to
decide.
\end{enumerate}
A channel supplied as a theorem constructor is how a link whose source or
target is not a matrix is expressed at all, and it is the only way an
obligation enters this module: no claim here discharges one.

\emph{(d) The declared descriptor family.} $\mathcal{D}_{\Xi}$, the
\termdef{term-declared-descriptor-family}{declared descriptor family}, is the family of source descriptors the
witness's identity is asserted for. It contains the declared source
descriptor; for a link whose channel is supplied as data it is that
descriptor alone; for a constructor link it is the family the named claim
proves the identity for. The link is \termdef{term-sound}{sound} when
$\Gamma_{\Xi}(\sigma^{\mathrm{in}}_x)=\sigma^{\mathrm{out}}_x$ holds for
every $x\in\Xset$ at every member of $\mathcal{D}_{\Xi}$; for a data link
this is what (S3) decides, and for a constructor link it is what the named
claim concludes and the obligation record names. A link whose family does not
contain its own declared descriptor asserts nothing about the descriptor it
carries, and is refused. For a constructor link that membership is not a
decidable condition: the family is the extension of a hypothesis of a claim,
the emitting module asserts that its descriptors satisfy it, and what a
consumer decides instead is the triple \eqref{eq:obligation} against the
registry and the provenance of the artefact
(Definition~\ref{def:source-segment}\,(A1) to (A3)). The refusal of a family
that excludes its own descriptor therefore applies where membership is
decidable, to a data link, whose family is by definition its own descriptor,
so that it refuses a data link whose family slot excludes that descriptor;
elsewhere
membership is an obligation of the emitting module, and what the registry
bounds is which claims a link may name for it.

\emph{(e) The retained-register sector map.} $\Smap_{\Xi}$, the
\termdef{term-retained-register-sector-map}{retained-register sector map}, is optional: a link may emit a matched
pair of \termref{term-sector-structure}{sector structures} over one ordered index set, one on Alice's
retained system $A$ and one on the Alice-side descriptor \emph{the
certificate is written against}, which is the \emph{source} descriptor of the
segment's first link, the virtual pair
(Definition~\ref{def:source-segment}\,(S1)), in the sense the \termref{term-block-form}{block form} of
Definition~\ref{def:srcert} is written against. It is
not the descriptor the link produces: a link's target may be a system that
carries no sector structure at all, and for a constructor link whose target
is infinite dimensional it never carries one. The sector map is a
presentation datum: it says how the certificate, the marginal and Alice's
measurement are presented, and it is no statement about any joint state
(Definition~\ref{def:source-spec}, ``What the sector structure is not'').
Where the slot is empty the \termref{term-one-sector-structure}{one-sector structure} is meant.

\emph{(f) The transported constraint entries and the licence.}
$\CList_{\Xi}$ is the list of \termdef{term-transported-constraint-entries}{transported constraint entries}, entries
of a consumer's constraint list carried across the link, each
carrying, \emph{per entry}, the certificate case that discharges it --- the
statement, for that entry, that what the entry asserts of the link's target
is implied by what holds of its source --- and \textsf{licence} is the
link's licence record, the slot every link of a chain has for its transported
entries (\xref{reductions/def:link}), present exactly when
$\CList_{\Xi}\ne\varnothing$. \emph{A non-empty list is refused}: a
consumer's constraint operators live on the measured system, and no claim of
this repository transports an Alice-side entry. An admitted link therefore
carries an empty list and no licence record of its own.

\emph{(g) The status.}
$\textsf{status}\in\{\textsf{exact},\textsf{certified\_outer},
\textsf{exploratory\_numeric}\}$, the link statuses of
\xref{reductions/def:link}, with \textsf{exact} meaning that the
identity of (c1) holds with no error term --- for a constructor link, that
the named claim concludes it with none.

\emph{Alice-side descriptor discipline.} Both descriptors of a source link
are Alice-side: each names a transmitted system and an ensemble on it, and
no condition on a source link mentions the receiver's systems, its outcome
alphabets or its measurements. A source link changes what Alice emits and
nothing else.

\emph{The degenerate link.} A link that supplies its channel as data, whose
declared descriptor family is that one descriptor, whose sector-map slot is
empty and whose entry list is empty gains nothing from its slots (d) and
(e): they carry exactly what the identity of (c1) asserts and the one-sector
structure.
\end{definition}
\claimmeta{\metalabel{Used by}\mref{Def.}{def:source-segment}, \mref{Def.}{def:tagged-source}, \mref{Prop.}{prop:tagging}, \mref{Cor.}{cor:emitted-decoy}, \extref{../reductions/contract.pdf}{reductions/def:chain}, \extref{../reductions/contract.pdf}{reductions/def:link}
\metasep newly derived
\metasep implemented
\metasep \metalabel{Code}\coderef{https://github.com/fraunhofersg/certiqs-key/blob/v0.1.0/src/certiqs_key/transmitter/links.py}{certiqs_key.transmitter.links:SourceLink}}

% claim transmitter/def:source-segment
\claimsummary{transmitter/def:source-segment}{A source segment chains source links from the virtual ensemble to the specified one, preserving labels and weights, with the source identity decided exactly for data links. A link kind is admitted only with a soundness claim, and a consumer decides its registry entry, ledger and provenance.}
\claimbookmark{transmitter/def:source-segment}{Definition}{def:source-segment}{Source segment and its interface conditions}
\claimleadin{%
Where the source of a protocol is not declared directly as the ensemble the
security analysis is about, a chain of source links connects the two. The
definition fixes what a consumer decides of such a chain and which link kinds
may appear in it at all; Remark~\ref{rem:segment-scope} says what these
conditions leave to other claims. The conditions are read by the modules
that consume the artefact this module emits: the declared data with the
segment, a soundness registry of admitted link kinds, and an assumption
ledger listing the assumptions the artefact rests on.
}
\begin{definition}[Source segment and its interface conditions]
\label{def:source-segment}
The \termdef{term-source-segment}{source segment} of a reduction chain is an ordered list
$\Chsrc=(\Xi_1,\dots,\Xi_{L_{\mathrm s}})$ of source links
(Definition~\ref{def:source-link}), $L_{\mathrm s}\ge0$, whose links act on
the transmitted system and replace Alice's signal ensemble; it is presented
together with a \termref{term-source-specification}{source specification} $\Spec_{\mathrm{src}}$
(Definition~\ref{def:source-spec}). The segment satisfies the
\termdef{term-interface-conditions}{interface conditions} when the following hold.
\begin{enumerate}[label=(S\arabic*),leftmargin=3.0em]
\item \emph{(descriptor match)} the target descriptor of $\Xi_i$ equals the
source descriptor of $\Xi_{i+1}$ for $1\le i<L_{\mathrm s}$; the target
descriptor of $\Xi_{L_{\mathrm s}}$ is the specified pair, the transmitted
system $\Asig$ named by $\Spec_{\mathrm{src}}$ with the specified ensemble
$\{p_x,\rho_x\}$ on it; and the source descriptor of $\Xi_1$ is the virtual
pair, the \termdef{term-virtual-transmitted-system}{virtual transmitted system} $\Avirt$ with the
\termdef{term-virtual-ensemble}{virtual ensemble} $\{p_x,\rvirt_x\}_{x\in\Xset}$ on it, on which the
\termref{term-source-replacement-certificate}{source-replacement certificate} is written (Definition~\ref{def:srcert}).
Every descriptor of
this segment is an Alice-side descriptor: no condition of the segment
mentions the receiver's systems, its outcome alphabets or its measurements;
\item \emph{(index preservation)} each link carries the signal index $x$
through unchanged: it maps a family indexed by $\Xset$ to a family indexed
by $\Xset$, with the weights $p_x$ untouched;
\item \emph{(source identity, where checkable)} writing
$\Xi\coloneqq\Xi_{L_{\mathrm s}}\circ\dots\circ\Xi_1$ for the composite
when every link supplies its channel as data,
\begin{equation}\label{eq:srcmap}
  \Xi(\rvirt_x)=\rho_x\qquad\text{for every }x\in\Xset ,
\end{equation}
an \emph{exact} matrix identity, decided entry by entry. The composite
$\Xi$ is \emph{formed} in exact arithmetic as well, so that the identity
decided is the identity written; that the operands are binary64 is not what
makes the decision exact (the number model of
Definition~\ref{def:source-spec}). The channel runs from the \emph{virtual}
states to the \emph{specified} ones, and the object a consumer receives is
the virtual protocol, the protocol with the virtual ensemble in the place of
the specified one. \emph{Where a link supplies its channel as a theorem
constructor} (Definition~\ref{def:source-link}\,(c2)) the identity of
\emph{that} link is not decided on numbers: it is the conclusion of the claim
the link's obligation record names, at the member of that claim's declared
descriptor family which the link's descriptors are, and the identity of the
composite is the composition of the per-link identities --- each one either
decided, where its link supplies data, or concluded by its named claim, where
its link supplies a constructor. In a segment that mixes the two kinds of
witness, each data link's identity is therefore decided link by link. So
\eqref{eq:srcmap} is what the composite satisfies in every case, and what
changes with the kind of witness is
\emph{by what} it is established and therefore what a consumer may re-decide:
nothing of a constructor link's identity, which is why the obligation
travels;
\item \emph{(slots)} each link carries its constraint entries and its
licence slot, both possibly empty, and its status tag; a link whose channel
is a theorem constructor carries in addition its obligation record, the claim
and the assumptions it names, which are the strings (A2) below decides.
\end{enumerate}

\emph{Which kinds are admitted, and what a consumer decides.} \emph{No
source link kind is admitted unless the submodule that declares it supplies
its soundness claim}: a link of kind $k$ may appear in a segment only where a
submodule of this module declares $k$ and supplies, as a claim of this
repository, the statement that the link's witness satisfies the identity of
Definition~\ref{def:source-link}\,(c1) at every member of its declared
descriptor family --- decided there, for a data link, or proved there, for a
constructor link. A link of a kind for which no such claim exists is
\emph{refused}, and so is a link with a non-empty list of transported
constraint entries, whatever its kind: no claim of this repository transports
an Alice-side entry.

That much is a condition on the \emph{kind}, and a consumer cannot decide it
of a claim it does not read. What a consumer decides is the following, and it
is decided on data.
\begin{enumerate}[label=(A\arabic*),leftmargin=3.0em]
\item \emph{(the soundness registry)} this module emits, with every artefact,
a finite list --- the \termdef{term-soundness-registry}{soundness registry} --- with one entry per
admitted source-link kind, \emph{closed} in that a kind without an entry is
not admitted (Definition~\ref{def:source-link}\,(a)), each entry the triple
\begin{equation}\label{eq:registry}
  \big(\textsf{kind},\ \text{the id of its soundness claim},\
  \mathcal{A}(\textsf{kind})\big) ,
\end{equation}
written out in the module's interface and emitted verbatim. \emph{The third
component $\mathcal{A}(\textsf{kind})$ is the assumptions the obligation
carries}, and it has two sources, both of which a registry entry lists:
\begin{enumerate}[label=(\alph*),leftmargin=2.2em]
\item the \emph{assumption closure of the soundness claim}: the assumptions
in the transitive closure of that claim under its \emph{uses} and
\emph{assumes} edges, the claims and assumptions each claim lists as its
dependencies; it may be empty, and is empty if and only if no claim of that
closure, the soundness claim included, lists an assumption among its
dependencies;
\item the \emph{physical-source assumption the submodule declares for the
source model the claim is stated over}, where the claim is stated over one.
That assumption is not a premise of the soundness claim and must not be made
one: what it attests is that the model the claim is about describes the
physical device, which is what an obligation on a physical source owes and a
mathematical identity does not.
\end{enumerate}
The registry of an artefact is the registry of the version of this module that
emitted it, not a datum a request supplies.
\item \emph{(the triple is decided)} a presented chain's link is
\emph{refused} unless the triple (its \textsf{kind}, the claim id of its
obligation record, the assumption short ids of that record) appears in the
registry of the artefact that carries it --- a comparison of finitely many
declared strings, taken entry by entry --- and unless every assumption of that
entry is in the artefact's assumption ledger, the list of assumptions the
artefact rests on. \emph{Both components of
$\mathcal{A}(\textsf{kind})$ are compared}, so a link that names the soundness
claim's closure and omits the model's attestation is refused exactly as one
that names neither. A link supplying its channel as data carries no obligation
record and is decided by (S3) instead.
\item \emph{(provenance)} a non-empty source segment is accepted only inside
an artefact emitted by the module that admits its links. A segment assembled
from links a consumer received in any other way is refused, whatever its
triples: the registry says which kinds this repository has proved, not that a
particular link was built by the module that proved them.
\end{enumerate}
The admission rule on the kind is applied where a link is admitted, which is
this module and the submodule whose kind it is; it is not among the conditions
(S1) to (S4), which are the conditions a consumer decides on a presented
chain, and (A1) to (A3) are conditions on the artefact and its registry rather
than on the segment's descriptors.

\emph{The empty segment.} When $L_{\mathrm s}=0$ the segment is empty, $\Xi$
is the identity, $\Avirt=\Asig$ and $\rvirt_x=\rho_x$, so \eqref{eq:srcmap}
holds by inspection. That is the case of every specification whose ensemble
is declared as matrices and that carries no link. A specification whose
ensemble is declared by a source model is presented with a non-empty segment
(Definition~\ref{def:source-spec}, case (b)), which is the other case and the
one the conditions above are written for.
\end{definition}
\claimmeta{\metalabel{Uses}\mref{Def.}{def:source-spec}, \mref{Def.}{def:source-link}
\metasep \metalabel{Used by}\mref{Rem.}{rem:segment-scope}, \mref{Def.}{def:srcert}, \mref{Rem.}{rem:boundary}, \mref{Def.}{def:tagged-source}, \mref{Cor.}{cor:emitted-decoy}, \extref{../meat-core/contract.pdf}{meat-core/lem:product-marginal}, \extref{../protocol-compiler/contract.pdf}{protocol-compiler/ass:link}, \extref{../protocol-compiler/contract.pdf}{protocol-compiler/def:banda}, \extref{../protocol-compiler/contract.pdf}{protocol-compiler/def:chaincond}, \extref{../protocol-compiler/contract.pdf}{protocol-compiler/def:spec}, \extref{../protocol-compiler/contract.pdf}{protocol-compiler/req:da1-source-identity}, \extref{../protocol-compiler/contract.pdf}{protocol-compiler/thm:band-a}, \extref{../protocol-compiler/contract.pdf}{protocol-compiler/thm:chain-consumption}, \extref{../reductions/contract.pdf}{reductions/def:chain}, \extref{../reductions/contract.pdf}{reductions/def:link}, \extref{../reductions/contract.pdf}{reductions/lem:source-transfer}, \extref{../reductions/contract.pdf}{reductions/thm:two-sided}
\metasep newly derived
\metasep implemented
\metasep \metalabel{Code}\coderef{https://github.com/fraunhofersg/certiqs-key/blob/v0.1.0/src/certiqs_key/transmitter/links.py}{certiqs_key.transmitter.links:check_source_link}}

% claim transmitter/rem:segment-scope
\claimsummary{transmitter/rem:segment-scope}{The registry conditions decide that a link names a claim proved for its kind, carries that claim's assumptions and comes from the admitting module; they do not decide that the claim applies to these descriptors. Composing the segment with the receiver segment is stated elsewhere.}
\claimbookmark{transmitter/rem:segment-scope}{Remark}{rem:segment-scope}{Scope of the source-segment conditions}
\begin{remark}[Scope of the source-segment conditions]\label{rem:segment-scope}
\emph{What (A1) to (A3) decide.} The conditions (A1) to (A3) of
Definition~\ref{def:source-segment} decide that a link's obligation names a
claim of this repository that was proved for its kind, that the assumptions
that obligation carries --- the claim's closure and the attestation for its
model --- travel in the ledger, and that the segment came from the module that
admits it. They do not decide that the named claim's conclusion applies to
these descriptors: family membership is asserted by the emitting module and
recorded as the fourth component of the obligation record
(Definition~\ref{def:source-link}\,(c2), (d)), and no consumer re-proves a
theorem. A link's obligation thus has two halves, what a consumer decides and
what it takes from the emitting module. The second is exactly what rests on
the emitting module for every artefact this repository passes between
modules, and it is the smaller half: the registry closes the case of a link
naming a claim that proves nothing about its kind, which is the case no other
condition reaches.

\emph{What is decided elsewhere.} The definition admits a \emph{kind} as a
type, on the strength of a claim and a registry entry, and it decides of a
presented segment the conditions (S1) to (S4) and (A1) to (A3). It asserts of
no specification that it declares a non-empty segment: whether one does is
\emph{data} of the presented chain, read off the segment itself. Two things a
non-empty segment needs are not the definition's and are named rather than
restated: the rule by which the segment composes with a receiver segment,
which is \xref{reductions/lem:precomposition} and
\xref{reductions/thm:two-sided} with the admission rule of
\xref{reductions/def:chain}, in \termref{term-trusted-source-mode}{trusted-source mode}; and the consumer that
decides that rule on an assembled declaration, which is
\xref{protocol-compiler/def:chaincond}. That consumer also refuses a
non-empty segment in \termref{term-untrusted-source-mode}{untrusted-source mode}.
\end{remark}
\claimmeta{\metalabel{Uses}\mref{Def.}{def:source-segment}
\metasep newly derived
\metasep accepted}

\section{The source-replacement certificate}\label{group-transmitter/fact:sr}
\sectionintro{Source replacement describes a preparation as the measurement of a system Alice retains, so that a security proof can work with an entangled state. The imported construction is canonical, and \hyperlink{claim-transmitter/def:srcert}{Definition~\ref*{def:srcert}} abstracts the property that makes declared data such a description, as an operator on Alice's system and the \termref{term-virtual-transmitted-system}{virtual transmitted system}, the system of the \termref{term-virtual-ensemble}{virtual ensemble} from which the \termref{term-source-segment}{source segment} produces the specified one. It is declared as a pure projector or as positive sector blocks, or defined from the virtual ensemble alone. The section's main result, \hyperlink{claim-transmitter/lem:ensemble-presentation}{Lemma~\ref*{lem:ensemble-presentation}}, constructs the certificate an \termref{term-ensemble-presentation}{ensemble presentation} defines, so that no square root or quotient of a declared number is ever declared. Every decided condition is decided exactly on the declared binary64 entries, in the sense of the number model of \hyperlink{claim-transmitter/def:source-spec}{Definition~\ref*{def:source-spec}}; the retention condition is an attestation.}

% claim transmitter/fact:sr
\claimsummary{transmitter/fact:sr}{Preparing labelled signals is equivalent to preparing purifications with a shield system and measuring a retained label register. This is the entanglement-based description on which the trusted-source constraints of the module rest.}
\claimbookmark{transmitter/fact:sr}{Fact}{fact:sr}{Source replacement for prepare-and-measure protocols}
\begin{fact}[Source replacement for prepare-and-measure protocols]\label{fact:sr}
\emph{(\cite[Lemma~8.3, Eqs.~(83)--(84)]{tupkary2026decoy}, with the
protocol framing of its Def.~8.1 and Def.~8.2.)}
Consider a protocol in which Alice first prepares the global state
$\rho_{X_1^n(A')_1^n}=\bigotimes_{j=1}^{n}\rho^{(j)}_{X_jA_j'}$ with
$\rho^{(j)}_{X_jA_j'}=\sum_{k=1}^{d_A}p^{(j)}_k\dyad{k}_{X_j}\otimes
(\rho^{(j)}_k)_{A_j'}$, where $\{\ket{k}\}_{X_j}$ is the classical basis of
$X_j$; she keeps $X_1^n$ and sends the $(A')_1^n$ systems to the adversary.
Then the state $\rho_{X_1^n(A')_1^n}$ can equivalently be obtained by
\emph{(1)} preparing the global state
$\rho_{\Abar_1^n\Ahat_1^n(A')_1^n}=\bigotimes_j\rho^{(j)}_{\Abar_j\Ahat_jA_j'}$
with
$\rho^{(j)}_{\Abar_j\Ahat_jA_j'}=\sum_{k,i=1}^{d_A}
\sqrt{p^{(j)}_ip^{(j)}_k}\,\ketbra{i}{k}_{\Abar}\otimes
\ketbra{\psi^{(j)}_i}{\psi^{(j)}_k}_{\Ahat_jA_j'}$, where
$\ket{\psi^{(j)}_k}_{\Ahat_jA_j'}$ is a purification of
$(\rho^{(j)}_k)_{A_j'}$, $\{\ket{k}\}$ is an orthonormal basis of $\Abar$
and $\Ahat_j$ is the \emph{shield system}; \emph{(2)} sending $(A')_1^n$ to
the adversary; and \emph{(3)} for each $j$, measuring $\Abar_j\Ahat_j$ with
the POVM $\{\dyad{k}_{\Abar_j}\otimes\idon{\Ahat_j}\}_{k\in\{1,\dots,d_A\}}$
and storing the result in $X_j$.

\emph{Symbol map and hypothesis check.} The source's $X_j$ is this module's
outcome register with alphabet $\Xset$ and $k\leftrightarrow x$; its $A_j'$
is the \termref{term-virtual-transmitted-system}{transmitted system} carrying the ensemble the lemma is applied to,
and its $\{p^{(j)}_k,\rho^{(j)}_k\}$ is that ensemble. That ensemble is the
virtual pair $(\Avirt,\{p_x,\rvirt_x\})$ on which the source-replacement
certificate is stated, which is the specified pair $(\Asig,\{p_x,\rho_x\})$
of Definition~\ref{def:source-spec} exactly when the \termref{term-source-segment}{source segment} is empty.
This module describes one round; that every round is described by the
same data is the consumer's round-homogeneity assumption and is not
asserted here, and it is under that assumption that the superscript $j$ is
dropped in what follows. The source's $\Abar_j\Ahat_j$ is this module's
retained system $A$, and its measurement POVM is $\{M_x\}$ with
$M_x=\dyad{x}_{\Abar}\otimes\idon{\Ahat}$. The $d_A$ of the cited lemma and
the $d_A$ of this module are \emph{not} the same number, and the
identification does not claim they
are: the published index runs $k\in\{1,\dots,d_A\}$ over an orthonormal
basis of $\Abar$, so there $d_A=\dim\Hs{\Abar}=|\Xset|$, whereas throughout
this module $d_A=\dim\Hs{A}=\dim\Hs{\Abar}\cdot\dim\Hs{\Ahat}$; the two
agree exactly when the shield system is trivial. The lemma's one hypothesis
is that Alice prepares the classical-quantum source state of its
Eq.~(83) --- this is Assumption~\ref{ass:ens} for the physical device, and
the declared source data for the specified object. Its only timing element is
that the global state is prepared first; it places no condition on the
sequential structure of later operations, being an identity between the state
of its Eq.~(83) and the measured state of its Eq.~(84). The cited work
applies it within the protocol framing of its Def.~8.1 and Def.~8.2, and this
module applies it to protocols with the sequential structure of
Assumption~\ref{ass:seq}, which supplies that structure: the assumption is the
hypothesis under which this module applies the fact, and not a hypothesis of
the cited lemma. The
published statement quantifies over an
arbitrary purification $\ket{\psi_k}$, so the shield system is trivial
exactly when every state of that ensemble is pure --- the case, for a pure
\emph{specified} ensemble, of \cite[Eqs.~(37)--(39)]{winick2018reliable},
which writes the source-replaced vector as
$\ket{\psi}_{\Abar\Asig}=\sum_i\sqrt{p_i}\ket{i}_{\Abar}\ket{\varphi_i}_{\Asig}$
and records that the resulting marginal
$\rho_{\Abar}=\tr_{\Asig}(\dyad{\psi})$ is fixed independently of the
adversary's attack.
\end{fact}
\claimmeta{\metalabel{Uses}\mref{Def.}{def:source-spec}
\metasep \metalabel{Assumes}\aref{A-ENS}{ass:ens}, \aref{A-SEQ}{ass:seq}
\metasep \metalabel{Used by}\mref{Rem.}{rem:import}, \mref{Cor.}{cor:noncanonical}
\metasep published in \cite[{Lemma 8.3, Eqs. (83)-(84)}]{tupkary2026decoy}; \cite[{Def. 8.1, Def. 8.2}]{tupkary2026decoy}; \cite[{Eqs. (37)-(39)}]{winick2018reliable}
\metasep accepted
\metasep \metalabel{Oracles}\coderef{https://github.com/fraunhofersg/certiqs-key/blob/v0.1.0/tests/oracles/registry_transmitter.py}{registry_transmitter.py}}

% claim transmitter/rem:import
\claimsummary{transmitter/rem:import}{The imported construction gives Alice's retained label register one dimension per letter, with a singular marginal whenever the signal states are linearly dependent, as in BB84. \hyperlink{claim-transmitter/thm:certificate-sufficiency}{Theorem~\ref*{thm:certificate-sufficiency}} removes that constraint; the import says nothing about granting the adversary a purification.}
\claimbookmark{transmitter/rem:import}{Remark}{rem:import}{Scope of the imported source replacement}
\begin{remark}[Scope of the imported source replacement]\label{rem:import}
Fact~\ref{fact:sr} supplies the \emph{canonical} source replacement: one
particular purification and one particular measurement, with $\Abar$ of
dimension $|\Xset|$. A consumer that admitted only that form would be
forced to declare $d_A=|\Xset|$ with a singular marginal in every protocol
whose signal states are not linearly independent --- including BB84, where
the canonical $\tau_{\Abar}$ has rank two in a four-dimensional space.
Theorem~\ref{thm:certificate-sufficiency} removes that constraint: it states the property that
makes a declared triple $(\Psi,\tau_A,\{M_x\})$ a source replacement
(Definition~\ref{def:srcert}), proves that any triple with that property
describes the same protocol, and thereby admits the qubit form that
Corollary~\ref{cor:noncanonical} exhibits for BB84. The import supplies no
statement about the adversary being granted a purification of Alice's
retained system; that step belongs to the consuming engine's round model and
is not claimed here.
\end{remark}
\claimmeta{\metalabel{Uses}\mref{Fact}{fact:sr}
\metasep newly derived
\metasep accepted
\metasep \metalabel{Oracles}\coderef{https://github.com/fraunhofersg/certiqs-key/blob/v0.1.0/tests/oracles/registry_transmitter.py}{registry_transmitter.py}}

% claim transmitter/def:srcert
\claimsummary{transmitter/def:srcert}{An operator on Alice's system and the virtual signal whose partial trace over the signal is her marginal and whose partial traces against her measurement elements are the weighted virtual members. It is declared in pure or block form or defined by an ensemble presentation; retention is attested.}
\claimbookmark{transmitter/def:srcert}{Definition}{def:srcert}{Source-replacement certificate}
\claimleadin{%
The imported source replacement is one canonical way of describing a
preparation as the measurement of a system Alice retains. The certificate
abstracts the property that makes declared data such a description. Its
algebraic conditions are either decided exactly on declared numbers, in the
sense of the number model of Definition~\ref{def:source-spec}, or, in an
ensemble presentation, theorems of the declared data; its retention condition
is an attestation. It comes in three forms because what is exactly declarable
depends on the source. It is written against the
\termref{term-virtual-ensemble}{virtual ensemble} on $\Avirt$, from which the links of the source segment, of
length $L_{\mathrm s}$, produce the specified ensemble on $\Asig$
(Definition~\ref{def:source-segment}); with no links the two coincide.
}
\begin{definition}[Source-replacement certificate]\label{def:srcert}
Let $\{p_x,\rvirt_x\}_{x\in\Xset}$ be the \emph{virtual} signal ensemble on
the \termref{term-virtual-transmitted-system}{virtual transmitted system} $\Avirt$ --- the pair the source segment
starts from, Definition~\ref{def:source-segment}\,(S1), (S3) --- with
\termref{term-weighted-members}{weighted members} $\sigma^{\mathrm v}_x\coloneqq p_x\rvirt_x\in\Pos{\Avirt}$,
$p_x=\tr\sigma^{\mathrm v}_x\ge0$ and $\sum_xp_x=1$. These conditions are
decided \emph{exactly}, as every declared condition is, the weight sum in
the declared order and rounding nowhere. A \termdef{term-source-replacement-certificate}{source-replacement
certificate} for the declared data $(A,\mathbb{T},\tau_A,\{M_x\})$ of a
\termref{term-source-specification}{source specification} (Definition~\ref{def:source-spec}) is an operator
$\Psi\in\Lin{A\Avirt}$ such that
\begin{enumerate}[label=(SR\arabic*),leftmargin=3.0em]
\item $\tr_{\Avirt}\Psi=\tau_A$;
\item $\tr_{A}\!\big[(M_x\otimes\idon{\Avirt})\,\Psi\big]=p_x\,\rvirt_x$
for every $x\in\Xset$; and
\item \emph{(retention, an attestation)} $A$ is retained in Alice's
laboratory and never transmitted; no operation of the protocol and no
operation of the adversary acts on $A$, except Alice's declared measurement
$\{M_x\}$.
\end{enumerate}
It is given in one of the three forms below. In the pure form and in the
block form $\Psi$ is declared, together with the declared computational basis
of $A$. In an ensemble presentation neither $\Psi$ nor a basis of $A$ is
declared: $\Psi$ is defined from a declared ensemble by \eqref{eq:enspres}.

\emph{Pure form.} In the \termdef{term-pure-form}{pure form} the \termref{term-sector-structure}{sector structure} of $A$ is the
\termref{term-one-sector-structure}{one-sector structure}, and $\Psi$ is a \emph{rank-one projector} --- $\Psi$
Hermitian, $\Psi^2=\Psi$ and $\tr\Psi=1$, each decided exactly and the trace
summed without rounding, so that $\Psi=\dyad{\psi}$ for a unit vector
$\ket{\psi}\in\Hs{A\Avirt}$.

\emph{Block form.} Let $\Hs{A}=\bigoplus_{t\in\mathbb{T}}\Hs{A_t}$ be the
sector structure of $A$ declared in Definition~\ref{def:source-spec}, along
which $\tau_A=\bigoplus_t\tau^{(t)}_A$ and $M_x=\bigoplus_tM^{(t)}_x$ are
presented as block lists. In the \termdef{term-block-form}{block form} the virtual transmitted
system carries a declared sector structure \emph{matched} to that of $A$: the
same ordered index set $\mathbb{T}$ and a decomposition
$\Hs{\Avirt}=\bigoplus_{t\in\mathbb{T}}\Hs{S_t}$ into subspaces each spanned
by vectors of the declared basis of $\Avirt$, with projectors
$\Pi^{\mathrm v}_t$. The certificate is a block list
$(\Psi_t)_{t\in\mathbb{T}}$ with $\Psi_t\in\Lin{A_tS_t}$, each block
Hermitian and positive semidefinite, decided exactly, and
$\sum_{t}\tr\Psi_t=1$, the traces summed in the declared order without
rounding. It denotes the operator $\Psi=\bigoplus_t\Psi_t$ on the
\termdef{term-matched-subspace}{matched subspace} $\bigoplus_t\Hs{A_t}\otimes\Hs{S_t}
\subseteq\Hs{A\Avirt}$, extended by zero to its orthogonal complement. No
purity condition is imposed, on $\Psi$ or on any block. The weighted virtual
members are required to be \emph{block diagonal} along $(S_t)_t$:
$\Pi^{\mathrm v}_t\sigma^{\mathrm v}_x\Pi^{\mathrm v}_{t'}=0$ for all
$t\ne t'$ and every $x$, that is, every entry of $\sigma^{\mathrm v}_x$ whose
row lies in one sector and whose column lies in another is zero, decided entry
by entry on the declared entries of $\sigma^{\mathrm v}_x$ (a virtual ensemble
declared as block lists along $(S_t)_t$ meets it by its presentation). Write
$\sigma^{\mathrm v,(t)}_x\coloneqq\Pi^{\mathrm v}_t\sigma^{\mathrm v}_x
\Pi^{\mathrm v}_t$, read as an operator on $S_t$, for its diagonal blocks.
Against a block list the two identities read block by block,
\begin{equation}\label{eq:srblock}
  \tr_{S_t}\Psi_t=\tau^{(t)}_A ,\qquad
  \tr_{A_t}\!\big[(M^{(t)}_x\otimes\idon{S_t})\,\Psi_t\big]
  =\sigma^{\mathrm v,(t)}_x
  \qquad\text{for every }t\in\mathbb{T},\ x\in\Xset ,
\end{equation}
and \eqref{eq:srblock} together with the block diagonality of every
$\sigma^{\mathrm v}_x$ is equivalent to (SR1) and (SR2) of the operator
$\Psi$. Indeed $\Psi$ vanishes off the matched subspace and $M_x$ is block
diagonal, so the partial traces of $\Psi$ are the direct sums of those of its
blocks: $\tr_{\Avirt}\Psi=\bigoplus_t\tr_{S_t}\Psi_t$, which is
$\tau_A=\bigoplus_t\tau^{(t)}_A$ exactly when the first identity of
\eqref{eq:srblock} holds; and $\tr_A[(M_x\otimes\idon{\Avirt})\Psi]=\bigoplus_t
\tr_{A_t}[(M^{(t)}_x\otimes\idon{S_t})\Psi_t]$, which is block diagonal along
$(S_t)_t$ and so equals $\sigma^{\mathrm v}_x$ exactly when
$\sigma^{\mathrm v}_x$ is block diagonal and its diagonal blocks satisfy the
second identity. Without the block diagonality, \eqref{eq:srblock} is (SR2)
against the pinched members $\bigoplus_t\sigma^{\mathrm v,(t)}_x$ and not
against $\sigma^{\mathrm v}_x$.

\emph{Ensemble presentation.} A block-form certificate need not be declared
at all. In an \termdef{term-ensemble-presentation}{ensemble presentation} the declared data are the matched
sector structure $(S_t)_{t\in\mathbb{T}}$ on $\Avirt$ and the weighted
virtual ensemble as block lists,
$\sigma^{\mathrm v}_x=\bigoplus_t\sigma^{\mathrm v,(t)}_x$ with
$\sigma^{\mathrm v,(t)}_x\in\Lin{S_t}$, and \emph{nothing else}: no $\Psi$,
no $\{M_x\}$, no $\tau_A$ and no basis of $A$. The conditions are
\begin{enumerate}[label=(E\arabic*),leftmargin=3.0em]
\item every block $\sigma^{\mathrm v,(t)}_x$ is Hermitian and positive
semidefinite, decided exactly;
\item $p_x=\sum_t\tr\sigma^{\mathrm v,(t)}_x$ for every $x$ and
$\sum_xp_x=1$, the sums formed in the declared orders and rounding nowhere;
and
\item for every $t$ the sum $\sum_{x\in\Xset}\sigma^{\mathrm v,(t)}_x$,
accumulated entry by entry in the declared order of $\Xset$, rounds nowhere,
so that its entries are binary64 numbers.
\end{enumerate}
The system $A$, its sector structure, the marginal $\tau_A$, the POVM
$\{M_x\}$ and the blocks $\Psi_t$ are then \emph{defined} by
\eqref{eq:enspres}: $A_t$ is a copy of $S_t$, $\tau^{(t)}_A$ is the
transpose of the sum in (E3), and $\Psi_t$ and $M^{(t)}_x$ are the canonical
purification and its square-root measurement. By
Lemma~\ref{lem:ensemble-presentation} they are a block-form certificate
with (SR1) and (SR2) holding as identities, with no error term, against the
declared ensemble itself, and with every non-zero block a multiple of a
rank-one projector.
(E1) to (E3) are all that is decided. (SR1), (SR2), positivity of the
blocks, the POVM condition and the purity of the blocks are \emph{theorems}
of the presentation and are decided by nothing, there being no second
declared datum they could be compared with; (SR3) is the attestation it is
in every form. Of the defined objects only $\tau_A$ is ever formed as
numbers, and by (E3) it is admissible declared data. $\Psi_t$ and
$M^{(t)}_x$ carry square roots and quotients of the declared numbers, are
formed by no one, and are never rounded.

\emph{What a consumer reads of a presentation.} What a consumer can read of
Alice's measurement without $\Psi_t$ and $M^{(t)}_x$ is the family
$N^{(t)}_x=(\sigma^{\mathrm v,(t)}_x)^{\transpose}$ of
Lemma~\ref{lem:ensemble-presentation}\,(v). \emph{A consumer whose
declaration carries Alice's measurement as the POVM $\{M_x\}$ cannot receive
an ensemble presentation and refuses it}; the form of declaration that
carries $\{N_x\}$ in its place is that consumer's to define, and the
consuming modules define it: the protocol compiler and the key-rate engine
receive a presentation through the ensemble form of their declarations,
which carries $\{N_x\}$ in the place of $\{M_x\}$
(\xref{protocol-compiler/def:declaration},
\xref{meat-core/def:declaration}), admitted beside a presentation by
\xref{protocol-compiler/def:eligible}. The mode, the alphabet, the source
segment and the attestations of a specification are what they are in the
other forms.

\emph{What is decided.} (SR1) and (SR2) are \emph{exact matrix identities}
on declared data, decided entry by entry. The conditions a consumer decides
--- in pure form the three on $\Psi$, (SR1) and (SR2); in block form
Hermiticity and positive semidefiniteness of every block, the trace, the block
diagonality of every $\sigma^{\mathrm v}_x$ along $(S_t)_t$, and
\eqref{eq:srblock}, which together are (SR1) and (SR2); in an ensemble
presentation (E1) to (E3) and no others
--- are stated in these terms and in no others: none of them is stated on a
vector, which is not a declared datum. (SR3) carries no number and is
decided by nothing; it is the attestation of Assumption~\ref{ass:src}. A
rank-one projector of unit trace is positive semidefinite of unit trace, so
a pure-form certificate satisfies every condition of the block form at the
one-sector structure, and one condition more.

\emph{One reading at every $L_{\mathrm s}$.} The pair
$(\Avirt,\{p_x,\rvirt_x\})$ is the only pair a certificate is ever written
against, at every length of the source segment, because the object a
consumer receives is the virtual protocol
(Definition~\ref{def:source-segment}\,(S3)): the specification's $\tau_A$
and $\{M_x\}$ describe the preparation of $\{p_x,\rvirt_x\}$, and it is the
source identity \eqref{eq:srcmap}, not the certificate, that carries that
description back to the specified ensemble $\{p_x,\rho_x\}$ on $\Asig$. At
$L_{\mathrm s}=0$ the two pairs are literally equal, $\Avirt=\Asig$ and
$\rvirt_x=\rho_x$. So (SR1)--(SR2) have exactly one reading whatever
$L_{\mathrm s}$ is.
\end{definition}
\claimmeta{\metalabel{Uses}\mref{Def.}{def:source-spec}, \mref{Def.}{def:source-segment}
\metasep \metalabel{Used by}\mref{Lem.}{lem:ensemble-presentation}, \mref{Thm.}{thm:certificate-sufficiency}, \mref{Cor.}{cor:noncanonical}, \mref{Rem.}{rem:certificate-forms}, \mref{Lem.}{lem:matched-pinching}, \mref{Rem.}{rem:boundary}, \mref{Rem.}{rem:da1}, \mref{Ass.}{ass:src}, \mref{Req.}{req:licence-scope}, \mref{Req.}{req:failing-certificate}, \mref{Prop.}{prop:decoy-cert}, \mref{Cor.}{cor:emitted-decoy}, \mref{Req.}{req:exact-decoy-blocks}, \mref{Req.}{req:decoy-ensemble-conditions}, \extref{../meat-core/contract.pdf}{meat-core/lem:product-marginal}, \extref{../protocol-compiler/contract.pdf}{protocol-compiler/def:eligible}, \extref{../protocol-compiler/contract.pdf}{protocol-compiler/def:spec}, \extref{../protocol-compiler/contract.pdf}{protocol-compiler/req:da1-source-identity}, \extref{../reductions/contract.pdf}{reductions/lem:source-transfer}
\metasep newly derived
\metasep implemented
\metasep \metalabel{Code}\coderef{https://github.com/fraunhofersg/certiqs-key/blob/v0.1.0/src/certiqs_key/transmitter/certificate.py}{certiqs_key.transmitter.certificate:SourceReplacementCertificate}
\metasep \metalabel{Oracles}\coderef{https://github.com/fraunhofersg/certiqs-key/blob/v0.1.0/tests/oracles/registry_instances_bb84.py}{registry_instances_bb84.py}, \coderef{https://github.com/fraunhofersg/certiqs-key/blob/v0.1.0/tests/oracles/registry_protocol_compiler.py}{registry_protocol_compiler.py}, \coderef{https://github.com/fraunhofersg/certiqs-key/blob/v0.1.0/tests/oracles/registry_transmitter.py}{registry_transmitter.py}}

% claim transmitter/lem:ensemble-presentation
\claimsummary{transmitter/lem:ensemble-presentation}{A virtual ensemble in sector blocks defines, by canonical purification and a square-root measurement, a block certificate with pure blocks that exactly reproduces Alice's marginal and, under that measurement, the ensemble. Only her marginal, a transposed sum of declared blocks, is formed as numbers.}
\claimbookmark{transmitter/lem:ensemble-presentation}{Lemma}{lem:ensemble-presentation}{Certificate defined by an ensemble presentation}
\claimleadin{%
The ensemble presentation of Definition~\ref{def:srcert} declares only the
weighted virtual ensemble, in blocks along a sector structure of the virtual
transmitted system. The lemma supplies the objects it defines and shows that
they form a block-form certificate, so that no square root or quotient of a
declared number ever has to be declared; of the defined objects only Alice's
marginal is formed as numbers, that is, computed entry by entry.
}
\begin{mainclaim}
\begin{lemma}[Certificate defined by an ensemble presentation]\label{lem:ensemble-presentation}
Let the virtual transmitted system carry a sector structure
$\Hs{\Avirt}\supseteq\bigoplus_{t\in\mathbb{T}}\Hs{S_t}$, each $S_t$ spanned
by vectors $(\ket{i}_{S_t})_i$ of the declared basis, and let
$(\sigma^{\mathrm v,(t)}_x)_{t\in\mathbb{T}}$, $x\in\Xset$, be block lists
with $\sigma^{\mathrm v,(t)}_x\in\Pos{S_t}$ and
$\sum_{x}\sum_t\tr\sigma^{\mathrm v,(t)}_x=1$. For every $t$ let $A_t$ be a
copy of $S_t$ with basis $(\ket{i}_{A_t})_i$ in the same order, let
transposition $X\mapsto X^{\transpose}$ be taken in these bases, and put
\begin{equation}\label{eq:enspres}
\begin{gathered}
  \tau^{(t)}_{S}\coloneqq\sum_{x\in\Xset}\sigma^{\mathrm v,(t)}_x ,\qquad
  \tau^{(t)}_A\coloneqq\big(\tau^{(t)}_{S}\big)^{\transpose} ,\\
  \ket{\psi_t}\coloneqq\sum_i\ket{i}_{A_t}\otimes
     \big(\tau^{(t)}_{S}\big)^{1/2}\ket{i}_{S_t} ,\qquad
  \Psi_t\coloneqq\dyad{\psi_t} ,\\
  M^{(t)}_x\coloneqq\big(\tau^{(t)}_A\big)^{-1/2}
     \big(\sigma^{\mathrm v,(t)}_x\big)^{\transpose}
     \big(\tau^{(t)}_A\big)^{-1/2}
     +\delta_{x,x_0}\big(\idon{A_t}-P_t\big) ,
\end{gathered}
\end{equation}
where $P_t$ is the projector onto the support of $\tau^{(t)}_A$, the inverse
square root is taken on that support and is zero off it, and $x_0$ is the
first letter of $\Xset$ in its declared order. Then, with
$A\coloneqq\bigoplus_tA_t$:
\begin{enumerate}[label=(\roman*),leftmargin=2.2em]
\item \emph{(a certificate in block form)} every $\Psi_t$ is positive
semidefinite, $\sum_t\tr\Psi_t=1$, $\{M_x\}_x$ with
$M_x=\bigoplus_tM^{(t)}_x$ is a POVM presented as a block list,
$\tau_A=\bigoplus_t\tau^{(t)}_A$ is a state presented as a block list, and
both identities of \eqref{eq:srblock} hold: $(\Psi_t)_t$ satisfies every
condition Definition~\ref{def:srcert} places on a block-form certificate
for $(A,\mathbb{T},\tau_A,\{M_x\})$ against the virtual ensemble with
weighted members $\sigma^{\mathrm v}_x=\bigoplus_t\sigma^{\mathrm v,(t)}_x$,
(SR3) apart, which is an attestation;
\item \emph{(pure blocks)} every non-zero $\Psi_t$ is
$g_t\coloneqq\tr\tau^{(t)}_S$ times a rank-one projector, so the converse of
Theorem~\ref{thm:certificate-sufficiency}\,(ii) holds of this certificate;
\item \emph{(what is formed as numbers)} the entry $(i,j)$ of
$\tau^{(t)}_A$ is the entry $(j,i)$ of the sum $\tau^{(t)}_S$ of the
declared blocks. No
other object of \eqref{eq:enspres} enters (i) or (ii) through its entries:
$\Psi_t$ and $M^{(t)}_x$ contain square roots and quotients of the declared
numbers and are, in general, not declarable;
\item \emph{(the completion is immaterial)} the letter $x_0$ enters the
POVM condition and nothing else: (SR1), (SR2), (ii) and (v) hold with the
term $\idon{A_t}-P_t$ assigned to any letter, or split among the letters by
any positive semidefinite operators on the kernel of $\tau^{(t)}_A$ that
sum to $\idon{A_t}-P_t$;
\item \emph{(Alice's side without a square root)}
$\big(\tau^{(t)}_A\big)^{1/2}M^{(t)}_x\big(\tau^{(t)}_A\big)^{1/2}
=\big(\sigma^{\mathrm v,(t)}_x\big)^{\transpose}$ for every $x$ and $t$: the
operators $N^{(t)}_x\coloneqq(\sigma^{\mathrm v,(t)}_x)^{\transpose}$ are
entrywise the declared numbers, and
$\sum_xN^{(t)}_x=\tau^{(t)}_A$.
\end{enumerate}
\end{lemma}
\end{mainclaim}
\claimmeta{\metalabel{Uses}\mref{Def.}{def:source-spec}, \mref{Def.}{def:srcert}
\metasep \metalabel{Used by}\mref{Thm.}{thm:certificate-sufficiency}, \mref{Prop.}{prop:decoy-cert}, \mref{Cor.}{cor:emitted-decoy}
\metasep newly derived
\metasep implemented
\metasep \metalabel{Code}\coderef{https://github.com/fraunhofersg/certiqs-key/blob/v0.1.0/src/certiqs_key/transmitter/presentation.py}{certiqs_key.transmitter.presentation:alice_marginal_blocks}}

\begin{proof}
\emph{Idea.} Each block is the canonical purification of the summed
ensemble of its sector, written through the maximally entangled vector; three
facts about transposition reduce every claim to an identity between the
declared blocks and their sum.

\emph{Preliminaries.} Fix $t$ and drop it from the notation; write
$K\coloneqq\tau_S^{1/2}$,
$\ket{\Omega}\coloneqq\sum_i\ket{i}_A\ket{i}_S$, so that
$\ket{\psi}=(\idon{A}\otimes K)\ket{\Omega}$, and let $Q$ be the projector
onto the support of $\tau_S$. Three facts are used. (a)
$\tr_A\big[(Y\otimes\idon{S})\dyad{\Omega}\big]=Y^{\transpose}$ for
$Y\in\Lin{A}$: with $\dyad{\Omega}=\sum_{ij}\ket{i}\!\bra{j}\otimes
\ket{i}\!\bra{j}$ the left-hand side is $\sum_{ij}\bra{j}Y\ket{i}\,
\ket{i}\!\bra{j}$, whose $(i,j)$ entry is $Y_{ji}$. (b)
$(\idon{}\otimes Z)\ket{\Omega}=(Z^{\transpose}\otimes\idon{})\ket{\Omega}$,
by comparing coefficients. (c) Transposition is linear, reverses products,
commutes with the adjoint and maps positive semidefinite operators to
positive semidefinite operators with the same spectrum; hence
$f(X^{\transpose})=f(X)^{\transpose}$ for every real function $f$ on the
spectrum of a Hermitian $X$. In particular $(\tau_A)^{1/2}=K^{\transpose}$,
$(\tau_A)^{-1/2}=(\tau_S^{-1/2})^{\transpose}$ and $P=Q^{\transpose}$.

\emph{(SR1).} By (b),
$\Psi=(K^{\transpose}\otimes\idon{})\dyad{\Omega}
(K^{\transpose}\otimes\idon{})^{\dagger}$, and
$\tr_S\dyad{\Omega}=\idon{A}$, so
$\tr_S\Psi=K^{\transpose}(K^{\transpose})^{\dagger}=(K^{\transpose})^2
=(K^2)^{\transpose}=\tau_S^{\transpose}=\tau_A$, the operator
$K^{\transpose}$ being Hermitian by (c).

\emph{(SR2).} $\tr_A\big[(M_x\otimes\idon{})\Psi\big]
=K\,\tr_A\big[(M_x\otimes\idon{})\dyad{\Omega}\big]K=K\,M_x^{\transpose}K$
by (a), $K$ acting on $S$ alone. By (c),
$M_x^{\transpose}=\tau_S^{-1/2}\sigma^{\mathrm v}_x\tau_S^{-1/2}
+\delta_{x,x_0}(\idon{S}-Q)$, and $K\tau_S^{-1/2}=Q$, $K(\idon{}-Q)=0$, so
$K\,M_x^{\transpose}K=Q\,\sigma^{\mathrm v}_xQ$. Finally
$0\preceq\sigma^{\mathrm v}_x\preceq\tau_S$ puts the support of
$\sigma^{\mathrm v}_x$ inside that of $\tau_S$, so
$Q\,\sigma^{\mathrm v}_xQ=\sigma^{\mathrm v}_x$.

\emph{The remaining conditions of (i).} $\Psi=\dyad{\psi}\succeq0$ and
$\tr\Psi=\tr\tr_S\Psi=\tr\tau_A=\tr\tau_S$, whose sum over $t$ is
$\sum_x\sum_t\tr\sigma^{\mathrm v,(t)}_x=1$. Each $M_x$ is positive
semidefinite: $(\sigma^{\mathrm v}_x)^{\transpose}\succeq0$ by (c), a
conjugation by the Hermitian operator $\tau_A^{-1/2}$ preserves positivity,
and $\idon{}-P\succeq0$. Their sum is
$\tau_A^{-1/2}\,\tau_S^{\transpose}\,\tau_A^{-1/2}+(\idon{}-P)
=P+(\idon{}-P)=\idon{A}$, by linearity of transposition. $\tau_A$ is
positive semidefinite by (c), and $\sum_t\tr\tau^{(t)}_A=1$ was just shown.
The weighted members $\bigoplus_t\sigma^{\mathrm v,(t)}_x$ of the virtual
ensemble are block diagonal along $(S_t)_t$ by their presentation as block
lists, with the declared blocks as their diagonal blocks, which is the
condition the block form places on the virtual ensemble; together with
\eqref{eq:srblock}, shown above sector by sector, this is (SR1) and (SR2) of
$\Psi=\bigoplus_t\Psi_t$ by Definition~\ref{def:srcert}.

(ii) $\Psi=\dyad{\psi}$ with $\braket{\psi}{\psi}=\tr\Psi=g$; where $g>0$,
$\Psi/g$ is a rank-one projector. Theorem~\ref{thm:certificate-sufficiency}\,(ii) states its
converse for exactly such block lists.

(iii) is the definition of the transpose; (i) and (ii) were proved from the
operator identities above and read no entry of $\Psi$ or of $M_x$. That
these are not declarable in general is seen on $S$ a qubit:
$\tau_S=\mathrm{diag}(\tfrac12,\tfrac12)$ gives $\Psi$ with every non-zero
entry equal to $\tfrac12$, while $\tau_S=\mathrm{diag}(\tfrac34,\tfrac14)$
gives $\Psi$ the entry $\sqrt{3/16}$, which is irrational.

(iv) In the proof of (SR2) the completion term contributed
$K(\idon{}-Q)K=0$, and the same holds of any operator supported on the
kernel of $\tau_S$; in (SR1) and (ii) $M_x$ does not occur; in (v) below it
contributes $\tau_A^{1/2}(\idon{}-P)\tau_A^{1/2}=0$.

(v) $\tau_A^{1/2}\tau_A^{-1/2}=P$, so
$\tau_A^{1/2}M_x\tau_A^{1/2}=P(\sigma^{\mathrm v}_x)^{\transpose}P
=\big(Q\,\sigma^{\mathrm v}_xQ\big)^{\transpose}
=(\sigma^{\mathrm v}_x)^{\transpose}$, using (c) and the support inclusion
above. The sum is $\tau_S^{\transpose}=\tau_A$.
\end{proof}

\section{Sufficiency of the certificate and its scope}\label{group-transmitter/thm:certificate-sufficiency}
\sectionintro{The section's first main result, \hyperlink{claim-transmitter/thm:certificate-sufficiency}{Theorem~\ref*{thm:certificate-sufficiency}}, shows that the certificate's algebraic conditions are, in either declared form, finitely many exact decisions and that Alice's measurement may be moved across every attack that acts only on the \termref{term-transmitted-system}{transmitted system} and the adversary's registers, starting from a register in a product state with the source; retention and the sequential round structure make the attack of a round such a channel. \hyperlink{claim-transmitter/rem:xreg}{Remark~\ref*{rem:xreg}} shows that the hypothesis on the attack cannot be dropped, \hyperlink{claim-transmitter/cor:noncanonical}{Corollary~\ref*{cor:noncanonical}} that BB84 admits a certificate on a qubit, and \hyperlink{claim-transmitter/rem:certificate-forms}{Remark~\ref*{rem:certificate-forms}} what purity adds. \hyperlink{claim-transmitter/lem:matched-pinching}{Lemma~\ref*{lem:matched-pinching}}, the second main result, states what a block certificate implies for the states an attack reaches; the closing remarks bound what a verified certificate buys, down to a one-dimensional retained system.}

% claim transmitter/thm:certificate-sufficiency
\claimsummary{transmitter/thm:certificate-sufficiency}{In any form (pure projector, sector blocks, ensemble presentation) the certificate's algebraic conditions are finitely many exact decisions making Alice's measurement reproduce the virtual ensemble. Under retention and sequential rounds she may measure before or after any attack: source replacement.}
\claimbookmark{transmitter/thm:certificate-sufficiency}{Theorem}{thm:certificate-sufficiency}{Checkability and sufficiency of the source-replacement certificate}
\claimleadin{%
The theorem licenses source replacement at the boundary of this module.
Recall the conditions of Definition~\ref{def:srcert}: (SR1)
$\tr_{\Avirt}\Psi=\tau_A$, (SR2)
$\tr_A[(M_x\otimes\idon{\Avirt})\Psi]=p_x\rvirt_x$, and the attestation (SR3)
that nothing but Alice's measurement acts on $A$; an \termref{term-ensemble-presentation}{ensemble presentation}
has instead the decided conditions (E1) to (E3), positive declared blocks,
normalised weights and per-sector sums that round nowhere. Part (i) is about
declared data before its conditions are known to hold, and says that the
conditions other than the attestation (SR3) can be decided. Part (ii) relates
a certificate to Alice's measurement and is stated for any Hermitian family,
since its converse concerns a family that fails the POVM condition on the
support of $\tau_A$, that is, $\Pi(\sum_xM_x)\Pi\ne\Pi$ for the projector
$\Pi$ onto that support, which no \termref{term-source-specification}{source specification} declares, a
mismatch the \termref{term-pure-form}{pure form} detects (Remark~\ref{rem:certificate-forms}). Part
(iii) shows that her measurement may be moved across any attack that leaves
her retained system and her outcome register alone and starts from an
adversary register in a product state with the source. Its identities hold
for every channel of that form; the attestation (SR3) of
Assumption~\ref{ass:src} and Assumption~\ref{ass:seq} are what make the
attack of a round such a channel, and Remark~\ref{rem:xreg} shows that the
restriction on the attack is necessary and not a weakening of the intended
reading.
}
\begin{mainclaim}
\begin{theorem}[Checkability and sufficiency of the source-replacement certificate]
\label{thm:certificate-sufficiency}
Let a \termref{term-source-specification}{source specification} (Definition~\ref{def:source-spec}) carry a
\termref{term-source-replacement-certificate}{source-replacement certificate} $\Psi$ for $(A,\mathbb{T},\tau_A,\{M_x\})$
against the \termref{term-virtual-ensemble}{virtual ensemble} $\{p_x,\rvirt_x\}$, in any form of
Definition~\ref{def:srcert}. In (ii) and (iii) and their proofs an ensemble
presentation is read as the block-form certificate it defines
(Lemma~\ref{lem:ensemble-presentation}\,(i)), so that ``either form'' below
covers it.
\begin{enumerate}[label=(\roman*),leftmargin=2.2em]
\item \emph{(checkability)} the conditions on $\Psi$ of its declared form
--- in pure form Hermiticity, $\Psi^2=\Psi$ and $\tr\Psi=1$; in \termref{term-block-form}{block form}
Hermiticity and positive semidefiniteness of every block and
$\sum_t\tr\Psi_t=1$ --- together with (SR1) and (SR2), are finitely many
exact identities and exact decisions on the declared data, decided by
finitely many arithmetic operations. In block form (SR1) and (SR2) are decided
as the block identities \eqref{eq:srblock} together with the block diagonality
of every weighted virtual member $\sigma^{\mathrm v}_x$ along $(S_t)_t$,
decided entry by entry, which by Definition~\ref{def:srcert} is their
conjunction; \eqref{eq:srblock} alone decides (SR2) only against the pinched
members $\bigoplus_t\Pi^{\mathrm v}_t\sigma^{\mathrm v}_x\Pi^{\mathrm v}_t$.
For a certificate given by an ensemble
presentation the decided conditions are (E1) to (E3) of
Definition~\ref{def:srcert}, equally finite and exact, and the conditions
of the block form hold of the defined certificate by
Lemma~\ref{lem:ensemble-presentation}\,(i) without being decided. (SR3) is
an attestation and is decided by no function of the declared data. (SR1)
moreover implies $\tau_A\in\Seq{A}$, $\tr_{\Avirt}\Psi$ being positive
semidefinite with trace $\tr\Psi=1$. In an ensemble presentation $\tau_A$ is
not declared: it is formed as the transposed sum of (E3), and that it is a
state is Lemma~\ref{lem:ensemble-presentation}\,(i).
\item \emph{(consistency)} more generally than in a source specification,
let $\{M_x\}$ be any family of Hermitian operators on $A$, presented as block
lists against the \termref{term-sector-structure}{sector structure}, of which $\sum_xM_x=\idon{A}$ is not
assumed; let $\Psi$ meet the conditions on $\Psi$ of its form listed in (i)
and (SR1); and let $p_x\rvirt_x$ stand for
$\tr_A[(M_x\otimes\idon{\Avirt})\Psi]$, so that (SR2) holds. Writing
$U=\operatorname{supp}\tau_A$ with projector $\Pi_U$, the certificate's
signal-side marginal and the virtual average signal state agree,
\begin{equation}\label{eq:srconsist}
  \sum_xp_x\rvirt_x=\tr_A\Psi ,
\end{equation}
\emph{if} $\Pi_U\big(\sum_xM_x\big)\Pi_U=\Pi_U$, in either form; in
particular they agree for the certificate of a source specification, whose
$\{M_x\}$ is a POVM by Definition~\ref{def:source-spec}. The converse,
\emph{only if}, holds in pure form, and in block form whenever every non-zero
block $\Psi_t$ is a positive multiple of a rank-one projector; it fails for a
general positive block.
\item \emph{(sufficiency at the module boundary)} Assume (SR1), (SR2), the
attestation (SR3) of Assumption~\ref{ass:src}, and
Assumption~\ref{ass:seq}; in block form (SR2) is the identity of operators on
$\Avirt$ against the weighted virtual members $p_x\rvirt_x$ themselves, and
not against their diagonal blocks alone: for a declared block form (i) decides
it with their block diagonality, and for an ensemble presentation it is
proved by Lemma~\ref{lem:ensemble-presentation}\,(i). Let $\pcNmsm$ be Alice's measurement map, from
her retained system $A$ to the classical outcome register $X$,
\begin{equation}\label{eq:Nmsm}
  \pcNmsm(\xi)\coloneqq\sum_{x\in\Xset}\tr[M_x\,\xi]\;\dyad{x}_{X},
  \qquad \xi\in\Lin{A} ,
\end{equation}
and let $\mathcal{E}$ be \emph{any} channel from $\Avirt E$ to $E'$, where
$E$ is the adversary's initial register and $E'$ collects her output
registers including whatever is delivered to the receiver. Thus
$\mathcal{E}$ models the adversary's attack and the remainder of the round
\emph{on the \termref{term-virtual-transmitted-system}{transmitted system} and the adversary's registers}, and it
acts on neither $A$ nor $X$: it is not permitted to depend on Alice's
outcome. Write
$\rho_{X\Avirt}\coloneqq\sum_{x\in\Xset}p_x\,\dyad{x}_X\otimes\rvirt_x$ for
the virtual prepare-and-measure source state. Then:
\begin{enumerate}[label=(\alph*),leftmargin=2.2em]
\item \emph{(the interchange)}
\begin{equation}\label{eq:srmain}
  (\idmap_{X}\otimes\mathcal{E})\circ(\pcNmsm\otimes\idmap_{\Avirt E})
  \;=\;\pcNmsm\otimes\mathcal{E}\;=\;
  (\pcNmsm\otimes\idmap_{E'})\circ(\idmap_{A}\otimes\mathcal{E})
\end{equation}
as linear maps on $\Lin{A\Avirt E}$, hence in particular on
$\Psi\otimes\sigma$ for every state $\sigma$ on $E$: Alice's measurement may
be performed before the attack or after it, indifferently. This uses no
property of the certificate, only that the two maps act on complementary
registers;
\item \emph{(the source identification)} for every state $\sigma$ on $E$,
\begin{equation}\label{eq:srident}
  (\pcNmsm\otimes\mathcal{E})\big(\Psi\otimes\sigma\big)
  \;=\;(\idmap_{X}\otimes\mathcal{E})\big(\rho_{X\Avirt}\otimes\sigma\big) ,
\end{equation}
and both sides are states on $XE'$. Together with (a): on every register
other than $A$, the declared entanglement-based description and the
virtual prepare-and-measure description produce the same state, against
every attack of this form. The states $\rvirt_x$ in $\rho_{X\Avirt}$ are the
declared virtual states, in block form as in every form, and not their
pinching along $(S_t)_t$.
\end{enumerate}
\end{enumerate}
\end{theorem}
\end{mainclaim}
\claimmeta{\metalabel{Uses}\mref{Def.}{def:source-spec}, \mref{Def.}{def:srcert}, \mref{Lem.}{lem:ensemble-presentation}
\metasep \metalabel{Assumes}\aref{A-SEQ}{ass:seq}, \aref{A-SRC}{ass:src}
\metasep \metalabel{Used by}\mref{Rem.}{rem:xreg}, \mref{Cor.}{cor:noncanonical}, \mref{Rem.}{rem:boundary}, \mref{Req.}{req:marginal-decided}, \mref{Req.}{req:attestation-record}, \mref{Cor.}{cor:emitted-decoy}, \extref{../meat-core/contract.pdf}{meat-core/lem:product-marginal}, \extref{../protocol-compiler/contract.pdf}{protocol-compiler/def:eligible}, \extref{../protocol-compiler/contract.pdf}{protocol-compiler/thm:gates}, \extref{../reductions/contract.pdf}{reductions/lem:source-transfer}
\metasep newly derived
\metasep implemented (effective: accepted)
\metasep \metalabel{Code}\coderef{https://github.com/fraunhofersg/certiqs-key/blob/v0.1.0/src/certiqs_key/transmitter/certificate.py}{certiqs_key.transmitter.certificate:verify_certificate}
\metasep \metalabel{Oracles}\coderef{https://github.com/fraunhofersg/certiqs-key/blob/v0.1.0/tests/oracles/registry_instances_bb84.py}{registry_instances_bb84.py}, \coderef{https://github.com/fraunhofersg/certiqs-key/blob/v0.1.0/tests/oracles/registry_protocol_compiler.py}{registry_protocol_compiler.py}, \coderef{https://github.com/fraunhofersg/certiqs-key/blob/v0.1.0/tests/oracles/registry_transmitter.py}{registry_transmitter.py}}

\begin{proof}
\emph{Idea.} Part (i) counts the operations each condition needs and
observes that they are formed exactly. Part (ii) reduces the consistency
identity to the vanishing of the compression $\Pi_UY\Pi_U$ of the defect
$Y=\sum_xM_x-\idon{A}$, using that $\Psi$ lives on the support of $\tau_A$.
Part (iii) is the interchange law for a tensor product of maps on
complementary registers, followed by (SR2).

(i) In pure form the conditions on $\Psi$ compare $(d_A\dim\Hs{\Avirt})^2$
entries each --- $\Psi=\Psi^{\dagger}$ and $\Psi^2=\Psi$ --- and one
accumulated trace; in block form they are, per block, one comparison of
$(d_{A_t}\dim\Hs{S_t})^2$ entries for Hermiticity and one exact decision of
positive semidefiniteness --- a symmetric elimination over the Gaussian
rationals in which a negative pivot refuses the block and a zero pivot
requires its whole row and column to vanish, the block being refused
otherwise, which decides that every principal minor is non-negative --- and
one accumulated trace over all blocks; in an ensemble presentation the same
decision per declared block, the traces of (E2) and the entrywise sums of
(E3); in pure form (SR1) compares $d_A^2$ entries after a partial trace and
(SR2) compares $|\Xset|$ matrices of size $\dim\Hs{\Avirt}$, each obtained
from $\Psi$ by a multiplication and a partial trace; in block form
\eqref{eq:srblock} compares, per sector, $d_{A_t}^2$ entries and $|\Xset|$
matrices of size $\dim\Hs{S_t}$, each obtained from $\Psi_t$ in the same way,
and the block diagonality compares with zero the entries of every
$\sigma^{\mathrm v}_x$ outside its diagonal blocks, at most
$|\Xset|(\dim\Hs{\Avirt})^2$ of them. By Definition~\ref{def:srcert} the block
identities and the block diagonality together are (SR1) and (SR2) of $\Psi$:
the partial traces of $\Psi$ are the direct sums of those of its blocks, and
the one of (SR2) is block diagonal along $(S_t)_t$, so it equals
$\sigma^{\mathrm v}_x$ exactly when $\sigma^{\mathrm v}_x$ has no entry
outside its diagonal blocks and those blocks satisfy the second identity of
\eqref{eq:srblock}. Each is a finite comparison, and each
is \emph{formed as well as taken} in exact arithmetic on the declared
dyadic entries (the number model of Definition~\ref{def:source-spec}): the
products entering $\Psi^2$ and (SR2), and the additions entering the
partial traces and the accumulated trace, are exact, so what is compared is
the quantity written and not a rounded image of it. That the operands are
exactly represented does not on its own deliver that; the exact formation
does. For the implication: $\tr_{\Avirt}\Psi$ is positive semidefinite as
the partial trace of a positive semidefinite operator --- in pure form
$\Psi$ is a projector, hence $\Psi=\Psi^{\dagger}\Psi\succeq0$; in block
form $\Psi$ is a direct sum of positive semidefinite blocks --- and its
trace is $\tr\Psi=1$; (SR1) is a matrix identity, so $\tau_A\in\Seq{A}$.
(SR3) is a statement about the apparatus, not about matrices, and no matrix
function decides it; it is the attestation of Assumption~\ref{ass:src}.

(ii) Put $Y\coloneqq\sum_xM_x-\idon{A}$, which is Hermitian, each $M_x$
being Hermitian by hypothesis and a sum of Hermitian matrices less
$\idon{A}$ being Hermitian; $Y$ is block diagonal along the sector structure,
every $M_x$ being presented as a block list. Summing (SR2) over $x$ gives
\begin{equation*}
  \tr_A\!\Big[\Big(\big(\textstyle\sum_xM_x\big)\otimes\idon{}\Big)\Psi\Big]
  =\sum_xp_x\rvirt_x ,
\end{equation*}
and the left-hand side is $\tr_A[(Y\otimes\idon{})\Psi]+\tr_A\Psi$ by
linearity of the partial trace, so \eqref{eq:srconsist} holds if and only
if $\tr_A[(Y\otimes\idon{})\Psi]=0$.

\emph{If, in either form.} $\Psi\succeq0$ in either form. Let $\ket{\phi}$
be a unit vector of $\Hs{A}$ orthogonal to $U$ and put
$P\coloneqq\dyad{\phi}\otimes\idon{\Avirt}$. By (SR1),
$\tr[P\Psi]=\bra{\phi}\tau_A\ket{\phi}=0$, so
$\Psi^{1/2}P\Psi^{1/2}$ is a positive semidefinite operator of trace zero,
hence zero, hence $P\Psi^{1/2}=0$ and $P\Psi=0$. Summing over an orthonormal
basis of the orthogonal complement of $U$ gives
$\Psi=(\Pi_U\otimes\idon{})\Psi(\Pi_U\otimes\idon{})$, and then
\begin{equation*}
  \tr_A\!\big[(Y\otimes\idon{})\,\Psi\big]
  =\tr_A\!\big[(Y\Pi_U\otimes\idon{})\,\Psi\,(\Pi_U\otimes\idon{})\big]
  =\tr_A\!\big[(\Pi_UY\Pi_U\otimes\idon{})\,\Psi\big] ,
\end{equation*}
the partial trace over $A$ being cyclic under operators that act on $A$
alone. This vanishes when $\Pi_UY\Pi_U=0$, i.e.\ when
$\Pi_U(\sum_xM_x)\Pi_U=\Pi_U$. If $\{M_x\}$ is a POVM then $Y=0$ and the
condition holds outright, and Definition~\ref{def:source-spec} delivers
exactly that: $\sum_xM^{(t)}_x=\idon{A_t}$ for every $t$ as an identity of
the declared numbers, its declared-order sum decided exactly.

\emph{Only if, in pure form.} $\Psi=\dyad{\psi}$ for a unit vector
$\ket{\psi}$; take a Schmidt decomposition
$\ket{\psi}=\sum_i\sqrt{\mu_i}\ket{e_i}_A\ket{f_i}_{\Avirt}$ with every
$\mu_i>0$ and $\{\ket{e_i}\}$, $\{\ket{f_i}\}$ orthonormal. By (SR1),
$\tau_A=\tr_{\Avirt}\Psi=\sum_i\mu_i\dyad{e_i}$, which is a spectral
decomposition, so $\{\ket{e_i}\}$ is an orthonormal basis of
$U=\operatorname{supp}\tau_A$. Then
\begin{equation*}
  \tr_A\!\big[(Y\otimes\idon{})\,\Psi\big]
  =\sum_{i,j}\sqrt{\mu_i\mu_j}\,\bra{e_j}Y\ket{e_i}\,\ketbra{f_i}{f_j},
\end{equation*}
and since the $\ket{f_i}$ are orthonormal and every $\mu_i>0$, this
operator vanishes if and only if $\bra{e_j}Y\ket{e_i}=0$ for all $i,j$,
i.e.\ if and only if $\Pi_UY\Pi_U=0$.

\emph{Only if, in block form with rank-one blocks.} Let
$\Psi_t=w_t\dyad{\psi_t}$ with $w_t>0$ and $\ket{\psi_t}$ a unit vector of
$\Hs{A_tS_t}$ for every non-zero block. $Y=\bigoplus_tY_t$ is block diagonal
because every $M_x$ is, so
$\tr_A[(Y\otimes\idon{})\Psi]=\bigoplus_tw_t\,
\tr_{A_t}[(Y_t\otimes\idon{S_t})\dyad{\psi_t}]$, a direct sum along
$(S_t)_t$, which vanishes if and only if every summand does. The pure-form
argument, applied to $\ket{\psi_t}$ with $\tau^{(t)}_A/w_t$ in the place of
$\tau_A$, shows that the $t$-th summand vanishes if and only if
$\Pi_{U,t}Y_t\Pi_{U,t}=0$, where $\Pi_{U,t}$ projects onto
$\operatorname{supp}\tau^{(t)}_A$; and
$\Pi_U=\bigoplus_t\Pi_{U,t}$, so the conjunction over $t$ is
$\Pi_UY\Pi_U=0$.

\emph{The converse fails for a general positive block.} Take one sector,
$\Psi=\tau_A\otimes\varsigma$ for a state $\varsigma$ on $\Avirt$, and any
Hermitian $Y$ with $\tr[Y\tau_A]=0$ and $\Pi_UY\Pi_U\ne0$; such a $Y$
exists whenever $\dim U\ge2$, and scaled so that $\idon{A}+Y\succeq0$ it is
realised by the one-letter family $M=\idon{A}+Y$, Hermitian and positive
semidefinite but not a POVM, which the hypothesis of (ii) admits, with
$p\rvirt=\tr_A[(M\otimes\idon{})\Psi]$ as in the statement. The block
$\Psi$ is positive semidefinite of unit trace and satisfies (SR1). Then
$\tr_A[(Y\otimes\idon{})\Psi]=\tr[Y\tau_A]\,\varsigma=0$.

(iii)(a) \emph{The interchange.} The registers do the work, and no
numerical hypothesis enters. The input operator lives on $A\Avirt E$. By
hypothesis $\pcNmsm$ is a linear map out of $A$ alone and $\mathcal{E}$ a
channel out of $\Avirt E$ alone, and $A$ and $\Avirt E$ are complementary
tensor factors of the input; the outputs $X$ and $E'$ are likewise
distinct registers. Hence both composites equal the tensor product
$\pcNmsm\otimes\mathcal{E}$ of the two maps, which is \eqref{eq:srmain};
this is the elementary interchange law for a tensor product of linear maps
and needs nothing about the particular $\pcNmsm$, $\mathcal{E}$ or input
operator. That $\mathcal{E}$ does not act on $A$ is (SR3); that it does not
act on $X$ is a hypothesis of the statement, and Remark~\ref{rem:xreg}
shows it cannot be dropped.

(iii)(b) \emph{The source identification.} Applying
$\pcNmsm\otimes\idmap_{\Avirt E}$ to $\Psi\otimes\sigma$ gives, by
\eqref{eq:Nmsm} and (SR2),
\begin{equation*}
  \sum_x\dyad{x}_X\otimes
  \tr_A\!\big[(M_x\otimes\idon{\Avirt})\,\Psi\big]\otimes\sigma
  =\sum_xp_x\,\dyad{x}_X\otimes\rvirt_x\otimes\sigma
  =\rho_{X\Avirt}\otimes\sigma ,
\end{equation*}
the right-hand side being exactly the classical-quantum source state from
which Fact~\ref{fact:sr} starts, under the symbol map stated with that
fact. Applying the linear map $\idmap_X\otimes\mathcal{E}$ to both
sides gives \eqref{eq:srident}, the left-hand side being
$(\pcNmsm\otimes\mathcal{E})(\Psi\otimes\sigma)$ by (a). That both sides
are states: $\{M_x\}$ is a POVM by Definition~\ref{def:source-spec}, so
$\pcNmsm$ is a channel; $\Psi$ is a state, being positive semidefinite of
unit trace in either form; and $\mathcal{E}$ is a channel by hypothesis.
Nothing in (iii) reads $\Psi^2=\Psi$.
\end{proof}

% claim transmitter/rem:xreg
\claimsummary{transmitter/rem:xreg}{Without the hypothesis that the attack leaves Alice's outcome register alone, part (iii) of \hyperlink{claim-transmitter/thm:certificate-sufficiency}{Theorem~\ref*{thm:certificate-sufficiency}} is ill-typed and false. The restriction costs nothing, because A-SEQ places the announcements after both parties' measurements.}
\claimbookmark{transmitter/rem:xreg}{Remark}{rem:xreg}{Necessity and cost of the outcome-register hypothesis}
\begin{remark}[Necessity and cost of the outcome-register hypothesis]\label{rem:xreg}
The hypothesis of Theorem~\ref{thm:certificate-sufficiency}\,(iii) that $\mathcal{E}$ acts
trivially on the outcome register cannot be dropped: without it the
statement is both ill-typed and false. \emph{Ill-typed}, because $X$ does
not exist until $\pcNmsm$ creates it: in the right-hand composite of
\eqref{eq:srmain} the channel $\mathcal{E}$ acts first, on an operator on
$A\Avirt E$ carrying no such register. \emph{False}, because a channel
that did act on $X$ could destroy exactly the information the equality is
about: let $\mathcal{E}_{\mathrm{ow}}$ overwrite $X$ with a fixed $\dyad{x_0}$ and act
as the identity elsewhere; then $\mathcal{E}_{\mathrm{ow}}\circ(\pcNmsm\otimes\idmap)$
produces the point distribution at $x_0$, while any composition in which
$\mathcal{E}_{\mathrm{ow}}$ precedes $\pcNmsm$ leaves the outcome distribution
$(p_x)_x$ intact. One overwriting channel is enough to refute a
commutation asserted for \emph{every} channel.

What the restriction excludes is exactly the operations that read Alice's
outcome --- in a QKD round, the public announcements. Nothing is lost,
because Assumption~\ref{ass:seq} places those \emph{after} both parties'
measurements: the round is preparation and transmission, then measurement
by both parties, then announcements. So the channel that models ``the
adversary's attack and the rest of the round up to the announcements'' is
of the admitted form, and the delayed-measurement reading, in which Alice
measures after the attack, survives in the form that is used --- Alice's
measurement may be moved across the attack, not across the announcements. The same restriction, in the same place, is
what the published framework this module imports from relies on: its
security reduction postpones Alice's measurement on the ground that her
measurements commute with all operations except the public announcements,
which by design occur only after the receiver's measurement in each round
\cite[Thm.~8.3]{tupkary2026decoy}. That corroboration is recorded, not
used: no step of the proof of Theorem~\ref{thm:certificate-sufficiency} appeals to it.
\end{remark}
\claimmeta{\metalabel{Uses}\mref{Thm.}{thm:certificate-sufficiency}
\metasep \metalabel{Assumes}\aref{A-SEQ}{ass:seq}
\metasep newly derived
\metasep accepted
\metasep \metalabel{Oracles}\coderef{https://github.com/fraunhofersg/certiqs-key/blob/v0.1.0/tests/oracles/registry_instances_bb84.py}{registry_instances_bb84.py}, \coderef{https://github.com/fraunhofersg/certiqs-key/blob/v0.1.0/tests/oracles/registry_protocol_compiler.py}{registry_protocol_compiler.py}, \coderef{https://github.com/fraunhofersg/certiqs-key/blob/v0.1.0/tests/oracles/registry_transmitter.py}{registry_transmitter.py}}

% claim transmitter/cor:noncanonical
\claimsummary{transmitter/cor:noncanonical}{The canonical source replacement satisfies the pure-form certificate conditions, but it is not the only data that do. Four-state BB84 admits a pure certificate with a two-dimensional retained system, where the canonical form needs four dimensions.}
\claimbookmark{transmitter/cor:noncanonical}{Corollary}{cor:noncanonical}{Certificates beyond the canonical form}
\claimleadin{%
The imported construction fixes Alice's label register $\Abar$ at dimension
$|\Xset|$. The corollary shows that the certificate is strictly more
permissive already in its pure form, the form in which both the canonical data
and the witness below are declared.
}
\begin{corollary}[Certificates beyond the canonical form]
\label{cor:noncanonical}
Throughout, the certificate is in the \emph{pure} form of
Definition~\ref{def:srcert}, against the \termref{term-one-sector-structure}{one-sector structure} on $A$ that
Definition~\ref{def:source-spec} makes the default.
\begin{enumerate}[label=(\roman*),leftmargin=2.2em]
\item Fact~\ref{fact:sr}'s canonical data --- $A=\Abar\Ahat$ with
$\dim\Hs{\Abar}=|\Xset|$, $M_x=\dyad{x}_{\Abar}\otimes\idon{\Ahat}$, and
$\Psi=\dyad{\psi}$ with
$\ket{\psi}=\sum_x\sqrt{p_x}\ket{x}_{\Abar}\ket{\psi_x}_{\Ahat\Avirt}$ and
$\ket{\psi_x}$ a purification of $\rvirt_x$ --- satisfy (SR1) and (SR2).
\item They are not the only such data: Theorem~\ref{thm:certificate-sufficiency}\,(iii)
applies to every triple $(\Psi,\tau_A,\{M_x\})$ satisfying (SR1)--(SR3), and
there are source
specifications with $d_A<|\Xset|$ carrying a certificate. In particular the
four-state BB84 ensemble on a qubit, $\Xset=\{(b,k):b,k\in\{0,1\}\}$ with
$p_{(b,k)}=\tfrac14$ and $\rvirt_{(b,k)}=\dyad{k_b}$ for the computational
basis ($b=0$) and the Hadamard basis ($b=1$), admits the certificate
$A=\Avirt=\mathbb{C}^2$, $M_{(b,k)}=\tfrac12\dyad{k_b}$,
$\tau_A=\tfrac12\idon{A}$, $\Psi=\dyad{\Phip}$ with
$\ket{\Phip}=\tfrac1{\sqrt2}(\ket{00}+\ket{11})$: $d_A=2$ where the
canonical form has $d_A=4$.
\end{enumerate}
\end{corollary}
\claimmeta{\metalabel{Uses}\mref{Def.}{def:source-spec}, \mref{Fact}{fact:sr}, \mref{Def.}{def:srcert}, \mref{Thm.}{thm:certificate-sufficiency}
\metasep newly derived
\metasep accepted
\metasep \metalabel{Oracles}\coderef{https://github.com/fraunhofersg/certiqs-key/blob/v0.1.0/tests/oracles/registry_instances_bb84.py}{registry_instances_bb84.py}, \coderef{https://github.com/fraunhofersg/certiqs-key/blob/v0.1.0/tests/oracles/registry_protocol_compiler.py}{registry_protocol_compiler.py}, \coderef{https://github.com/fraunhofersg/certiqs-key/blob/v0.1.0/tests/oracles/registry_transmitter.py}{registry_transmitter.py}}

\begin{proof}
(i) \emph{The canonical data.} $\Psi$ is a rank-one projector of unit
trace, $\ket{\psi}$ being a unit vector --- the $\ket{\psi_x}$ being unit
vectors and the $p_x$ summing to one. (SR1) is the definition of $\tau_A$
for the canonical data. For (SR2),
$(M_x\otimes\idon{\Avirt})\ket{\psi}
=\sqrt{p_x}\ket{x}_{\Abar}\ket{\psi_x}_{\Ahat\Avirt}$, since
$\dyad{x}\ket{x'}=\delta_{xx'}\ket{x}$ and $\idon{\Ahat}$ acts trivially;
$M_x$ is a projector, so
$\tr_A[(M_x\otimes\idon{})\,\Psi]
=\tr_A[(M_x\otimes\idon{})\,\Psi\,(M_x\otimes\idon{})]
=p_x\tr_{\Ahat}\dyad{\psi_x}=p_x\rvirt_x$, the last step because
$\ket{\psi_x}$ purifies $\rvirt_x$.

(ii) \emph{The qubit witness.} $\{M_{(b,k)}\}$ is a POVM on $\mathbb{C}^2$:
each element is positive semidefinite and
$\sum_{b,k}\tfrac12\dyad{k_b}=\tfrac12\idon{}+\tfrac12\idon{}=\idon{}$,
each basis resolving the identity; so the data are a source specification
(Definition~\ref{def:source-spec}). $\dyad{\Phip}$ is a rank-one projector
of unit trace with every entry in $\{0,\tfrac12\}$. (SR1):
$\tr_{\Avirt}\dyad{\Phip}=\tfrac12\idon{A}=\tau_A$. (SR2): both bases are
real, so $(\bra{k_b}\otimes\idon{})\ket{\Phip}=\tfrac1{\sqrt2}\ket{k_b}$,
whence $\tr_A[(M_{(b,k)}\otimes\idon{})\dyad{\Phip}]
=\tfrac12\cdot\tfrac12\dyad{k_b}=p_{(b,k)}\rvirt_{(b,k)}$. Four vectors
in a two-dimensional space are not linearly independent, so the canonical
form has $\dim\Hs{\Abar}=|\Xset|=4$ with a $\tau_{\Abar}$ of rank two; the
witness has $d_A=2$.
\end{proof}

% claim transmitter/rem:certificate-forms
\claimsummary{transmitter/rem:certificate-forms}{Only the converse in part (ii) of \hyperlink{claim-transmitter/thm:certificate-sufficiency}{Theorem~\ref*{thm:certificate-sufficiency}} depends on purity. A certificate vector, a globally pure operator over sectors and a pure block over non-orthogonal states can all need irrational entries, which is why the certificate is an operator with a block form and an ensemble form.}
\claimbookmark{transmitter/rem:certificate-forms}{Remark}{rem:certificate-forms}{Purity and exact declarability of the certificate forms}
\begin{remark}[Purity and exact declarability of the certificate forms]\label{rem:certificate-forms}
\emph{What purity adds.} A pure-form certificate satisfies every condition of
the block form at the one-sector structure and one condition more
(Definition~\ref{def:srcert}). Every statement this module proves of the
block form therefore holds of the pure form, and
Theorem~\ref{thm:certificate-sufficiency}\,(ii) is the only statement whose \emph{conclusion}
depends on the additional condition: in pure form a source/POVM mismatch of
the kind (ii) describes is a detected defect and not a silent one, while in
block form it is the POVM condition of Definition~\ref{def:source-spec},
decided on its own, that excludes it. Part (i) of the theorem lists purity as
a condition of the pure form, and a consumer that demands purity of every
certificate re-decides it and therefore refuses a block form. Everything
Corollary~\ref{cor:noncanonical} exhibits is admitted in block form as well.
Not every block-form certificate is pure: already at one sector the block
form admits a positive semidefinite operator of unit trace and rank above
one, which the pure form refuses.

\emph{Why the certificate is an operator and not a vector.} No binary64
number $a$ satisfies $a^2=\tfrac12$, so a declared vector with entries
$\mathrm{fl}(1/\sqrt2)$ yields $\dyad{\psi}$ with entries
$\mathrm{fl}(1/\sqrt2)^2\ne\tfrac12$, and (SR1) fails exactly \emph{for the
certificate as an instance declares it} --- $\ket{\Phip}$ in that
instance's own computational basis, whose amplitudes are $1/\sqrt2$. The
\emph{operator} $\dyad{\Phip}$, by contrast, has every entry equal to
$\tfrac12$, which is a binary64 number; and every quantity (SR1), (SR2)
and Theorem~\ref{thm:certificate-sufficiency} read off the certificate is a function of
$\Psi$ and never of a vector. In pure form the purity conditions ---
$\Psi^2=\Psi$ with $\tr\Psi=1$ --- are what recover the vector where an
argument needs one (Theorem~\ref{thm:certificate-sufficiency}\,(ii)), and they are two exact
identities on the declared matrix. Whether a given declaration of the
canonical data of Corollary~\ref{cor:noncanonical} is admissible is decided,
not assumed: the canonical $\Psi$ must be declared as a binary64 rank-one
projector, and a declaration whose entries do not satisfy $\Psi^2=\Psi$
exactly is refused.

\emph{Why a block form and an ensemble presentation.} Exact purity is not
declarable for a source that draws a sector and then prepares a state in it.
A globally pure $\Psi$ over sectors of weights $g_t$ carries cross-sector
amplitudes $\sqrt{g_tg_{t'}}$, which need not be dyadic rationals, and a
pure block whose marginal has the spectrum $(\tfrac12,\tfrac14,\tfrac14)$
carries the entry $\sqrt{1/8}$, which is not one. The direct sum of positive
blocks is exactly declarable, it is the state of the source that first draws
$t$ with probability $\tr\Psi_t$ and then prepares $\Psi_t/\tr\Psi_t$, and
the sufficiency statement, Theorem~\ref{thm:certificate-sufficiency}\,(iii), reads purity
nowhere. Where even the blocks are not declarable --- a pure block over
non-orthogonal signal states carries square roots of their weights --- the
ensemble presentation is the form to use: it has the purity
(Lemma~\ref{lem:ensemble-presentation}\,(ii)) and declares none of the
numbers, which is what it exists for.
\end{remark}
\claimmeta{\metalabel{Uses}\mref{Def.}{def:srcert}
\metasep newly derived
\metasep accepted}

% claim transmitter/lem:matched-pinching
\claimsummary{transmitter/lem:matched-pinching}{A block certificate is unchanged by pinching Alice's sectors or the virtual signal's sectors paired with them, so every state an admitted attack reaches from it is block diagonal with Alice's declared marginal. An operator merely constrained to have that marginal need not be block diagonal.}
\claimbookmark{transmitter/lem:matched-pinching}{Lemma}{lem:matched-pinching}{Matched pinching of a block certificate}
\claimleadin{%
A block certificate declares its sectors on Alice's side and on the virtual
signal at once. The lemma shows that this matching survives every admitted
attack, for the states the attack actually reaches; it is about those states
only, and part (iii) shows why it does not restrict any optimisation's
feasible set.
}
\begin{mainclaim}
\begin{lemma}[Matched pinching of a block certificate]\label{lem:matched-pinching}
Let a source specification carry a certificate in block form
(Definition~\ref{def:srcert}), with sector projectors $\Pi^{A}_t$ on $A$ and
$\Pi^{\mathrm v}_t$ on $\Avirt$, $t\in\mathbb{T}$, which sum to the identity
on $A$ and on $\Avirt$ respectively (Definitions~\ref{def:source-spec}
and~\ref{def:srcert}), and write
\begin{equation*}
  \mathrm{Pin}_{A}(\xi)\coloneqq\sum_{t\in\mathbb{T}}\Pi^{A}_t\,\xi\,\Pi^{A}_t ,
  \qquad
  \mathrm{Pin}_{\mathrm v}(\zeta)\coloneqq
  \sum_{t\in\mathbb{T}}\Pi^{\mathrm v}_t\,\zeta\,\Pi^{\mathrm v}_t
\end{equation*}
for the two pinching channels.
\begin{enumerate}[label=(\roman*),leftmargin=2.2em]
\item \emph{(matched structure)} for every $t\in\mathbb{T}$,
\begin{equation}\label{eq:matched}
  (\Pi^{A}_t\otimes\idon{\Avirt})\,\Psi
  =\Psi\,(\Pi^{A}_t\otimes\idon{\Avirt})
  =(\idon{A}\otimes\Pi^{\mathrm v}_t)\,\Psi
  =\Psi\,(\idon{A}\otimes\Pi^{\mathrm v}_t)
  =\Psi_t ,
\end{equation}
the block $\Psi_t$ read as an operator on $\Hs{A\Avirt}$ extended by zero.
Hence
$(\mathrm{Pin}_{A}\otimes\idmap_{\Avirt})(\Psi)=\Psi
=(\idmap_{A}\otimes\mathrm{Pin}_{\mathrm v})(\Psi)$.
\item \emph{(reached states)} assume (SR3), and let $\mathcal{E}$ be any
channel from $\Avirt E$ to $E'$ and $\sigma$ any state on $E$, as in
Theorem~\ref{thm:certificate-sufficiency}\,(iii). The state
$\omega\coloneqq(\idmap_{A}\otimes\mathcal{E})(\Psi\otimes\sigma)$ on $AE'$
satisfies
\begin{equation}\label{eq:reached}
  \tr_{E'}\omega=\tau_A ,\qquad
  (\mathrm{Pin}_{A}\otimes\idmap_{E'})(\omega)=\omega ,\qquad
  \omega=(\idmap_{A}\otimes\mathcal{E}')(\Psi\otimes\sigma)
  \ \text{ with }\
  \mathcal{E}'\coloneqq\mathcal{E}\circ(\mathrm{Pin}_{\mathrm v}\otimes\idmap_{E}) .
\end{equation}
So $\omega=\bigoplus_t\omega_t$ with $\omega_t\in\Pos{A_tE'}$ and
$\tr_{E'}\omega_t=\tau^{(t)}_A$; and the attack $\mathcal{E}$ and the attack
$\mathcal{E}'$ that first measures the sector of the virtual signal reach
the same state.
\item \emph{(operators that are not reached)} the block diagonality of (ii)
does not extend to operators merely constrained to have the marginal: in
general a positive operator on $AE'$ with marginal
$\tau_A=\bigoplus_t\tau^{(t)}_A$ need not be block diagonal along
$\mathbb{T}$, and this happens already for two one-dimensional sectors,
$\tau_A=\tfrac12\idon{A}$ and $E'$ a qubit.
\end{enumerate}
\end{lemma}
\end{mainclaim}
\claimmeta{\metalabel{Uses}\mref{Def.}{def:source-spec}, \mref{Def.}{def:srcert}
\metasep \metalabel{Assumes}\aref{A-SRC}{ass:src}
\metasep \metalabel{Used by}\mref{Rem.}{rem:boundary}, \mref{Cor.}{cor:emitted-decoy}
\metasep newly derived
\metasep accepted}

\begin{proof}
(i) By Definition~\ref{def:srcert}, $\Psi=\bigoplus_t\Psi_t$ on the \termref{term-matched-subspace}{matched
subspace} $\bigoplus_t\Hs{A_t}\otimes\Hs{S_t}$ and zero on its orthogonal
complement. The projector $\Pi^{A}_t\otimes\idon{\Avirt}$ has range
$\bigoplus_{t'}\Hs{A_t}\otimes\Hs{S_{t'}}$, whose intersection with the
matched subspace is $\Hs{A_t}\otimes\Hs{S_t}$, and it commutes with the
projector onto the matched subspace, both being diagonal in the product of
the declared bases. So $(\Pi^{A}_t\otimes\idon{})\Psi$ is the part of
$\Psi$ with range in $\Hs{A_t}\otimes\Hs{S_t}$, which is $\Psi_t$; taking
adjoints, $\Psi$ and $\Psi_t$ being Hermitian, gives the second equality,
and the same argument with $\idon{A}\otimes\Pi^{\mathrm v}_t$, whose range
meets the matched subspace in the same summand, gives the third and fourth.
Summing $(\Pi^{A}_t\otimes\idon{})\Psi(\Pi^{A}_t\otimes\idon{})
=\Psi_t(\Pi^{A}_t\otimes\idon{})=\Psi_t$ over $t$ gives
$(\mathrm{Pin}_{A}\otimes\idmap)(\Psi)=\Psi$, and likewise on the other
factor.

(ii) $\mathcal{E}$ is trace preserving and does not act on $A$, the latter
being (SR3); so $\tr_{E'}\omega=\tr_{\Avirt E}[\Psi\otimes\sigma]
=\tr_{\Avirt}\Psi=\tau_A$ by (SR1). The maps $\mathrm{Pin}_{A}$ and
$\mathcal{E}$ act on complementary tensor factors, so they commute, and by
(i)
\begin{equation*}
  (\mathrm{Pin}_{A}\otimes\idmap_{E'})(\omega)
  =(\idmap_{A}\otimes\mathcal{E})
   \big((\mathrm{Pin}_{A}\otimes\idmap_{\Avirt})(\Psi)\otimes\sigma\big)
  =(\idmap_{A}\otimes\mathcal{E})(\Psi\otimes\sigma)=\omega .
\end{equation*}
Replacing $(\mathrm{Pin}_{A}\otimes\idmap)(\Psi)$ by
$(\idmap\otimes\mathrm{Pin}_{\mathrm v})(\Psi)$, equal to it by (i), in the
middle expression gives
$\omega=(\idmap_{A}\otimes\mathcal{E}')(\Psi\otimes\sigma)$. An operator
fixed by $\mathrm{Pin}_{A}\otimes\idmap_{E'}$ is the direct sum of its
compressions $\omega_t=(\Pi^{A}_t\otimes\idon{})\omega(\Pi^{A}_t\otimes\idon{})$,
each positive semidefinite as a compression of one, and
$\tr_{E'}\omega_t=\Pi^{A}_t\tau_A\Pi^{A}_t=\tau^{(t)}_A$.

(iii) With two one-dimensional sectors spanned by $\ket{0}$ and $\ket{1}$,
$\tau_A=\tfrac12\idon{A}$ and $E'$ a qubit, the projector $\omega$ onto
$(\ket{0}\ket{0}+\ket{1}\ket{1})/\sqrt2$ is positive semidefinite with
marginal $\tau_A$ on $A$, and its compression
$(\Pi^{A}_0\otimes\idon{})\,\omega\,(\Pi^{A}_1\otimes\idon{})
=\tfrac12\ketbra{0}{1}\otimes\ketbra{0}{1}$ is not zero.
\end{proof}

% claim transmitter/rem:boundary
\claimsummary{transmitter/rem:boundary}{A verified certificate buys the interchange and the source identification of \hyperlink{claim-transmitter/thm:certificate-sufficiency}{Theorem~\ref*{thm:certificate-sufficiency}} outright, and in block form the reached-state statement of \hyperlink{claim-transmitter/lem:matched-pinching}{Lemma~\ref*{lem:matched-pinching}}. It is no security reduction, restricts no consumer's feasible set and decides nothing about retention.}
\claimbookmark{transmitter/rem:boundary}{Remark}{rem:boundary}{Scope of a verified certificate}
\begin{remark}[Scope of a verified certificate]
\label{rem:boundary}
Theorem~\ref{thm:certificate-sufficiency}\,(iii) is a statement about the \emph{source}: it
says that Alice's measurement commutes with every attack of the admitted
form --- every channel on the \termref{term-virtual-transmitted-system}{transmitted system} and the adversary's
registers, which by Assumption~\ref{ass:seq} is every operation of the
round up to the announcements (Remark~\ref{rem:xreg}) --- and that the
declared trusted-source data describe the virtual preparation (the
specified one exactly when the \termref{term-source-segment}{source segment} is empty) against every such
attack. It is not a security reduction, and it does not say that the
consuming engine's optimization is sound; that argument --- granting the
adversary a purification of Alice's retained system and the
\termref{term-virtual-transmitted-system}{transmitted system}, restricting to the declared feasible set, accumulating over rounds
--- lives in the engine's claims, which import Assumption~\ref{ass:src} for
it. What this module owes depends on the form of the certificate
(Theorem~\ref{thm:certificate-sufficiency}\,(i)). In pure and in block form it is the
decided conditions of the declared certificate, those on $\Psi$ of its form
together with (SR1) and (SR2), and the attestation (SR3). The conditions on $\Psi$
are Hermiticity, $\Psi^2=\Psi$ and $\tr\Psi=1$ in pure form, and Hermiticity
and positive semidefiniteness of every block and $\sum_t\tr\Psi_t=1$ in block
form. For an ensemble presentation it is exactly (E1) to (E3), decided on
the declared blocks, and the attestation (SR3); (SR1) and (SR2) then hold of
the certificate the presentation defines as theorems and are decided by
nothing.

\emph{What a verified certificate buys.} A verified certificate, one whose
decided conditions (Theorem~\ref{thm:certificate-sufficiency}\,(i)) hold, buys
\eqref{eq:srmain} and \eqref{eq:srident} outright: the commutation, and the
identification of the declared trusted-source data with the virtual source,
as identities for every channel of the admitted form, with no residue and no
qualification. That is what the exact reading of (SR1)--(SR2) is for, and it
is what a consumer of the licence record, the record of the verified
certificate that the artefact carries (Requirement~\ref{req:attestation-record}),
is entitled to read the record as asserting. What the certificate does
\emph{not} buy is (SR3), which no function of the declared data decides
(Assumption~\ref{ass:src}); it is (SR3), with Assumption~\ref{ass:seq}, that
makes the attack of a round a channel of the admitted form.

\emph{What the block form adds.} A certificate declared in block form
(Definition~\ref{def:srcert}) buys one thing more, and it is a statement about
\emph{reached} states: Lemma~\ref{lem:matched-pinching}, that every state an
admitted attack reaches from that certificate is block diagonal along the
declared sectors with $A$-marginal $\tau_A$, and that pinching $A$ after the
attack is the same operation as pinching the virtual signal before it. That is
what a consumer may read the record as asserting of the declared \termref{term-sector-structure}{sector
structure}, and it is all of it. \emph{It does not buy a restriction of any
consumer's feasible set.} A feasible set is a set of operators constrained to
have the marginal $\tau_A$, not a set of states reached from a certificate,
and a positive operator with that marginal need not be block diagonal
(Lemma~\ref{lem:matched-pinching}\,(iii)). A consumer that minimises over
block-diagonal operators alone does so under a theorem of its own module, with
its own preconditions on declared data, and neither this module's licence
record nor the presentation of $\tau_A$ and $\{M_x\}$ as block lists is that
theorem (Definition~\ref{def:source-spec}, ``What the sector structure is
not'').

\emph{Links supplied by a theorem.} The same boundary, in the other
direction, governs a \termref{term-source-link}{source link} whose channel is supplied by a named claim
of a submodule rather than as data (a \termref{term-theorem-constructor}{theorem constructor},
Definition~\ref{def:source-link}\,(c2)). For such a link, what the artefact
carries about its source identity \eqref{eq:srcmap} is the link's \termref{term-obligation-record}{obligation
record}, which travels with the link in the source segment
(Definition~\ref{def:source-segment}\,(S4)) and not in the licence record: a
named claim and the assumptions the obligation carries, never a verdict this
module decided, and there is nothing of that identity a consumer could
re-decide.
\end{remark}
\claimmeta{\metalabel{Uses}\mref{Def.}{def:source-segment}, \mref{Def.}{def:srcert}, \mref{Thm.}{thm:certificate-sufficiency}, \mref{Lem.}{lem:matched-pinching}
\metasep \metalabel{Used by}\mref{Req.}{req:attestation-record}
\metasep newly derived
\metasep accepted
\metasep \metalabel{Oracles}\coderef{https://github.com/fraunhofersg/certiqs-key/blob/v0.1.0/tests/oracles/registry_instances_bb84.py}{registry_instances_bb84.py}, \coderef{https://github.com/fraunhofersg/certiqs-key/blob/v0.1.0/tests/oracles/registry_protocol_compiler.py}{registry_protocol_compiler.py}, \coderef{https://github.com/fraunhofersg/certiqs-key/blob/v0.1.0/tests/oracles/registry_transmitter.py}{registry_transmitter.py}}

% claim transmitter/rem:da1
\claimsummary{transmitter/rem:da1}{For a one-dimensional retained system the two source modes impose the same condition and the licence excludes no attack. A trusted non-empty source segment there for which no certificate exists rests on the supplier obligation A-LINK, a case no artefact reaches.}
\claimbookmark{transmitter/rem:da1}{Remark}{rem:da1}{The one-dimensional retained system}
\begin{remark}[The one-dimensional retained system]
\label{rem:da1}
Suppose $d_A=1$. Then $\Seq{A}=\{[1]\}$, so the only marginal a
trusted-mode specification can carry is $\tau_A=[1]$; and for any state
$\omega$ on $A$ and the \termref{term-virtual-transmitted-system}{transmitted system} the $A$-marginal $\omega_A$ is
the $1\times1$ matrix $\tr\omega$, so the conditions $\omega_A=\tau_A$ and
$\tr\omega=1$ are the same condition on the same number. The two source
modes, \termref{term-trusted-source-mode}{trusted-source mode} and \termref{term-untrusted-source-mode}{untrusted-source mode}, therefore impose the same
condition on every consumer, and what the source-replacement certificate
licenses at $d_A=1$ is the condition $\tr\omega=1$, which every normalized
state satisfies: it excludes nothing, and Assumption~\ref{ass:src} is
\emph{vacuous} --- it deletes no attack. That is why the certificate of
Definition~\ref{def:srcert} is demanded only at $d_A\ge2$
(Requirement~\ref{req:licence-scope}). A specification at $d_A=1$ is also
the shape in which a protocol with no trusted source is declared: the
consuming declaration folds Alice's measured system into the receiver-side
systems and names, in its partition of the outcome among the parties, which
components are hers (\xref{meat-core/rem:modes-da1}). That shape is the
\emph{safe} direction (Remark~\ref{rem:source-mode}): it removes an assumption
rather than adding one.

\emph{A non-empty source segment at $d_A=1$.} Every such specification that
an artefact of this module can carry is covered by theorems, and the one case
left to an assumption is reached by no artefact. In detail, the theorem route
applies to a non-empty segment whose consumer declares the ensemble form
(Definition~\ref{def:srcert}, ``What a consumer reads of a presentation'')
and which passes every condition that admits it: the admission rule on the
kind, (S1) to (S4) and (A1) to (A3) of Definition~\ref{def:source-segment},
the refusals of Definition~\ref{def:source-link}\,(d) and (f), and the
admission rule of \xref{reductions/def:chain}, which admits the segment only
in trusted-source mode and only where every entry of the constraint list is
licensed, as \xref{protocol-compiler/def:chaincond} decides it. There the
security statement is carried across the segment by the source identity
\eqref{eq:srcmap} at the descriptors of the emitted links, by
Theorem~\ref{thm:certificate-sufficiency}\,(iii) with (SR1) to (SR3) for the source-replacement
recasting, and by \xref{reductions/thm:two-sided} for the composition with
the receiver segment; at $d_A\ge2$ Assumption~\ref{ass:src} supplies the
attestation. The degenerate declaration of the decoy submodule
(Definition~\ref{def:tagged-source}) is covered by the same route: with one
polarisation, intensity set $\{0\}$, its \termref{term-entropy-cutoff}{entropy cutoff} and
\termref{term-source-cutoff}{source cutoff} both zero and
all weight on the zero-photon sector it has one one-dimensional sector, hence
$d_A=1$; its ensemble presentation supplies (SR1) and (SR2) as theorems
(Lemma~\ref{lem:ensemble-presentation}\,(i)), its defined source state and
marginal are both $[1]$, and (SR3) is vacuous, a trace-preserving map being
unable to act nontrivially on a one-dimensional retained register. The
hypotheses of \xref{reductions/lem:source-transfer}\,(ii), (iii) and
\xref{reductions/thm:two-sided} therefore hold there as well.

\emph{The case left to an assumption.} A trusted specification at $d_A=1$
that declares no certificate, with a non-empty segment and a virtual
ensemble that is not constant over the letters with $p_x>0$, has none to
declare: with a one-dimensional retained register (SR2) would force all
positively weighted virtual states to coincide. The theorem route does not
cover that case. Its transfer across the segment is the source clause of the
protocol compiler's supplier obligation A-LINK
(\xref{protocol-compiler/ass:link}), declared on that module's
chain-condition definition, \xref{protocol-compiler/def:chaincond}, and on
\xref{protocol-compiler/thm:chain-consumption}\,(v$'$), never discharged, and emitted in
its assumption ledger. No artefact reaches the case: (A3) of
Definition~\ref{def:source-segment} accepts a segment only inside an artefact
of the admitting module, and the one admitted constructor kind emits an
ensemble presentation, which supplies the certificate described above.
\end{remark}
\claimmeta{\metalabel{Uses}\mref{Def.}{def:source-spec}, \mref{Def.}{def:srcert}
\metasep \metalabel{Used by}\mref{Req.}{req:licence-scope}, \mref{Req.}{req:marginal-fill}, \mref{Req.}{req:failing-certificate}
\metasep newly derived
\metasep accepted}

\section{Assumptions}\label{group-transmitter/ass:ens}
\sectionintro{Three assumptions carry what declared matrices cannot establish, that the source emits the declared ensemble in every round and independently of every other round, that each round runs preparation, measurement and then announcements, and, in \termref{term-trusted-source-mode}{trusted-source mode}, that Alice's retained system never reaches the adversary. Each is an attestation about the devices and how the protocol is run; the module records each in the artefact's assumption ledger, the list of assumptions the artefact rests on, and never verifies it. A specification declared by the decoy submodule's \termref{term-source-model}{source model} carries that submodule's premise, \hyperlink{claim-transmitter/decoy/ass:wcp}{A-WCP}, in the place of \hyperlink{claim-transmitter/ass:ens}{A-ENS}.}

% claim transmitter/ass:ens
\claimsummary{transmitter/ass:ens}{The physical source draws each round's letter independently with the declared weights and emits the product of the declared states, no round depending on another. It ties a source specification declared as matrices to the device it describes.}
\claimbookmark{transmitter/ass:ens}{Assumption}{ass:ens}{Emission of the specified ensemble}
\claimleadin{%
The assumption is about the source as it runs: in every round it emits
exactly the declared ensemble, not an approximation of it, and it does so
afresh, independently of every other round. The published prepare-and-measure
protocol makes the same premise, Alice preparing each round's state
``independently for each round'' \cite[Protocol~2, step~1(a)]{tupkary2026decoy};
the product it gives the $n$ rounds is the state from which the
source-replacement lemma starts. A specification declared by a \termref{term-source-model}{source
model} carries the premise its submodule declares in the place of this one. In
\termref{term-untrusted-source-mode}{untrusted-source mode} the protocol compiler's assumption on Alice's
measurement governs the measurement by which Alice obtains a letter from a
received system.
}
\SetShortId{A-ENS}
\begin{assumption}[Emission of the specified ensemble]
\label{ass:ens}
For the declared signal ensemble $\{p_x,\rho_x\}$ of
Definition~\ref{def:source-spec} on the \termref{term-transmitted-system}{transmitted system}, the
letters the physical source draws and the systems it emits satisfy the
following; a letter that Alice obtains in untrusted-source mode by measuring a
received system is not drawn by the source, and nothing below concerns it:
(i) the letters $x_1,\dots,x_n$ it draws in the $n$ rounds are drawn
independently, each with the weights $p_x$;
(ii) given those letters, the systems it emits in the $n$ rounds are in the
state $\rho_{x_1}\otimes\dots\otimes\rho_{x_n}$;
(iii) neither the letter it draws in a round nor the state it emits in it
depends on what happens in the other rounds, their outcomes and public
messages included.
\end{assumption}
\claimverification{The engine does not verify this assumption; a run that rests on it records it as unverified.}
\claimreading{Your source prepares, in every round, one of the declared states with the declared probabilities. Each round's letter is drawn independently with weights p\_x, and the states emitted over the n rounds form a product: no round's letter or state depends on another round, its detection results or its public messages. A source with memory or drift between rounds (for example pattern effects in a modulator) violates this. Where your source is a pair source inside your laboratory with one half measured there to herald or fix the letter, those local detectors belong to the declared ensemble: their efficiency mismatch must be carried by the declared weights and states, and their memory across rounds (for example afterpulsing or dead time) violates this like any other dependence between rounds.}
\claimmeta{\metalabel{Uses}\mref{Def.}{def:source-spec}
\metasep \metalabel{Category}source
\metasep \metalabel{Affects}secrecy
\metasep \metalabel{Assumed by}\mref{Fact}{fact:sr}, \extref{../meat-core/contract.pdf}{meat-core/lem:product-marginal}, \extref{../meat-core/contract.pdf}{meat-core/thm:accumulated-entropy}, \extref{../meat-core/contract.pdf}{meat-core/thm:keylength}
\metasep \metalabel{Machine key}\iface{source_ensemble}
\metasep newly derived
\metasep accepted}

% claim transmitter/ass:seq
\claimsummary{transmitter/ass:seq}{Each round runs the delivery of its quantum systems, then every party's measurement, then the round's public messages, then sifting; each party measures in round order, post-processing follows all rounds and every message is public. It is the timing premise of the security analysis.}
\claimbookmark{transmitter/ass:seq}{Assumption}{ass:seq}{Sequential round structure}
\claimleadin{%
With one measuring party this is the structure of
\cite[Def.~8.1, Def.~8.2]{tupkary2026decoy}, whose state evolution is its
Lemma~8.2, and it is the hypothesis under which this module applies the
imported source-replacement lemma, whose own statement places no condition on
the sequential structure of the operations after the preparation. The cited definitions are that case and
no more: they have one measuring party and one announcement block per round,
and the reading for several measuring parties is this contract's own. Of the
messages, the assumption fixes only that all of them follow all measurements
of the round; a message may depend on earlier messages of the same round,
which is a property of the declared messages and not of this assumption. The
clauses have two consequences worth naming. Where one party
measures, as in a prepare-and-measure protocol, every measurement of the round
is that party's. Where two parties measure systems delivered by a source
neither of them prepared, neither party's message precedes the other party's
measurement. Clause (v) is the measurement part of the timing condition of
the published Protocol~2 (Bob's measurement times increasing there), read for
each measuring party \cite[Sec.~8.1.1, Protocol~2]{tupkary2026decoy}; the
monotone order of preparations and of announcements is not adopted, and the
analysis does not need it; with (i) it puts every
measurement of a round before every public message of a later round.
}
\SetShortId{A-SEQ}
\begin{assumption}[Sequential round structure]\label{ass:seq}
The physical protocol has the sequential structure the declared protocol
describes:
(i) each round runs, in this order, state preparation and transmission; the
measurements of the round, by every party that measures in it; the public
messages of the round, in their declared order, \emph{each of them sent after
every measurement of the round is complete}, whichever party sends it and
whichever party measured; and announcement-determined sifting;
(ii) the classical post-processing is applied after the $n$ rounds;
(iii) every message is public;
(iv) no operation on a quantum system of round $j$ depends on a public
message of round $j$;
(v) every party that measures makes its measurements in round order: its
measurement of round $j$ is complete before its measurement of round $j+1$
is.
\end{assumption}
\claimverification{The engine does not verify this assumption; a run that rests on it records it as unverified.}
\claimreading{In every round the quantum systems come first: where your source is trusted, it prepares and sends that round's signal, and otherwise the systems of the round are delivered to the measuring stations. Then every party that measures completes its measurement, and only then are that round's public messages sent, in the declared order. Sifting is then determined by the announcements, and no operation on a quantum system of a round depends on a public message of that round. All classical post-processing (for example error correction, verification and privacy amplification) runs after all n rounds, and every classical message is public. A message sent before a measurement of the same round has finished (for example an early basis announcement) violates this. Each party that measures completes its measurement of each round before it completes that of the next round.}
\claimmeta{\metalabel{Category}protocol structure
\metasep \metalabel{Affects}secrecy
\metasep \metalabel{Assumed by}\mref{Fact}{fact:sr}, \mref{Thm.}{thm:certificate-sufficiency}, \mref{Rem.}{rem:xreg}, \mref{Cor.}{cor:emitted-decoy}, \extref{../meat-core/contract.pdf}{meat-core/lem:honest-attacks}, \extref{../meat-core/contract.pdf}{meat-core/lem:product-marginal}, \extref{../meat-core/contract.pdf}{meat-core/thm:accumulated-entropy}, \extref{../meat-core/contract.pdf}{meat-core/thm:keylength}, \extref{../reductions/contract.pdf}{reductions/lem:source-transfer}, \extref{../reductions/contract.pdf}{reductions/thm:factorization}
\metasep \metalabel{Machine key}\iface{sequential_round_model}
\metasep newly derived
\metasep accepted}

% claim transmitter/ass:src
\claimsummary{transmitter/ass:src}{In trusted-source mode the retained system A stays in the laboratory of the party that holds it and is acted on only by that party's declared measurement, condition (SR3) of the certificate. It is what lets a declared marginal delete attacks.}
\claimbookmark{transmitter/ass:src}{Assumption}{ass:src}{Retention of the system A in trusted-source mode}
\claimleadin{%
The attestation below is the premise the published trusted-source analyses
state in words: \cite[Eqs.~(36)--(39)]{winick2018reliable} adds the marginal
constraints to the entanglement-based statistics constraints for
prepare-and-measure protocols on the stated ground that ``system $A$ is
inaccessible to Eve'', and \cite[Rem.~1]{kamin2025renyi} imposes the
$n$-round product form ``because we assume Eve has no access to Alice's
lab''. Under it, after the source-replacement recasting, the register $A$ is
\emph{inaccessible to the adversary} in every round and its marginal is
$\tau_A$ whatever the adversary does; together with the independence of the
source's rounds, or the premise a \termref{term-source-model}{source model} declares in its place,
it licenses that $n$-round product form. For a prepare-and-measure source the
register $A$ exists only in the source-replacement description; what is
attested of the apparatus is then what those analyses assume in words, that
the adversary has no access to Alice's laboratory. (SR3) is the one part of a
\termref{term-source-replacement-certificate}{source-replacement certificate} that no declared number carries, and it
is demanded only at $d_A\ge2$: at $d_A=1$ the marginal excludes nothing.
}
\SetShortId{A-SRC}
\begin{assumption}[Retention of the system A in trusted-source mode]
\label{ass:src}
For a \termref{term-source-specification}{source specification} (Definition~\ref{def:source-spec}) in
\termref{term-trusted-source-mode}{trusted-source mode}, which fills the marginal slot with
$\tau_A\in\Seq{A}$, condition~(SR3) of Definition~\ref{def:srcert} holds of
the apparatus:
(i) the system $A$ is retained in the laboratory of the party whose label
the party partition of the consuming declaration
(\xref{protocol-compiler/def:declaration}) gives $A$, and is never
transmitted;
(ii) no operation of the protocol or of the adversary acts on $A$ except the
declared measurement $\{M_x\}$ of that party.
\end{assumption}
\claimverification{The engine does not verify this assumption; a run that rests on it records it as unverified.}
\claimreading{The source's retained half, system A, never leaves the laboratory of the party the declaration assigns it to, and nothing acts on it except that party's own declared measurement: no other step of the protocol and no action of the adversary. In a prepare-and-measure setup A exists only in the source-replacement description, and this is the premise that the adversary has no access to that party's laboratory. Only then may you declare the source's reduced state \ensuremath{\tau}\_A (trusted-source mode), which removes attacks from the analysis when A has dimension at least two. If you cannot vouch for this, use untrusted-source mode: the key rate can only drop, and the result stays valid.}
\claimmeta{\metalabel{Uses}\mref{Def.}{def:source-spec}, \mref{Def.}{def:srcert}
\metasep \metalabel{Category}source
\metasep \metalabel{Affects}secrecy
\metasep \metalabel{Assumed by}\mref{Thm.}{thm:certificate-sufficiency}, \mref{Lem.}{lem:matched-pinching}, \mref{Cor.}{cor:emitted-decoy}, \extref{../meat-core/contract.pdf}{meat-core/ex:fss}, \extref{../meat-core/contract.pdf}{meat-core/lem:product-marginal}, \extref{../meat-core/contract.pdf}{meat-core/lem:range}, \extref{../meat-core/contract.pdf}{meat-core/rem:cgate}, \extref{../meat-core/contract.pdf}{meat-core/rem:modes-da1}, \extref{../meat-core/contract.pdf}{meat-core/rem:not-licensed}, \extref{../meat-core/contract.pdf}{meat-core/thm:accumulated-entropy}
\metasep \metalabel{Machine key}\iface{source_replacement_license}
\metasep newly derived
\metasep accepted}

\clearpage
\part{Engine requirements}
\sectionintro{What an implementation of this module must do, record or refuse; each requirement names the result that motivates it.}

\section{Certificate and licence requirements}\label{group-transmitter/req:licence-scope}
\sectionintro{These requirements bind the engine, here the implementation of this module, which answers a request, a submitted \termref{term-source-specification}{source specification} whose marginal slot may be empty, with an artefact or a refusal. The licence is what a verified certificate with its retention attestation gives a consumer, the right to impose the declared marginal on the states it optimises over, and the licence record is its data form. The requirements fix when a certificate is demanded, how the one-dimensional case is served, when a certificate is refused, how a declared marginal is decided, and how the retention attestation is recorded.}

% claim transmitter/req:licence-scope
\claimsummary{transmitter/req:licence-scope}{A trusted-source request must carry the certificate, and with it the retention attestation, only when Alice's retained system has dimension at least two. At dimension one the licensed marginal excludes no state, so demanding them there would refuse data for nothing.}
\claimbookmark{transmitter/req:licence-scope}{Requirement}{req:licence-scope}{Certificate and attestation demanded only at d\_A≥2}
\begin{requirement}[Certificate and attestation demanded only at $d_A\ge2$]\label{req:licence-scope}
The engine demands the \termref{term-source-replacement-certificate}{source-replacement certificate} of
Definition~\ref{def:srcert} of a request in \termref{term-trusted-source-mode}{trusted-source mode} only at
$d_A\ge2$, and there refuses such a request that carries none
(\iface{trusted_mode_without_certificate}). At $d_A=1$, where the licensed
condition $\omega_A=\tau_A$ excludes no state (Remark~\ref{rem:da1}),
neither the engine nor a consumer of its artefact demands the certificate or
the attestation of Assumption~\ref{ass:src}.
\end{requirement}
\claimmeta{\metalabel{Motivated by}\mref{Def.}{def:srcert}, \mref{Rem.}{rem:da1}, \mref{Ass.}{ass:src}
\metasep \metalabel{Used by}\mref{Req.}{req:attestation-record}, \extref{../protocol-compiler/contract.pdf}{protocol-compiler/req:gate-outcomes}
\metasep accepted
\metasep \metalabel{Code}\coderef{https://github.com/fraunhofersg/certiqs-key/blob/v0.1.0/src/certiqs_key/transmitter/artifact.py}{certiqs_key.transmitter.artifact:emit_artifact}}

% claim transmitter/req:marginal-fill
\claimsummary{transmitter/req:marginal-fill}{With a one-dimensional retained system a trusted-source request may leave the marginal empty; the engine fills it with the only state there and records the fill. The record keeps a filled value distinguishable from a declared one, and no other slot is ever filled.}
\claimbookmark{transmitter/req:marginal-fill}{Requirement}{req:marginal-fill}{The unique marginal served at d\_A=1}
\begin{requirement}[The unique marginal served at $d_A=1$]\label{req:marginal-fill}
At $d_A=1$ a request in trusted-source mode may leave the marginal slot
empty. The engine then fills the slot with $\tau_A=[1]$, the unique element
of $\Seq{A}$ (Remark~\ref{rem:da1}); the emitted specification is in
trusted-source mode, trusted-source and \termref{term-untrusted-source-mode}{untrusted-source mode} imposing the
same condition there; and the emitted artefact records that the slot was
filled. This is the only slot the engine ever fills.
\end{requirement}
\claimmeta{\metalabel{Motivated by}\mref{Rem.}{rem:da1}
\metasep accepted
\metasep \metalabel{Code}\coderef{https://github.com/fraunhofersg/certiqs-key/blob/v0.1.0/src/certiqs_key/transmitter/artifact.py}{certiqs_key.transmitter.artifact:emit_artifact}}

% claim transmitter/req:failing-certificate
\claimsummary{transmitter/req:failing-certificate}{A request whose certificate is present and fails a condition decided of its form, pure, block or ensemble presentation, is refused at every dimension, including dimension one, where none is demanded. A certificate that does not verify is a defect in the supplied data, not a weaker claim.}
\claimbookmark{transmitter/req:failing-certificate}{Requirement}{req:failing-certificate}{Refusal of a certificate that fails a decided condition}
\begin{requirement}[Refusal of a certificate that fails a decided condition]\label{req:failing-certificate}
The engine refuses a request whose certificate is present and fails one of
the conditions Definition~\ref{def:srcert} decides of its form --- in \termref{term-pure-form}{pure
form} the three on $\Psi$, (SR1) and (SR2); in \termref{term-block-form}{block form} Hermiticity and
positive semidefiniteness of every block, the trace, the block diagonality of
every weighted virtual member along the matched sectors, and
\eqref{eq:srblock}; in an \termref{term-ensemble-presentation}{ensemble presentation} (E1) to (E3) --- at every
$d_A$, including $d_A=1$, where no certificate is demanded
(Requirement~\ref{req:licence-scope}): a certificate that does not verify is
a defect in the supplied data, not a weaker claim.
\end{requirement}
\claimmeta{\metalabel{Motivated by}\mref{Def.}{def:srcert}, \mref{Rem.}{rem:da1}, \mref{Req.}{req:licence-scope}
\metasep accepted
\metasep \metalabel{Code}\coderef{https://github.com/fraunhofersg/certiqs-key/blob/v0.1.0/src/certiqs_key/transmitter/artifact.py}{certiqs_key.transmitter.artifact:emit_artifact}}

% claim transmitter/req:marginal-decided
\claimsummary{transmitter/req:marginal-decided}{A declared marginal is checked to be a state on the declared matrix itself, even where a certificate satisfies (SR1). An implication proved of the certificate's objects is no check of the datum, so it is never used.}
\claimbookmark{transmitter/req:marginal-decided}{Requirement}{req:marginal-decided}{Direct decision of the declared marginal}
\claimleadin{%
Theorem~\ref{thm:certificate-sufficiency}\,(i) shows that (SR1) implies that $\tau_A$ is a
state. That implication is proved of the objects the declared data describe;
a check of the declared matrix is a check of the datum itself, and it is the
check the engine makes.
}
\begin{requirement}[Direct decision of the declared marginal]\label{req:marginal-decided}
Where $\tau_A$ is declared, the engine decides $\tau_A\in\Seq{A}$ on the
declared matrix directly, also where a certificate is present and satisfies
(SR1), and never infers it from the implication
(SR1)\,$\Rightarrow\tau_A\in\Seq{A}$ of Theorem~\ref{thm:certificate-sufficiency}\,(i); it
refuses a request whose declared $\tau_A$ is not a state
(\iface{alice_marginal_not_admissible}). Consumers of the declared data
decide it in the same way under their own rules; this requirement binds the
engine.
\end{requirement}
\claimmeta{\metalabel{Motivated by}\mref{Thm.}{thm:certificate-sufficiency}
\metasep implemented (effective: accepted)
\metasep \metalabel{Code}\coderef{https://github.com/fraunhofersg/certiqs-key/blob/v0.1.0/src/certiqs_key/transmitter/artifact.py}{certiqs_key.transmitter.artifact:emit_artifact}}

% claim transmitter/req:attestation-record
\claimsummary{transmitter/req:attestation-record}{Where a licence record is emitted, in trusted-source mode at dimension at least two, the retention attestation (SR3) is recorded in it exactly as supplied, never verified and never asserted of a physical apparatus. No function of the declared data decides it, so the engine can only record it.}
\claimbookmark{transmitter/req:attestation-record}{Requirement}{req:attestation-record}{Pass-through of the retention attestation}
\begin{requirement}[Pass-through of the retention attestation]\label{req:attestation-record}
Where the engine emits the licence record \iface{source_replacement_license},
which it does only in \termref{term-trusted-source-mode}{trusted-source mode} at $d_A\ge2$, the scope in
which the certificate is demanded (Requirement~\ref{req:licence-scope}), it
records the (SR3) attestation it receives (Definition~\ref{def:srcert}) in that
record, and passes it through: it never
verifies it and never asserts it of any physical apparatus
(Assumption~\ref{ass:src}; Remark~\ref{rem:boundary}). No consumer can
verify it either, since no function of the declared data decides it
(Theorem~\ref{thm:certificate-sufficiency}\,(i)).
\end{requirement}
\claimmeta{\metalabel{Motivated by}\mref{Def.}{def:srcert}, \mref{Thm.}{thm:certificate-sufficiency}, \mref{Rem.}{rem:boundary}, \mref{Ass.}{ass:src}, \mref{Req.}{req:licence-scope}
\metasep implemented (effective: accepted)
\metasep \metalabel{Code}\coderef{https://github.com/fraunhofersg/certiqs-key/blob/v0.1.0/src/certiqs_key/transmitter/certificate.py}{certiqs_key.transmitter.certificate:SourceReplacementLicense}}

\clearpage
\part{Decoy-state sources and the tagging channel}
\sectionintro{The decoy submodule declares a phase-randomised weak-coherent-pulse source and turns it into a finite virtual source that the engine can read. The physical states live on an infinite-dimensional two-mode Fock space, so they are declared as a source model, with the attestation \hyperlink{claim-transmitter/decoy/ass:wcp}{A-WCP} in the place of \hyperlink{claim-transmitter/ass:ens}{A-ENS}. Up to an entropy cutoff the photon-number sectors stay quantum; above it and up to a source cutoff each photon number and polarisation carries one flag shared by every intensity, which ties a photon number's weight across intensities; the remaining weight sits on one tail flag per polarisation and intensity. Its main results are \hyperlink{claim-transmitter/decoy/prop:tagging}{Proposition~\ref*{prop:tagging}}, one channel carrying this virtual ensemble to the physical states for any sector weights at or below the Poisson weights, and \hyperlink{claim-transmitter/decoy/cor:emitted-decoy}{Corollary~\ref*{cor:emitted-decoy}}, by which the emitted declaration, an ensemble presentation, has an exact source-replacement certificate wherever its declared blocks meet the conditions of that form, and a tagging link whose source identity holds at the emitted descriptors. Every number is used as declared: the Poisson weights are met by a certified inequality and never by an equality, and no square root or quotient is ever declared. The submodule states no key rate and compares none, and its emission, the step that outputs the declaration, refuses entropy cutoffs above two, quantum sectors whose signal dyads it cannot form exactly, flag families finer than the tagged declaration's and formed entries that are not binary64 numbers.}

\section{The source model and the tagged declaration}\label{group-transmitter/decoy/def:wcp-model}
\sectionintro{The \termref{term-weak-coherent-pulse-model}{weak-coherent-pulse model} fixes the declared labels, intensities and phase-averaged coherent states on the two-mode Fock space, which no finite matrix can carry. The tagged declaration \hyperlink{claim-transmitter/decoy/def:tagged-source}{Definition~\ref*{def:tagged-source}} replaces it by a finite virtual source in three bands: quantum photon-number sectors up to the \termref{term-entropy-cutoff}{entropy cutoff}; above it and up to the \termref{term-source-cutoff}{source cutoff}, one one-dimensional \termref{term-polarisation-flag}{polarisation flag} per photon number and polarisation, shared by all intensities; and one \termref{term-tail-flag}{tail flag} per polarisation and intensity for the rest. Sharing a flag across intensities is what relates a photon number's weight at different intensities, the decoy content of the declaration. Its sector weights are declared numbers bounded by the \termref{term-poisson-weight}{Poisson weights}, and its sector list is read off them; \hyperlink{claim-transmitter/decoy/lem:poisson-enclosure}{Lemma~\ref*{lem:poisson-enclosure}} encloses each Poisson weight between two rationals formed from the declared intensity, which is how that bound is decided. \hyperlink{claim-transmitter/decoy/rem:cutoffs}{Remark~\ref*{rem:cutoffs}} separates its two cutoffs from the receiver's.}

% claim transmitter/decoy/def:wcp-model
\claimsummary{transmitter/decoy/def:wcp-model}{A weak-coherent-pulse model declares polarisations, intensities, weighted labels and the phase-averaged coherent states on the two-mode Fock space, as mathematical objects rather than as matrices. It is the specified ensemble over which the tagged declaration is built.}
\claimbookmark{transmitter/decoy/def:wcp-model}{Definition}{def:wcp-model}{Phase-randomised weak-coherent-pulse source model}
\claimleadin{%
A phase-randomised weak-coherent-pulse source prepares, for a polarisation $a$
and an intensity $\mu$, the coherent state of amplitude $\sqrt{\mu}$ in the
mode $\ket{a}$ averaged over its global phase. That average is diagonal in the
photon number, which is why full phase randomisation is one of the conditions
Assumption~\ref{ass:wcp} attests; without it the emitted state is not the one
below. The state is \cite[Eq.~(144)]{tupkary2026decoy}, whose
$\sigma_{(a,\mu,z)}$, $\ket{N}_a$ and setting choice $(a,\mu,z)$ are $\rho_x$,
$\ket{\mathsf{n},a}$ and the letter $x$ here. Its \termref{term-transmitted-system}{transmitted system} is
an infinite-dimensional Fock space, so the model fixes the states as
mathematical objects and not as matrices. It fixes only the
\termref{term-specified-signal-ensemble}{specified ensemble} $\{p_x,\rho_x\}$, the
weighted states the source is declared to prepare:
Alice's retained system, her measurement, her marginal and the certificate come
from the tagged declaration built over it (Definition~\ref{def:tagged-source}).
We write $\mathsf{n}$ for a photon number, $n$ being the number of rounds.
}
\begin{definition}[Phase-randomised weak-coherent-pulse source model]\label{def:wcp-model}
A \termdef{term-weak-coherent-pulse-model}{weak-coherent-pulse model}
is a \termref{term-source-model}{source model}, the case (b) of
Definition~\ref{def:source-spec} in which the specified ensemble
$\{p_x,\rho_x\}$ is fixed as mathematical objects rather than as matrices, on
a transmitted system whose Hilbert space is separable. It consists of the
following declared data.
\begin{enumerate}[label=(W\arabic*),leftmargin=3.0em]
\item \emph{(modes and polarisations)} two declared optical modes with a fixed
orthonormal basis of the one-photon space $\mathbb{C}^2$, and a finite
non-empty set $\mathrm{Pol}$ of
\termdef{term-polarisations}{polarisations},
declared labels each standing for a unit vector $\ket{a}\in\mathbb{C}^2$,
$a\in\mathrm{Pol}$, with creation operator $c^{\dagger}_a$. The vector
$\ket{a}$ is a mathematical object of the model, as the states of (W4) are,
and not an array of declared numbers: a declaration names the vector a label
stands for and carries none of its coordinates, which need not be binary64
numbers ($(\ket{H}+\ket{V})/\sqrt2$ has the coordinates $1/\sqrt2$). The
transmitted system $\Asig$ is the Fock space over the two modes, with vacuum
$\ket{\vac}$; its Hilbert space is separable and infinite dimensional.
\item \emph{(intensities and labels)} a finite non-empty set
$\mathrm{Int}\subset[0,\infty)$ of
\termdef{term-intensities}{intensities};
the signal alphabet
$\Xset\subseteq\mathrm{Pol}\times\mathrm{Int}\times\{\tst,\gen\}$ with its
\exttermref{../foundations/contract.pdf}{term-declared-order}{declared order}, an element of
$\mathrm{Pol}\times\mathrm{Int}\times\{\tst,\gen\}$ being a label and a label in
$\Xset$ a letter, written $x=(a_x,\mu_x,z_x)$, whose third component is the
\termdef{term-round-type}{round type};
and weights $p_x>0$ for every $x\in\Xset$ with $\sum_{x\in\Xset}p_x=1$ in the
declared order, \exttermref{../foundations/contract.pdf}{term-rounds-nowhere}{rounding nowhere}. A label
of weight zero is not a letter of $\Xset$.
\item \emph{(number states)} for $a\in\mathrm{Pol}$ and an integer
$\mathsf{n}\ge0$, the unit vector
\begin{equation}\label{eq:numberstate}
  \ket{\mathsf{n},a}\;\coloneqq\;
  \frac{(c^{\dagger}_a)^{\mathsf{n}}}{\sqrt{\mathsf{n}!}}\,\ket{\vac} ,
\end{equation}
and the subspace
$S_{\mathsf{n}}\coloneqq\operatorname{span}\{\ket{\mathsf{n},a}:a\in\mathrm{Pol}\}$
of the $\mathsf{n}$-photon subspace. Since the $\mathsf{n}$-photon subspace of
a two-mode Fock space has dimension $\mathsf{n}+1$ and $S_{\mathsf{n}}$ is
spanned by $|\mathrm{Pol}|$ vectors,
$\dim S_{\mathsf{n}}\le\min(\mathsf{n}+1,|\mathrm{Pol}|)$; and
$S_{\mathsf{n}}\perp S_{\mathsf{n}'}$ for $\mathsf{n}\ne\mathsf{n}'$, the two
subspaces lying in different photon-number eigenspaces.
$S_0=\mathbb{C}\ket{\vac}$ is one dimensional, every polarisation sharing the
vacuum.
\item \emph{(the specified states)} with the
\termdef{term-poisson-weight}{Poisson weight}
$P_{\mu}(\mathsf{n})\coloneqq e^{-\mu}\mu^{\mathsf{n}}/\mathsf{n}!$, the state
prepared at the letter $x=(a,\mu,z)$ is
\begin{equation}\label{eq:wcpstate}
  \rho_x\;=\;\sum_{\mathsf{n}\ge0}P_{\mu}(\mathsf{n})\,
  \dyad{\mathsf{n},a} ,
\end{equation}
which does not depend on the round type $z$. It is a state:
$\sum_{\mathsf{n}\ge0}P_{\mu}(\mathsf{n})=1$ and each $\dyad{\mathsf{n},a}$
is a rank-one projector.
\end{enumerate}
The physical-source assumption of the model, which asserts that the device
emits \eqref{eq:wcpstate}, is Assumption~\ref{ass:wcp}; the model declares it
in the place of Assumption~\ref{ass:ens}.
\end{definition}
\claimmeta{\metalabel{Uses}\mref{Def.}{def:source-spec}
\metasep \metalabel{Used by}\mref{Def.}{def:tagged-source}, \mref{Lem.}{lem:poisson-enclosure}, \mref{Prop.}{prop:tagging}, \mref{Prop.}{prop:decoy-cert}, \mref{Cor.}{cor:emitted-decoy}, \mref{Ass.}{ass:wcp}, \mref{Rem.}{rem:wcp-virtual-model}, \mref{Req.}{req:signal-dyads}
\metasep newly derived
\metasep implemented
\metasep \metalabel{Code}\coderef{https://github.com/fraunhofersg/certiqs-key/blob/v0.1.0/src/certiqs_key/transmitter/decoy/model.py}{certiqs_key.transmitter.decoy.model:WcpModel}}

% claim transmitter/decoy/def:tagged-source
\claimsummary{transmitter/decoy/def:tagged-source}{A tagged declaration is a finite virtual source on the sectors a letter weights: quantum to the entropy cutoff, then to the source cutoff a flag per photon number and polarisation shared by all intensities, then a tail flag per polarisation and intensity. Its weights are at most the Poisson ones.}
\claimbookmark{transmitter/decoy/def:tagged-source}{Definition}{def:tagged-source}{Tagged virtual source: cutoffs, flags and sector map}
\claimleadin{%
The model's states live on an infinite-dimensional Fock space, which no
declaration can carry as matrices. The tagged declaration replaces them by a
finite virtual source with three bands of photon numbers: up to the entropy
cutoff each photon-number sector stays quantum; above it, up to the source
cutoff, each photon number and polarisation keeps a one-dimensional flag; and
the remaining weight goes to one tail flag per polarisation and intensity. Only
the sectors that some letter weights are declared. A polarisation flag is
shared by every intensity, so that the weight of one photon number at different
intensities sits on one sector; that relation is the decoy content of the
declaration (Remark~\ref{rem:no-intensity-split} gives the detail). The
Poisson weights $P_\mu(\mathsf{n})$ the sector weights are bounded by are not
declared numbers, so the inequality of (T2) is not decided against them: the
emission decides it against a rational lower bound on $P_\mu(\mathsf{n})$
that it forms from the declared intensity
(Lemma~\ref{lem:poisson-enclosure}, Requirement~\ref{req:sector-weight-bound}).
The polarisation vectors, the number states and the basis of each quantum
sector are objects of the model as well, so the entries of the signal dyads in
that basis are not declared either: the emission forms them by a construction
(Requirement~\ref{req:signal-dyads}). The
construction generalises the virtual source of
\cite[Eq.~(145)]{tupkary2026decoy}, which is the case
$N_{\mathrm{src}}=N_{\mathrm{ent}}$ with the exact Poisson weights
(Proposition~\ref{prop:tagging}\,(iv)): here the declared weights need only lie
at or below the Poisson weights, and the polarisation flags between the two
cutoffs and the sector map are new. Remark~\ref{rem:cutoffs} compares the two
cutoffs with the receiver's.
}
\begin{definition}[Tagged virtual source: cutoffs, flags and sector map]\label{def:tagged-source}
Let a weak-coherent-pulse model be declared (Definition~\ref{def:wcp-model}). A
\termdef{term-tagged-source-declaration}{tagged source declaration}
over it consists of the following.
\begin{enumerate}[label=(T\arabic*),leftmargin=3.0em]
\item \emph{(the two cutoffs)} integers $N_{\mathrm{ent}}$ and
$N_{\mathrm{src}}$ with $0\le N_{\mathrm{ent}}\le N_{\mathrm{src}}$: the
\termdef{term-entropy-cutoff}{entropy cutoff},
the largest photon number whose sector stays quantum, and the
\termdef{term-source-cutoff}{source cutoff},
the largest photon number that keeps a sector of its own.
\item \emph{(the sector weights)} declared numbers $\tilde p(\mathsf{n}\mid\mu)$
for $0\le\mathsf{n}\le N_{\mathrm{src}}$ and $\mu\in\mathrm{Int}$, depending on
the intensity alone and not on the polarisation or the round type, satisfying
\begin{equation}\label{eq:weightineq}
  0\;\le\;\tilde p(\mathsf{n}\mid\mu)\;\le\;P_{\mu}(\mathsf{n})
  \qquad\text{for every }\mathsf{n}\le N_{\mathrm{src}},\ \mu\in\mathrm{Int} ,
\end{equation}
and the
\termdef{term-tail-weight}{tail weight}
\begin{equation}\label{eq:tailweight}
  \tilde w_{\mu}\;\coloneqq\;
  1-\sum_{\mathsf{n}=0}^{N_{\mathrm{src}}}\tilde p(\mathsf{n}\mid\mu) ,
\end{equation}
the sum formed in increasing $\mathsf{n}$ and rounding nowhere. Then
$\tilde w_{\mu}\ge0$, because
$\sum_{\mathsf{n}\le N_{\mathrm{src}}}\tilde p(\mathsf{n}\mid\mu)
\le\sum_{\mathsf{n}\le N_{\mathrm{src}}}P_{\mu}(\mathsf{n})\le1$ by
\eqref{eq:weightineq}.
\item \emph{(the sector index set and the \termref{term-virtual-transmitted-system}{virtual transmitted system})} a
sector that no letter of $\Xset$ puts positive weight on is not declared: the
\termdef{term-sector-index-set}{sector index set}
carries the sectors the declaration uses and is read off the weights,
\begin{equation}\label{eq:tset}
\begin{aligned}
  \mathbb{T}\;=\;
  &\{\mathsf{n}\le N_{\mathrm{ent}}:\
     \tilde p(\mathsf{n}\mid\mu_x)>0\text{ for some }x\in\Xset\}\\
  \cup\;&\{(\mathsf{n},a):N_{\mathrm{ent}}<\mathsf{n}\le N_{\mathrm{src}},\
     \tilde p(\mathsf{n}\mid\mu_x)>0\text{ for some }x\in\Xset
     \text{ with }a_x=a\}\\
  \cup\;&\{(a,\mu):\tilde w_{\mu}>0\text{ and }(a,\mu)=(a_x,\mu_x)
     \text{ for some }x\in\Xset\} ,
\end{aligned}
\end{equation}
in that declared order: the quantum sectors in increasing photon number, then
the polarisation flags in the declared order of $(\mathsf{n},a)$, then the tail
flags in the declared order of $(a,\mu)$. The finite-dimensional virtual
transmitted system is the orthogonal direct sum over it,
\begin{equation}\label{eq:avirt}
  \Avirt\;=\;\bigoplus_{t\in\mathbb{T}}S_t ,\qquad
  S_t=\begin{cases}
    S_{\mathsf{n}} & t=\mathsf{n}\le N_{\mathrm{ent}},\\
    \mathbb{C}f_{\mathsf{n},a} & t=(\mathsf{n},a),\\
    \mathbb{C}f_{a,\mu} & t=(a,\mu),
  \end{cases}
\end{equation}
with the basis of each summand described next; it is the declared basis of
$\Avirt$ in the sense of Definition~\ref{def:srcert}, the basis in which
the blocks of the declaration are written. At a quantum sector
$t=\mathsf{n}$ the summand is the space $S_{\mathsf{n}}$ of
Definition~\ref{def:wcp-model}\,(W3), taken into the direct sum with its
vectors $\ket{\mathsf{n},a}$ and an orthonormal basis of $S_{\mathsf{n}}$;
like those vectors, that basis is a mathematical object of the model and not
an array of declared numbers (Proposition~\ref{prop:tagging} returns
$S_{\mathsf{n}}$ to the Fock space by an isometry $V_{\mathsf{n}}$). There is
one
\termdef{term-polarisation-flag}{polarisation flag}
$f_{\mathsf{n},a}$ per photon number above the entropy cutoff and per
polarisation of a letter that weights it, shared across intensities and round
types; and one
\termdef{term-tail-flag}{tail flag}
$f_{a,\mu}$ per pair a letter carries with
$\tilde w_{\mu}>0$, shared across round types. The flag vectors are declared
orthonormal, the summands of \eqref{eq:avirt} are mutually orthogonal in
$\Avirt$, and $\Pi^{\mathrm v}_t$ is the projector onto $S_t$. Every flag
sector is one dimensional, and so is the quantum sector $\mathsf{n}=0$. The
index set is a function of the declared weights.
\item \emph{(the \termref{term-virtual-ensemble}{virtual ensemble})} for $x=(a,\mu,z)\in\Xset$,
\begin{equation}\label{eq:rvirt}
  \rvirt_x\;=\;
  \sum_{\mathsf{n}\le N_{\mathrm{ent}}}
    \tilde p(\mathsf{n}\mid\mu)\,\dyad{\mathsf{n},a}
  \;+\;
  \sum_{N_{\mathrm{ent}}<\mathsf{n}\le N_{\mathrm{src}}}
    \tilde p(\mathsf{n}\mid\mu)\,\dyad{f_{\mathsf{n},a}}
  \;+\;
  \tilde w_{\mu}\,\dyad{f_{a,\mu}} ,
\end{equation}
every sum running over the sectors of the index set \eqref{eq:tset} only: a
quantum sector or a polarisation flag that no letter puts positive weight on is
not a sector of $\Avirt$ by (T3) and contributes no term, as the tail flag
contributes none where $\tilde w_{\mu}=0$. Each $\rvirt_x$ is a state on
$\Avirt$ that is block diagonal along $\mathbb{T}$. Its component in a quantum
sector $\mathsf{n}\le N_{\mathrm{ent}}$ is
$\tilde p(\mathsf{n}\mid\mu)\dyad{\mathsf{n},a}$; in the polarisation flag
sector $(\mathsf{n},a')$ it is $\tilde p(\mathsf{n}\mid\mu)$ if $a'=a$ and zero
otherwise; in the tail flag sector $(a',\mu')$ it is $\tilde w_{\mu}$ if
$(a',\mu')=(a,\mu)$ and zero otherwise. It is positive semidefinite, being a
sum of non-negative multiples of rank-one projectors, and its trace is
$\sum_{\mathsf{n}\le N_{\mathrm{src}}}\tilde p(\mathsf{n}\mid\mu)
+\tilde w_{\mu}=1$ by \eqref{eq:tailweight}, the absent terms contributing
nothing to it: a sector that \eqref{eq:tset} leaves out is one on which
$\tilde p(\mathsf{n}\mid\mu)$ vanishes for every letter, this one included, or
a tail flag with $\tilde w_{\mu}=0$. Two labels of the same polarisation and
different intensities put weight on the same polarisation flag, and two labels
of the same polarisation and intensity and different round types put weight on
the same tail flag.
\item \emph{(the matched sector map)} the retained system and its
\termref{term-sector-structure}{sector structure} are not declared: a tagged source is presented as an
\termref{term-ensemble-presentation}{ensemble presentation} (Definition~\ref{def:srcert}: a certificate defined from declared
weighted blocks rather than declared), in which $A$, its structure, $\tau_A$
and $\{M_x\}$ are the objects \eqref{eq:enspres} defines, $A_t$ being a copy of
$S_t$ sector by sector and in the same order. In particular
$\dim\Hs{A_t}=\dim\Hs{S_t}$ whatever the rank of the sector's marginal
(Proposition~\ref{prop:decoy-cert}\,(iii)). The matched pair
$(\{\Pi^{A}_t\},\{\Pi^{\mathrm v}_t\})$, with $\Pi^{A}_t$ the projectors onto
Alice's sectors (Definition~\ref{def:source-spec}), is the structure of (T3)
read on both sides. It is the \termref{term-retained-register-sector-map}{retained-register sector map} of the link of
(T6), the slot (e) of Definition~\ref{def:source-link}, which fixes how the
certificate, $\tau_A$ and $\{M_x\}$ are presented, sector by sector. It is a
statement about those three objects only: it does not assert that any joint
state of $A$ and another system is block diagonal
(Definition~\ref{def:source-spec}).
\item \emph{(the link)} the channel $\Xi_{\mathrm{tag}}$ that
Proposition~\ref{prop:tagging} constructs and proves to carry each $\rvirt_x$
to $\rho_x$, recorded as one \termref{term-source-link}{source link}
(Definition~\ref{def:source-link}) with the following slots: the kind
\textsf{tagging} (slot (a)); the \termref{term-source-descriptor}{source descriptor}
$(\Avirt,\{p_x,\rvirt_x\})$ of (T3) and (T4) and the \termref{term-target-descriptor}{target descriptor}
$(\Asig,\{p_x,\rho_x\})$ of the model (slot (b)); the channel witness supplied
as a \termref{term-theorem-constructor}{theorem constructor}, namely Proposition~\ref{prop:tagging}, with an
\termref{term-obligation-record}{obligation record} that names Assumption~\ref{ass:wcp} (slot (c2)); the
\termref{term-declared-descriptor-family}{declared descriptor family}, every tagged source declaration over every
weak-coherent-pulse model (slot (d)); the retained-register sector map, the
matched pair of (T5) (slot (e)); an empty list of
\termref{term-transported-constraint-entries}{transported constraint entries}, so that no constraint is carried
across the link, and hence no licence record, a licence being present only
beside a non-empty list (slot (f)); and the status \textsf{exact}, the identity
holding with no error term (slot (g)).
\end{enumerate}
The cutoffs of (T1), the weights $\tilde p(\mathsf{n}\mid\mu)$ of (T2), and the
model's alphabet, polarisation labels, intensities and label weights are
\emph{declared}. The polarisation vectors $\ket{a}$, the number states
$\ket{\mathsf{n},a}$ and the orthonormal basis of each quantum sector
$S_{\mathsf{n}}$ of (T3) are objects of the model and are not declared
numbers. The tail weight \eqref{eq:tailweight}, the index set \eqref{eq:tset}
and the \termref{term-weighted-members}{weighted members} $p_x\rvirt_x$ of \eqref{eq:rvirt} are
\emph{formed}: the tail weight and the index set from the declared numbers in
exact arithmetic, and each entry of a weighted member, in the basis of its
sector, as the product of declared numbers with an entry of the dyad of the
sector's signal vector in that basis, a number fixed by the model and the
basis. They are the declaration's own data: read sector by sector, the
weighted members are the blocks that the ensemble presentation of (T5) takes as
its data. Everything on Alice's side --- $A$, $\tau_A$, $\{M_x\}$ and the
certificate --- is defined from those blocks and declared by no one.
\end{definition}
\claimmeta{\metalabel{Uses}\mref{Def.}{def:source-link}, \mref{Def.}{def:source-segment}, \mref{Def.}{def:wcp-model}
\metasep \metalabel{Used by}\mref{Rem.}{rem:cutoffs}, \mref{Prop.}{prop:tagging}, \mref{Rem.}{rem:no-intensity-split}, \mref{Prop.}{prop:decoy-cert}, \mref{Cor.}{cor:emitted-decoy}, \mref{Ass.}{ass:wcp}, \mref{Rem.}{rem:wcp-virtual-model}, \mref{Req.}{req:decoy-emission}, \mref{Req.}{req:cutoff-order}, \mref{Req.}{req:signal-dyads}, \mref{Req.}{req:sector-weight-bound}, \mref{Req.}{req:shared-flags}
\metasep newly derived
\metasep implemented
\metasep \metalabel{Code}\coderef{https://github.com/fraunhofersg/certiqs-key/blob/v0.1.0/src/certiqs_key/transmitter/decoy/tagged.py}{certiqs_key.transmitter.decoy.tagged:TaggedSourceDeclaration}}

% claim transmitter/decoy/lem:poisson-enclosure
\claimsummary{transmitter/decoy/lem:poisson-enclosure}{For a rational intensity \ensuremath{\mu} \ensuremath{\ge} 0 and an even integer K \ensuremath{\ge} \ensuremath{\mu}, two consecutive partial sums of the exponential series, scaled by \ensuremath{\mu}\textasciicircum{}n/n!, enclose the Poisson weight between rationals. The enclosure narrows monotonically to zero width as K grows, so the bound on the sector weights is decided exactly.}
\claimbookmark{transmitter/decoy/lem:poisson-enclosure}{Lemma}{lem:poisson-enclosure}{Rational enclosure of the Poisson weights}
\claimleadin{%
The inequality \eqref{eq:weightineq} of Definition~\ref{def:tagged-source}\,(T2)
compares a declared number with $P_\mu(\mathsf{n})$, which at a positive
intensity is transcendental. A declared intensity is a binary64 number, hence a
dyadic rational, and the exponential series at a negative rational argument is
alternating with terms that decrease from some index on, so two consecutive
partial sums bracket $e^{-\mu}$. Multiplied by the rational factor
$\mu^{\mathsf{n}}/\mathsf{n}!$ they bracket the Poisson weight, and every
number involved is a rational formed from $\mu$ in finitely many steps.
}
\begin{technicalclaim}
\begin{lemma}[Rational enclosure of the Poisson weights]\label{lem:poisson-enclosure}
Let $\mu\ge0$ be a rational number. For integers $k\ge0$ put
$t_k\coloneqq\mu^k/k!$, with $\mu^0\coloneqq1$, and for integers $K\ge0$
\begin{equation}\label{eq:expsum}
  S_K\;\coloneqq\;\sum_{k=0}^{K}(-1)^k\,t_k .
\end{equation}
Let $K$ be an even integer with $K\ge\mu$, and let $\mathsf{n}\ge0$ be an
integer. Then:
\begin{enumerate}[label=(\roman*),leftmargin=2.2em]
\item \emph{(the enclosure)} with the \termref{term-poisson-weight}{Poisson weight}
$P_\mu(\mathsf{n})=e^{-\mu}\mu^{\mathsf{n}}/\mathsf{n}!$ of
Definition~\ref{def:wcp-model}\,(W4),
\begin{equation}\label{eq:poisson-enclosure}
  \max\big(0,S_{K+1}\big)\,t_{\mathsf{n}}
  \;\le\;P_\mu(\mathsf{n})\;\le\;S_K\,t_{\mathsf{n}} ;
\end{equation}
\item \emph{(rational ends)} both ends of \eqref{eq:poisson-enclosure} are
rational numbers, formed from $\mu$ by finitely many additions, subtractions,
multiplications and divisions by positive integers, and a comparison of
either end with a binary64 number is an exact comparison of two rationals;
\item \emph{(the width)} $S_K-S_{K+1}=t_{K+1}$, so the two ends of
\eqref{eq:poisson-enclosure} differ by at most $t_{K+1}t_{\mathsf{n}}$, and
$t_{K+1}\to0$ as $K\to\infty$ through the even integers. At $\mu=0$ both ends
equal $P_0(\mathsf{n})$, which is $1$ for $\mathsf{n}=0$ and $0$ for
$\mathsf{n}\ge1$.
\item \emph{(monotone in $K$)} over the even integers $K\ge\mu$ the lower end
$\max(0,S_{K+1})\,t_{\mathsf{n}}$ is non-decreasing and the upper end
$S_K\,t_{\mathsf{n}}$ is non-increasing: for even $K\le K'$ with $K\ge\mu$,
the enclosure at $K'$ lies inside the enclosure at $K$.
\end{enumerate}
\end{lemma}
\end{technicalclaim}
\claimmeta{\metalabel{Uses}\mref{Def.}{def:wcp-model}
\metasep \metalabel{Used by}\mref{Req.}{req:sector-weight-bound}
\metasep newly derived
\metasep implemented
\metasep \metalabel{Code}\coderef{https://github.com/fraunhofersg/certiqs-key/blob/v0.1.0/src/certiqs_key/transmitter/decoy/enclosure.py}{certiqs_key.transmitter.decoy.enclosure:poisson_enclosures}}

\begin{proof}
\emph{Idea.} The tail of the exponential series beyond $K$ is alternating with
non-increasing terms, so grouping it in consecutive pairs fixes its sign.

\emph{Step 1 (the series).} The exponential series
$e^{-\mu}=\sum_{k\ge0}(-\mu)^k/k!=\sum_{k\ge0}(-1)^kt_k$ converges absolutely,
$\sum_{k\ge0}t_k=e^{\mu}$ being finite. Hence for every $K'\ge0$
\begin{equation*}
  e^{-\mu}-S_{K'}=\sum_{k>K'}(-1)^kt_k ,
\end{equation*}
and this tail may be summed in consecutive pairs of terms: the partial sums of
the paired series are partial sums of the tail, which converge to its value.

\emph{Step 2 (the tail decreases).} For $k\ge K+1$,
$t_{k+1}=t_k\,\mu/(k+1)$ with $0\le\mu\le K<k+1$, so $0\le t_{k+1}\le t_k$.

\emph{Step 3 (the upper end).} $K$ is even, so the first term of the tail of
$S_K$ is $(-1)^{K+1}t_{K+1}=-t_{K+1}$, and pairing gives
\begin{equation*}
  e^{-\mu}-S_K=-\sum_{j\ge0}\big(t_{K+1+2j}-t_{K+2+2j}\big)\le0 ,
\end{equation*}
every bracket being non-negative by Step~2. So $e^{-\mu}\le S_K$.

\emph{Step 4 (the lower end).} $K+2$ is even, so the first term of the tail of
$S_{K+1}$ is $t_{K+2}$, and pairing gives
\begin{equation*}
  e^{-\mu}-S_{K+1}=\sum_{j\ge0}\big(t_{K+2+2j}-t_{K+3+2j}\big)\ge0 ,
\end{equation*}
again by Step~2. So $S_{K+1}\le e^{-\mu}$, and since $e^{-\mu}>0$ also
$\max(0,S_{K+1})\le e^{-\mu}$.

\emph{Step 5 (the Poisson weight).} $P_\mu(\mathsf{n})=e^{-\mu}t_{\mathsf{n}}$
with $t_{\mathsf{n}}\ge0$, and multiplying the inequalities of Steps 3 and 4 by
the non-negative number $t_{\mathsf{n}}$ preserves them. This is (i).

(ii) $t_0=1$ and $t_k=t_{k-1}\,\mu/k$ for $k\ge1$, so every $t_k$ is a rational
formed from $\mu$ by multiplications and divisions by positive integers; $S_K$
and $S_{K+1}$ are finite sums and differences of them, and the maximum of two
rationals is one of them. A binary64 number is a dyadic rational, so the
comparison is one of two rationals.

(iii) $S_K-S_{K+1}=-(-1)^{K+1}t_{K+1}=t_{K+1}$, $K+1$ being odd. The upper end
minus the lower end is
$\big(S_K-\max(0,S_{K+1})\big)t_{\mathsf{n}}\le(S_K-S_{K+1})t_{\mathsf{n}}
=t_{K+1}t_{\mathsf{n}}$. The $t_k$ are the terms of the convergent series of
Step~1, so $t_k\to0$. At $\mu=0$, $t_k=0$ for $k\ge1$, so the width
$t_{K+1}t_{\mathsf{n}}$ is zero and both ends equal $P_0(\mathsf{n})$; and
$S_K=S_{K+1}=1$, so the common value is $t_{\mathsf{n}}$, which is $1$ for
$\mathsf{n}=0$ and $0$ for $\mathsf{n}\ge1$.

(iv) It suffices to compare $K$ with $K+2$, both even and at least $\mu$. By
\eqref{eq:expsum}, $S_{K+3}-S_{K+1}=t_{K+2}-t_{K+3}\ge0$ and
$S_{K+2}-S_K=-t_{K+1}+t_{K+2}\le0$, both by Step~2, which applies from the
index $K+1$ on. So $S_{K+1}$ does not decrease and $S_K$ does not increase
from $K$ to $K+2$; $\max(0,\cdot)$ preserves the first, and multiplying by
$t_{\mathsf{n}}\ge0$ preserves both.
\end{proof}

% claim transmitter/decoy/rem:cutoffs
\claimsummary{transmitter/decoy/rem:cutoffs}{The entropy cutoff bounds the quantum sectors and the source cutoff the flagged ones, and raising the second never requires raising the first. The receiver's cutoff is a different object, which no condition of this submodule reads.}
\claimbookmark{transmitter/decoy/rem:cutoffs}{Remark}{rem:cutoffs}{Entropy, source and receiver cutoffs}
\begin{remark}[Entropy, source and receiver cutoffs]\label{rem:cutoffs}
A reader of a protocol with a decoy source meets three cutoffs that are easily
conflated, and only two of them are data of this submodule.

\emph{The entropy cutoff $N_{\mathrm{ent}}$} is the largest photon number whose
sector of $\Avirt$ stays quantum: at $\mathsf{n}\le N_{\mathrm{ent}}$ the sector
is $S_{\mathsf{n}}$ and carries the signal states $\ket{\mathsf{n},a}$
themselves, one sector for all polarisations of that photon number. Above it the
polarisation is tagged: the sector is one dimensional and there is one per
polarisation, so the declaration separates what a quantum sector keeps
together. Whether keeping a photon number quantum or tagging it per
polarisation gives the higher key rate is a question for the consuming program,
settled by computation and not here.

\emph{The source cutoff $N_{\mathrm{src}}$} is the largest photon number that
keeps a sector of its own at all. Raising it by one adds at most
$|\mathrm{Pol}|$ one-dimensional sectors, one per polarisation some letter
carries at a positive weight, the index set of
Definition~\ref{def:tagged-source}\,(T3) being read off the weights; and it
moves the weight $\tilde p(\mathsf{n}\mid\mu)$ out of the tail flag into them,
so that the weight of that photon number sits on a flag shared by every
intensity instead of on a flag carried by one intensity alone. Raising
$N_{\mathrm{src}}$ never requires raising $N_{\mathrm{ent}}$: (T1) demands only
$N_{\mathrm{ent}}\le N_{\mathrm{src}}$, and the sectors between the two cutoffs
are one dimensional whatever $N_{\mathrm{src}}$ is.

\emph{The receiver cutoff}, the photon number at which a receiver-side
reduction stops retaining blocks and starts flagging, is a datum of the
receiver and of the reduction that consumes it, and not of this submodule. No
condition of this submodule mentions it, no claim here reads it, and nothing
ties it to the entropy or the source cutoff: the transmitter's data and the
receiver's meet only where the protocol compiler composes them. A declaration
with $N_{\mathrm{ent}}=1$ and a receiver cutoff of two is as well formed as any
other.

Every claim of this submodule is stated for arbitrary
$N_{\mathrm{ent}}\le N_{\mathrm{src}}$. The bound $N_{\mathrm{ent}}\le2$ is a
requirement of the emission, the step that outputs the declaration
(Requirement~\ref{req:entropy-cutoff}), and concerns which numbers a
declaration can carry exactly, not any of the mathematics here: for the four
polarisations $\Dset$, at one and at two photons,
Corollary~\ref{cor:emitted-decoy}\,(iv)(b) exhibits a declared basis in which
the signal dyads have the entries $0$ and $\pm\tfrac12$ (and at one photon
also $1$); at three photons and above no such basis is given here.
\end{remark}
\claimmeta{\metalabel{Uses}\mref{Def.}{def:tagged-source}
\metasep newly derived
\metasep accepted}

\section{The tagging channel}\label{group-transmitter/decoy/prop:tagging}
\sectionintro{The virtual and the physical source are linked by one channel, the tagging channel, which is fixed before a letter is drawn. The main result of the section, \hyperlink{claim-transmitter/decoy/prop:tagging}{Proposition~\ref*{prop:tagging}}, constructs it for any sector weights at or below the \termref{term-poisson-weight}{Poisson weights}: it keeps each quantum sector, prepares from each \termref{term-polarisation-flag}{polarisation flag} the number state it stands for, and prepares from each \termref{term-tail-flag}{tail flag} the normalised remainder of the Poisson distribution. So a declared weight need only satisfy a certified inequality, never equal a Poisson weight. \hyperlink{claim-transmitter/decoy/rem:no-intensity-split}{Remark~\ref*{rem:no-intensity-split}} explains why a polarisation flag is shared across intensities: a finer family would still admit such a channel but would lose the relation across intensities that is the decoy content of the declaration.}

% claim transmitter/decoy/prop:tagging
\claimsummary{transmitter/decoy/prop:tagging}{One channel, independent of the label, maps each tagged virtual state to its phase-randomised weak coherent state whenever the declared sector weights lie at or below the Poisson weights. The weights therefore need only a certified inequality, never an exact Poisson weight.}
\claimbookmark{transmitter/decoy/prop:tagging}{Proposition}{prop:tagging}{Tagging channel for weights at or below the Poisson weights}
\claimleadin{%
Recall the virtual source of Definition~\ref{def:tagged-source}\,(T3) and
(T4): $\Avirt$ is the direct sum, over the sectors that some letter weights, of
the quantum sectors $S_{\mathsf{n}}$ for $\mathsf{n}\le N_{\mathrm{ent}}$, the
one-dimensional \termref{term-polarisation-flag}{polarisation flags}
$\mathbb{C}f_{\mathsf{n},a}$ for $N_{\mathrm{ent}}<\mathsf{n}\le N_{\mathrm{src}}$
and the one-dimensional \termref{term-tail-flag}{tail flags} $\mathbb{C}f_{a,\mu}$; the
virtual state $\rvirt_x$ of a letter $x=(a,\mu,z)$ puts
$\tilde p(\mathsf{n}\mid\mu)\dyad{\mathsf{n},a}$ on each quantum sector, the
weight $\tilde p(\mathsf{n}\mid\mu)$ on the polarisation flag $f_{\mathsf{n},a}$
of its own polarisation, and the \termref{term-tail-weight}{tail weight} $\tilde w_\mu$ on its tail
flag $f_{a,\mu}$. The tail flag absorbs whatever weight the declared sectors
leave out, so the construction needs the sector weights only to lie at or below
the \termref{term-poisson-weight}{Poisson weights}: a declared weight may be any number
certified to lie in $[0,P_\mu(\mathsf{n})]$, and none need equal a Poisson
weight. Lowering $\tilde p(\mathsf{n}\mid\mu)$ moves the
difference into $\tilde w_\mu$, and so into the tail flag of every polarisation
at that intensity; Corollary~\ref{cor:emitted-decoy}\,(v) says what else
changes. Each flag, polarisation or tail, is sent to a state that depends on
the flag alone, so one channel, fixed before a letter is drawn, serves every
letter, which is what a \termref{term-source-link}{source link} requires of its channel witness
(Definition~\ref{def:source-link}). That holds for any flag family, the tail
flags indexed by $\mu$ included; Remark~\ref{rem:no-intensity-split} says why
the polarisation flags are nonetheless shared across intensities.
}
\begin{mainclaim}
\begin{proposition}[Tagging channel for weights at or below the Poisson weights]\label{prop:tagging}
Let a \termref{term-weak-coherent-pulse-model}{weak-coherent-pulse model} be given
(Definition~\ref{def:wcp-model}): labels $x=(a,\mu,z)$ with $z$ the
\termref{term-round-type}{round type}, unit polarisation vectors
$\ket{a}$, Poisson weights
$P_\mu(\mathsf{n})=e^{-\mu}\mu^{\mathsf{n}}/\mathsf{n}!$ and model states
$\rho_x=\sum_{\mathsf{n}\ge0}P_\mu(\mathsf{n})\,\dyad{\mathsf{n},a}$ on the
two-mode Fock space. Let a tagged source be declared on it
(Definition~\ref{def:tagged-source}) with cutoffs
$0\le N_{\mathrm{ent}}\le N_{\mathrm{src}}$, virtual states $\rvirt_x$ of
\eqref{eq:rvirt} and declared weights $\tilde p(\mathsf{n}\,|\,\mu)$,
$\mathsf{n}\le N_{\mathrm{src}}$, satisfying the inequality
\eqref{eq:weightineq} of Definition~\ref{def:tagged-source}\,(T2), which
reads
\begin{equation}\label{eq:tag-ineq}
  0\le\tilde p(\mathsf{n}\,|\,\mu)\le P_\mu(\mathsf{n})
  \qquad\text{for every intensity }\mu\text{ and every }
  \mathsf{n}\le N_{\mathrm{src}} ,
\end{equation}
with tail weights
$\tilde w_\mu=1-\sum_{\mathsf{n}\le N_{\mathrm{src}}}\tilde p(\mathsf{n}\,|\,\mu)$.
Then:
\begin{enumerate}[label=(\roman*),leftmargin=2.2em]
\item \emph{(the tail states)} for every polarisation $a$ and every intensity
$\mu$ with $\tilde w_\mu>0$,
\begin{equation}\label{eq:tailstate}
  \vartheta_{a,\mu}\coloneqq\frac{1}{\tilde w_\mu}\Big[
  \sum_{\mathsf{n}\le N_{\mathrm{src}}}
  \big(P_\mu(\mathsf{n})-\tilde p(\mathsf{n}\,|\,\mu)\big)\dyad{\mathsf{n},a}
  +\sum_{\mathsf{n}>N_{\mathrm{src}}}P_\mu(\mathsf{n})\,\dyad{\mathsf{n},a}
  \Big]
\end{equation}
is a state on the Fock space. Where $\tilde w_\mu=0$ the declaration carries no
tail flag for $(a,\mu)$ (Definition~\ref{def:tagged-source}) and no tail state
is defined.
\item \emph{(one channel)} the map
\begin{equation}\label{eq:xitag}
\begin{aligned}
  \Xi_{\mathrm{tag}}(\zeta)\coloneqq{}&
  \sum_{\mathsf{n}\le N_{\mathrm{ent}}}
    V_{\mathsf{n}}\,\Pi^{\mathrm v}_{\mathsf{n}}\,\zeta\,
    \Pi^{\mathrm v}_{\mathsf{n}}V_{\mathsf{n}}^{\dagger}
  +\sum_{N_{\mathrm{ent}}<\mathsf{n}\le N_{\mathrm{src}}}\sum_{a}
    \bra{f_{\mathsf{n},a}}\zeta\ket{f_{\mathsf{n},a}}\,\dyad{\mathsf{n},a}\\
  &+\sum_{(a,\mu):\,\tilde w_\mu>0}
    \bra{f_{a,\mu}}\zeta\ket{f_{a,\mu}}\,\vartheta_{a,\mu} ,
\end{aligned}
\end{equation}
in which $\Pi^{\mathrm v}_{\mathsf{n}}$ projects $\Hs{\Avirt}$ onto the quantum
sector $S_{\mathsf{n}}$, $V_{\mathsf{n}}$ is the isometry that embeds that
sector into the Fock space as $\operatorname{span}\{\ket{\mathsf{n},a}\}_a$,
and every sum runs over the sectors of the declared index set $\mathbb{T}$
only, is a channel from $\Avirt$ to the Fock space, and
\begin{equation}\label{eq:tagid}
  \Xi_{\mathrm{tag}}(\rvirt_x)=\rho_x\qquad\text{for every label }x .
\end{equation}
A quantum sector or a polarisation flag that no letter puts weight on is not a
sector of $\Avirt$ by Definition~\ref{def:tagged-source}\,(T3) and contributes
no term to \eqref{eq:xitag}, exactly as for the tail flags.
The channel $\Xi_{\mathrm{tag}}$ depends on the declared weights and the
cutoffs and on no label.
\item \emph{(the role of the inequality)} \eqref{eq:tagid} holds for every
family of weights satisfying \eqref{eq:tag-ineq}. Where \eqref{eq:tag-ineq}
fails at an intensity $\mu$, its two halves fail differently.
\begin{enumerate}[label=(\alph*),leftmargin=2.2em]
\item If $\tilde p(\mathsf{n}\,|\,\mu)>P_\mu(\mathsf{n})$ for some
$\mathsf{n}\le N_{\mathrm{src}}$ and $\tilde w_\mu>0$, the operator
\eqref{eq:tailstate} is not positive semidefinite, for every polarisation $a$,
and the construction is not available: where a letter $x$ has intensity $\mu$,
the map \eqref{eq:xitag} sends the state $\dyad{f_{a_x,\mu}}$ to
$\vartheta_{a_x,\mu}$ and is not positive.
\item If $\tilde p(\mathsf{n}\,|\,\mu)<0$ for some
$\mathsf{n}\le N_{\mathrm{src}}$ while
$\tilde p(\mathsf{n}'\,|\,\mu)\le P_\mu(\mathsf{n}')$ for every
$\mathsf{n}'\le N_{\mathrm{src}}$, then $\tilde w_\mu>0$ and the operator
\eqref{eq:tailstate} is a state for every polarisation $a$. What fails is the
virtual state: for every letter $x$ of intensity $\mu$ the operator $\rvirt_x$
of \eqref{eq:rvirt} is not positive semidefinite or has trace greater than
one, so it is not a state.
\end{enumerate}
\item \emph{(the published case)} for
$\tilde p(\mathsf{n}\,|\,\mu)=P_\mu(\mathsf{n})$ and
$N_{\mathrm{ent}}=N_{\mathrm{src}}$ the virtual states are those of
\cite[Eq.~(145)]{tupkary2026decoy} and \eqref{eq:tagid} is
\cite[Lemma~9.2]{tupkary2026decoy}.
\end{enumerate}
\end{proposition}
\end{mainclaim}
\claimmeta{\metalabel{Uses}\mref{Def.}{def:source-link}, \mref{Def.}{def:wcp-model}, \mref{Def.}{def:tagged-source}
\metasep \metalabel{Used by}\mref{Cor.}{cor:emitted-decoy}, \mref{Req.}{req:sector-weight-bound}
\metasep newly derived
\metasep accepted}

\begin{proof}
\emph{Idea.} The channel keeps each quantum sector, prepares from each
polarisation flag the number state it stands for, and prepares from each tail
flag the normalised remainder of the Poisson distribution, which the
inequality makes a state.

(i) Every coefficient of \eqref{eq:tailstate} is non-negative, the first family
by \eqref{eq:tag-ineq}, and the vectors $\ket{\mathsf{n},a}$, $\mathsf{n}\ge0$,
are orthonormal for a fixed $a$, being number states of the one mode of
polarisation $a$ with the orthogonal mode in its vacuum. So the bracket is a
positive trace-class operator with trace
$\sum_{\mathsf{n}\ge0}P_\mu(\mathsf{n})
-\sum_{\mathsf{n}\le N_{\mathrm{src}}}\tilde p(\mathsf{n}\,|\,\mu)
=1-(1-\tilde w_\mu)=\tilde w_\mu$, and $\vartheta_{a,\mu}$ has unit trace.

(ii) \emph{A channel.} The sector projectors $\Pi^{\mathrm v}_{\mathsf{n}}$,
$\dyad{f_{\mathsf{n},a}}$ and $\dyad{f_{a,\mu}}$ are mutually orthogonal and sum
to $\idon{\Avirt}$ (Definition~\ref{def:tagged-source}). Each of the three sums
of \eqref{eq:xitag} is completely positive: the first is a sum of conjugations
by $V_{\mathsf{n}}\Pi^{\mathrm v}_{\mathsf{n}}$; the second and third are
measure-and-prepare maps, a rank-one projective measurement followed by the
preparation of a fixed state, the states $\vartheta_{a,\mu}$ being states by
(i). For the trace, $V_{\mathsf{n}}$ is an isometry on $S_{\mathsf{n}}$ and
every prepared operator has unit trace, so
$\tr\Xi_{\mathrm{tag}}(\zeta)=\sum_{\mathsf{n}\le N_{\mathrm{ent}}}
\tr[\Pi^{\mathrm v}_{\mathsf{n}}\zeta]
+\sum_{\mathsf{n},a}\bra{f_{\mathsf{n},a}}\zeta\ket{f_{\mathsf{n},a}}
+\sum_{(a,\mu):\,\tilde w_\mu>0}\bra{f_{a,\mu}}\zeta\ket{f_{a,\mu}}
=\tr\zeta$, the projectors summed being all the sector projectors of the
declaration. The input is finite dimensional and every output operator is trace
class, so $\Xi_{\mathrm{tag}}$ is a completely positive trace-preserving map
into the trace class of the Fock space.

\emph{The identity.} Fix $x=(a,\mu,z)$. By Definition~\ref{def:tagged-source},
\begin{equation*}
  \rvirt_x=\sum_{\mathsf{n}\le N_{\mathrm{ent}}}
  \tilde p(\mathsf{n}\,|\,\mu)\,\dyad{\mathsf{n},a}_{S_{\mathsf{n}}}
  +\sum_{N_{\mathrm{ent}}<\mathsf{n}\le N_{\mathrm{src}}}
  \tilde p(\mathsf{n}\,|\,\mu)\,\dyad{f_{\mathsf{n},a}}
  +\tilde w_\mu\,\dyad{f_{a,\mu}} ,
\end{equation*}
which is block diagonal along the sectors. The first sum of \eqref{eq:xitag}
maps its first sum to $\sum_{\mathsf{n}\le N_{\mathrm{ent}}}
\tilde p(\mathsf{n}\,|\,\mu)\dyad{\mathsf{n},a}$, the embedding carrying the
vector $\ket{\mathsf{n},a}$ of $S_{\mathsf{n}}$ to the Fock vector of the same
name. The second maps the flag $f_{\mathsf{n},a}$, which carries weight
$\tilde p(\mathsf{n}\,|\,\mu)$ in $\rvirt_x$ while the flags
$f_{\mathsf{n},a'}$ of the other polarisations $a'\ne a$ carry none, to
$\tilde p(\mathsf{n}\,|\,\mu)\dyad{\mathsf{n},a}$.
The third maps the one tail flag of $\rvirt_x$, where $\tilde w_\mu>0$, to
$\tilde w_\mu\vartheta_{a,\mu}$, which by \eqref{eq:tailstate} is the bracket
there. When $\tilde w_\mu=0$ there is no tail flag and $\rvirt_x$ has no third
term; then \eqref{eq:tag-ineq} and $\sum_{\mathsf{n}}P_\mu(\mathsf{n})=1$ force
$\tilde p(\mathsf{n}\,|\,\mu)=P_\mu(\mathsf{n})$ for
$\mathsf{n}\le N_{\mathrm{src}}$ and $P_\mu(\mathsf{n})=0$ beyond, so the
bracket vanishes as well. Adding the three,
\begin{equation*}
  \Xi_{\mathrm{tag}}(\rvirt_x)
  =\sum_{\mathsf{n}\le N_{\mathrm{src}}}\tilde p(\mathsf{n}\,|\,\mu)
   \dyad{\mathsf{n},a}
  +\sum_{\mathsf{n}\le N_{\mathrm{src}}}
   \big(P_\mu(\mathsf{n})-\tilde p(\mathsf{n}\,|\,\mu)\big)\dyad{\mathsf{n},a}
  +\sum_{\mathsf{n}>N_{\mathrm{src}}}P_\mu(\mathsf{n})\dyad{\mathsf{n},a}
  =\rho_x .
\end{equation*}
Nothing in \eqref{eq:xitag} refers to $x$: the state prepared from a flag, a
polarisation flag indexed by $(\mathsf{n},a)$ or a tail flag indexed by
$(a,\mu)$, is a function of the flag alone.

(iii) The first sentence is (ii), whose only hypothesis on the weights is
\eqref{eq:tag-ineq}.

(a) For every $a$ the vectors $\ket{\mathsf{n}',a}$, $\mathsf{n}'\ge0$, are
orthonormal (proof of (i)), so
$\bra{\mathsf{n},a}\vartheta_{a,\mu}\ket{\mathsf{n},a}
=\big(P_\mu(\mathsf{n})-\tilde p(\mathsf{n}\,|\,\mu)\big)/\tilde w_\mu<0$ and
\eqref{eq:tailstate} is not positive semidefinite. A letter $x$ of intensity
$\mu$ carries the tail flag $f_{a_x,\mu}$, which is a sector of $\mathbb{T}$ by
\eqref{eq:tset} because $\tilde w_\mu>0$. On $\dyad{f_{a_x,\mu}}$ the first two
sums of \eqref{eq:xitag} vanish, the sectors being orthogonal, and the third
gives $\vartheta_{a_x,\mu}$; so \eqref{eq:xitag} maps a positive operator to
one that is not positive semidefinite.

(b) Here
$\sum_{\mathsf{n}'\le N_{\mathrm{src}}}\tilde p(\mathsf{n}'\,|\,\mu)
<\sum_{\mathsf{n}'\le N_{\mathrm{src}}}P_\mu(\mathsf{n}')\le1$, the inequality
being strict at $\mathsf{n}'=\mathsf{n}$, where
$\tilde p(\mathsf{n}\,|\,\mu)<0\le P_\mu(\mathsf{n})$, so $\tilde w_\mu>0$; and
the proof of
(i), which reads only the upper half of \eqref{eq:tag-ineq}, makes
$\vartheta_{a,\mu}$ a state for every $a$. Let $x$ be a letter with
$\mu_x=\mu$, write $a=a_x$, and let $t$ be the sector that carries the term of
$\tilde p(\mathsf{n}\,|\,\mu)$ in \eqref{eq:rvirt}: $t=\mathsf{n}$ if
$\mathsf{n}\le N_{\mathrm{ent}}$ and $t=(\mathsf{n},a)$ otherwise. If
$t\in\mathbb{T}$, the component of $\rvirt_x$ in $S_t$ is
$\tilde p(\mathsf{n}\,|\,\mu)$ times the dyad of the unit vector
$\ket{\mathsf{n},a}$, resp.\ $f_{\mathsf{n},a}$, so the expectation of
$\rvirt_x$ in that vector is $\tilde p(\mathsf{n}\,|\,\mu)<0$ and $\rvirt_x$ is
not positive semidefinite. If $t\notin\mathbb{T}$, \eqref{eq:rvirt} has no term
in $t$. By \eqref{eq:tset} every sector on which $x$ puts a positive weight lies
in $\mathbb{T}$, the tail flag $f_{a,\mu}$ among them since $\tilde w_\mu>0$;
so every term that \eqref{eq:rvirt} omits for $x$ has a weight
$\tilde p(\mathsf{n}'\,|\,\mu)\le0$, the term at $\mathsf{n}$ among them, and by
\eqref{eq:tailweight}
\begin{equation*}
  \tr\rvirt_x
  =1-\sum_{\mathsf{n}'\ \text{omitted}}\tilde p(\mathsf{n}'\,|\,\mu)
  \;\ge\;1-\tilde p(\mathsf{n}\,|\,\mu)\;>\;1 .
\end{equation*}

(iv) With $\tilde p=P_\mu$ and $N_{\mathrm{ent}}=N_{\mathrm{src}}$ the middle
sum of $\rvirt_x$ is empty, the low sectors carry the exact Poisson weights on
$\dyad{\mathsf{n},a}$, and the tail flag carries
$1-\sum_{\mathsf{n}\le N_{\mathrm{src}}}P_\mu(\mathsf{n})$. This is the
displayed state of the cited equation with its photon cutoff equal to
$N_{\mathrm{src}}$ and its flag $\dyad{a,\mu}$ equal to $\dyad{f_{a,\mu}}$; the
cited lemma asserts a channel with \eqref{eq:tagid} for every choice of $a$,
$\mu$ and $z$.
\end{proof}

% claim transmitter/decoy/rem:no-intensity-split
\claimsummary{transmitter/decoy/rem:no-intensity-split}{A flag family split by intensity would still admit a label-independent tagging channel, but it would drop the relation between one photon number's behaviour at different intensities, which is the decoy content of the declaration. The emission therefore refuses it with a named refusal.}
\claimbookmark{transmitter/decoy/rem:no-intensity-split}{Remark}{rem:no-intensity-split}{Polarisation flags shared across intensities}
\begin{remark}[Polarisation flags shared across intensities]
\label{rem:no-intensity-split}
Definition~\ref{def:tagged-source}\,(T3) declares \emph{one} polarisation flag
$f_{\mathsf{n},a}$ per photon number $N_{\mathrm{ent}}<\mathsf{n}\le
N_{\mathrm{src}}$ and per polarisation, shared by every intensity and both
round types. A declarer might instead declare a finer family
$f_{\mathsf{n},a,\mu}$, one flag per intensity as well; this remark says what
that would and would not be.

\emph{The source identity would still hold.} The finer declaration admits a
channel of the shape Proposition~\ref{prop:tagging} constructs, and with it the
source identity, for this link \eqref{eq:tagid}: send $f_{\mathsf{n},a,\mu}$ to
$\dyad{\mathsf{n},a}$, which is independent of $\mu$, and leave the rest as it
is. So splitting is not excluded by the source identity and is not a security
error of the kind this submodule is about.

\emph{The decoy content would be lost.} With one shared flag, the letters
$(a,\mu,z)$ and $(a,\mu',z')$ of the same polarisation put their weight on the
same one-dimensional sector of $\Avirt$, so one sector of the declaration
carries the whole of that photon number's weight across intensities and round
types. With one flag per intensity they are different sectors, and nothing in
the declaration relates what happens in the sector of $\mu$ to what happens in
the sector of $\mu'$: a tagged photon number of one intensity and the same
photon number of another intensity become unrelated parts of the declaration.
That relation is the decoy content of the declaration, the reason for declaring
intensities at all being that the same photon number behaves the same way
whichever intensity produced it; so a split declaration is a declaration of a
different protocol, one with intensities and without decoy. The same applies to
the round type: the shared flag is why a tagged photon number of a test round
and of a generation round are one sector, which is not an assumption added here
but a consequence of declaring one flag.

\emph{The direction of the loss.} Splitting removes a relation from the
declaration and adds none, so a consumer of a split declaration is told less
about the source than a consumer of the shared one. The loss would not show as
an error but as a rate that does not improve when an intensity is added, which
is why the emission, the step that outputs the declaration, refuses a split
family with a named refusal (Requirement~\ref{req:shared-flags}). No claim of
this repository compares the key lengths of two declarations, and this remark
asserts no such comparison: it asserts which relation is present in which
declaration. The reverse, a coarser flag than
Definition~\ref{def:tagged-source}\,(T3) declares (for instance one flag for
two polarisations of the same photon number), is excluded for a stricter
reason: it is not a \termref{term-tagged-source-declaration}{tagged source declaration} at all. Which vectors the flags
are and which letters put weight on them is fixed by
Definition~\ref{def:tagged-source}\,(T4), and Proposition~\ref{prop:tagging}
proves the source identity for those declarations only.
\end{remark}
\claimmeta{\metalabel{Uses}\mref{Def.}{def:tagged-source}
\metasep newly derived
\metasep accepted}

\section{The certificate and the emitted declaration}\label{group-transmitter/decoy/prop:decoy-cert}
\sectionintro{The emission, the step of this submodule that outputs a tagged declaration as a decoy declaration, declares the weighted virtual blocks sector by sector and nothing on Alice's side. \hyperlink{claim-transmitter/decoy/prop:decoy-cert}{Proposition~\ref*{prop:decoy-cert}} shows that these blocks satisfy the positivity and weight conditions (E1) and (E2) of an \termref{term-ensemble-presentation}{ensemble presentation}, which gives an exact certificate, and computes the marginal and the numbers it forms; whether those numbers are binary64, as (E3) requires, is decided on each declaration and not there. The main result of the section, \hyperlink{claim-transmitter/decoy/cor:emitted-decoy}{Corollary~\ref*{cor:emitted-decoy}}, states what the emitted declaration satisfies: an exact certificate against the tagged ensemble where its blocks meet (E1) to (E3), the source identity at the emitted descriptors as a conclusion of \hyperlink{claim-transmitter/decoy/prop:tagging}{Proposition~\ref*{prop:tagging}}, block-diagonal states after any attack, the entries on which exactness is decided, and that lowering a weight gives another tagged declaration. \hyperlink{claim-transmitter/decoy/rem:decoy-consumers}{Remark~\ref*{rem:decoy-consumers}} says how the consuming modules receive it.}

% claim transmitter/decoy/prop:decoy-cert
\claimsummary{transmitter/decoy/prop:decoy-cert}{Read sector by sector, the tagged virtual ensemble meets (E1) and (E2) of an ensemble presentation, giving an exact certificate with pure blocks whose sector marginals are transposed sums of signal dyads. Whether the formed numbers are binary64, as (E3) requires, is not decided here.}
\claimbookmark{transmitter/decoy/prop:decoy-cert}{Proposition}{prop:decoy-cert}{Ensemble presentation of the tagged source}
\claimleadin{%
The certificate of a tagged source is never declared: by the
\termref{term-ensemble-presentation}{ensemble presentation} of Definition~\ref{def:srcert} it suffices to
declare the weighted \termref{term-virtual-ensemble}{virtual ensemble} sector by sector. A presentation
has three conditions: (E1) every block is Hermitian and positive semidefinite;
(E2) the traces of a letter's blocks sum to $p_x$, and the $p_x$ to one; (E3)
for every sector the sum of the blocks over the letters, accumulated entry by
entry in the \exttermref{../foundations/contract.pdf}{term-declared-order}{declared order} of $\Xset$, \exttermref{../foundations/contract.pdf}{term-rounds-nowhere}{rounds nowhere}, so that
its entries are \exttermref{../foundations/contract.pdf}{term-binary64-numbers}{binary64 numbers}. The
certificate it defines satisfies (SR1), its partial trace over the virtual
system is Alice's marginal $\tau_A$, and (SR2), Alice's outcome $x$ leaves
$p_x\rvirt_x$ on the virtual system; (SR3), that Alice retains her system, is
an attestation. The sectors are those of
Definition~\ref{def:tagged-source}\,(T3), the ones that some letter weights: a
quantum sector $t=\mathsf{n}$, with
$S_{\mathsf{n}}=\operatorname{span}\{\ket{\mathsf{n},a}:a\in\mathrm{Pol}\}$, for
each $\mathsf{n}\le N_{\mathrm{ent}}$ (the \termref{term-entropy-cutoff}{entropy cutoff}); a
\termref{term-polarisation-flag}{polarisation flag} $t=(\mathsf{n},a)$ for
$N_{\mathrm{ent}}<\mathsf{n}\le N_{\mathrm{src}}$ (the \termref{term-source-cutoff}{source cutoff}); and
a \termref{term-tail-flag}{tail flag} $t=(a,\mu)$ for each pair a letter carries with
$\tilde w_\mu>0$. The proposition checks (E1) and (E2), which every tagged
declaration meets, computes the sector marginals and their ranks (a rank below
$\dim S_t$ is where the completion term $\idon{A_t}-P_t$ of \eqref{eq:enspres}
is non-zero), and records which numbers the presentation forms. Whether those
products and sums are binary64 numbers, as the number model and (E3) require,
is a property of a particular declaration, decided on it and not asserted here
(Corollary~\ref{cor:emitted-decoy}\,(iv)).
}
\begin{technicalclaim}
\begin{proposition}[Ensemble presentation of the tagged source]
\label{prop:decoy-cert}
Let a tagged source be declared on a weak-coherent-pulse \termref{term-source-model}{source model}
(Definitions~\ref{def:wcp-model} and~\ref{def:tagged-source}), with labels
$x=(a_x,\mu_x,z_x)$, $z_x$ the \termref{term-round-type}{round type}, of weights $p_x>0$,
\termref{term-sector-index-set}{sector index set} $\mathbb{T}$ and
virtual sectors $S_t\subseteq\Hs{\Avirt}$. For $t\in\mathbb{T}$ and a label $x$
write $w_{x,t}\ge0$ and $\ket{\varphi_{x,t}}\in\Hs{S_t}$ for the weight and the
unit vector that $\rvirt_x$ carries in sector $t$:
$w_{x,t}=\tilde p(\mathsf{n}\,|\,\mu_x)$ and
$\ket{\varphi_{x,t}}=\ket{\mathsf{n},a_x}$ on a quantum sector $t=\mathsf{n}$;
$w_{x,t}=\tilde p(\mathsf{n}\,|\,\mu_x)$ if $a_x=a$ and $0$ otherwise, with
$\ket{\varphi_{x,t}}=\ket{f_{\mathsf{n},a}}$, on a polarisation flag
$t=(\mathsf{n},a)$; and $w_{x,t}=\tilde w_{\mu}$ if $(a_x,\mu_x)=(a,\mu)$ and
$0$ otherwise, with $\ket{\varphi_{x,t}}=\ket{f_{a,\mu}}$, on a tail flag
$t=(a,\mu)$. So $\rvirt_x=\sum_tw_{x,t}\dyad{\varphi_{x,t}}$ and
$\sum_tw_{x,t}=1$. Put $\Xset_t\coloneqq\{x:w_{x,t}>0\}$, let the
\termref{term-sector-structure}{sector structure} carry exactly the sectors with $\Xset_t\ne\varnothing$, and let
\begin{equation}\label{eq:decoy-blocks}
  \sigma^{\mathrm v,(t)}_x\;\coloneqq\;p_x\,w_{x,t}\,\dyad{\varphi_{x,t}}
  \;\in\;\Pos{S_t},\qquad x\in\Xset,\ t\in\mathbb{T} ,
\end{equation}
the block of the \termref{term-weighted-members}{weighted member} $p_x\rvirt_x$ in
sector $t$.
\begin{enumerate}[label=(\roman*),leftmargin=2.2em]
\item \emph{(an ensemble presentation)} the block lists
\eqref{eq:decoy-blocks} satisfy conditions (E1) and (E2) of
Definition~\ref{def:srcert} as identities of real numbers, with the weights
$p_x$ of Definition~\ref{def:wcp-model}. Hence, by
Lemma~\ref{lem:ensemble-presentation}, the objects \eqref{eq:enspres} formed
from them --- a retained system $A=\bigoplus_tA_t$ with $A_t$ a copy of $S_t$,
the marginal $\tau_A$, a POVM $\{M_x\}$ and blocks $(\Psi_t)_t$ --- are a
block-form \termref{term-source-replacement-certificate}{source-replacement certificate}, (SR1) and (SR2) holding exactly
against the tagged virtual ensemble $\{p_x,\rvirt_x\}$ of
Definition~\ref{def:tagged-source}\,(T4) itself, and every non-zero $\Psi_t$ is
a multiple of a rank-one projector.
\item \emph{(the marginal)} in the declared basis of $S_t$, copied to $A_t$,
\begin{equation}\label{eq:decoy-marginal}
  \tau^{(t)}_A
  =\Big(\sum_{x\in\Xset_t}p_xw_{x,t}\,\dyad{\varphi_{x,t}}\Big)^{\transpose} ,
  \qquad
  g_t\coloneqq\tr\tau^{(t)}_A=\sum_{x\in\Xset_t}p_xw_{x,t} ,
\end{equation}
$g_t$ being the probability of sector $t$ and $\sum_tg_t=1$. On the vacuum
sector and on every flag sector $A_t$ is one dimensional and $\tau^{(t)}_A$ is
the number $g_t$.
\item \emph{(ranks)} $r_t\coloneqq\operatorname{rank}\tau^{(t)}_A
=\dim\operatorname{span}\{\ket{\varphi_{x,t}}:x\in\Xset_t\}$. Hence $r_t=1$ on
every polarisation flag, on every tail flag and on the vacuum sector
$\mathsf{n}=0$, and on a quantum sector $\mathsf{n}\ge1$
\begin{equation}\label{eq:rankbound}
  r_{\mathsf{n}}=\dim\operatorname{span}\{\ket{\mathsf{n},a_x}:x\in\Xset_{\mathsf{n}}\}
  \le\min\big(\dim\Hs{S_{\mathsf{n}}},\ |\{a_x:x\in\Xset_{\mathsf{n}}\}|\big) ,
\end{equation}
while $\dim\Hs{A_{\mathsf{n}}}=\dim\Hs{S_{\mathsf{n}}}$ whatever the rank.
\item \emph{(the formed numbers)} every entry of \eqref{eq:decoy-blocks} is a
product of the declared numbers $p_x$ and $w_{x,t}$ with an entry of
$\dyad{\varphi_{x,t}}$ in the declared basis, and every entry of
$\tau^{(t)}_A$ is a sum of such products: no square root and no quotient of a
declared number occurs in anything a declaration of this presentation carries.
\end{enumerate}
\end{proposition}
\end{technicalclaim}
\claimmeta{\metalabel{Uses}\mref{Def.}{def:srcert}, \mref{Lem.}{lem:ensemble-presentation}, \mref{Def.}{def:wcp-model}, \mref{Def.}{def:tagged-source}
\metasep \metalabel{Used by}\mref{Cor.}{cor:emitted-decoy}
\metasep newly derived
\metasep implemented
\metasep \metalabel{Code}\coderef{https://github.com/fraunhofersg/certiqs-key/blob/v0.1.0/src/certiqs_key/transmitter/decoy/emit.py}{certiqs_key.transmitter.decoy.emit:DecoySourceDeclaration}}

\begin{proof}
(i) (E1): $\sigma^{\mathrm v,(t)}_x$ is a non-negative multiple of a dyad, hence
Hermitian and positive semidefinite. (E2):
$\sum_t\tr\sigma^{\mathrm v,(t)}_x=p_x\sum_tw_{x,t}=p_x$, the
$\ket{\varphi_{x,t}}$ being unit vectors, and $\sum_xp_x=1$ by
Definition~\ref{def:wcp-model}; a sector with $\Xset_t=\varnothing$ carries the
zero block for every letter, so leaving it out of the structure changes neither
sum. That the blocks are those of $p_x\rvirt_x$ is the decomposition
$\rvirt_x=\sum_tw_{x,t}\dyad{\varphi_{x,t}}$, which is \eqref{eq:rvirt} read
sector by sector, the sectors being mutually orthogonal. The hypotheses of
Lemma~\ref{lem:ensemble-presentation} are (E1) and the total trace
$\sum_x\sum_t\tr\sigma^{\mathrm v,(t)}_x=\sum_xp_x=1$, so its parts (i) and
(ii) give the statement; the virtual ensemble of its part (i) has the weighted
members $\bigoplus_t\sigma^{\mathrm v,(t)}_x=p_x\rvirt_x$.

(ii) is \eqref{eq:enspres} with the blocks \eqref{eq:decoy-blocks}, the terms
with $w_{x,t}=0$ vanishing; the trace of a transpose is the trace, and
$\sum_tg_t=\sum_xp_x\sum_tw_{x,t}=1$. The vacuum sector is spanned by the one
vector $\ket{0,a}$, which does not depend on $a$, and a flag sector by its flag
vector, so there $\dim\Hs{S_t}=1$.

(iii) Transposition preserves the rank, and the range of
$\sum_{x\in\Xset_t}p_xw_{x,t}\dyad{\varphi_{x,t}}$, a sum of positive multiples
of dyads, is the span of the vectors $\ket{\varphi_{x,t}}$, $x\in\Xset_t$. On a
flag sector and on the vacuum that span is one dimensional. On a quantum sector
it is spanned by the vectors $\ket{\mathsf{n},a}$ with $a$ ranging over the
polarisations of $\Xset_{\mathsf{n}}$ and lies in $\Hs{S_{\mathsf{n}}}$, which
gives both bounds. $A_{\mathsf{n}}$ is a copy of $S_{\mathsf{n}}$ by
construction.

(iv) is read off \eqref{eq:decoy-blocks} and \eqref{eq:decoy-marginal}.
\end{proof}

% claim transmitter/decoy/cor:emitted-decoy
\claimsummary{transmitter/decoy/cor:emitted-decoy}{The emitted decoy declaration has an exact certificate against the tagged ensemble, its source identity by the tagging channel, and block-diagonal attacked states. Exactness is decided on products of declared weights and formed dyad entries; lowering a weight gives another tagged declaration.}
\claimbookmark{transmitter/decoy/cor:emitted-decoy}{Corollary}{cor:emitted-decoy}{The emitted decoy declaration}
\claimleadin{%
The emission is the step of this submodule that outputs a
\termref{term-tagged-source-declaration}{tagged source declaration} as a decoy declaration; the emitted
declaration is that output, a record and not a pulse. It declares the tagged
virtual ensemble itself, sector by sector, and defines everything on Alice's
side from it. The corollary collects what such a declaration satisfies: a
certificate and a \termref{term-source-specification}{source specification} (i), the source identity at the emitted descriptors and the
admission of its segment (ii), what the certificate reaches (iii), the entries
on which exactness is decided for a particular declaration (iv), and the effect
of lowering a weight (v), which gives a different declaration and not a
perturbation of the first one. The inequality on the weights is decided by
the emission against the rational lower bound on $P_\mu(\mathsf{n})$ of
Lemma~\ref{lem:poisson-enclosure} (Requirement~\ref{req:sector-weight-bound}),
and the entries of the signal dyads are formed by the construction of (iv)(b)
(Requirement~\ref{req:signal-dyads}). A-SEQ and A-SRC enter through
Theorem~\ref{thm:certificate-sufficiency} and Lemma~\ref{lem:matched-pinching}, in (i) and (iii).
A-WCP is the physical-source assumption that a specification declared by this
source model carries (Definition~\ref{def:source-spec}, case (b)) and that the
\termref{term-obligation-record}{obligation record} of (ii) names; it is no hypothesis of
Proposition~\ref{prop:tagging}. The artefact the emission outputs lists A-WCP
in its assumption ledger (Requirement~\ref{req:decoy-emission}), so the
admission of its segment in (ii) applies there. Part (iv) matters because
the emission requires every formed entry to be a binary64 number and refuses one that is not
(Requirement~\ref{req:exact-decoy-blocks}). Remark~\ref{rem:decoy-consumers}
says what the consuming modules do with the declaration.
}
\begin{mainclaim}
\begin{corollary}[The emitted decoy declaration]\label{cor:emitted-decoy}
Let a \termref{term-weak-coherent-pulse-model}{weak-coherent-pulse model} (Definition~\ref{def:wcp-model}) and a tagged
source declaration over it (Definition~\ref{def:tagged-source}) be presented,
with the certified inequality \eqref{eq:weightineq},
$0\le\tilde p(\mathsf{n}\mid\mu)\le P_\mu(\mathsf{n})$, discharged for every
$\mathsf{n}\le N_{\mathrm{src}}$ and every intensity $\mu\in\mathrm{Int}$. Its
\emph{decoy declaration} is the ensemble presentation, in the sense of
Definition~\ref{def:srcert}, whose declared data are exactly
\begin{itemize}[leftmargin=1.6em]
\item the alphabet $\Xset$ with its declared order and the weights $p_x$;
\item the matched sector structure, matched meaning that Alice's system and
  $\Avirt$ carry one ordered index set: the sectors $S_t$ of
  Definition~\ref{def:tagged-source}\,(T3) that some letter puts positive
  weight on, in the declared order of $\mathbb{T}$, each with its orthonormal
  basis of (T3);
\item the weighted blocks $\sigma^{\mathrm v,(t)}_x=p_xw_{x,t}
  \dyad{\varphi_{x,t}}$ of \eqref{eq:decoy-blocks}, as block lists;
\item the specified ensemble $\{p_x,\rho_x\}$ of the model as the source model,
  the \termref{term-source-segment}{source segment} with the one tagging link of
  Definition~\ref{def:tagged-source}\,(T6) and its obligation record, and the
  assumption ledger;
\end{itemize}
so that no retained system, no measurement, no marginal and no certificate is
among them. Below, \emph{emitted} refers to this decoy declaration: the emitted
\termref{term-source-descriptor}{source descriptor}, segment and assumption ledger are the ones it
declares. Then
the following hold.
\begin{enumerate}[label=(\roman*),leftmargin=3.0em]
\item \emph{(a certificate, and an admissible specification)} where (E1) to
(E3) of Definition~\ref{def:srcert} hold of the declared blocks, the objects
\eqref{eq:enspres} --- the retained system $A=\bigoplus_tA_t$ with $A_t$ a copy
of $S_t$, the marginal $\tau_A$, the POVM $\{M_x\}$ and the blocks
$(\Psi_t)_t$ --- are a source-replacement certificate in \termref{term-block-form}{block form}
with every non-zero block a multiple of a rank-one projector, satisfying (SR1)
and (SR2) exactly against the tagged virtual ensemble of
Definition~\ref{def:tagged-source}\,(T4); the data are a source specification
in the sense of Definition~\ref{def:source-spec} in
\termref{term-trusted-source-mode}{trusted-source mode}; and
Theorem~\ref{thm:certificate-sufficiency} holds of it, its part (i) reading (E1) to (E3) as the
decided conditions.
\item \emph{(the segment, and the source identity at the emitted descriptors)}
\begin{enumerate}[label=(\alph*),leftmargin=2.2em]
\item \emph{(the source identity)} The emitted source descriptor is
$(\Avirt,\{\sigma^{\mathrm v}_x\})$ with $\sigma^{\mathrm v}_x=p_x\rvirt_x$ for
the tagged ensemble of Definition~\ref{def:tagged-source}\,(T4) with the
declared weights, hence a member of the \termref{term-declared-descriptor-family}{declared descriptor family} of
Proposition~\ref{prop:tagging}, every tagged source declaration over every
weak-coherent-pulse model; so the source identity \eqref{eq:srcmap}, here
$\Xi_{\mathrm{tag}}(\rvirt_x)=\rho_x$, holds exactly at the emitted
descriptors, as that proposition's conclusion.
\item \emph{(the segment record)} The emitted segment satisfies the
\termref{term-interface-conditions}{interface conditions} (S1) (descriptor match), (S2) (index preservation)
and (S4) (slots) of Definition~\ref{def:source-segment}, and its (S3) (source
identity) is the conclusion of (a). Its list of
\termref{term-transported-constraint-entries}{transported constraint entries} is empty and it carries no licence
record. Its obligation record, the triple a consumer compares with the
\termref{term-soundness-registry}{soundness registry} to admit the link, is
$\big(\textsf{tagging},\ \text{the id of Proposition~\ref{prop:tagging}},\
\{\mathrm{A\textnormal{-}WCP}\}\big)$, together with the member of the declared
descriptor family that its descriptors are
(Definition~\ref{def:source-link}\,(c2)). The third component is the union of
the assumption closure of Proposition~\ref{prop:tagging}, which is empty, and
the physical-source assumption A-WCP that this submodule declares for the model
that proposition is stated over (Definition~\ref{def:source-segment}\,(A1)).
Inside an artefact that this submodule emits and whose assumption ledger lists
A-WCP, the segment meets the admission conditions (A1) to (A3) of
Definition~\ref{def:source-segment}: the registry entry, the decided triple
and the provenance.
\end{enumerate}
\item \emph{(what the certificate reaches)} every state reached from the
certificate of (i) by an attack on the
\termref{term-virtual-transmitted-system}{transmitted system} and the adversary's
registers is block diagonal along $\mathbb{T}$ with $A$-marginal $\tau_A$, and
pinching $A$ along $\mathbb{T}$ after the attack is the same operation as
pinching the virtual signal before it (Lemma~\ref{lem:matched-pinching});
Theorem~\ref{thm:certificate-sufficiency}\,(iii) holds against the tagged ensemble of
Definition~\ref{def:tagged-source}\,(T4) itself;
and a positive operator constrained only to have the marginal $\tau_A$ need not
be block diagonal.
\item \emph{(the entries of the declared blocks)} (E1) to (E3) are conditions on
the declared blocks, and every entry of a block is a product
$p_x\,w_{x,t}\,\big(\dyad{\varphi_{x,t}}\big)_{ij}$ of the declared weight
$p_x$, the weight $w_{x,t}$, which is a declared sector weight or a
\termref{term-tail-weight}{tail weight} formed from them, and an entry of the dyad of $\ket{\varphi_{x,t}}$ in the
orthonormal basis of $S_t$ of Definition~\ref{def:tagged-source}\,(T3), a
number fixed by the model and that basis and not a declared number.
\begin{enumerate}[label=(\alph*),leftmargin=2.2em]
\item \emph{(the weights)} The products $p_xw_{x,t}$ are formed from the
declared $p_x$, the declared $\tilde p(\mathsf{n}\mid\mu)$ and the
\termref{term-tail-weight}{tail weight} \eqref{eq:tailweight}, which is itself formed exactly and rounds nowhere.
Choosing the $p_x$ and the $\tilde p(\mathsf{n}\mid\mu)$ with few enough
significant bits makes every product exact, and lowering a $\tilde p$ to a
value that is still non-negative is permitted by \eqref{eq:weightineq} and
yields again a tagged source declaration by (v).
\item \emph{(the signal dyads)} The entries of $\dyad{\varphi_{x,t}}$ are the
number $1$ on every flag sector and on the vacuum, and on a quantum sector
$\mathsf{n}\ge1$ the entries of $\dyad{\mathsf{n},a}$ in the basis of
$S_{\mathsf{n}}$. For every unit vector $\ket{b}\in\mathbb{C}^2$, in
$\mathrm{Pol}$ or not, write
$\ket{\mathsf{n},b}\coloneqq(c^{\dagger}_b)^{\mathsf{n}}\ket{\vac}/
\sqrt{\mathsf{n}!}$, which for $b\in\mathrm{Pol}$ is the vector of
\eqref{eq:numberstate}; then
$\big|\braket{\mathsf{n},a}{\mathsf{n},b}\big|^2=|\braket{a}{b}|^{2\mathsf{n}}$
for all unit vectors $a,b$. Let every polarisation be one of the four
vectors of $\Dset$, with $\ket{H},\ket{V}$ the fixed orthonormal basis of the
one-photon space of Definition~\ref{def:wcp-model}\,(W1) and
$\ket{\pm}=(\ket{H}\pm\ket{V})/\sqrt2$, distinct polarisations being distinct
vectors, and put $v_H=(1,1,0)$, $v_V=(1,-1,0)$ and $v_{\pm}=(1,0,\pm1)$. If
$|\mathrm{Pol}|\ge2$, then $S_1$ is the one-photon space,
$\{\ket{1,H},\ket{1,V}\}$ is an orthonormal basis of it, and in that basis
$\dyad{1,H}$, $\dyad{1,V}$ and $\dyad{1,\pm}$ are
$\big(\begin{smallmatrix}1&0\\0&0\end{smallmatrix}\big)$,
$\big(\begin{smallmatrix}0&0\\0&1\end{smallmatrix}\big)$ and
$\tfrac12\big(\begin{smallmatrix}1&\pm1\\\pm1&1\end{smallmatrix}\big)$. If
$|\mathrm{Pol}|\ge3$, then $S_2$ is the two-photon space, and it has an
orthonormal basis in which $\dyad{2,a}=\tfrac12\,v_av_a^{\transpose}$ for every
$a\in\Dset$. So every entry of these dyads is $0$, $1$ or $\pm\tfrac12$ at one
photon and $0$ or $\pm\tfrac12$ at two.
\item \emph{(the sector sums)} With the entries of (a) and (b) representable,
(E3) is a condition on the sums $\sum_x\sigma^{\mathrm v,(t)}_x$, accumulated
entry by entry in the declared order, and on nothing else.
\end{enumerate}
\item \emph{(lowered weights)} if
$0\le\tilde p'(\mathsf{n}\mid\mu)\le\tilde p(\mathsf{n}\mid\mu)$ for every
$\mathsf{n}\le N_{\mathrm{src}}$ and every $\mu$, then the weights $\tilde p'$
are again a tagged source declaration over the same model and
Proposition~\ref{prop:tagging} holds of it, the difference passing into the
tail weight. Its weights, its virtual ensemble, its tail states and its sector
list are those formed from $\tilde p'$: the sector list may gain a tail flag,
where $\tilde w_{\mu}$ was zero and $\tilde w'_{\mu}$ is not, and it may lose a
polarisation flag or a quantum sector, where a weight is lowered to zero at
every letter that carried it, so that neither sector list need contain the
other.
\end{enumerate}
\end{corollary}
\end{mainclaim}
\claimmeta{\metalabel{Uses}\mref{Def.}{def:source-spec}, \mref{Def.}{def:source-link}, \mref{Def.}{def:source-segment}, \mref{Def.}{def:srcert}, \mref{Lem.}{lem:ensemble-presentation}, \mref{Thm.}{thm:certificate-sufficiency}, \mref{Lem.}{lem:matched-pinching}, \mref{Def.}{def:wcp-model}, \mref{Def.}{def:tagged-source}, \mref{Prop.}{prop:tagging}, \mref{Prop.}{prop:decoy-cert}
\metasep \metalabel{Assumes}\aref{A-WCP}{ass:wcp}, \aref{A-SRC}{ass:src}, \aref{A-SEQ}{ass:seq}
\metasep \metalabel{Used by}\mref{Rem.}{rem:decoy-consumers}, \mref{Req.}{req:decoy-emission}, \mref{Req.}{req:entropy-cutoff}, \mref{Req.}{req:signal-dyads}, \mref{Req.}{req:exact-decoy-blocks}, \mref{Req.}{req:decoy-ensemble-conditions}
\metasep newly derived
\metasep implemented (effective: accepted)
\metasep \metalabel{Code}\coderef{https://github.com/fraunhofersg/certiqs-key/blob/v0.1.0/src/certiqs_key/transmitter/decoy/emit.py}{certiqs_key.transmitter.decoy.emit:emit_decoy_declaration}}

\begin{proof}
(i) Proposition~\ref{prop:decoy-cert}\,(i) is exactly this statement for the
blocks \eqref{eq:decoy-blocks}: (E1) and (E2) hold as identities of real
numbers, and Lemma~\ref{lem:ensemble-presentation} makes the objects
\eqref{eq:enspres} a block-form certificate with pure blocks against the
weighted members $\bigoplus_t\sigma^{\mathrm v,(t)}_x=p_x\rvirt_x$, which are
those of Definition~\ref{def:tagged-source}\,(T4). (E3) is the one condition of the presentation that is not an
identity of real numbers but a property of the declared numbers, and (iv) says
what it is a condition on. That the data are a source specification is
Definition~\ref{def:source-spec}: the alphabet and its order are declared, the
specified ensemble is declared by a source model, which is that definition's
case (b) and comes with a non-empty source segment, and in trusted-source mode
with more than one sector an ensemble presentation is one of the two forms it
admits, under which $A$, $\{M_x\}$ and $\tau_A$ are not declared but defined,
the conditions on them holding by
Lemma~\ref{lem:ensemble-presentation}\,(i). The dimension of Alice's system is
$d_A=\sum_t\dim\Hs{S_t}\ge2$ whenever the structure carries two sectors, and
Definition~\ref{def:source-spec} demands a certificate in trusted-source mode
at $d_A\ge2$; so the certificate is demanded, and it is supplied. Theorem~\ref{thm:certificate-sufficiency} applies to a certificate in either form, and
its part (i) says what an ensemble presentation decides.

(ii)(a) By Definition~\ref{def:tagged-source}\,(T4) the emitted weighted
members are $p_x\rvirt_x$ with the declared
$\tilde p(\mathsf{n}\mid\mu)$ and $\tilde w_{\mu}$, which is a tagged source
declaration over the given model: it is the one presented, no number having
been changed. The declared descriptor family of the link is every tagged source
declaration over every weak-coherent-pulse model
(Definition~\ref{def:tagged-source}\,(T6)), so the emitted descriptor is a
member of it and Proposition~\ref{prop:tagging}\,(ii) concludes
$\Xi_{\mathrm{tag}}(\rvirt_x)=\rho_x$ for every letter, which is
\eqref{eq:srcmap} for the one-link composite.

(ii)(b) (S1) holds by Definition~\ref{def:tagged-source}\,(T6): the source
descriptor is the virtual pair and the \termref{term-target-descriptor}{target descriptor} is the specified pair
of the model. (S2) holds because both are indexed by $\Xset$ and
$\tr\sigma^{\mathrm v}_x=p_x\sum_tw_{x,t}=p_x$ by
Proposition~\ref{prop:decoy-cert}\,(i), the same $p_x$ the model declares. (S4)
holds because the link carries its empty entry list, its empty licence slot,
its status and its obligation record. The triple of that record is the registry
entry of the kind \textsf{tagging} (Definition~\ref{def:source-segment}\,(A1)),
which is what a consumer decides against. The third component is as that
definition specifies it. Its first half, the assumption closure of
Proposition~\ref{prop:tagging}, is empty by the criterion of (A1)(a), no claim
of that closure listing an assumption among its dependencies: the proposition
lists none and uses Definitions~\ref{def:wcp-model}, \ref{def:tagged-source}
and~\ref{def:source-link}; what these use, directly or in turn, lies among
them and Definitions~\ref{def:source-segment} and~\ref{def:source-spec}; and
none of these definitions lists an assumption. Its second half is the
physical-source assumption this submodule declares for the model, A-WCP; so
the union is $\{\mathrm{A\textnormal{-}WCP}\}$. A-WCP is not a
hypothesis of Proposition~\ref{prop:tagging}, and nothing here makes it one.
For (A1) to (A3), let the segment be carried by an artefact that this
submodule emits and whose assumption ledger lists A-WCP. The registry entry of
the kind \textsf{tagging} is the triple above, which gives (A1) and the
comparison of triples of (A2); the one assumption of that entry, A-WCP, is in
the artefact's ledger by that hypothesis, which is the rest of (A2); and the
artefact is emitted by this submodule, the module that admits the kind, which
is (A3).

(iii) The certificate of (i) is in block form against the matched pair of
Definition~\ref{def:tagged-source}\,(T5), which is the hypothesis of
Lemma~\ref{lem:matched-pinching}, and the attacks quantified over are those of
Theorem~\ref{thm:certificate-sufficiency}\,(iii), whose conclusion is stated against the virtual
ensemble the certificate is written against: the tagged one, by (i). The
operator of the last clause is the one exhibited in the proof of that lemma.

(iv) That every entry of a block is such a product is
Proposition~\ref{prop:decoy-cert}\,(iv). (a) is the definition of $w_{x,t}$ in
that proposition and \eqref{eq:tailweight}; that lowering a $\tilde p$ leaves a
tagged declaration is (v). (b) On a flag sector and on the
vacuum $S_t$ is one dimensional (Proposition~\ref{prop:decoy-cert}\,(ii)) and
the entry is $\braket{\varphi_{x,t}}{\varphi_{x,t}}=1$.

\emph{The overlaps.} For unit vectors $\ket{a},\ket{b}\in\mathbb{C}^2$, in
$\mathrm{Pol}$ or not, write $\ket{a}=\sum_ia_i\ket{i}$ in the fixed basis of
(W1) and $c^{\dagger}_a=\sum_ia_ic^{\dagger}_i$ for the mode operators $c_i$
of that basis, so that $[c_a,c^{\dagger}_b]=\sum_i\bar a_ib_i
=\braket{a}{b}$ and $c_a\ket{\vac}=0$. The commutator is a number, so
$[c_a,(c^{\dagger}_b)^{\mathsf{n}}]=\mathsf{n}\braket{a}{b}
(c^{\dagger}_b)^{\mathsf{n}-1}$ by induction on $\mathsf{n}$, and hence
$c_a(c^{\dagger}_b)^{\mathsf{n}}\ket{\vac}=\mathsf{n}\braket{a}{b}
(c^{\dagger}_b)^{\mathsf{n}-1}\ket{\vac}$. Applying this $\mathsf{n}$ times,
$\bra{\vac}c_a^{\mathsf{n}}(c^{\dagger}_b)^{\mathsf{n}}\ket{\vac}
=\mathsf{n}!\,\braket{a}{b}^{\mathsf{n}}$, and the definition of
$\ket{\mathsf{n},b}$ in (b), which is \eqref{eq:numberstate} for
$b\in\mathrm{Pol}$, gives
$\braket{\mathsf{n},a}{\mathsf{n},b}=\braket{a}{b}^{\mathsf{n}}$, whence
$\big|\braket{\mathsf{n},a}{\mathsf{n},b}\big|^2=|\braket{a}{b}|^{2\mathsf{n}}$.

\emph{One photon.} The map $\ket{a}\mapsto c^{\dagger}_a\ket{\vac}=\ket{1,a}$
is linear and isometric from $\mathbb{C}^2$ onto the one-photon space, by the
overlaps at $\mathsf{n}=1$. If $|\mathrm{Pol}|\ge2$, $\mathrm{Pol}$ holds two
distinct vectors of $\Dset$, and no vector of $\Dset$ is a multiple of another,
their overlaps having modulus $0$ or $1/\sqrt2$; so two vectors of
$\mathrm{Pol}$ span $\mathbb{C}^2$, and $S_1$, the span of the $\ket{1,a}$,
$a\in\mathrm{Pol}$, is the whole one-photon space. $\ket{H}$ and $\ket{V}$ are
the fixed basis of (W1), whether or not they lie in $\mathrm{Pol}$, so
$\{\ket{1,H},\ket{1,V}\}$ is an orthonormal basis of it, and
$\ket{1,\pm}=(\ket{1,H}\pm\ket{1,V})/\sqrt2$. The four dyads are therefore
$\dyad{1,H}$, $\dyad{1,V}$ and $\tfrac12\big(\dyad{1,H}+\dyad{1,V}
\pm\ketbra{1,H}{1,V}\pm\ketbra{1,V}{1,H}\big)$, which are the stated matrices.

\emph{Two photons.} Write $\ket{2,H},\ket{1_H1_V},\ket{2,V}$ for the
orthonormal number states of the two modes, which span the two-photon space,
three dimensional. Expanding
$(c^{\dagger}_{\pm})^2=\tfrac12\big((c^{\dagger}_H)^2
\pm2c^{\dagger}_Hc^{\dagger}_V+(c^{\dagger}_V)^2\big)$ gives
$\ket{2,\pm}=\tfrac12\ket{2,H}\pm\tfrac1{\sqrt2}\ket{1_H1_V}
+\tfrac12\ket{2,V}$, so the four vectors $\ket{2,a}$, $a\in\Dset$, have the
Gram matrix with diagonal $1$, $\braket{2,H}{2,V}=0$, $\braket{2,+}{2,-}=0$
and $\braket{2,H}{2,\pm}=\braket{2,V}{2,\pm}=\tfrac12$. The four vectors
$r_a\coloneqq v_a/\sqrt2$,
\begin{equation*}
  r_H=\tfrac1{\sqrt2}(1,1,0),\quad r_V=\tfrac1{\sqrt2}(1,-1,0),\quad
  r_+=\tfrac1{\sqrt2}(1,0,1),\quad r_-=\tfrac1{\sqrt2}(1,0,-1)
\end{equation*}
of $\mathbb{C}^3$ are unit vectors with the same Gram matrix --- each inner
product is a sum of three products of coordinates, and
$r_H^{\dagger}r_V=\tfrac12(1-1+0)=0$, $r_+^{\dagger}r_-=\tfrac12(1+0-1)=0$,
$r_H^{\dagger}r_{\pm}=r_V^{\dagger}r_{\pm}=\tfrac12$ --- and any three of them
are linearly independent, the determinant of any three of $v_H,v_V,v_+,v_-$
being $\pm2$. Equal Gram matrices force the same linear relations, a vector
$c$ of coefficients annihilating one family exactly when $c^{\dagger}Gc=0$ for
the common Gram matrix $G$, so the map $\ket{2,a}\mapsto r_a$ extends to a
well-defined linear map $U$ from the span of the $\ket{2,a}$, $a\in\Dset$,
onto $\mathbb{C}^3$ that preserves inner products; by the same relations any
three of the $\ket{2,a}$ are linearly independent, so that span is the whole
two-photon space and $U$ is unitary. If $|\mathrm{Pol}|\ge3$, $S_2$ contains
three of the $\ket{2,a}$ and is therefore the whole two-photon space. The
inverse image under $U$ of the standard basis is then an orthonormal basis of
$S_2$, in which $\dyad{2,a}$ has the entries
$(r_a)_i\overline{(r_a)_j}=\tfrac12(v_a)_i(v_a)_j\in\{0,\pm\tfrac12\}$, the
$v_a$ being real. (c) is (E3) read on the entries of (a) and (b).

(v) Let $0\le\tilde p'\le\tilde p$ pointwise. Then
$0\le\tilde p'(\mathsf{n}\mid\mu)\le\tilde p(\mathsf{n}\mid\mu)\le
P_{\mu}(\mathsf{n})$, so \eqref{eq:weightineq} holds of $\tilde p'$, and
$\tilde w'_{\mu}=1-\sum_{\mathsf{n}\le N_{\mathrm{src}}}
\tilde p'(\mathsf{n}\mid\mu)\ge\tilde w_{\mu}\ge0$ by \eqref{eq:tailweight}.
Every item of Definition~\ref{def:tagged-source} is then satisfied by the primed
data, the sector structure (T3) included, which is formed from the weights: a
tail flag exists exactly where $\tilde w'_{\mu}>0$, a polarisation flag or a
quantum sector exactly where some letter puts positive weight on it, and both
sets are read off $\tilde p'$. So the primed declaration is a tagged source
declaration and Proposition~\ref{prop:tagging} holds of it. Neither inclusion
between the two sector lists holds in general: raising a tail weight from zero
adds a tail flag, and lowering $\tilde p(\mathsf{n}\mid\mu)$ to zero for every
intensity that a polarisation $a$ occurs with removes the sector
$(\mathsf{n},a)$.
\end{proof}

% claim transmitter/decoy/rem:decoy-consumers
\claimsummary{transmitter/decoy/rem:decoy-consumers}{The ensemble form of the consuming declarations reads the decoy declaration, and exactly trusted-source mode, the rule against general constraint entries, which refuses nothing it emits, and the emission's own refusals stand in between. It licenses no restriction of a consumer's feasible set.}
\claimbookmark{transmitter/decoy/rem:decoy-consumers}{Remark}{rem:decoy-consumers}{Consumers of the emitted decoy declaration}
\begin{remark}[Consumers of the emitted decoy declaration]\label{rem:decoy-consumers}
By Corollary~\ref{cor:emitted-decoy}\,(i) and (ii) the decoy declaration is a
source specification in trusted-source mode given by an ensemble presentation,
and its source segment meets the admission conditions of the transmitter inside
the artefact the emission outputs (Requirement~\ref{req:decoy-emission}). Item
(a) says which consumer reads it; what stands between the decoy declaration and
a consumer of that form is exactly (b) to (d), each a named refusal.
\begin{enumerate}[label=(\alph*),leftmargin=2.2em]
\item \emph{The consumer's form.} The protocol compiler's declaration
(\xref{protocol-compiler/def:declaration}) and the MEAT core's
(\xref{meat-core/def:declaration}) each have an ensemble form, which carries
$\{N_x\}$ in the place of Alice's POVM, and
\xref{protocol-compiler/def:eligible} admits an ensemble presentation beside a
declaration of that form, so the decoy declaration has a consumer. A consumer
whose declaration carries $\{M_x\}$ refuses it (Definition~\ref{def:srcert}),
which is a refusal of that consumer's form and not a condition on the decoy
declaration.
\item \emph{Trusted-source mode only.} A non-empty source segment is admitted in
trusted-source mode alone, by the admission rule of \xref{reductions/def:chain},
and the two-sided reduction theorem (\xref{reductions/thm:two-sided}) is stated
for that mode alone; an ensemble form is refused in
\termref{term-untrusted-source-mode}{untrusted-source mode}, the
mode without a marginal or a certificate (Definition~\ref{def:source-spec}), by
the consuming module.
\item \emph{Scalar constraint entries only.} Beside a non-empty source segment
the protocol compiler requires every entry of a consumer's constraint list
$\CList$ to have a real scalar as its right-hand side, and refuses an entry
declared general (the licensed-only rule of
\xref{protocol-compiler/def:chaincond}). The decoy declaration carries no
constraint entry (Corollary~\ref{cor:emitted-decoy}\,(ii)), so the rule refuses
nothing it emits.
\item \emph{The emission's own refusals.} The requirements of this submodule
refuse a declaration before any consumer reads it; the bound
$N_{\mathrm{ent}}\le2$ on the entropy cutoff
(Requirement~\ref{req:entropy-cutoff}) and the exactness of the formed entries
(Requirement~\ref{req:exact-decoy-blocks}) are the two of them a declarer of a
physically interesting source meets first.
\end{enumerate}
One further point is not a refusal but the absence of a licence: the emission
licenses no restriction of any consumer's feasible set. The states reached from
its certificate are block diagonal, but a positive operator constrained only to
have the marginal $\tau_A$ need not be (Corollary~\ref{cor:emitted-decoy}\,(iii)),
and a consumer that restricts its own optimisation does so under a theorem of
its own module (Remark~\ref{rem:boundary}).
\end{remark}
\claimmeta{\metalabel{Uses}\mref{Cor.}{cor:emitted-decoy}, \mref{Req.}{req:decoy-emission}
\metasep newly derived
\metasep accepted}

\section{The physical-source assumption}\label{group-transmitter/decoy/ass:wcp}
\sectionintro{One attestation carries the model to the device: \hyperlink{claim-transmitter/decoy/ass:wcp}{A-WCP} asserts that the physical source emits the declared phase-randomised weak coherent states, exactly and independently in every round. It refines \hyperlink{claim-transmitter/ass:ens}{A-ENS} and takes its place for a specification declared by this model, and an artefact emitted from a tagged declaration carries A-WCP and not A-ENS in its assumption ledger (\hyperlink{claim-transmitter/decoy/req:decoy-emission}{Requirement~\ref*{req:decoy-emission}}). No function of the declared data decides it. \hyperlink{claim-transmitter/decoy/rem:wcp-virtual-model}{Remark~\ref*{rem:wcp-virtual-model}} reads it as making the tagged declaration, with its tagging link and at the declared numbers exactly, a virtual model of the device.}

% claim transmitter/decoy/ass:wcp
\claimsummary{transmitter/decoy/ass:wcp}{The physical source emits exactly the declared phase-randomised weak coherent state of each round's letter, drawn with its declared weight, in the declared modes only and independently across rounds. It ties the tagged declaration to the device, in the place of \hyperlink{claim-transmitter/ass:ens}{A-ENS}.}
\claimbookmark{transmitter/decoy/ass:wcp}{Assumption}{ass:wcp}{Emission of the declared weak-coherent-pulse states}
\claimleadin{%
The model declares its states as mathematical objects, and this assumption
ties the device to them. It refines Assumption~\ref{ass:ens}: for a
specification whose specified ensemble is declared by this \termref{term-source-model}{source
model}, A-WCP is declared in the place of A-ENS; it asserts the emission A-ENS
asserts and adds the conditions this model needs, so it is never the weaker of
the two. Full phase randomisation makes the emitted state diagonal in the
photon number. That no degree of freedom outside the two declared modes
depends on the letter excludes side channels: no timing, spectral or spatial
correlation with the label and no leakage of the label beyond the declared
modes.
}
\SetShortId{A-WCP}
\begin{assumption}[Emission of the declared weak-coherent-pulse states]\label{ass:wcp}
For the \termref{term-tagged-source-declaration}{tagged source declaration} of Definition~\ref{def:tagged-source},
with the cutoffs, the weights and the polarisations it declares, the physical
source satisfies:
(i) in every round it emits the specified state $\rho_x$ of
Definition~\ref{def:wcp-model} of the letter $x=(a,\mu,z)$ drawn in that
round, and nothing else;
(ii) the intensity is the declared $\mu$ exactly, with no fluctuation and no
uncertainty interval;
(iii) the phase randomisation is full, so that the emitted state is the phase
average \eqref{eq:wcpstate};
(iv) the emission is in the two optical modes declared in
Definition~\ref{def:wcp-model}\,(W1), and the polarisation is the declared
$\ket{a}$;
(v) no degree of freedom other than those two modes carries any dependence on
$x$;
(vi) the \termref{term-round-type}{round type} $z$ changes nothing the source emits;
(vii) the letters $x_1,\dots,x_n$ of the $n$ rounds are drawn independently,
each with the weights $p_x$ of Definition~\ref{def:wcp-model}\,(W2);
(viii) given those letters, the systems the source emits in the $n$ rounds
are in the product state $\rho_{x_1}\otimes\dots\otimes\rho_{x_n}$ of the
per-round states;
(ix) neither the letter of a round nor the state emitted in it depends on
what happens in the other rounds, their outcomes and public messages
included.
\end{assumption}
\claimverification{The engine does not verify this assumption; a run that rests on it records it as unverified.}
\claimreading{Your weak-coherent-pulse source emits, in every round, a fully phase-randomised coherent pulse at exactly the declared intensity \ensuremath{\mu} and polarisation, in the two declared optical modes only, and nothing else; the numbers are those of your tagged declaration, exactly as declared. No other property of the pulse (timing, spectrum, spatial shape) depends on the chosen bit, basis, intensity or round type, and the round type (test or generation) changes nothing emitted. Intensities do not fluctuate. Each round's letter is drawn independently with the declared probabilities, and no round's letter or emitted pulse depends on another round, its detection results or its public messages. Intensity fluctuations, imperfect phase randomisation or pattern effects violate this.}
\claimmeta{\metalabel{Uses}\mref{Def.}{def:wcp-model}, \mref{Def.}{def:tagged-source}
\metasep \metalabel{Category}source
\metasep \metalabel{Affects}secrecy
\metasep \metalabel{Refines}\aref{A-ENS}{ass:ens}
\metasep \metalabel{Assumed by}\mref{Cor.}{cor:emitted-decoy}
\metasep \metalabel{Machine key}\iface{wcp_source_model}
\metasep newly derived
\metasep accepted}

% claim transmitter/decoy/rem:wcp-virtual-model
\claimsummary{transmitter/decoy/rem:wcp-virtual-model}{Under \hyperlink{claim-transmitter/decoy/ass:wcp}{A-WCP} the tagged declaration with its tagging link is a virtual model of the device: the states the tagging channel produces are the ones it emits. The assumption carries no arithmetic identity, so it says nothing of a declaration whose numbers were altered.}
\claimbookmark{transmitter/decoy/rem:wcp-virtual-model}{Remark}{rem:wcp-virtual-model}{The tagged declaration as a virtual model of the device}
\begin{remark}[The tagged declaration as a virtual model of the device]\label{rem:wcp-virtual-model}
Assumption~\ref{ass:wcp} is asserted of the \termref{term-tagged-source-declaration}{tagged source declaration}
of Definition~\ref{def:tagged-source}, with the cutoffs, the weights and the
polarisations it declares, which is what the emitted decoy declaration of
Corollary~\ref{cor:emitted-decoy} presents. That declared virtual source,
together with its tagging link, is then a virtual model of the physical
source: the states $\Xi_{\mathrm{tag}}(\rvirt_x)=\rho_x$ that the link's
channel is proved to produce (Proposition~\ref{prop:tagging}\,(ii)) are the
states the device emits, the device being the one
Assumption~\ref{ass:wcp} describes at those intensities and polarisations.
The assumption asserts no arithmetic identity: the source identity of the
link is a theorem about the declared numbers and holds of them exactly
(Corollary~\ref{cor:emitted-decoy}\,(ii)), and if a declared number is
altered in any way, rounded or shifted included, the assumption says nothing
about the result. This contract asserts Assumption~\ref{ass:wcp} of no
particular apparatus.
\end{remark}
\claimmeta{\metalabel{Uses}\mref{Def.}{def:wcp-model}, \mref{Def.}{def:tagged-source}
\metasep newly derived
\metasep accepted}

\section{Emission and refusals}\label{group-transmitter/decoy/req:decoy-emission}
\sectionintro{These requirements bind the engine, the implementation, where it turns a tagged declaration into a decoy declaration. Recall the three bands: photon numbers up to the \termref{term-entropy-cutoff}{entropy cutoff} keep quantum sectors, those up to the \termref{term-source-cutoff}{source cutoff} one flag per polarisation, and the rest sits on the \termref{term-tail-flag}{tail flags}. What the emission outputs is the tagged ensemble as declared, with no number altered and with the entries of its signal dyads formed by the engine. It refuses, each with a named reason, cutoffs that are not integers or are out of order, entropy cutoffs above two, quantum sectors whose signal dyads the construction it carries does not cover, sector weights that are negative or exceed the rational lower bound it forms on their \termref{term-poisson-weight}{Poisson weights} or are presented as equalities, flag families finer than the tagged declaration's, formed entries that are not binary64 numbers, and blocks or weight sums that fail the conditions of an \termref{term-ensemble-presentation}{ensemble presentation}.}

% claim transmitter/decoy/req:decoy-emission
\claimsummary{transmitter/decoy/req:decoy-emission}{What the engine emits from a tagged declaration is the decoy declaration of \hyperlink{claim-transmitter/decoy/cor:emitted-decoy}{Corollary~\ref*{cor:emitted-decoy}} and nothing else, with every declared number unchanged. Its ledger carries A-WCP in the place of A-ENS.}
\claimbookmark{transmitter/decoy/req:decoy-emission}{Requirement}{req:decoy-emission}{Emission of the tagged ensemble as declared}
\claimleadin{%
The source identity holds at the emitted descriptors because the emitted
ensemble is the tagged one with the declared weights
(Corollary~\ref{cor:emitted-decoy}\,(ii)).
}
\begin{requirement}[Emission of the tagged ensemble as declared]\label{req:decoy-emission}
When the engine emits an artefact from a \termref{term-tagged-source-declaration}{tagged source declaration} over
a \termref{term-weak-coherent-pulse-model}{weak-coherent-pulse model}, the artefact is that declaration's decoy
declaration (Corollary~\ref{cor:emitted-decoy}) and nothing else: the engine
declares no retained system, no measurement, no marginal and no certificate,
and it rounds or otherwise alters no declared number. The matched
\termref{term-sector-structure}{sector structure} among the declared data of that corollary is emitted as
the \termref{term-retained-register-sector-map}{retained-register sector map} of
Definition~\ref{def:tagged-source}\,(T5), with
the block kind \iface{proof_sector}, the tag that records a decomposition
licensed by a proof, here source replacement, rather than by the device's own
mode structure;
it names Alice's sectors only as copies of the declared sectors $S_t$, so it
declares no retained system. An artefact emitted from a tagged declaration
carries A-WCP and not A-ENS in its assumption ledger
(Assumption~\ref{ass:wcp}).
\end{requirement}
\claimmeta{\metalabel{Motivated by}\mref{Def.}{def:tagged-source}, \mref{Cor.}{cor:emitted-decoy}, \mref{Ass.}{ass:wcp}
\metasep \metalabel{Used by}\mref{Rem.}{rem:decoy-consumers}
\metasep implemented (effective: accepted)
\metasep \metalabel{Code}\coderef{https://github.com/fraunhofersg/certiqs-key/blob/v0.1.0/src/certiqs_key/transmitter/decoy/emit.py}{certiqs_key.transmitter.decoy.emit:emit_decoy_declaration}}

% claim transmitter/decoy/req:cutoff-order
\claimsummary{transmitter/decoy/req:cutoff-order}{The engine refuses cutoffs unless both are integers and the entropy cutoff lies between zero and the source cutoff, the first condition of a tagged source declaration.}
\claimbookmark{transmitter/decoy/req:cutoff-order}{Requirement}{req:cutoff-order}{Order of the entropy and source cutoffs}
\begin{requirement}[Order of the entropy and source cutoffs]\label{req:cutoff-order}
The engine refuses, with no artefact, a declaration whose cutoffs are not
integers with $0\le N_{\mathrm{ent}}\le N_{\mathrm{src}}$, which is (T1) of
Definition~\ref{def:tagged-source} (\iface{cutoffs_out_of_order}, the one
refusal for a cutoff that is not an integer and for cutoffs out of order).
\end{requirement}
\claimmeta{\metalabel{Motivated by}\mref{Def.}{def:tagged-source}
\metasep implemented
\metasep \metalabel{Code}\coderef{https://github.com/fraunhofersg/certiqs-key/blob/v0.1.0/src/certiqs_key/transmitter/decoy/emit.py}{certiqs_key.transmitter.decoy.emit:emit_decoy_declaration}}

% claim transmitter/decoy/req:entropy-cutoff
\claimsummary{transmitter/decoy/req:entropy-cutoff}{The engine refuses an entropy cutoff of three or more. This is a scope limit and no consequence of the claims, which hold for every entropy cutoff; a basis with exactly declarable dyad entries is exhibited only for the polarisations H, V, + and \ensuremath{-}, at one and at two photons.}
\claimbookmark{transmitter/decoy/req:entropy-cutoff}{Requirement}{req:entropy-cutoff}{Entropy cutoff of at most two photons}
\begin{requirement}[Entropy cutoff of at most two photons]\label{req:entropy-cutoff}
The engine refuses, with no artefact, a declaration with
$N_{\mathrm{ent}}\ge3$ (\iface{entropy_cutoff_above_two}). The refusal is a
scope decision and not a consequence of anything the claims state: every claim
of this submodule is stated for arbitrary $N_{\mathrm{ent}}\le N_{\mathrm{src}}$
and nothing in them fails at $N_{\mathrm{ent}}\ge3$. What is missing is a
construction of the kind Corollary~\ref{cor:emitted-decoy}\,(iv)(b) supplies
for the four polarisations $\Dset$ at one and at two photons: a declared
basis in which the signal dyads have exactly declarable entries. This
repository has no such construction at three photons and above. Admitting such
a sector would mean deciding exact conditions on dyad entries that differ from
the entries of the signal states, so until the construction exists the sector
is refused rather than declared with entries that are not the numbers written.
\end{requirement}
\claimmeta{\metalabel{Motivated by}\mref{Cor.}{cor:emitted-decoy}
\metasep implemented (effective: accepted)
\metasep \metalabel{Code}\coderef{https://github.com/fraunhofersg/certiqs-key/blob/v0.1.0/src/certiqs_key/transmitter/decoy/emit.py}{certiqs_key.transmitter.decoy.emit:emit_decoy_declaration}}

% claim transmitter/decoy/req:signal-dyads
\claimsummary{transmitter/decoy/req:signal-dyads}{The engine forms the entries of the quantum-sector signal dyads itself, by the construction of \hyperlink{claim-transmitter/decoy/cor:emitted-decoy}{Corollary~\ref*{cor:emitted-decoy}} (iv)(b), and takes no dyad or basis from a request. It refuses a quantum sector above the vacuum that the construction does not cover.}
\claimbookmark{transmitter/decoy/req:signal-dyads}{Requirement}{req:signal-dyads}{Signal dyads formed by the exhibited construction}
\claimleadin{%
The source identity at the emitted descriptors
(Corollary~\ref{cor:emitted-decoy}\,(ii)) holds because the emitted blocks are
the tagged ensemble of the model: on a quantum sector their dyads are the
matrices of the projectors $\dyad{\mathsf{n},a}$ in an orthonormal basis of
$S_{\mathsf{n}}$. The polarisation vectors and that basis are objects of the
model whose coordinates are in general not binary64 numbers, so those entries
cannot be declared; they exist as numbers only through a construction, and
Corollary~\ref{cor:emitted-decoy}\,(iv)(b) gives one for polarisations among
the four of $\Dset$ at one and at two photons.
}
\begin{requirement}[Signal dyads formed by the exhibited construction]\label{req:signal-dyads}
The engine forms the entries of the signal dyads of a decoy declaration
itself, and takes neither a dyad nor a basis of a sector from a request. On
the vacuum sector and on every flag sector it forms the number $1$. On a
quantum sector $\mathsf{n}\in\{1,2\}$ of the index set it forms them only where
every polarisation of $\mathrm{Pol}$ is declared as one of the four vectors of
$\Dset$, distinct polarisations as distinct vectors, and
$|\mathrm{Pol}|\ge\mathsf{n}+1$; it then forms the dyad of $\ket{\mathsf{n},a}$
with the entries Corollary~\ref{cor:emitted-decoy}\,(iv)(b) states for it. It
decides of the formed dyads, exactly, that each is a rank-one projector on
$S_{\mathsf{n}}$ (Hermitian, idempotent and of unit trace) and that
$\tr\big[\dyad{\mathsf{n},a}\,\dyad{\mathsf{n},b}\big]
=|\braket{a}{b}|^{2\mathsf{n}}$ for every pair of polarisations of
$\mathrm{Pol}$, the value that Definition~\ref{def:wcp-model}\,(W3) gives by
that corollary and that is $0$, $2^{-\mathsf{n}}$ or $1$ for two vectors of
$\Dset$. It refuses, with no artefact, a declaration whose index set carries
a quantum sector $\mathsf{n}\ge1$ outside this construction, or at which one of
these decisions fails (\iface{quantum_sector_outside_construction}). A
declaration with $N_{\mathrm{ent}}\ge3$ is refused under its own reason
(Requirement~\ref{req:entropy-cutoff}).
\end{requirement}
\claimmeta{\metalabel{Motivated by}\mref{Def.}{def:wcp-model}, \mref{Def.}{def:tagged-source}, \mref{Cor.}{cor:emitted-decoy}, \mref{Req.}{req:entropy-cutoff}
\metasep implemented (effective: accepted)
\metasep \metalabel{Code}\coderef{https://github.com/fraunhofersg/certiqs-key/blob/v0.1.0/src/certiqs_key/transmitter/decoy/emit.py}{certiqs_key.transmitter.decoy.emit:emit_decoy_declaration}}

% claim transmitter/decoy/req:sector-weight-bound
\claimsummary{transmitter/decoy/req:sector-weight-bound}{Each sector weight is decided against a rational lower bound on its Poisson weight that the engine forms from the declared intensity, and a negative weight, one above that bound, or one presented as an equality is refused. No Poisson weight is compared for equality or at a tolerance.}
\claimbookmark{transmitter/decoy/req:sector-weight-bound}{Requirement}{req:sector-weight-bound}{Certified bound on the sector weights}
\claimleadin{%
Proposition~\ref{prop:tagging} needs the inequality \eqref{eq:weightineq} and
nothing more: its identity holds for every family of weights at or below the
\termref{term-poisson-weight}{Poisson weights}, the tail absorbing the difference. The
right-hand side is not a number the engine can compare with, so it compares
with the lower end of the rational enclosure of
Lemma~\ref{lem:poisson-enclosure}, which lies at or below the Poisson weight.
A weight between that lower end and the Poisson weight is refused as well. The
engine fixes the summation index by one rule, so a declarer can compute the
same lower end and declare at or below it; by parts (iii) and (iv) of that
lemma the rule's index gives a lower end at least as large as any smaller
index would, within $2^{-90}\mu^{\mathsf{n}}/\mathsf{n}!$ of the Poisson
weight.
}
\begin{requirement}[Certified bound on the sector weights]\label{req:sector-weight-bound}
The obligation \eqref{eq:weightineq} of Definition~\ref{def:tagged-source}\,(T2),
$0\le\tilde p(\mathsf{n}\mid\mu)\le P_\mu(\mathsf{n})$ for every
$\mathsf{n}\le N_{\mathrm{src}}$ and every intensity $\mu\in\mathrm{Int}$ of
Definition~\ref{def:wcp-model}\,(W2), is an inequality and never an equality:
its right-hand side is not a declared datum, and at a positive intensity,
declared as a binary64 number, it is a transcendental number (at $\mu=0$ it is
$1$ for $\mathsf{n}=0$ and $0$ for $\mathsf{n}\ge1$). For every such
$(\mathsf{n},\mu)$ the engine forms, from the declared intensity $\mu$ and in
exact rational arithmetic, the lower end
$\underline{P}_\mu(\mathsf{n})\coloneqq\max(0,S_{K+1})\,\mu^{\mathsf{n}}/\mathsf{n}!$
of the enclosure of Lemma~\ref{lem:poisson-enclosure}, with $S_{K+1}$ the
partial sum \eqref{eq:expsum} at $K=K_\mu$, the least even integer
$K\ge\mu$ with $t_{K+1}=\mu^{K+1}/(K+1)!\le2^{-90}$, which exists because
$t_{K+1}\to0$. It records $K_\mu$ for every intensity in the artefact, and it
reads no enclosure of $P_\mu(\mathsf{n})$ from the request. By part (iv) of
that lemma no even $K$ with $\mu\le K<K_\mu$ gives a larger lower end, so the
rule refuses no weight that a smaller admissible index would admit. It
refuses, with no artefact, a declaration in which
$\tilde p(\mathsf{n}\mid\mu)<0$ or
$\tilde p(\mathsf{n}\mid\mu)>\underline{P}_\mu(\mathsf{n})$ at some
$(\mathsf{n},\mu)$ (\iface{sector_weight_above_poisson}). Since
$\underline{P}_\mu(\mathsf{n})\le P_\mu(\mathsf{n})$ by that lemma, every
declaration violating \eqref{eq:weightineq} is refused, and every declaration
it admits satisfies \eqref{eq:weightineq}. It refuses, with no artefact, a
declaration presenting the obligation as an equality
$\tilde p(\mathsf{n}\mid\mu)=P_\mu(\mathsf{n})$ to be decided
(\iface{sector_weight_as_equality}). No condition of this submodule compares
$\tilde p(\mathsf{n}\mid\mu)$ with $P_{\mu}(\mathsf{n})$ for equality and none
is decided at a tolerance.
\end{requirement}
\claimmeta{\metalabel{Motivated by}\mref{Def.}{def:wcp-model}, \mref{Def.}{def:tagged-source}, \mref{Lem.}{lem:poisson-enclosure}, \mref{Prop.}{prop:tagging}
\metasep implemented (effective: accepted)
\metasep \metalabel{Code}\coderef{https://github.com/fraunhofersg/certiqs-key/blob/v0.1.0/src/certiqs_key/transmitter/decoy/emit.py}{certiqs_key.transmitter.decoy.emit:emit_decoy_declaration}}

% claim transmitter/decoy/req:shared-flags
\claimsummary{transmitter/decoy/req:shared-flags}{The engine refuses a flag family finer than the one the tagged declaration declares, in particular polarisation flags split by intensity. For that split the refusal keeps the loss of the decoy content from being silent.}
\claimbookmark{transmitter/decoy/req:shared-flags}{Requirement}{req:shared-flags}{Shared polarisation and tail flags}
\begin{requirement}[Shared polarisation and tail flags]\label{req:shared-flags}
The engine refuses, with no artefact, a declaration carrying a flag family finer
than the one Definition~\ref{def:tagged-source}\,(T3) declares
(\iface{flag_family_split_by_intensity}). That family is one
\termref{term-polarisation-flag}{polarisation flag} $f_{\mathsf{n},a}$ for each sector $(\mathsf{n},a)$ of
the index set of (T3), $N_{\mathrm{ent}}<\mathsf{n}\le N_{\mathrm{src}}$, shared
by every intensity and both \termref{term-round-type}{round types}, and one
\termref{term-tail-flag}{tail flag} $f_{a,\mu}$ for each sector $(a,\mu)$ of
that index set, shared by both round types; a finer family is one with separate
flags where (T3) declares one. In particular the engine refuses a family whose
polarisation flags are split by intensity at a photon number
$N_{\mathrm{ent}}<\mathsf{n}\le N_{\mathrm{src}}$. A family split by intensity
still satisfies the source identity but gives up the relation across
intensities that is the decoy content of the declaration, one photon number's
weight at every intensity sitting on one flag
(Remark~\ref{rem:no-intensity-split}), and without the refusal that loss would
be silent.
\end{requirement}
\claimmeta{\metalabel{Motivated by}\mref{Def.}{def:tagged-source}, \mref{Rem.}{rem:no-intensity-split}
\metasep implemented
\metasep \metalabel{Code}\coderef{https://github.com/fraunhofersg/certiqs-key/blob/v0.1.0/src/certiqs_key/transmitter/decoy/emit.py}{certiqs_key.transmitter.decoy.emit:emit_decoy_declaration}}

% claim transmitter/decoy/req:exact-decoy-blocks
\claimsummary{transmitter/decoy/req:exact-decoy-blocks}{Every entry of an emitted block and of each sector sum must be a binary64 number, and the engine refuses rather than rounds one that is not. Rounding would change which ensemble is declared, and so what the source identity is asserted of.}
\claimbookmark{transmitter/decoy/req:exact-decoy-blocks}{Requirement}{req:exact-decoy-blocks}{Binary64 entries of the emitted blocks and sector sums}
\claimleadin{%
The three conditions of an \termref{term-ensemble-presentation}{ensemble presentation}
(Definition~\ref{def:srcert}) are: (E1) every block is Hermitian and positive
semidefinite; (E2) the traces of a letter's blocks sum to $p_x$, and the $p_x$
to one; (E3) for every sector the sum of the blocks over the letters,
accumulated entry by entry in the \exttermref{../foundations/contract.pdf}{term-declared-order}{declared order} of $\Xset$,
\exttermref{../foundations/contract.pdf}{term-rounds-nowhere}{rounds nowhere}. A declarer for which a product is
not representable removes the refusal by declaring on a coarser dyadic grid
(Corollary~\ref{cor:emitted-decoy}\,(iv)(a)).
}
\begin{requirement}[Binary64 entries of the emitted blocks and sector sums]\label{req:exact-decoy-blocks}
The engine decides (E3) of Definition~\ref{def:srcert} on the sums of the
declared blocks of a decoy declaration, accumulated entry by entry in the
declared order. Each entry of a block, a product of declared numbers, the
\termref{term-tail-weight}{tail weight} formed from them and a formed entry of a signal dyad
(Corollary~\ref{cor:emitted-decoy}\,(iv)), is
required to be a binary64 number and is refused, with no artefact, when it is
not (\iface{formed_entry_not_declarable}). Those sums are required to round
nowhere, so that their entries are binary64 numbers, and a declaration at a
sector of which the accumulation rounds fails (E3) and is refused, with no
artefact (\iface{ensemble_presentation_not_admissible}). No entry is ever
rounded, because rounding an entry of a declared block would change which
ensemble is declared and so which ensemble the source identity of
Corollary~\ref{cor:emitted-decoy}\,(ii) is asserted of.
\end{requirement}
\claimmeta{\metalabel{Motivated by}\mref{Def.}{def:srcert}, \mref{Cor.}{cor:emitted-decoy}
\metasep implemented (effective: accepted)
\metasep \metalabel{Code}\coderef{https://github.com/fraunhofersg/certiqs-key/blob/v0.1.0/src/certiqs_key/transmitter/decoy/emit.py}{certiqs_key.transmitter.decoy.emit:emit_decoy_declaration}}

% claim transmitter/decoy/req:decoy-ensemble-conditions
\claimsummary{transmitter/decoy/req:decoy-ensemble-conditions}{The engine decides on the declared blocks that each is positive semidefinite and that a letter's block traces sum to its weight, the weights to one. It refuses a declaration in which either fails.}
\claimbookmark{transmitter/decoy/req:decoy-ensemble-conditions}{Requirement}{req:decoy-ensemble-conditions}{Positivity and weight sums of the emitted blocks}
\begin{requirement}[Positivity and weight sums of the emitted blocks]\label{req:decoy-ensemble-conditions}
The engine decides (E1) and (E2) of Definition~\ref{def:srcert} on the declared
blocks of a decoy declaration, as it decides (E3)
(Requirement~\ref{req:exact-decoy-blocks}). It refuses, with no artefact, a
declared block failing (E1) or a weight sum failing (E2)
(\iface{ensemble_presentation_not_admissible}).
\end{requirement}
\claimmeta{\metalabel{Motivated by}\mref{Def.}{def:srcert}, \mref{Cor.}{cor:emitted-decoy}, \mref{Req.}{req:exact-decoy-blocks}
\metasep implemented (effective: accepted)
\metasep \metalabel{Code}\coderef{https://github.com/fraunhofersg/certiqs-key/blob/v0.1.0/src/certiqs_key/transmitter/decoy/emit.py}{certiqs_key.transmitter.decoy.emit:emit_decoy_declaration}}
\clearpage
\bookmarksetupnext{level=part}
\section*{Interface}\addcontentsline{toc}{section}{Interface}
\sectionintro{The declared inputs and outputs of the module and its submodules: each key with its symbol, when it is present and the claims that fix it. The domains, types and refusal rules are in the interface files themselves.}
\subsection*{Inputs: Transmitter source model (transmitter)}
\noindent{\small From \href{interface.yaml}{\nolinkurl{interface.yaml}}.}\par\smallskip
{\footnotesize\begin{xltabular}{\linewidth}{@{}>{\raggedright\arraybackslash}p{4.4cm}>{\raggedright\arraybackslash}X>{\raggedright\arraybackslash}p{2.9cm}@{}}
\toprule
Key and symbol & Presence & Fixed by\\
\midrule
\endhead
\iface{signal_alphabet}\newline X & always & \mref{Def.}{def:source-spec}\\
\addlinespace
\iface{alice_system}\newline A, d\_A = dim H(A) & always, EXCEPT in an ensemble presentation, where A is not declared but defined (the `ensemble\_presentation` row) & \mref{Def.}{def:source-spec}, \mref{Def.}{def:srcert}\\
\addlinespace
\iface{alice_sector_structure}\newline T, \{Pi\textasciicircum{}A\_t\}, d\_\{A\_t\} & always; THE ONE-SECTOR STRUCTURE IS THE DEFAULT & \mref{Def.}{def:source-spec}, \mref{Def.}{def:tagged-source}\\
\addlinespace
\iface{alice_povm_spec}\newline \{M\_x\}\_\{x in X\} & always, EXCEPT in an ensemble presentation, where \{M\_x\} is not declared but defined and the POVM condition is a theorem (the `ensemble\_presentation` row) & \mref{Def.}{def:source-spec}\\
\addlinespace
\iface{specified_ensemble}\newline \{p\_x, rho\_x\}\_\{x in X\} on A' & always & \mref{Def.}{def:source-spec}, \mref{Def.}{def:wcp-model}\\
\addlinespace
\iface{source_mode_request} & always & \mref{Def.}{def:source-spec}, \mref{Rem.}{rem:da1}\\
\addlinespace
\iface{alice_marginal}\newline tau\_A & trusted mode only; null in untrusted mode; NOT DECLARED in an ensemble presentation, where tau\_A is the transpose of the declared sum (the `ensemble\_presentation` row) & \mref{Def.}{def:source-spec}, \mref{Rem.}{rem:source-mode}, \mref{Rem.}{rem:da1}\\
\addlinespace
\iface{source_replacement_certificate}\newline (tau\_A, Psi, \{p\_x, rho\textasciicircum{}v\_x\}) & trusted mode; demanded only at d\_A \textgreater{}= 2 & \mref{Def.}{def:srcert}, \mref{Def.}{def:source-spec}, \mref{Thm.}{thm:certificate-sufficiency}\\
\addlinespace
\iface{ensemble_presentation}\newline (S\_t)\_t and \{sigma\textasciicircum{}v\_x\} as block lists & the third form of a certificate; trusted mode & \mref{Def.}{def:srcert}, \mref{Lem.}{lem:ensemble-presentation}, \mref{Def.}{def:source-spec}\\
\addlinespace
\iface{source_segment}\newline Ch\_src = (Xi\_1, ..., Xi\_\{L\_s\}) & always & \mref{Def.}{def:source-segment}, \mref{Def.}{def:source-link}\\
\bottomrule
\end{xltabular}}
\subsection*{Outputs: Transmitter source model (transmitter)}
\noindent{\small From \href{interface.yaml}{\nolinkurl{interface.yaml}}.}\par\smallskip
{\footnotesize\begin{xltabular}{\linewidth}{@{}>{\raggedright\arraybackslash}p{4.4cm}>{\raggedright\arraybackslash}X>{\raggedright\arraybackslash}p{2.9cm}@{}}
\toprule
Key and symbol & Presence & Fixed by\\
\midrule
\endhead
\iface{signal_alphabet}\newline X & always & \mref{Def.}{def:source-spec}\\
\addlinespace
\iface{alice_dim}\newline d\_A & always & \mref{Def.}{def:source-spec}\\
\addlinespace
\iface{alice_marginal}\newline tau\_A, or none & always & \mref{Def.}{def:source-spec}, \mref{Rem.}{rem:source-mode}\\
\addlinespace
\iface{alice_marginal_slot_filled} & always & \mref{Req.}{req:marginal-fill}\\
\addlinespace
\iface{alice_povm}\newline \{M\_x\}\_\{x in X\} & always; an empty mapping with an ensemble presentation, where \{M\_x\} is defined by eq:enspres and formed by no one. A consumer in ensemble form reads `alice\_ensemble\_operators` instead, and one whose declaration carries \{M\_x\} refuses the presentation (def:srcert, 'What a consumer reads of a presentation'). & \mref{Def.}{def:source-spec}\\
\addlinespace
\iface{alice_ensemble_operators}\newline \{N\textasciicircum{}(t)\_x\} & with an ensemble presentation only & \mref{Lem.}{lem:ensemble-presentation}, \mref{Def.}{def:srcert}\\
\addlinespace
\iface{virtual_transmitted_system}\newline A\textasciicircum{}v & always & \mref{Def.}{def:source-segment}\\
\addlinespace
\iface{soundness_registry} & always & \mref{Def.}{def:source-segment}, \mref{Def.}{def:source-link}\\
\addlinespace
\iface{virtual_ensemble}\newline \{p\_x, rho\textasciicircum{}v\_x\}\_\{x in X\} & always; with an ensemble presentation its `states` are an empty mapping, rho\textasciicircum{}v\_x = sigma\textasciicircum{}v\_x / p\_x being a quotient of declared numbers that need not be binary64 and is formed by no one (`virtual\_ensemble\_weighted\_members`). & \mref{Def.}{def:source-segment}, \mref{Def.}{def:srcert}\\
\addlinespace
\iface{virtual_ensemble_weighted_members}\newline \{sigma\textasciicircum{}v\_x\}, per virtual sector & with a block-form certificate or an ensemble presentation only & \mref{Def.}{def:srcert}, \mref{Lem.}{lem:ensemble-presentation}\\
\addlinespace
\iface{alice_sector_map}\newline T, \{Pi\textasciicircum{}A\_t\}, and the matched \{Pi\textasciicircum{}v\_t\} where a link emitted one & always; the one-sector structure by default & \mref{Def.}{def:source-spec}, \mref{Def.}{def:source-link}, \mref{Lem.}{lem:matched-pinching}\\
\addlinespace
\iface{source_segment}\newline Ch\_src & always; EMPTY unless a submodule supplied an admitted link & \mref{Def.}{def:source-segment}, \mref{Def.}{def:source-link}\\
\addlinespace
\iface{source_replacement_license} & trusted mode at d\_A \textgreater{}= 2 only & \mref{Def.}{def:srcert}, \mref{Ass.}{ass:src}, \mref{Rem.}{rem:da1}, \mref{Req.}{req:licence-scope}, \mref{Req.}{req:attestation-record}, \mref{Rem.}{rem:boundary}\\
\addlinespace
\iface{certificate_form} & on every record but the degenerate one; an absent key means `pure` & \mref{Def.}{def:srcert}\\
\addlinespace
\iface{assumption_ledger} & always & \mref{Ass.}{ass:ens}, \mref{Ass.}{ass:seq}, \mref{Ass.}{ass:src}, \mref{Rem.}{rem:da1}, \mref{Def.}{def:source-spec}, \mref{Def.}{def:source-link}, \mref{Def.}{def:source-segment}\\
\addlinespace
\iface{assumption_bindings} & always & \mref{Ass.}{ass:ens}, \mref{Ass.}{ass:seq}, \mref{Ass.}{ass:src}, \mref{Ass.}{ass:wcp}\\
\addlinespace
\iface{mathematical_spec_version} & always & ---\\
\bottomrule
\end{xltabular}}
\subsection*{Inputs: Decoy-state sources and the tagging channel}
\noindent{\small From \href{decoy/interface.yaml}{\nolinkurl{decoy/interface.yaml}}.}\par\smallskip
{\footnotesize\begin{xltabular}{\linewidth}{@{}>{\raggedright\arraybackslash}p{4.4cm}>{\raggedright\arraybackslash}X>{\raggedright\arraybackslash}p{2.9cm}@{}}
\toprule
Key and symbol & Presence & Fixed by\\
\midrule
\endhead
\iface{polarisation_set}\newline Pol, with \textbar{}a\textgreater{} in C\textasciicircum{}2 per a & always & \mref{Def.}{def:wcp-model}, \mref{Cor.}{cor:emitted-decoy}, \mref{Req.}{req:signal-dyads}\\
\addlinespace
\iface{intensity_set}\newline Int & always & \mref{Def.}{def:wcp-model}\\
\addlinespace
\iface{signal_alphabet}\newline X, x = (a, mu, z) & always & \mref{Def.}{def:wcp-model}, \mref{Def.}{def:source-spec}\\
\addlinespace
\iface{entropy_cutoff}\newline N\_ent & always & \mref{Def.}{def:tagged-source}, \mref{Req.}{req:entropy-cutoff}, \mref{Rem.}{rem:cutoffs}\\
\addlinespace
\iface{source_cutoff}\newline N\_src & always & \mref{Def.}{def:tagged-source}, \mref{Rem.}{rem:cutoffs}\\
\addlinespace
\iface{sector_weights}\newline ptilde(n \textbar{} mu), 0 \textless{}= n \textless{}= N\_src & always & \mref{Def.}{def:tagged-source}, \mref{Lem.}{lem:poisson-enclosure}, \mref{Req.}{req:sector-weight-bound}\\
\addlinespace
\iface{weighted_blocks}\newline sigma\textasciicircum{}\{v,(t)\}\_x = p\_x w\_\{x,t\} \textbar{}phi\_\{x,t\}\textgreater{}\textless{}phi\_\{x,t\}\textbar{} & always & \mref{Prop.}{prop:decoy-cert}, \mref{Cor.}{cor:emitted-decoy}, \mref{Req.}{req:signal-dyads}, \mref{Def.}{def:srcert}\\
\addlinespace
\iface{sector_structure}\newline T, \{Pi\textasciicircum{}v\_t\}, the declared basis of each S\_t & always & \mref{Def.}{def:tagged-source}, \mref{Prop.}{prop:decoy-cert}, \mref{Req.}{req:decoy-emission}\\
\bottomrule
\end{xltabular}}
\subsection*{Outputs: Decoy-state sources and the tagging channel}
\noindent{\small From \href{decoy/interface.yaml}{\nolinkurl{decoy/interface.yaml}}.}\par\smallskip
{\footnotesize\begin{xltabular}{\linewidth}{@{}>{\raggedright\arraybackslash}p{4.4cm}>{\raggedright\arraybackslash}X>{\raggedright\arraybackslash}p{2.9cm}@{}}
\toprule
Key and symbol & Presence & Fixed by\\
\midrule
\endhead
\iface{source_specification}\newline S\_src & always & \mref{Cor.}{cor:emitted-decoy}, \mref{Def.}{def:source-spec}, \mref{Def.}{def:srcert}\\
\addlinespace
\iface{virtual_ensemble}\newline \{sigma\textasciicircum{}v\_x\} on A\textasciicircum{}v, by weighted blocks & always & \mref{Def.}{def:tagged-source}, \mref{Prop.}{prop:decoy-cert}, \mref{Def.}{def:srcert}\\
\addlinespace
\iface{source_segment}\newline Ch\_src = (Xi\_tag) & always & \mref{Def.}{def:tagged-source}, \mref{Cor.}{cor:emitted-decoy}\\
\addlinespace
\iface{link_obligation_record} & always & \mref{Def.}{def:source-link}, \mref{Def.}{def:source-segment}, \mref{Cor.}{cor:emitted-decoy}\\
\addlinespace
\iface{alice_sector_map}\newline (\{Pi\textasciicircum{}A\_t\}, \{Pi\textasciicircum{}v\_t\}) & always & \mref{Def.}{def:tagged-source}, \mref{Cor.}{cor:emitted-decoy}\\
\addlinespace
\iface{assumption_ledger} & always & \mref{Ass.}{ass:wcp}, \mref{Req.}{req:decoy-emission}, \mref{Def.}{def:source-spec}, \mref{Cor.}{cor:emitted-decoy}\\
\addlinespace
\iface{mathematical_spec_version} & always & ---\\
\bottomrule
\end{xltabular}}
\clearpage
\bookmarksetupnext{level=part}
\section*{Claim index}\addcontentsline{toc}{section}{Claim index}
{\small\begin{xltabular}{\linewidth}{@{}>{\raggedright\arraybackslash}p{2.7cm}>{\raggedright\arraybackslash}Xr@{}}
\toprule
Claim & Title and identifier & Page\\
\midrule
\endhead
\mref{Definition}{def:source-spec} & Source specification\newline{\footnotesize\color{metaGrey}\claimid{transmitter/def:source-spec}} & \pageref{def:source-spec}\\
\addlinespace[2pt]
\mref{Remark}{rem:source-mode} & The two source modes\newline{\footnotesize\color{metaGrey}\claimid{transmitter/rem:source-mode}} & \pageref{rem:source-mode}\\
\addlinespace[2pt]
\mref{Definition}{def:source-link} & Source link\newline{\footnotesize\color{metaGrey}\claimid{transmitter/def:source-link}} & \pageref{def:source-link}\\
\addlinespace[2pt]
\mref{Definition}{def:source-segment} & Source segment and its interface conditions\newline{\footnotesize\color{metaGrey}\claimid{transmitter/def:source-segment}} & \pageref{def:source-segment}\\
\addlinespace[2pt]
\mref{Remark}{rem:segment-scope} & Scope of the source-segment conditions\newline{\footnotesize\color{metaGrey}\claimid{transmitter/rem:segment-scope}} & \pageref{rem:segment-scope}\\
\addlinespace[2pt]
\mref{Fact}{fact:sr} & Source replacement for prepare-and-measure protocols\newline{\footnotesize\color{metaGrey}\claimid{transmitter/fact:sr}} & \pageref{fact:sr}\\
\addlinespace[2pt]
\mref{Remark}{rem:import} & Scope of the imported source replacement\newline{\footnotesize\color{metaGrey}\claimid{transmitter/rem:import}} & \pageref{rem:import}\\
\addlinespace[2pt]
\mref{Definition}{def:srcert} & Source-replacement certificate\newline{\footnotesize\color{metaGrey}\claimid{transmitter/def:srcert}} & \pageref{def:srcert}\\
\addlinespace[2pt]
\mref{Lemma}{lem:ensemble-presentation} & Certificate defined by an ensemble presentation\newline{\footnotesize\color{metaGrey}\claimid{transmitter/lem:ensemble-presentation}} & \pageref{lem:ensemble-presentation}\\
\addlinespace[2pt]
\mref{Theorem}{thm:certificate-sufficiency} & Checkability and sufficiency of the source-replacement certificate\newline{\footnotesize\color{metaGrey}\claimid{transmitter/thm:certificate-sufficiency}} & \pageref{thm:certificate-sufficiency}\\
\addlinespace[2pt]
\mref{Remark}{rem:xreg} & Necessity and cost of the outcome-register hypothesis\newline{\footnotesize\color{metaGrey}\claimid{transmitter/rem:xreg}} & \pageref{rem:xreg}\\
\addlinespace[2pt]
\mref{Corollary}{cor:noncanonical} & Certificates beyond the canonical form\newline{\footnotesize\color{metaGrey}\claimid{transmitter/cor:noncanonical}} & \pageref{cor:noncanonical}\\
\addlinespace[2pt]
\mref{Remark}{rem:certificate-forms} & Purity and exact declarability of the certificate forms\newline{\footnotesize\color{metaGrey}\claimid{transmitter/rem:certificate-forms}} & \pageref{rem:certificate-forms}\\
\addlinespace[2pt]
\mref{Lemma}{lem:matched-pinching} & Matched pinching of a block certificate\newline{\footnotesize\color{metaGrey}\claimid{transmitter/lem:matched-pinching}} & \pageref{lem:matched-pinching}\\
\addlinespace[2pt]
\mref{Remark}{rem:boundary} & Scope of a verified certificate\newline{\footnotesize\color{metaGrey}\claimid{transmitter/rem:boundary}} & \pageref{rem:boundary}\\
\addlinespace[2pt]
\mref{Remark}{rem:da1} & The one-dimensional retained system\newline{\footnotesize\color{metaGrey}\claimid{transmitter/rem:da1}} & \pageref{rem:da1}\\
\addlinespace[2pt]
\mref{Assumption}{ass:ens} & Emission of the specified ensemble\newline{\footnotesize\color{metaGrey}\claimid{transmitter/ass:ens}} & \pageref{ass:ens}\\
\addlinespace[2pt]
\mref{Assumption}{ass:seq} & Sequential round structure\newline{\footnotesize\color{metaGrey}\claimid{transmitter/ass:seq}} & \pageref{ass:seq}\\
\addlinespace[2pt]
\mref{Assumption}{ass:src} & Retention of the system A in trusted-source mode\newline{\footnotesize\color{metaGrey}\claimid{transmitter/ass:src}} & \pageref{ass:src}\\
\addlinespace[2pt]
\mref{Requirement}{req:licence-scope} & Certificate and attestation demanded only at $d_A\ge2$\newline{\footnotesize\color{metaGrey}\claimid{transmitter/req:licence-scope}} & \pageref{req:licence-scope}\\
\addlinespace[2pt]
\mref{Requirement}{req:marginal-fill} & The unique marginal served at $d_A=1$\newline{\footnotesize\color{metaGrey}\claimid{transmitter/req:marginal-fill}} & \pageref{req:marginal-fill}\\
\addlinespace[2pt]
\mref{Requirement}{req:failing-certificate} & Refusal of a certificate that fails a decided condition\newline{\footnotesize\color{metaGrey}\claimid{transmitter/req:failing-certificate}} & \pageref{req:failing-certificate}\\
\addlinespace[2pt]
\mref{Requirement}{req:marginal-decided} & Direct decision of the declared marginal\newline{\footnotesize\color{metaGrey}\claimid{transmitter/req:marginal-decided}} & \pageref{req:marginal-decided}\\
\addlinespace[2pt]
\mref{Requirement}{req:attestation-record} & Pass-through of the retention attestation\newline{\footnotesize\color{metaGrey}\claimid{transmitter/req:attestation-record}} & \pageref{req:attestation-record}\\
\addlinespace[2pt]
\mref{Definition}{def:wcp-model} & Phase-randomised weak-coherent-pulse source model\newline{\footnotesize\color{metaGrey}\claimid{transmitter/decoy/def:wcp-model}} & \pageref{def:wcp-model}\\
\addlinespace[2pt]
\mref{Definition}{def:tagged-source} & Tagged virtual source: cutoffs, flags and sector map\newline{\footnotesize\color{metaGrey}\claimid{transmitter/decoy/def:tagged-source}} & \pageref{def:tagged-source}\\
\addlinespace[2pt]
\mref{Lemma}{lem:poisson-enclosure} & Rational enclosure of the Poisson weights\newline{\footnotesize\color{metaGrey}\claimid{transmitter/decoy/lem:poisson-enclosure}} & \pageref{lem:poisson-enclosure}\\
\addlinespace[2pt]
\mref{Remark}{rem:cutoffs} & Entropy, source and receiver cutoffs\newline{\footnotesize\color{metaGrey}\claimid{transmitter/decoy/rem:cutoffs}} & \pageref{rem:cutoffs}\\
\addlinespace[2pt]
\mref{Proposition}{prop:tagging} & Tagging channel for weights at or below the Poisson weights\newline{\footnotesize\color{metaGrey}\claimid{transmitter/decoy/prop:tagging}} & \pageref{prop:tagging}\\
\addlinespace[2pt]
\mref{Remark}{rem:no-intensity-split} & Polarisation flags shared across intensities\newline{\footnotesize\color{metaGrey}\claimid{transmitter/decoy/rem:no-intensity-split}} & \pageref{rem:no-intensity-split}\\
\addlinespace[2pt]
\mref{Proposition}{prop:decoy-cert} & Ensemble presentation of the tagged source\newline{\footnotesize\color{metaGrey}\claimid{transmitter/decoy/prop:decoy-cert}} & \pageref{prop:decoy-cert}\\
\addlinespace[2pt]
\mref{Corollary}{cor:emitted-decoy} & The emitted decoy declaration\newline{\footnotesize\color{metaGrey}\claimid{transmitter/decoy/cor:emitted-decoy}} & \pageref{cor:emitted-decoy}\\
\addlinespace[2pt]
\mref{Remark}{rem:decoy-consumers} & Consumers of the emitted decoy declaration\newline{\footnotesize\color{metaGrey}\claimid{transmitter/decoy/rem:decoy-consumers}} & \pageref{rem:decoy-consumers}\\
\addlinespace[2pt]
\mref{Assumption}{ass:wcp} & Emission of the declared weak-coherent-pulse states\newline{\footnotesize\color{metaGrey}\claimid{transmitter/decoy/ass:wcp}} & \pageref{ass:wcp}\\
\addlinespace[2pt]
\mref{Remark}{rem:wcp-virtual-model} & The tagged declaration as a virtual model of the device\newline{\footnotesize\color{metaGrey}\claimid{transmitter/decoy/rem:wcp-virtual-model}} & \pageref{rem:wcp-virtual-model}\\
\addlinespace[2pt]
\mref{Requirement}{req:decoy-emission} & Emission of the tagged ensemble as declared\newline{\footnotesize\color{metaGrey}\claimid{transmitter/decoy/req:decoy-emission}} & \pageref{req:decoy-emission}\\
\addlinespace[2pt]
\mref{Requirement}{req:cutoff-order} & Order of the entropy and source cutoffs\newline{\footnotesize\color{metaGrey}\claimid{transmitter/decoy/req:cutoff-order}} & \pageref{req:cutoff-order}\\
\addlinespace[2pt]
\mref{Requirement}{req:entropy-cutoff} & Entropy cutoff of at most two photons\newline{\footnotesize\color{metaGrey}\claimid{transmitter/decoy/req:entropy-cutoff}} & \pageref{req:entropy-cutoff}\\
\addlinespace[2pt]
\mref{Requirement}{req:signal-dyads} & Signal dyads formed by the exhibited construction\newline{\footnotesize\color{metaGrey}\claimid{transmitter/decoy/req:signal-dyads}} & \pageref{req:signal-dyads}\\
\addlinespace[2pt]
\mref{Requirement}{req:sector-weight-bound} & Certified bound on the sector weights\newline{\footnotesize\color{metaGrey}\claimid{transmitter/decoy/req:sector-weight-bound}} & \pageref{req:sector-weight-bound}\\
\addlinespace[2pt]
\mref{Requirement}{req:shared-flags} & Shared polarisation and tail flags\newline{\footnotesize\color{metaGrey}\claimid{transmitter/decoy/req:shared-flags}} & \pageref{req:shared-flags}\\
\addlinespace[2pt]
\mref{Requirement}{req:exact-decoy-blocks} & Binary64 entries of the emitted blocks and sector sums\newline{\footnotesize\color{metaGrey}\claimid{transmitter/decoy/req:exact-decoy-blocks}} & \pageref{req:exact-decoy-blocks}\\
\addlinespace[2pt]
\mref{Requirement}{req:decoy-ensemble-conditions} & Positivity and weight sums of the emitted blocks\newline{\footnotesize\color{metaGrey}\claimid{transmitter/decoy/req:decoy-ensemble-conditions}} & \pageref{req:decoy-ensemble-conditions}\\
\bottomrule
\end{xltabular}}

\clearpage
\bookmarksetupnext{level=part}
\printindex

\clearpage\phantomsection\bookmarksetupnext{level=part}%
\addcontentsline{toc}{section}{References}
\bibliographystyle{plain}
\bibliography{../../../refs/references}

\end{document}
